\documentclass[11pt]{article}
\usepackage[T1]{fontenc}
\usepackage{lmodern,amsmath,amssymb,amsthm,microtype,hyperref}
\usepackage[margin=0.95in]{geometry}
\hypersetup{colorlinks=true,linkcolor=black,urlcolor=blue,citecolor=black}
\newtheorem{theorem}{Theorem}
\newtheorem{corollary}[theorem]{Corollary}
\newtheorem{proposition}[theorem]{Proposition}
\newtheorem{lemma}[theorem]{Lemma}
\title{Analytic Estimates and Source-Level Certificates\\for a Terminal-Core Navier--Stokes Expansion}
\author{Tony Newton\\Newton Astro Labs\thanks{Newton Astro Labs is the trading name under which Tony Newton conducts independent computational research in the United Kingdom; it is not a limited company.}\\London, UK\\\texttt{tony.newton79@gmail.com}}
\date{}
\begin{document}
\maketitle
\begin{abstract}
We reassess the proposed unforced terminal-core Navier--Stokes expansion
against the source's pulse residuals and forced global theorem. Existing
algebra and interval-comparison methods reduce parts of the problem,
but no complete source-derived pulse intervals have been produced.
We prove an exact centerline precision requirement and a polynomial
cancellation formula suitable for high-accuracy interval evaluation.
A strengthened fail-closed verifier distinguishes actual file-hash
replay from proof that the intervals enclose the prescribed smooth
fields. The full pulse stress cone, Appendix-B contraction, unforced
correction hierarchy and physical zero-force residual remain open.
The proposed unforced Theorem~1.1 is not established.
\end{abstract}

\section{The common weighted space}
The central analytic question in \cite{newton2026} is whether all orders of a formal inner correction hierarchy can be estimated on one terminal similarity core. For $R,\rho>0$, let
\begin{equation}\label{eq:weight}
 w_{a,b}=\frac{R^{-a}\rho^{-b}\,b!\binom{2a+b}{b}}{(a+1)^2(b+1)^2},
 \qquad a,b\in\mathbb N_0.
\end{equation}
For $F(X,\eta)=\sum_{a\ge0}F_a(\eta)X^a$ with analytic coefficient functions on a fixed parameter domain $E$, put
\begin{equation}\label{eq:norm}
 \|F\|_{R,\rho;E}=\sup_{a,b\ge0}\ \sup_{\eta\in E}
 \frac{|\partial_\eta^b F_a(\eta)|}{w_{a,b}},
\end{equation}
whenever the right-hand side is finite. The estimate below applies equally to an algebra of complex analytic germs provided the chosen coefficient norm is finite. The radial regular-singular inverse is defined coefficientwise by
\begin{equation}\label{eq:radinv}
 X(J_\nu F)_{XX}+\nu(J_\nu F)_X=F,\quad (J_\nu F)(0,\eta)=0,
 \qquad (J_\nu F)_{a+1}=\frac{F_a}{(a+1)(a+\nu)},
\end{equation}
with $\nu=1$ or $2$. We use the ideal $I_m=\{F:F_a=0\text{ for }a<m\}$ with the \emph{induced} norm $\|F\|_{I_m}:=\|F\|_{R,\rho;E}$.

\section{Three inverse channels}
For $k\in\{0,1,2\}$ define the normalized weight ratio
\begin{equation}\label{eq:ratio_def}
 T_k(a,b,\nu):=\frac{w_{a,b+k}}{(a+1)(a+\nu)w_{a+1,b}}.
\end{equation}
Direct factorial cancellation yields, for $a,b\in\mathbb N_0$ and $\nu\ge1$,
\begin{align}
 T_0&=R\frac{2(a+2)^2(2a+1)}{(a+1)^2(a+\nu)(2a+b+1)(2a+b+2)},\label{eq:t0}\\
 T_1&=\frac R\rho\frac{2(a+2)^2(2a+1)(b+1)^2}{(a+1)^2(a+\nu)(2a+b+2)(b+2)^2},\label{eq:t1}\\
 T_2&=\frac R{\rho^2}\frac{2(a+2)^2(2a+1)(b+1)^2}{(a+1)^2(a+\nu)(b+3)^2}.\label{eq:t2}
\end{align}
These identities match the two-derivative multiplier in equation (45) of \cite{newton2026}; the zero- and one-derivative ratios are included here because the linear current-order channels contain at most one $\eta$-derivative.

\begin{theorem}[Radial valuation gains]\label{thm:ratios}
For all integers $a,b\ge0$, all $\nu\ge1$, and $R,\rho>0$,
\begin{equation}\label{eq:bounds}
 0\le T_0\le \frac{4R}{(a+1)^2},\qquad
 0\le T_1\le \frac{8R}{\rho(a+1)},\qquad
 0\le T_2\le \frac{8R}{\rho^2}.
\end{equation}
The constant $8$ in the last inequality is optimal as a uniform bound over $a,b\in\mathbb N_0$ and $\nu\in\{1,2\}$.
\end{theorem}
\begin{proof}
For the first estimate, both factors in $(2a+b+1)(2a+b+2)$ are nondecreasing in $b$. Hence
\[
 T_0\le R\frac{(a+2)^2}{(a+1)^3(a+\nu)}
 \le R\frac{(a+2)^2}{(a+1)^4}\le\frac{4R}{(a+1)^2},
\]
where $a+\nu\ge a+1$ and $a+2\le2(a+1)$. For the second, we use $(2a+1)/(2a+b+2)\le1$, $(b+1)^2/(b+2)^2\le1$, and $(a+2)^2/(a+1)^2\le4$; thus
\[
 T_1\le\frac{8R}{\rho(a+\nu)}\le\frac{8R}{\rho(a+1)}.
\]
For the third, the polynomial identity
\[
 4(a+1)^3-(a+2)^2(2a+1)=a^2(2a+3)\ge0
\]
gives
\[
 T_2\le\frac{8R}{\rho^2}\frac{(b+1)^2}{(b+3)^2}
 \le\frac{8R}{\rho^2}.
\]
At $a=0$, $\nu=1$ this last ratio is exactly $8R(b+1)^2/[\rho^2(b+3)^2]$, which approaches $8R/\rho^2$ as $b\to\infty$.
\end{proof}

\begin{corollary}[Ambient weighted inverse estimates]\label{cor:ambient}
For each $m\ge0$, for $F\in I_m$ on which the indicated operations are defined,
\begin{align}
 \|J_\nu F\|_{I_{m+1}}&\le\frac{4R}{(m+1)^2}\|F\|_{I_m},\label{eq:zero}\\
 \|J_\nu\partial_\eta F\|_{I_{m+1}}&\le\frac{8R}{\rho(m+1)}\|F\|_{I_m},\label{eq:one}\\
 \|J_\nu\partial_\eta^2F\|_{I_{m+1}}&\le\frac{8R}{\rho^2}\|F\|_{I_m}.\label{eq:two}
\end{align}
In particular, with $D_X=X\partial_X$ and $\lambda_n=2hn$ for $h>0$, one has for $F\in I_n$,
\begin{equation}\label{eq:shift}
 \big\|J_\nu(\lambda_n F+D_XF+\partial_\eta F)\big\|_{I_{n+1}}
 \le \frac{8hR+4R+8R/\rho}{n+1}\,\|F\|_{I_n}.
\end{equation}
\end{corollary}
\begin{proof}
For each output pair $(a+1,b)$ the corresponding weighted coefficient is bounded by $T_k(a,b,\nu)\|F\|_{I_m}$; taking suprema proves \eqref{eq:zero}--\eqref{eq:two}. The multiplier $D_X$ contributes a factor $a$ to the $a$th coefficient. Thus $aT_0\le4Ra/(a+1)^2\le4R/(n+1)$ for $a\ge n$. Also $\lambda_nT_0\le8hRn/(n+1)^2\le8hR/(n+1)$. Applying the triangle inequality with \eqref{eq:one} gives \eqref{eq:shift}.
\end{proof}

\section{Why the ideal norm must be specified}
Writing $I_m=X^m\mathcal B$ leaves open two inequivalent choices of norm. The \emph{induced} norm is the restriction of \eqref{eq:norm} to the series having no radial terms of degree less than $m$, and is the choice used in Corollary~\ref{cor:ambient}. The alternative \emph{quotient} convention identifies $X^m f$ with $f\in\mathcal B$ and assigns $\|X^mf\|_{I_m}^{\mathrm{quot}}=\|f\|_{\mathcal B}$.

\begin{proposition}[Failure of two-derivative control in quotient ideal norms]\label{prop:quotient}
Fix $m\ge0$, $R,\rho>0$ and $\nu>0$. Under the quotient convention, $J_\nu\partial_\eta^2:I_m\to I_{m+1}$ is unbounded.
\end{proposition}
\begin{proof}
For each $k>0$, let $f_k(\eta)=e^{k\eta}$ on a fixed compact parameter interval and put $F_k(X,\eta)=X^m f_k(\eta)$. The factorial growth in \eqref{eq:weight} makes $\|f_k\|_{\mathcal B}$ finite and nonzero for every fixed $k$. Directly from \eqref{eq:radinv},
\[
 J_\nu\partial_\eta^2F_k
 =\frac{k^2}{(m+1)(m+\nu)}X^{m+1}f_k.
\]
The ratio of output and input \emph{quotient} norms is therefore $k^2/[(m+1)(m+\nu)]$, which tends to infinity with $k$.
\end{proof}

\noindent This result does not contradict the induced-norm estimate \eqref{eq:two}. Its importance is expository and structural: bounds used to pass from the radial valuation ladder to an analytic space must specify which of these norms is intended. Moreover, \eqref{eq:two} cannot be strengthened to a decay bound of order $1/m$ without additional structure; in the normalized coefficient calculation \eqref{eq:t2}, taking $b$ large and then $a\to\infty$ with $\nu$ fixed leaves a positive limiting multiplier $4R/\rho^2$. The delayed two-derivative channel must consequently remain separate from the one-derivative current feedback.

\section{Relation to the coupled current-order problem}
The coefficient equations of \cite{newton2026} write the current pair $w_n=(\varphi_n,U_n)$ as $w_n=S_nf_n+K_nw_n$, where $K_n$ consists of the four operators in equations (47)--(50). Estimates \eqref{eq:bounds}--\eqref{eq:shift} quantify the scalar radial inverse in those expressions. They do not by themselves bound the map $U_n\mapsto V_n[U_n]$, multiplication by the leading profiles, or the pressure primitive and its shifted axial derivative. Those maps must be constructed and estimated in the same induced weighted norm before a bound of the form $\|K_n\|\le C/n$ can be used.

For the sum norm $\|(\varphi,U)\|=\|\varphi\|+\|U\|$, a sufficient operator-norm estimate follows from four concrete block bounds. If
\[
 \|K_n^{ij}\|\le\frac{c_{ij}}{n+1}\qquad (i,j\in\{\varphi,U\}),
\]
then the triangle inequality gives
\begin{equation}\label{eq:matrix}
 \|K_n\|\le\frac{1}{n+1}
 \max\{c_{\varphi\varphi}+c_{U\varphi},\
 c_{\varphi U}+c_{UU}\}.
\end{equation}
The estimate \eqref{eq:matrix} explains the exact remaining verification task: the coefficients $c_{ij}$ must be derived from the four actual PDE channels, not assigned as hypotheses. Only then do Neumann inversion, the quadratic coefficient majorant, normal convergence, and physical PDE reconstruction become available as subsequent steps.

The pointwise weight identity and sharp two-derivative inequality have already been checked in Lean 4. The one-derivative inequality and the integrated Banach-space consequences require separate kernel certificates; the rational regression program checks the three ratio formulas and their bounds on an extensive finite grid but is not a proof for all orders. The reproducibility materials are maintained at \url{https://github.com/tonynewton79-web/unforced-navier-stokes-lean/pull/1}.


\section{Composite derivatives and the four current-order blocks}
The estimates of Theorem~\ref{thm:ratios} must be applied to complete
operators, not to naked derivatives. Indeed, in the induced coefficient
norm \eqref{eq:norm}, the series $F_m=w_{m,0}X^m$ has norm one, whereas
$\|D_XF_m\|=m$. Likewise $\partial_\eta$ is unbounded, as shown by
coefficients $e^{ik\eta}$. Consequently a proof of the current-order
feedback bound cannot treat $T_{a,n}$ or $\partial_\eta A_X$ as bounded
before the regular-singular inverse is applied.

Write $\mathcal B$ for the analytic coefficient space with the induced
norm and $I_m\subset\mathcal B$ for its radial-valuation ideals.
We use the bounded multiplication estimate
$\|fg\|\le Q\|f\|\|g\|$, $Q=16\pi^4/9$, from the original coefficient
algebra argument. The commutators
\[
 g\partial_\eta F=\partial_\eta(gF)-(\partial_\eta g)F,\qquad
 gD_XF=D_X(gF)-(D_Xg)F
\]
and Corollary~\ref{cor:ambient} give, when the indicated multipliers
and derivatives belong to $\mathcal B$,
\begin{align}
 \|J_\nu(gF)\|&\le\frac{4RQ\|g\|}{(m+1)^2}\|F\|,\label{eq:cmp0}\\
 \|J_\nu(g\partial_\eta F)\|&\le RQ
 \left(\frac{8\|g\|}{\rho(m+1)}
 +\frac{4\|\partial_\eta g\|}{(m+1)^2}\right)\|F\|,\label{eq:cmp1}\\
 \|J_\nu(gD_XF)\|&\le RQ
 \left(\frac{4\|g\|}{m+1}
 +\frac{4\|D_Xg\|}{(m+1)^2}\right)\|F\|.\label{eq:cmpX}
\end{align}
Direct coefficient division proves
\[
 \|A_XF\|_{I_m}\le\frac{\|F\|_{I_m}}{m+1},\qquad
 \left\|\int_0^X F(x,\eta)\,dx\right\|_{I_{m+1}}
 \le\frac{4R}{m+1}\|F\|_{I_m}.
\]
These are statements about a concrete induced analytic norm, not an
unquantified order ledger.

Use the leading profiles $\varphi_0,U_0,V_0$ and the notation
$A=1/2+h$, $D=1/2-h$, $d=1-\eta^2$, $L=1-2h\eta^2$,
$\lambda_n=2hn$, $b=-A-1/2$, $c=-A$. Define fixed multipliers
\begin{equation}\label{eq:mult}
 v=\frac{V_0}{X},\quad g=\frac{-1+2\eta U_0}{L},\quad
 H=\frac{D\eta+dU_0}{L},\quad
 k=\frac{1-2\eta U_0}{L}+v,\quad
 m_\varphi=\frac{D_X\varphi_0+\varphi_0}{L},\quad
 m_U=\frac{D_XU_0}{L}.
\end{equation}
Also put $q_\varphi=Z_{b,0}\varphi_0$ and $q_U=Z_{c,0}U_0$.
Substitution of the exact $T$ and $Z$ operators into the
current swirl equation yields
\begin{equation}\label{eq:phiphi}
 K_n^{\varphi\varphi}F=\frac12 J_2
 [(b+\lambda_n)gF+vF+H\partial_\eta F+kD_XF].
\end{equation}
Using \eqref{eq:cmp0}--\eqref{eq:cmpX} proves
\begin{equation}\label{eq:cpp}
 \|K_n^{\varphi\varphi}\|\le\frac{C_{\varphi\varphi}}{n+1},\quad
 C_{\varphi\varphi}=2RQ[(|b|+2h)\|g\|+\|v\|
 +\|\partial_\eta H\|+\|k\|+\|D_Xk\|]
 +\frac{4RQ}{\rho}\|H\|.
\end{equation}

For either $\nu=1$ or $2$, define the complete radial-transport
channel
\[
 \mathcal V_{n,\nu}^{m,q}F=\frac12J_\nu
 [(2\eta m+q)F-(D+\lambda_n)(2\eta m)A_XF
 -(dm)\partial_\eta A_XF].
\]
Commuting $\partial_\eta$ through the full product $(dm)A_XF$
and using the radial average estimate gives
\begin{equation}\label{eq:CV}
 \|\mathcal V_{n,\nu}^{m,q}\|\le\frac{C_V(m,q)}{(n+1)^2},\qquad
 C_V(m,q)=2RQ[\|2\eta m+q\|+
 (|D|+2h)\|2\eta m\|+\|\partial_\eta(dm)\|]
 +\frac{4RQ}{\rho}\|dm\|.
\end{equation}
This proof does not require the generally false assertion that
$\partial_\eta A_X$ is bounded by itself on the analytic coefficient
space. The two velocity blocks are exactly
\begin{align}
 K_n^{\varphi U}F&=\mathcal V_{n,2}^{m_\varphi,q_\varphi}F,
 \label{eq:offphi}\\
 K_n^{UU}F&=\frac12 J_1[(c+\lambda_n)gF+H\partial_\eta F+
 kD_XF+q_UF]+\mathcal V_{n,1}^{m_U,0}F.
 \label{eq:uu}
\end{align}
Consequently
\[
 \|K_n^{\varphi U}\|\le\frac{C_V(m_\varphi,q_\varphi)}
 {(n+1)^2},\quad
 \|K_n^{UU}\|\le\frac{C_{UU}}{n+1},
\]
where
\[
 C_{UU}=2RQ[(|c|+2h)\|g\|+\|q_U\|+
 \|\partial_\eta H\|+\|k\|+\|D_Xk\|]
 +\frac{4RQ}{\rho}\|H\|+C_V(m_U,0).
\]

For pressure set $\gamma=2C^{-2}$, using $C$ for the source
azimuthal normalization, and let
$P_F=\gamma\int_0^X\varphi_0(x,\eta)F(x,\eta)\,dx$.
With $C_P=4R\gamma Q\|\varphi_0\|$ one has
$\|P_F\|\le C_P\|F\|/(n+1)$. Put
$s_1=2\eta/L$, $s_2=d/L$, and $s_3=-2\eta/L$.
The fourth block is precisely
\[
 K_n^{U\varphi}F=\frac12J_1[
 (-2A+\lambda_n)s_1P_F+s_2\partial_\eta P_F+s_3D_XP_F].
\]
The same commutator estimates, now applied at ideal index $n+1$,
give
\begin{align}
 \|K_n^{U\varphi}\|&\le\frac{C_{U\varphi}}{(n+1)^2},
 \label{eq:offpress}\\
 C_{U\varphi}&=2RQC_P[
 (2|A|+2h)\|s_1\|+\tfrac{2}{\rho}\|s_2\|
 +\|\partial_\eta s_2\|+\|s_3\|+\|D_Xs_3\|].\notag
\end{align}
In fact $D_Xs_3=0$.

\begin{theorem}[Four-block bound under explicit profile regularity]
\label{thm:fourblock}
Assume the analytic coefficient multiplication estimate above and
that all fixed leading multipliers and their derivatives appearing
in \eqref{eq:cpp}, \eqref{eq:CV}, and \eqref{eq:offpress}
belong to the \emph{same} space $\mathcal B$. Define the four blocks
by the exact equations (47)--(50) of \cite{newton2026} on the
induced-norm ideals $I_n$. Then, for all $n\ge1$,
\[
 \|K_n\|_{(I_n\oplus_1 I_n)\to(I_n\oplus_1 I_n)}
 \le\frac{C_K}{n+1},\qquad
 C_K=\max\{C_{\varphi\varphi}+C_{U\varphi},
 C_{UU}+C_V(m_\varphi,q_\varphi)\}.
\]
\end{theorem}
\begin{proof}
The four exact block rearrangements and the commutator bounds
establish \eqref{eq:cpp}, \eqref{eq:CV} and \eqref{eq:offpress}.
The sum-norm bound of a two-by-two block operator is at most the
maximum of its two column sums. Finally,
$(n+1)^{-2}\le(n+1)^{-1}$.
\end{proof}
This is a conditional concrete-space theorem, not yet a Lean proof of
the full PDE statement. A separate check must verify that the actual
imported leading profiles and all required multiplier derivatives lie
in one fixed analytic norm. Holomorphy on a strictly larger complex
polydisc is a sufficient route, but the precise source-to-norm
verification remains to be completed.

\section{A constructive scalar convergence majorant}
The lower-order forcing is bounded in the source manuscript by the
recurrence
\begin{equation}\label{eq:recur}
 A_n\le C_0(A_{n-1}+A_{n-2})+
 C_1\sum_{i=1}^{n-1}A_iA_{n-i}\quad(n\ge3).
\end{equation}
The scalar summation step can be made fully explicit.
\begin{proposition}[Constructive geometric majorant]\label{prop:maj}
Let $A_n,C_0,C_1\ge0$ satisfy \eqref{eq:recur}. Set
\[
 M=\frac{1}{16(1+C_1)},\qquad
 R\ge\max\{1,8C_0,4\sqrt{C_0},
 4A_1/M,3\sqrt{A_2/M}\}.
\]
Then $A_n\le MR^n/(n+1)^2$ for every $n\ge1$.
An arbitrary finite number of exceptional starting indices
can be absorbed by enlarging $R$.
\end{proposition}
\begin{proof}
For $x=i+1$, $y=n-i+1$, one has $x+y=n+2$ and
\[
 \frac{1}{x^2y^2}\le\frac{2}{(n+2)^2}
 \left(\frac1{x^2}+\frac1{y^2}\right).
\]
Summing and using $1/k^2\le1/[k(k-1)]$ for $k\ge2$
gives
\[
 \sum_{i=1}^{n-1}\frac1{(i+1)^2(n-i+1)^2}
 \le\frac4{(n+2)^2}.
\]
The claims at $n=1,2$ follow from the choice of $R$.
Assuming the conclusion below $n\ge3$ and applying the
recurrence shows
\[
 \frac{A_n}{MR^n/(n+1)^2}\le
 \frac{C_0}{R}\frac{(n+1)^2}{n^2}
 +\frac{C_0}{R^2}\frac{(n+1)^2}{(n-1)^2}
 +4C_1M
 \le \frac{2C_0}{R}+\frac{4C_0}{R^2}+4C_1M
 \le\frac34.
\]
This closes induction; a later recurrence starting index only
changes the finite base conditions.
\end{proof}
This proposition establishes exponential summability from
\eqref{eq:recur}; deriving that recurrence uniformly from the
\emph{complete} Navier--Stokes lower-order forcing and formalizing
the differentiated PDE reconstruction remain distinct tasks.


\section{A quantitative source-to-norm criterion}
\label{sec:cauchy-bridge}
The four-block estimates above assume that the leading-profile multipliers and a finite number of their derivatives lie in one fixed weighted coefficient space. This assumption can be reduced to a quantitative holomorphic-extension test. The reduction is particularly useful here because it avoids an order-dependent loss of analytic radius, although it does not by itself verify the extension for the particular imported profiles.

\begin{theorem}[Cauchy embedding in the doubled analytic weight]
\label{thm:cauchy-embed}
Let $E\subset\mathbb R$ be compact. Let
\(F(X,\eta)=\sum_{a\ge0}F_a(\eta)X^a\) be jointly holomorphic for $|X|<S$ and for every complex $\eta$ at distance less than $\delta$ from $E$. Suppose that $|F|\le M$ on this domain, where $S,\delta,M>0$. For any $R,\rho>0$ satisfying
\[
    R\le \frac S4,\qquad \rho\le \frac\delta4,
\]
the induced coefficient norm \eqref{eq:norm} satisfies
\[
                         \|F\|_{R,\rho;E}\le M.
\]
As usual, one can take slightly smaller $S$ and $\delta$ if the original holomorphic extension is given only on open domains without a uniform boundary bound.
\end{theorem}
\begin{proof}
Two applications of Cauchy's integral estimate, first in $X$ and then on each $\eta$-disc centered at a point of $E$, give
\[
    |\partial_\eta^bF_a(\eta)|\le M S^{-a}\delta^{-b}b!,
    \qquad \eta\in E.
\]
Consequently the coefficient divided by the manuscript weight \eqref{eq:weight} is at most
\[
 M\frac{(a+1)^2(b+1)^2}{\binom{2a+b}{b}}
     \left(\frac R S\right)^a
     \left(\frac\rho\delta\right)^b.
\]
For each integer $j\ge0$ and $0\le q\le1/4$ one has
\begin{equation}\label{eq:poly-geom}
                         (j+1)^2q^j\le1.
\end{equation}
Indeed the assertion is an equality at $j=0$, and the induction step follows from $(j+2)^2\le4(j+1)^2$. Since $\binom{2a+b}{b}\ge1$, applying \eqref{eq:poly-geom} to the radial and parameter indices proves the claim.
\end{proof}

\begin{corollary}[Concrete leading-multiplier criterion]
\label{cor:source-to-norm}
Suppose the leading profiles $\varphi_0,U_0,V_0$ of the imported system admit a jointly holomorphic extension to a complex product neighborhood of a compact terminal-core rectangle; suppose also that $V_0(0,\eta)=0$, and that $L(\eta)=1-2h\eta^2$ has no zero in that neighborhood. On a smaller product neighborhood, each fixed multiplier appearing in \eqref{eq:leadingmult}, the pressure multipliers $s_1,s_2,s_3$, and every finite-order derivative of those multipliers is uniformly bounded and holomorphic. For sufficiently small common $R,\rho>0$, all those multipliers therefore belong to the \emph{same} weighted space $\mathcal B$ required by Theorem~\ref{thm:fourblock}.
\end{corollary}
\begin{proof}
The quotient $V_0/X$ extends holomorphically across $X=0$ by its removable singularity. Every other multiplier is assembled from the leading profiles and their finite derivatives by addition, multiplication, and division by the nonvanishing analytic function $L$. All such functions are holomorphic on the assumed product neighborhood and bounded on any smaller compact product neighborhood. Apply Theorem~\ref{thm:cauchy-embed} separately to the finite list, taking a single pair $R,\rho$ small enough for every entry.
\end{proof}

For the paper's parameter range $0<h<1/100$, the specific denominator $L$ is easily controlled in a moderate complex axial disc. For instance $|\eta|\le2$ implies $|2h\eta^2|<2/25$ and hence $|L(\eta)|>23/25$. This explicitly removes one potential source of singularity, but does not establish that the \emph{leading profiles themselves} extend to the requisite complex neighborhood.

The polynomial-geometric bound \eqref{eq:poly-geom}, its two-index product, and the resulting normalized jet-array inequality have been encoded in \texttt{CauchyProfileBridge.lean}. These are elementary kernel-checkable prerequisites, not a Lean formalization of Cauchy's integral theorem or of the imported Navier--Stokes profiles. The full claim that the specific source leading profiles meet the extension hypothesis still requires an independently checked source theorem and a concrete domain/radius extraction. Once that step is supplied, the preceding corollary turns the profile regularity required by the four-block estimate into a consequence rather than an additional assumption.



\section{Uniform reconstruction from the radial valuation}
\label{sec:reconstruction}
The previous estimates control each coefficient in the analytic norm of
\eqref{eq:norm}. To pass from coefficient norms to the actual complex
similarity variables, one must also ensure that the radial Taylor series and
its parameter derivatives converge uniformly. The radial valuation gives a
quantitative improvement on a smaller fixed core.

\begin{theorem}[Uniform derivative tails for valuation ideals]
\label{thm:radialtails}
Let $R,\rho>0$ and $n\ge1$, and let
$F(X,\eta)=\sum_{a\ge n}F_a(\eta)X^a\in I_n$ have finite induced
coefficient norm $\|F\|_{R,\rho;E}\le A$. Suppose the coefficient functions
are analytic on a common neighborhood of the compact parameter set $E$.
For each integer $b\ge0$ and $\eta\in E$, if $|X|\le R/8$, then
\begin{equation}\label{eq:tailreconstruct}
 |\partial_\eta^b F(X,\eta)|\le
 2^{1-n}A\,\frac{b!}{(b+1)^2}\left(\frac2\rho\right)^b.
\end{equation}
The radial series defining each derivative converges absolutely
and uniformly on this closed core.
\end{theorem}
\begin{proof}
By the binomial theorem, for $a,b\ge0$,
$\binom{2a+b}{b}\le 2^{2a+b}=4^a2^b$. Hence the defining
coefficient norm and $|X|/R\le1/8$ imply
\[
 |\partial_\eta^b F_a(\eta)X^a|
 \le A\,\frac{b!}{(a+1)^2(b+1)^2}
 \left(\frac2\rho\right)^b
 \left(\frac{4|X|}{R}\right)^a.
\]
Sum over $a\ge n$. Because $4|X|/R\le1/2$ and
$\sum_{a\ge n}2^{-a}=2^{1-n}$, the claim follows from
the Weierstrass test.
\end{proof}

\begin{corollary}[Sufficient terminal-series convergence condition]
\label{cor:reconstruct}
Suppose $F_n\in I_n$ and
$\|F_n\|_{R,\rho;E}\le MT^n/(n+1)^2$ with $M,T>0$.
For each fixed $b\ge0$, the series
$\sum_{n\ge1}s^n\partial_\eta^b F_n$ converges normally
on $|X|\le R/8$, $\eta\in E$, whenever $|s|T<2$, and its
absolute sum is bounded by
\begin{equation}\label{eq:reconstructmaj}
 \frac{2Mb!}{(b+1)^2}\left(\frac2\rho\right)^b
 \sum_{n\ge1}\frac{(|s|T/2)^n}{(n+1)^2}.
\end{equation}
\end{corollary}
\begin{proof}
Apply Theorem~\ref{thm:radialtails} to $F_n$, multiply by
$|s|^n$, and sum. The right-hand series converges geometrically
for $|s|T<2$.
\end{proof}

\noindent The estimate applies to fixed derivatives in $\eta$.
Derivatives in $X$ introduce polynomial factors in radial degree,
which remain geometrically summable on a smaller core.
Physical $t,r,z$ derivatives need the exact similarity chain rules
and pressure reconstruction; those operations are not certified by
the present coefficient-space estimate. The all-index binomial
envelope and finite geometric-tail identity are submitted for Lean
verification in \texttt{AnalyticReconstruction.lean}. The actual
Navier--Stokes lower-order forcing recurrence remains to be proved.


\section{Delayed pressure: an exact four-jet source ledger}\label{sec:delayed-pressure}
The last missing step in the local convergence argument is not the abstract
quadratic majorant. It is the uniform estimate that places the complete
lower-order right-hand side of the coefficient equations into that majorant.
In particular, the pressure forcing contains a second-delay term which is
not covered by applying the two-derivative radial-inverse estimate once.
Here we isolate its highest-derivative coefficient and establish a sharp
bound in the original analytic weight. The remaining variable-coefficient
estimates are stated separately.

Write $\mathcal Z_{a,n}=Z_{a+\lambda_n}$ and
$\mathcal Z^{[2]}_{a,n}=Z_{a+\lambda_n-D}Z_{a+\lambda_n}$,
with $Z_a$ as in equation (18) of \cite{newton2026}. Extracting the
current-order terms (47)--(50) from equations (25)--(27) leaves the
following \emph{exact} lower-order source decomposition, for $n\ge2$:
\begin{align}
 f_n^\varphi={}&-\mathcal Z^{[2]}_{b,n-1}\varphi_{n-1}
 +\sum_{i=1}^{n-1}\left[
 V_i\left(\partial_X\varphi_{n-i}+\frac{\varphi_{n-i}}X\right)
 +U_i\mathcal Z_{b,n-i}\varphi_{n-i}\right],
 \label{eq:lower-swirl}\\
 f_n^U={}&-\mathcal Z^{[2]}_{c,n-1}U_{n-1}
 +\sum_{i=1}^{n-1}\left[V_i\partial_XU_{n-i}
 +U_i\mathcal Z_{c,n-i}U_{n-i}\right]
 +\mathcal Z_{-2A,n}\Pi_n^{\mathrm{low}},
 \label{eq:lower-axial}\\
 \Pi_n^{\mathrm{low}}={}&\int_0^X\left[
 C^{-2}\sum_{i=1}^{n-1}\varphi_i\varphi_{n-i}
 -\frac{\Omega_{n-1}}{2x}\right]dx.
 \label{eq:lower-pressure}
\end{align}
Here $\Omega_{n-1}$ is the \emph{full} expression (28), not an
independent pressure datum; in particular it includes the delayed
axial-viscosity term $-\mathcal Z^{[2]}_{0,n-2}V_{n-2}$.
The radial-velocity reconstruction is precisely
\begin{equation}\label{eq:V-delay}
 V_{n-2}=\frac X L\left[2\eta U_{n-2}
 -2\eta(D+\lambda_{n-2})A_XU_{n-2}
 -d\,\partial_\eta A_XU_{n-2}\right].
\end{equation}
Thus the path through the last term in \eqref{eq:V-delay},
$\mathcal Z^{[2]}_{0,n-2}$, \eqref{eq:lower-pressure},
$\mathcal Z_{-2A,n}$, and finally $J_1$ contains four axial
derivatives. It is exactly the path
$U_{n-2}\to V_{n-2}\to\Omega_{n-1}\to\Pi_n\to U_n$.
The first arrow and the final radial inverse each add one radial
degree. The radial average and pressure primitive provide the
factor $(a+1)^{-2}$; $J_1$ at input radial index $a+1$ adds
$(a+2)^{-2}$. This explains the \emph{principal} weight ratio below.

\begin{theorem}[Sharp fourth-jet, two-step radial compensation]
\label{thm:pressure-fourjet}
For the coefficient weights $w_{a,b}$ in \eqref{eq:weight}, and
integers $a,b\ge0$,
\begin{equation}\label{eq:fourjet-exact}
 \frac{w_{a,b+4}}{(a+1)^2(a+2)^2 w_{a+2,b}}
 =\frac{R^2}{\rho^4}
 \frac{(2a+1)(2a+2)(2a+3)(2a+4)(a+3)^2}
 {(a+1)^4(a+2)^2}
 \frac{(b+1)^2}{(b+5)^2}
 \le\frac{54R^2}{\rho^4}.
\end{equation}
The constant $54$ is optimal for this uniform coefficient ratio.
\end{theorem}
\begin{proof}
The factorial representation
\[
 w_{a,b}=R^{-a}\rho^{-b}
 \frac{(2a+b)!}{(2a)!}\frac1{(a+1)^2(b+1)^2}
\]
gives the equality in \eqref{eq:fourjet-exact}, since the
factorials with arguments $2a+b+4$ cancel exactly after a
radial lift by two degrees. For $a\ge0$, the polynomial identity
\begin{align*}
&54(a+1)^4(a+2)^2-(2a+1)(2a+2)(2a+3)(2a+4)(a+3)^2\\
&\hspace{5mm}=2a(a+1)(a+2)(19a^3+71a^2+69a+9)\ge0
\end{align*}
proves the radial bound. The remaining factor
$(b+1)^2/(b+5)^2$ is at most one. For $a=0$ it gives
$54(b+1)^2/(b+5)^2$, tending to $54$ as $b\to\infty$.
\end{proof}

This theorem is an exact estimate for the highest-jet scalar
multiplier along one source-native composite; it is \emph{not}
a norm bound for the complete variable-coefficient pressure
channel. In the latter the factor $(d/L)^4$, its derivatives,
the similarity shifts, the lower-order differential terms in
$Z^{[2]}$, and the other contributions to $\Omega_{n-1}$
produce commutators. Each of those terms must be estimated
on the same induced analytic space, with constants uniform
in the expansion order. One cannot replace this step by the
assertion that individual $\partial_\eta$ or $D_X$ operators
are bounded there.

There is also no automatic gain $1/n$ in the four-jet
term: at large $a$ and then large $b$ the dimensionless
factor in \eqref{eq:fourjet-exact} tends to $16$, not to zero.
That is consistent with the intended two-step linear forcing
term $C_0A_{n-2}$ in (53), but it prevents treating this
channel as a perturbative current-order block. The complete
recurrence requires a separate uniform bound for
\eqref{eq:lower-swirl}--\eqref{eq:lower-pressure} including
all pressure and viscosity contributions. Once that estimate
is certified, the already proved scalar majorant and
reconstruction lemmas can be applied to the \emph{actual}
coefficient sequence.

The factorization and its nonnegativity are encoded in the
new Lean file \texttt{DelayedPressureFourJet.lean}, submitted
for kernel checking. An independent SymPy and exact-rational
script verifies the analytic weight identity and inequality
for 9,333 nontrivial index and parameter triples and checks
the asymptotic sharpness. These computational checks do not
replace a proof of the full operator bound.


\subsection{The two remaining direct linear pressure terms}
\label{subsec:remaining-linear-pressure}

The pressure vorticity source $\Omega_{n-1}$ also contains
$T_{0,n-1}V_{n-1}-2X\partial_X^2V_{n-1}$. Unlike the delayed
viscosity source these terms act on $U_{n-1}$ and produce
\emph{different radial valuations}. Put $k=n-1$, $r=a+1$,
$a\ge n-1$, and define $B=2\eta/L$, $G=d/L$ and the coefficient
differential operators
\[
H_{a,k}=B(r-D-\lambda_k)-G\partial_\eta,\quad
Z_q^{(s)}=G\partial_\eta+B(q-s),\quad
T_q^{(s)}=\frac{D\eta}{L}\partial_\eta+\frac{s-q}{L}.
\]
Equation (24) reads $V_k[X^aF_a]=X^rH_{a,k}F_a/r$.
Taking each of the two sources through
$\Pi_n'=-\Omega_{n-1}/(2X)$ and $\frac12J_1Z_{-2A,n}$
gives the exact maps
\begin{align}
(\mathcal L_nF)_{a+1}
 &=\frac1{2r^2}Z_{\lambda_n-2A}^{(a)}H_{a,k}F_a,
\label{eq:radialpressure}\\
(\mathcal T_nF)_{a+2}
 &=-\frac1{4r^2(r+1)^2}
 Z_{\lambda_n-2A}^{(r)}T_{\lambda_k}^{(r)}H_{a,k}F_a.
\label{eq:timepressure}
\end{align}
The source $-2XV_k''$ produces the first map, from
$I_{n-1}$ to $I_n$, whereas $T_{0,k}V_k$ produces the
second map, from $I_{n-1}$ to $I_{n+1}$.

\begin{theorem}[Uniform bounds for the direct linear pressure sources]
For $0<h<1/100$, $R>0$, $0<\rho\le1/4$ and $E=[-1,1]$,
\[
\|\mathcal L_n\|_{I_{n-1}\to I_n}\le C_{\rm lap},\qquad
\|\mathcal T_n\|_{I_{n-1}\to I_{n+1}}\le C_{\rm time}/n,
\]
where
\begin{align*}
C_{\rm lap}&=\frac{1157}{2}\frac{80}{27}R
 \left(4+\frac8\rho+\frac8{\rho^2}\right),\\
C_{\rm time}&=\frac{9256}{4}\frac{80}{27}\,54R^2
 \sum_{j=0}^3\frac{25\,3^{4-j}}{(j+1)^2\rho^j}.
\end{align*}
In particular their combined action maps $I_{n-1}$ into
$I_n$ with a uniform operator bound.
\end{theorem}

\begin{proof}
On $|\eta|\le3$, the estimates $|L|>41/50$,
$|B|\le8$, $|G|\le13$, $|D\eta/L|\le2$ hold.
Since $r\ge n$, the zeroth- and first-order coefficient
bounds for $H_{a,k}$ are $13r$ and $13$.
For each shifted $Z$ factor they are $24r$ and $13$;
for $T_{\lambda_k}^{(r)}$ they are $2r$ and $2$.
Writing each composition as $\sum_j c_j\partial_\eta^j$,
apply Cauchy's estimate upon shrinking the complex disc
by $1/4$ at each stage. A $Z$ composition maps the bound
$|c_j|\le C_mr^{m-j}$ to
$|c_j^{\rm new}|\le89C_mr^{m+1-j}$, by estimating
$Gc_{j-1}$, $Gc'_j$ and $B(q-r)c_j$ separately.
A $T$ composition admits the factor $8$ instead.
Thus the Laplacian path has coefficient bounds
$1157r^{2-j}$, $j\le2$, on $|\eta|\le11/4$;
the time path has $9256r^{3-j}$, $j\le3$,
on $|\eta|\le5/2$. Derivatives of these coefficients
of order $\ell$ on $E$ gain at most $\ell!$.

The established one-radial-inverse inequalities give
\[
 \frac{r^{2-j}w_{a,b+j}}{r^2w_{a+1,b}}
 \le R(4,8/\rho,8/\rho^2)_j,\quad j=0,1,2.
\]
For the second source, the mixed four-jet estimate in
\eqref{eq:alljet} and $n\le r$ give
\[
 \frac{nr^{3-j}w_{a,b+j}}{r^2(r+1)^2w_{a+2,b}}
 \le54R^2\rho^{-j}\frac{25\,3^{4-j}}{(j+1)^2},
 \quad j\le3.
\]
Leibniz differentiation is absorbed by the coefficient-weight
factor \eqref{eq:commutatorratio}; its sum is at most $80/27$
when $\rho\le1/4$. Combining the prefactors $1/2$, $1/4$
and the indicated finite sums proves the theorem.
\end{proof}

\subsection{The remaining leading-profile and nonlinear pressure sources}
\label{subsec:nonlinear-pressure-ledger}

The full $\Omega_k$, $k\ge1$, may now be written as
\begin{align}
\Omega_k={}&T_{0,k}V_k-2XV_k''-Z_{0,k-1}^{[2]}V_{k-1}\notag\\
 &+V_0V_k'+V_kV_0'-\frac{V_0V_k}{X}
   +U_0Z_{0,k}V_k+U_kZ_{0,0}V_0\notag\\
 &+\sum_{\substack{i+j=k\\i,j\ge1}}
 \left(V_iV_j'-\frac{V_iV_j}{2X}+U_iZ_{0,j}V_j\right).
\label{eq:omega-complete-ledger}
\end{align}
The first line is now bounded term by term: the two estimates
just proved and the earlier complete delayed-viscosity pressure
estimate. The second line consists of \emph{linear}
leading-profile interactions and must be estimated separately
using concrete multiplier norms for $U_0,V_0$.
The final line contains all the genuinely positive-order
quadratic terms, for which the convolution structure is
visible but a full weighted pressure estimate is still needed.

The direct delays in the momentum equations, the interior
momentum products, and the swirl-square contribution in
$\Pi_n'$ remain additional proof obligations. Consequently
the lower-order forcing recurrence (53), normal convergence of
the actual Navier--Stokes hierarchy, and Theorem 1.1 are not
established by the present direct-pressure estimates.
The all-index Banach-space proof above is mathematical;
the accompanying Lean checks presently address its radial
integration factors and exact coefficient indexing only.


\subsection{The five leading-profile pressure interactions}
\label{subsec:leading-profile-pressure}

The remaining linear interactions in equation (28) appear at first to
involve an unbounded radial derivative of $V_k$. Their ordered sum has
a cancellation which is more useful than estimating those derivatives
separately. The reduction is exact and uses no assumptions about the
size or sign of the leading profiles.

Put $v=V_0/X$, let $F=U_k\in I_k$, and write $V_k=XW_k$.
Equation (24) gives, with $B=2\eta/L$, $G=d/L$,
\begin{equation}\label{eq:leading-W}
 W_k=BF-B(D+\lambda_k)A_XF-G\partial_\eta A_XF.
\end{equation}
In particular $W_k\in I_k$ coefficientwise whenever $F\in I_k$;
this is a valuation statement and not a claim that the unregularised
map $F\mapsto W_k$ is bounded in the analytic norm. Define
\begin{equation}\label{eq:leading-K0}
 K_0=G\partial_\eta v-B(v+D_Xv).
\end{equation}
The exact chain rule (18) yields
\[
 \frac{Z_{0,k}V_k}{X}
  =G\partial_\eta W_k+B(\lambda_k-1-D_X)W_k,
 \qquad \frac{Z_{0,0}V_0}{X}=K_0.
\]

\begin{proposition}[Radial cancellation and exact pressure primitive]
\label{prop:leading-cancellation}
For every $k\ge1$, the sum of the five leading-profile terms in
$\Omega_k$, namely
\[
 S_k=V_0V_k'+V_kV_0'-\frac{V_0V_k}{X}
     +U_0Z_{0,k}V_k+U_kZ_{0,0}V_0,
\]
admits the exact representation
\begin{equation}\label{eq:leading-source}
 \frac{S_k}{X}=(D_X+1)(vW_k)
     +U_0\bigl[G\partial_\eta+B(\lambda_k-1-D_X)\bigr]W_k
     +FK_0.
\end{equation}
Consequently its contribution to the order-$(k+1)$ pressure is
\begin{equation}\label{eq:leading-primitive}
 \Pi_{k+1}^{\mathrm{lead}}
 =-\frac12 XvW_k-\frac12\int_0^X
 \left\{U_0\bigl[G\partial_\eta+B(\lambda_k-1-D_X)\bigr]W_k
     +FK_0\right\}(x,\eta)\,dx.
\end{equation}
After the shifted axial operator and the axis-regular inverse, its
contribution to the next axial coefficient is precisely
\begin{equation}\label{eq:leading-pressure-operator}
 \mathcal P_{k+1}^{\mathrm{lead}}F
 =-\frac14 J_1 Z_{-2A,k+1}
 \left[XvW_k+\int_0^X
 \left\{U_0\bigl[G\partial_\eta+B(\lambda_k-1-D_X)\bigr]W_k
   +FK_0\right\}dx\right].
\end{equation}
If $v,U_0,K_0$ are regular at the axis, then
$\mathcal P_{k+1}^{\mathrm{lead}}(I_k)\subseteq I_{k+2}$
coefficientwise.
\end{proposition}
\begin{proof}
Since $V_0=Xv$ and $V_k=XW_k$, direct differentiation gives
\[
 V_0V_k'+V_kV_0'-\frac{V_0V_k}{X}
   =X(D_X+1)(vW_k)=X\partial_X(XvW_k).
\]
This proves the first term of \eqref{eq:leading-source}. The other
two terms follow from the displayed $Z$ identities. Equation (27)
contributes $-S_k/(2X)$ to $\Pi_{k+1}'$.
Integrating from the axis, where $XvW_k$ vanishes, gives
\eqref{eq:leading-primitive}. Equation (26), with its factor
$\frac12J_1 Z_{-2A,k+1}$, gives
\eqref{eq:leading-pressure-operator}.
Finally $W_k\in I_k$, and multiplication by axis-regular fixed
profiles preserves radial valuation. The bracket in
\eqref{eq:leading-pressure-operator} lies in $I_{k+1}$;
$J_1$ raises degree by one, proving the inclusion.
\end{proof}

The cancellation has a second useful consequence: the three source
families have different \emph{input} derivative orders after pressure
reconstruction. The $vW_k$ term has at most two derivatives of $F$;
the $U_0Z_{0,k}V_k$ term has at most three; and $FK_0$ has at most one.
Those statements count derivatives of the unknown $F$, not of the
leading profiles. Derivatives of $v,U_0,B,G$ must be included
explicitly when a bounded analytic operator is constructed.

\begin{lemma}[Two-radial-lift scalar derivative budget]
\label{lem:leading-weight-budget}
For $a,b\ge0$, $r=a+1$ and $j\in\{0,1,2,3\}$, the exact
coefficient weights \eqref{eq:weight} satisfy
\begin{equation}\label{eq:leading-mixed-bound}
 \frac{r^{3-j}w_{a,b+j}}{r^2(r+1)^2 w_{a+2,b}}
 \le \frac{54R^2}{r\rho^j}
       \frac{25\,3^{4-j}}{(j+1)^2}.
\end{equation}
\end{lemma}
\begin{proof}
Apply the all-index two-lift inequality
\eqref{eq:alljet} with the auxiliary index $n=a+2$.
It bounds the same weight ratio with the numerator
$(a+3)^{4-j}$ in place of $r^{4-j}$.
Since $r=a+1\le a+3$, division by $r$ gives
\eqref{eq:leading-mixed-bound}.
\end{proof}

\paragraph{What remains to prove.}
The coefficient budget \eqref{eq:leading-mixed-bound} is compatible
with a $1/(k+1)$ estimate for the regularised leading-profile
pressure operator. It is not that operator estimate: each analytic
profile contains a full radial Taylor series, and commuting its
multipliers through the parameter and radial derivatives produces
additional terms. These must be bounded after radial convolution in
the \emph{induced} norm with constants independent of $k$.
In particular, applying a boundedness claim to the separate map
$F\mapsto\partial_\eta A_XF$ would be invalid. The precise remaining
analytic statement is a uniform bound for the full
\eqref{eq:leading-pressure-operator}, with the common holomorphic
radius and multiplication constants of the actual leading profiles
specified. The positive-order quadratic pressure sum and the other
momentum forcing channels are separate subsequent steps; neither the
recurrence (53) nor Theorem~1.1 follows from Proposition
\ref{prop:leading-cancellation} alone.

The accompanying exact symbolic replay compares both sides of the
five-term source identity and its reconstructed pressure coefficient
on 72 independently specified nontrivial cases, and verifies 41,472
rational inequalities of the form \eqref{eq:leading-mixed-bound}.
The corresponding Lean source encodes the finite radial-cancellation,
primitive-factor and valuation lemmas; kernel certification is pending
successful compilation in the pinned GitHub workflow. The complete analytic
operator is not represented as Lean-certified.


\subsection{Ordered quadratic pressure symmetrization}
\label{subsec:quadratic-pressure}

The remaining positive-order velocity products in $\Omega_k$ have a
related cancellation. It relies on the source sum being over
\emph{ordered} pairs $i+j=k$; estimating individual summands would
hide it. Write $V_i=XW_i$ for every positive order and define
\[
 Q_k=\sum_{\substack{i+j=k\\i,j\ge1}}W_iW_j,
 \qquad H_k=\sum_{\substack{i+j=k\\i,j\ge1}}
 U_i\left[G\partial_\eta+B(\lambda_j-1-D_X)\right]W_j.
\]
Both sums vanish for $k=1$.

\begin{proposition}[Quadratic radial total derivative]
\label{prop:quad-pressure}
The full interior quadratic contribution to the pressure-vorticity
source (28) has the exact form
\begin{equation}\label{eq:quad-omega}
 \Omega_k^{\mathrm{quad}}=\frac X2(D_X+1)Q_k+XH_k.
\end{equation}
Its pressure primitive, fixed to vanish at the axis, is therefore
\begin{equation}\label{eq:quad-pi}
 \Pi_{k+1}^{\mathrm{quad}}=-\frac14XQ_k
                         -\frac12\int_0^XH_k(x,\eta)\,dx,
\end{equation}
and the corresponding axial forcing is
\begin{equation}\label{eq:quad-reconstruct}
 \mathcal P_{k+1}^{\mathrm{quad}}
 =-\frac18J_1Z_{-2A,k+1}
   \left[XQ_k+2\int_0^XH_k(x,\eta)\,dx\right].
\end{equation}
If $U_i,W_i\in I_i$ for every positive order, this contribution
belongs coefficientwise to $I_{k+2}$ after the final radial inverse.
\end{proposition}
\begin{proof}
For $V_i=XW_i$, differentiation gives
\[
 V_iV_j'-\frac{V_iV_j}{2X}
 =X\left(W_iD_XW_j+\frac12W_iW_j\right).
\]
The map $(i,j)\mapsto(j,i)$ is a bijection of the ordered-pair
index set. Applying the product rule and exchanging the indices,
\[
 \sum_{i+j=k}W_iD_XW_j
 =\frac12\sum_{i+j=k}D_X(W_iW_j)
 =\frac12D_XQ_k.
\]
The remaining product is
$U_iZ_{0,j}(XW_j)=XU_i[G\partial_\eta+
B(\lambda_j-1-D_X)]W_j$. This proves
\eqref{eq:quad-omega}. By (27), the quadratic part of
$\Pi_{k+1}'$ is $-\Omega_k^{\mathrm{quad}}/(2X)$.
The identity
$\int_0^X(D_X+1)Q_k\,dx=XQ_k$ gives
\eqref{eq:quad-pi}; applying the axial operator in (26) yields
\eqref{eq:quad-reconstruct}. The valuation statement follows
because every summand of $Q_k$ and $H_k$ has radial degree
at least $i+j=k$; the integral and $J_1$ each add one degree.
\end{proof}

The reduction removes a radial derivative and a pressure primitive
from the entire $V_iV_j$ quadratic source. It does \emph{not}
control the remaining $H_k$ as a bounded bilinear operator.
For that step the complete analytic norm must accommodate the
radial-average and parameter derivatives inside each $W_j$,
the $Z$-operator acting on $W_j$, derivatives of the rational
coefficients, and the convolution over $i+j=k$. A complete estimate
would have to imply, with a constant independent of $k$,
\[
 \bigl\|\mathcal P_{k+1}^{\mathrm{quad}}((U_i)_{i=1}^{k-1})\bigr\|_{I_{k+2}}
 \le C_{\mathrm{quad}}\sum_{\substack{i+j=k\\i,j\ge1}}
   \|U_i\|_{I_i}\|U_j\|_{I_j}.
\]
The scalar convolution inequality in the previous section would
then furnish the desired weighted summation factor.
We do not assert that the missing bilinear Banach-norm estimate
has been established. The exact source identities were independently
replayed in 120 symbolic tests, with 29,106 exact ordered-pair
Euler-coefficient comparisons. Lean encodes the finite ordered-pair
algebra; a kernel proof of the full analytic statement is still open.


\subsection{Radial convolution after two inverse lifts}
\label{subsec:two-lift-convolution}

The preceding pressure cancellations reduce the remaining estimates to products of analytic multipliers with regularized derivatives. A coefficientwise bound on an individual Taylor monomial is not enough for this purpose: the leading profiles have infinitely many radial coefficients, and differentiation distributes among them. We establish here an operator estimate that includes the complete radial convolution, without imposing a finite-degree cutoff. It also records which normal-form reduction is still needed to apply the estimate to the actual pressure source.

\begin{theorem}[A uniform multiplication constant]\label{thm:algebra64}
For the doubled coefficient weight \eqref{eq:weight}, the induced norm \eqref{eq:norm} satisfies
\begin{equation}\label{eq:algebra64}
 \|fg\|_{R,\rho;E}\le 64\|f\|_{R,\rho;E}\|g\|_{R,\rho;E}
\end{equation}
for all coefficient series for which these norms are finite. Moreover, $I_mI_p\subseteq I_{m+p}$, with the same constant. This improves the admissible coefficient-algebra constant in Appendix B.1 of \cite{newton2026}.
\end{theorem}
\begin{proof}
Write $N=i+j$ and $b=u+v$. Leibniz's rule and Vandermonde's inequality give
\[
 \binom b u\frac{w_{i,u}w_{j,v}}{w_{N,b}}
 \le \frac{(N+1)^2}{(i+1)^2(j+1)^2}
       \frac{(b+1)^2}{(u+1)^2(v+1)^2},
\]
since $\binom{2i+u}{u}\binom{2j+v}{v}\le\binom{2N+b}{b}$.
For each $N\ge0$,
\[
 (N+1)^2\sum_{i=0}^N\frac1{(i+1)^2(N-i+1)^2}
 \le4\sum_{r=1}^{N+1}\frac1{r^2}\le8.
\]
Indeed $(N+1)^2\le 2[(i+1)^2+(N-i+1)^2]$, while
$\sum_{r=1}^{N+1}r^{-2}\le1+\sum_{r=2}^{N+1}[r(r-1)]^{-1}<2$.
The same estimate applies to the Leibniz derivative index $b$. Their product is $64$, uniformly in both indices. The valuation assertion follows directly from coefficient convolution.
\end{proof}

The next result allows multiplication by a nonpolynomial analytic profile \emph{before} the radial inverses act. Let $D_X=X\partial_X$ and $n\ge1$. For $0\le j\le3$ and $F\in I_n$, define the two-lift jet operator
\begin{equation}\label{eq:regularized-jet}
 (\mathcal R_{n,j}F)_{a+2}
 =\frac{(a+n+2)^{3-j}}{(a+1)^2(a+2)^2}\,\partial_\eta^j F_a,
 \qquad a\ge n.
\end{equation}
For a multiplier $g=\sum_{c\ge0}g_c(\eta)X^c$, define the convolution-inverse operator coefficientwise by
\begin{equation}\label{eq:radial-convolution}
 (\mathcal C_{n,j}(g,F))_{N}
 =\sum_{\substack{a\ge n,\ c\ge0\\a+c+2=N}}
 \frac{(N+n)^{3-j}}{(N-1)^2N^2}\,
 g_c\,\partial_\eta^jF_a,
 \qquad N\ge n+2.
\end{equation}
No restriction is placed on the number of nonzero Taylor coefficients of $g$.

\begin{theorem}[Tame analytic radial-convolution smoothing]\label{thm:convolution-smoothing}
Let $0\le j\le3$, $n\ge1$, $R,\rho>0$, and put
\begin{equation}\label{eq:tame-constant}
 D_j(R,\rho)=\frac{675\,2^{4-j}R^2}{(j+1)^2\rho^j}.
\end{equation}
Then
\begin{equation}\label{eq:tame-basic}
 \|\mathcal R_{n,j}F\|_{R,\rho;E}
 \le\frac{D_j(R,\rho)}{n+1}\|F\|_{R,\rho;E}.
\end{equation}
If $g,D_Xg,\ldots,D_X^{3-j}g$ belong to the same weighted space, then $\mathcal C_{n,j}(g,F)\in I_{n+2}$ and
\begin{equation}\label{eq:tame-convolution}
 \|\mathcal C_{n,j}(g,F)\|_{R,\rho;E}
 \le\frac{64D_j(R,\rho)}{n+1}
 \left(\sum_{s=0}^{3-j}\binom{3-j}{s}\|D_X^sg\|_{R,\rho;E}\right)
 \|F\|_{R,\rho;E}.
\end{equation}
Both constants are independent of $n$ and of the radial Taylor support of $g$.
\end{theorem}
\begin{proof}
Set $t=a+n+2$. Since $a\ge n$, one has
$2(n+1)\le t\le2(a+1)$. The exact fourth-jet weight identity
\eqref{eq:fourjet-exact} implies, for $0\le j\le3$,
\[
 \frac{t^{4-j}w_{a,b+j}}{(a+1)^2(a+2)^2w_{a+2,b}}
 \le \frac{54R^2}{\rho^j}\frac{25\,2^{4-j}}{(j+1)^2}.
\]
To obtain this, divide the $j$-jet ratio by the fourth-jet ratio. The quotient is
$\rho^{4-j}(b+5)^2/(b+j+1)^2$ times the reciprocal of
$\prod_{m=j+1}^4(2a+b+m)$, where the first fraction is at most
$25\rho^{4-j}/(j+1)^2$ and each denominator is at least $a+1$.
Dividing by $t\ge2(n+1)$ gives \eqref{eq:tame-basic}, with
$54\cdot25/2=675$.

For the full convolution, write $p=3-j$, $N=a+c+2$, and
$t=a+n+2$. The elementary facts
\[
 (N+n)^p=(t+c)^p=\sum_{s=0}^p\binom p s t^{p-s}c^s,
 \qquad (N-1)^2N^2\ge(a+1)^2(a+2)^2,
\]
show that each coefficient and each parameter derivative of
$\mathcal C_{n,j}(g,F)$ is dominated term by term by the nonnegative
Leibniz majorants of
\[\sum_{s=0}^p\binom p s(D_X^sg)\,\mathcal R_{n,j}F.\]
Here we used $t\ge1$ to replace $t^{p-s}$ by $t^p$; no derivative is moved between $g$ and $F$ without the corresponding Leibniz terms. Apply the multiplication estimate
\eqref{eq:algebra64} to these majorants and then
\eqref{eq:tame-basic}. This proves \eqref{eq:tame-convolution}.
\end{proof}

\paragraph{The highest leading-profile pressure jet.}
The source-native formula \eqref{eq:leading-pressure-operator} allows the highest input derivative to be identified without freezing any leading-profile coefficient. Fix $k\ge1$, $a\ge k$ and $c\ge0$, set $r=a+1$, $m=a+c+1$, $N=m+1$, and let $u_c$ be the $X^c$ coefficient of $U_0$. For one input coefficient $F_a$, the radial-velocity identity \eqref{eq:leading-W} reads
\[
 W_{k,a}=w_0F_a+w_1\partial_\eta F_a,\quad
 w_0=B\frac{r-D-\lambda_k}{r},\qquad w_1=-\frac Gr.
\]
Only the $U_0Z_{0,k}V_k$ contribution can produce three derivatives after the final pressure reconstruction. The full source operator's third-jet coefficient is exactly
\begin{equation}\label{eq:highest-leading-jet}
 \bigl(\mathcal P_{k+1}^{\mathrm{lead}}F\bigr)_{a+c+2}\big|_{\partial_\eta^3F_a}
 =\frac{u_cG^3}{4r m N^2}\,\partial_\eta^3F_a.
\end{equation}
This is an identity of variable-coefficient operators, including the radial degree $c$ of the leading profile.

\begin{corollary}[Uniform leading-profile highest-jet channel]\label{cor:leading-highest}
If $g=\tfrac14G^3(U_0+D_XU_0)$ belongs to the weighted algebra, the complete third-derivative component of \eqref{eq:leading-pressure-operator}, summed over all $a,c$, satisfies
\[
 \|\mathcal P_{k+1}^{\mathrm{lead},(3)}F\|_{R,\rho;E}
 \le\frac{64D_3(R,\rho)}{k+1}\|g\|_{R,\rho;E}\|F\|_{R,\rho;E},\qquad F\in I_k.
\]
\end{corollary}
\begin{proof}
The coefficient of $\partial_\eta^2F_a$ in the integrand's $U_0Z_{0,k}V_k$ term is $u_cGw_1=-u_cG^2/r$. The radial primitive supplies $1/m$, and the outer $G\partial_\eta$ followed by $-J_1/4$ yields \eqref{eq:highest-leading-jet}. In the model operator \eqref{eq:radial-convolution} with $j=3$ the denominator is $m^2N^2$. Since $m/r=1+c/r\le1+c$,
\[
 \frac{|u_cG^3|}{4rmN^2}
 \le\frac{(1+c)|u_cG^3|}{4m^2N^2}.
\]
The multiplier on the right is precisely the $X^c$ coefficient of $g$. The same domination holds termwise after any $\eta$-derivatives, because $m/r$ is independent of $\eta$. Apply Theorem~\ref{thm:convolution-smoothing} with $j=3$.
\end{proof}

This result is directly applicable whenever a term of the source operator can be normal-ordered into a finite sum of the coefficients in \eqref{eq:radial-convolution}, with multipliers in the stated common analytic space. In particular, it removes the need to truncate a leading profile's radial series and controls the additional factors of $c$ generated by $D_X$ commutators. The specific highest input-jet channel in the leading-profile pressure source is covered by Corollary~\ref{cor:leading-highest}. The next subsection closes the three remaining leading-profile input-jet channels. This abstract theorem alone does \emph{not} prove that the complete pressure operators \eqref{eq:leading-pressure-operator} and \eqref{eq:quad-reconstruct} have yet been reduced to such finite sums with a uniform multiplier budget. Their exact symbolic normal forms include intermediate radial averages and, in the quadratic case, two separate unknown inputs. The required source-specific reduction and the bilinear estimate for $H_k$ remain the next proof obligations. The complete recurrence (53) and Theorem~1.1 therefore remain open.


\subsection{The complete leading-profile pressure operator}
\label{subsec:full-leading-normalform}
The highest-derivative estimate above treats only one channel of the five-term leading-profile pressure source. We now control the remaining three input derivative orders as well. The point is to expand the complete source before applying the analytic norm, retaining every radial index until both inverse lifts have acted.

Put $u=U_0$, $v=V_0/X$, $K=G\partial_\eta v-B(v+D_Xv)$, $B=2\eta/L$ and $G=(1-\eta^2)/L$. Fix $k\ge1$, $F=U_k\in I_k$, input degree $a\ge k$, and leading-profile radial degree $c\ge0$. Set
\begin{align*}
r&=a+1,\quad m=a+c+1,\quad N=m+1,\\
\alpha&=r-D-\lambda_k,\quad\epsilon=\lambda_k-r,\\
\theta&=\lambda_{k+1}-2A-m=2hk-1-m .
\end{align*}
Then $W_{k,a}=(\alpha B F_a-G\partial_\eta F_a)/r$.
All products of $u,v,K$ below are interpreted at radial degree $c$.
Introduce three finite collections of pairs $(p,g)$, with $p$ independent
of $\eta$ and $g$ an analytic multiplier:
\begin{align*}
\mathcal H_0={}&\{(m\alpha,vB),\,(\alpha,uGB'),\,
(\epsilon\alpha,uB^2),\,(r,K)\},\\
\mathcal H_1={}&\{(-m,vG),\,(\alpha,uGB),\,
(-1,uGG'),\,(-\epsilon,uBG)\},\\
\mathcal H_2={}&\{(-1,uG^2)\},\quad
\mathcal H_{-1}=\mathcal H_3=\varnothing.
\end{align*}
For $j=0,1,2,3$ put, retaining multiplicities,
\begin{equation}\label{eq:leading-pair-recursion}
\mathcal E_j=
\{(\theta p,Bg),(p,Gg'):(p,g)\in\mathcal H_j\}
\mathbin{\cup}
\{(p,Gg):(p,g)\in\mathcal H_{j-1}\}.
\end{equation}
The numbers of pairs are $8,12,6,1$, respectively.

\begin{proposition}[Exact complete leading-pressure normal form]
\label{prop:leading-all-jets}
For the full five-term leading-pressure operator in
\eqref{eq:leading-pressure-operator},
\begin{equation}\label{eq:complete-leading-normalform}
(\mathcal P_{k+1}^{\mathrm{lead}}F)_{a+c+2}
=-\frac1{4rmN^2}
\sum_{j=0}^3\sum_{(p,g)\in\mathcal E_j}
p\,g_c(\eta)\,\partial_\eta^jF_a(\eta).
\end{equation}
For every contributing pair, with $t=N+k=m+k+1$,
\begin{equation}\label{eq:leading-poly-bound}
|p|\le t^{3-j},\qquad m/r\le1+c .
\end{equation}
\end{proposition}
\begin{proof}
Write $H_\ell=\sum_{(p,g)\in\mathcal H_\ell}p g_c$
for $\ell=0,1,2$, with $H_{-1}=H_3=0$. Substituting
$W_{k,a}$ into the exact pressure primitive gives
$X^m(H_0F_a+H_1F_a'+H_2F_a'')/(rm)$.
The shifted operator $Z_{-2A,k+1}=
G\partial_\eta+B(\lambda_{k+1}-2A-D_X)$
sends its coefficient of $F_a^{(j)}$ to
$\theta BH_j+GH_j'+GH_{j-1}$. The final factor $-J_1/4$
supplies $1/(4N^2)$, yielding the stated normal form.
Since $a\ge k\ge1$, $r\ge k+1$ and $m=r+c$.
For $0<h<1/100$, $0<D<1$ and
$0<\lambda_k<k/50$, hence
$0\le\alpha\le r$, $-r\le\epsilon\le0$,
and $-(m+1)\le\theta\le0$.
Each $p$ in $\mathcal E_j$ is a signed monomial of degree
at most $3-j$ in $r,m,\alpha,\epsilon,\theta$.
All five have absolute value at most $t$, proving the first
bound. The second is $m/r=1+c/r\le1+c$.
\end{proof}

Define the finite multiplier budget
\begin{equation}\label{eq:leading-multiplier-budget}
\mathfrak M_j=\sum_{(p,g)\in\mathcal E_j}
\sum_{s=0}^{3-j}\binom{3-j}{s}
\|D_X^s(1+D_X)g\|_{R,\rho;E},\quad0\le j\le3.
\end{equation}
The 27 analytic multipliers are fully specified by
\eqref{eq:leading-pair-recursion}; none contains an unknown positive-order field.

\begin{theorem}[Uniform complete leading-profile pressure estimate]
\label{thm:full-leading-pressure}
Assume $0<h<1/100$, $R>0$, $0<\rho\le1/4$,
and finiteness of all four budgets
\eqref{eq:leading-multiplier-budget} on one common induced
analytic coefficient space. For every $k\ge1$ and $F\in I_k$,
\begin{equation}\label{eq:leading-full-bound}
\|\mathcal P_{k+1}^{\mathrm{lead}}F\|_{R,\rho;E}
\le\frac{C_{\mathrm{lead}}}{k+1}\|F\|_{R,\rho;E},
\quad C_{\mathrm{lead}}=
16\sum_{j=0}^3D_j(R,\rho)\mathfrak M_j,
\end{equation}
where $D_j(R,\rho)=675\,2^{4-j}R^2/
[(j+1)^2\rho^j]$ as in \eqref{eq:tame-constant}.
The complete five-term operator maps $I_k$ into $I_{k+2}$.
\end{theorem}
\begin{proof}
For $(p,g)\in\mathcal E_j$,
\[
\frac{|p|}{4rmN^2}|g_c\partial_\eta^jF_a|
\le\frac{t^{3-j}}{4m^2N^2}
|(1+c)g_c\partial_\eta^jF_a|.
\]
The right side is one quarter of the positive coefficient
majorant of the convolution-inverse operator
\eqref{eq:radial-convolution} with multiplier
$(1+D_X)g$ and $n=k$. The same domination applies
to all $\eta$-derivatives, since $p$ and $m/r$ are
$\eta$-independent, with each product derivative retained
by Leibniz's rule. The proof of Theorem
\ref{thm:convolution-smoothing} bounds precisely those
positive Leibniz majorants without cancellations and yields
for each pair
\[
\|\text{pair's contribution}\|
\le\frac{16D_j}{k+1}
\sum_{s=0}^{3-j}\binom{3-j}{s}
\|D_X^s(1+D_X)g\|\|F\|.
\]
Sum over the 27 pairs. The valuation follows
directly from \eqref{eq:complete-leading-normalform}.
\end{proof}

A sufficient, but still unverified, source condition is a
common bounded holomorphic extension of $U_0$ and $V_0/X$
to a larger complex product neighborhood where $L\ne0$.
By the Cauchy embedding theorem \ref{thm:cauchy-embed},
on a smaller common domain all the finitely many budgets
above are then finite after choosing $R,\rho$ small enough.
This does not certify the actual imported profiles.

The full normal form has been checked symbolically, and
45,360 exact-rational cases of the scalar polynomial
majorants passed independent replay. The Lean source
contains finite scalar/index proof gates, not the complete
Banach-space formalization. The bilinear $H_k$ pressure
source, other momentum sources, concrete imported profile
radii and the full recurrence (53) remain open. In
particular this theorem does not establish Theorem~1.1.


\paragraph{The highest jet of the ordered quadratic pressure operator.}
The preceding two-lift estimate also controls the most differentiated component of the genuinely bilinear pressure source. In the ordered source term $U_iZ_{0,j}V_j$ of \eqref{eq:quad-reconstruct}, fix $i+j=k$, and radial indices $a\ge i$, $b\ge j$. Put $N=a+b$ and $r=b+1$. The coefficient of $\partial_\eta^2(U_j)_b$ in
$U_i[G\partial_\eta+B(\lambda_j-1-D_X)]W_j$
is $-(U_i)_aG^2/r$. The outer shifted $Z$ gives the third derivative, the pressure primitive gives $1/(N+1)$, and the regular-singular inverse gives $1/(N+2)^2$. Hence the exact coefficient of $\partial_\eta^3(U_j)_b$ is
\begin{equation}\label{eq:quad-highest-jet}
\frac{G^3(U_i)_a}{4(b+1)(N+1)(N+2)^2}
\partial_\eta^3(U_j)_b.
\end{equation}
This is one derivative channel, not the full bilinear $H_k$.
Since $N\ge b$,
\[
(b+1)(N+1)(N+2)^2\ge(b+1)^2(b+2)^2.
\]
The inverse-denominator comparison is independent of $\eta$ and therefore persists under every axial derivative.

\begin{corollary}[Uniform highest-jet bilinear pressure bound]
\label{cor:quad-highest}
If $G^3$ has finite norm in the same weighted analytic algebra, the complete radial convolution of this $(0,3)$ jet component satisfies
\begin{equation}\label{eq:quad-highest-bound}
\|\mathcal Q_{ij}^{(0,3)}(U_i,U_j)\|_{R,\rho;E}
\le\frac{1024D_3(R,\rho)}{j+1}
\|G^3\|_{R,\rho;E}\|U_i\|_{R,\rho;E}\|U_j\|_{R,\rho;E}.
\end{equation}
\end{corollary}
\begin{proof}
The absolute coefficients and their Leibniz derivatives are dominated by those of
$\tfrac14(G^3U_i)\mathcal R_{j,3}U_j$.
Apply the multiplication theorem \eqref{eq:algebra64} twice and then
$\|\mathcal R_{j,3}U_j\|\le D_3\|U_j\|/(j+1)$.
The constant is $64^2/4=1024$.
\end{proof}
The bound fits the quadratic-convolution budget because $1/(j+1)\le1$.
Mixed input derivatives, the lower jets and the remaining quadratic
pressure/momentum contributions are not covered. The full coefficient
recurrence (53) remains open. Exact symbolic checks reproduce the
highest-jet coefficient, and 7,569 exact denominator comparisons passed.


\paragraph{The complete positive-order quadratic pressure operator.}
The exact symmetrization in Proposition~\ref{prop:quad-pressure} leaves two bilinear operators. The first comes from the axial coupling $H_k$, and the second from $Q_k=\sum_{i+j=k}W_iW_j$. Their estimates require more than the pure third-jet bound: the outer shifted $Z$ differentiates either input, and an operator norm cannot rely on cancellations between those Leibniz terms. We give a common majorant argument for all mixed derivative orders.

Throughout this paragraph $i,j\ge1$, $k=i+j$, $U_i\in I_i$, $U_j\in I_j$, and $0<h<1/100$, $E=[-1,1]$, $0<\rho\le1/4$. Put $a\ge i$, $b\ge j$, $N=a+b$, $r=b+1$, and $t=N+k+2$. Write
\[
 B=\frac{2\eta}{L},\qquad G=\frac{1-\eta^2}{L},\qquad
 \alpha=1-\frac{D+\lambda_j}{r},\quad
 e=\lambda_j-1-b,\quad\theta=\lambda_{k+1}-2A-(N+1).
\]
Since $j\le b$, $k\le N$ and $2h<1/50$, we have
\begin{equation}\label{eq:bilinear-index}
 0\le\alpha\le1,\quad r^{-1}\le1,\quad
 |e|\le t,\quad |\theta|\le t,\quad
 N+1\le t\le2(N+1).
\end{equation}
On the $b$-th radial coefficient, $W_j=\alpha B U_j-(G/r)\partial_\eta U_j$. Define three analytic coefficient functions
\begin{align}\label{eq:bilinear-c}
 c_0&=\alpha(GB'+eB^2),\\
 c_1&=\alpha BG-r^{-1}(GG'+eBG),\notag\\
 c_2&=-r^{-1}G^2.\notag
\end{align}
The exact expression $H_{ij}=U_i[G\partial_\eta+B(\lambda_j-1-D_X)]W_j$ is
$U_i(c_0U_j+c_1U_j'+c_2U_j'')$ coefficientwise. After the pressure primitive and the final shifted $Z$, its contribution is
\begin{equation}\label{eq:H-composite}
 (\mathcal H_{ij})_{N+2}=-\frac1{4(N+1)(N+2)^2}
 \sum_{(p,q)\in\mathcal E_H} h_{p,q}\,\partial_\eta^p(U_i)_a\,\partial_\eta^q(U_j)_b,
\end{equation}
where
\begin{align*}
\mathcal E_H&=\{(0,0),(1,0),(0,1),(1,1),(0,2),(1,2),(0,3)\},\\
h_{0,0}&=B\theta c_0+Gc_0',&h_{1,0}&=Gc_0,\\
h_{0,1}&=B\theta c_1+Gc_1'+Gc_0,&h_{1,1}&=Gc_1,\\
h_{0,2}&=B\theta c_2+Gc_2'+Gc_1,&h_{1,2}&=Gc_2,\\
h_{0,3}&=Gc_2.
\end{align*}
This is a complete seven-channel normal form, not merely a derivative-order ledger.

We first record the analytic smoothing estimate needed for its bilinear terms. For any $N\ge k\ge2$, $\beta\ge0$, and $d=0,1,2,3$, direct cancellation of the weight factorials gives
\begin{align}\label{eq:bilinear-weight-identity}
 \frac{w_{N,\beta+d}}{w_{N+2,\beta}}
 &=R^2\rho^{-d}\frac{(N+3)^2}{(N+1)^2}
  \frac{(\beta+1)^2}{(\beta+d+1)^2}
  \frac{\prod_{\ell=1}^{4}(2N+\ell)}
       {\prod_{\ell=d+1}^{4}(2N+\beta+\ell)}.
\end{align}
By $(N+3)/(N+1)\le3$, $(2N+4)/(N+1)\le4$ and $2N+\beta+\ell\ge N+1$, it follows that
\begin{align}\label{eq:bilinear-smoothing-bounds}
 \frac{t^{3-d}w_{N,\beta+d}}
 {(N+1)(N+2)^2w_{N+2,\beta}}
 &\le E_d:=\frac{2304\,2^{3-d}R^2}{\rho^d},\\
 \frac{t^{2-d}w_{N,\beta+d}}
 {(N+2)^2w_{N+2,\beta}}
 &\le F_d:=\frac{2304\,2^{\max(2-d,0)}R^2}{\rho^d},
 \qquad d=0,1,2,3.\notag
\end{align}
In the second inequality $t^{2-d}=t^{-1}$ when $d=3$. These are all-index bounds, uniform in $i$, $j$, $N$, and $\beta$.

To handle mixed derivatives, define the nonnegative Leibniz jet majorant
\[
 \mathcal M_{N,s}(U_i,U_j)=
 \sum_{a+b=N}\sum_{u+v=s}\binom{s}{u}
 |\partial_\eta^u(U_i)_a|\,|\partial_\eta^v(U_j)_b|.
\]
The weighted multiplication estimate \eqref{eq:algebra64} applies without cancellation and yields
\begin{equation}\label{eq:bilinear-unsigned}
 \mathcal M_{N,s}(U_i,U_j)
 \le64\|U_i\|\|U_j\|w_{N,s}.
\end{equation}
Moreover, for any $p,q\ge0$, $d=p+q$, $s\ge0$, and $0\le u\le s$,
$\binom{s}{u}\le\binom{s+d}{u+p}$. Thus every differentiated split product is dominated by the corresponding unsigned majorant at order $s+d$.

The rational multipliers can be bounded without a leading-profile hypothesis. On $|\eta|\le3$, $|L|>41/50$, and $|B|<300/41$, $|G|<500/41$. For $q=0,1,2$, Cauchy's estimate first bounds $B^{(q)}$ and $G^{(q)}$ on $|\eta|\le2$ by $q!\,300/41$ and $q!\,500/41$, using unit-radius discs. A second unit-radius estimate at $E$, together with $(\ell+1)^2\rho^\ell\le1$, yields
\begin{equation}\label{eq:bilinear-multiplier-M}
 M:=\max\{1,\|B\|,\|G\|,\|B'\|,\|G'\|,
 \|B''\|,\|G''\|\}\le25.
\end{equation}
These norms refer to the radial-degree-zero elements of $\mathcal B_{R,\rho;E}$. By the multiplication bound $64$, \eqref{eq:bilinear-index}--\eqref{eq:bilinear-c} give
\begin{align*}
 \|c_0\|&\le2\cdot64M^2t,& \|c_0'\|&\le4\cdot64M^2t,\\
 \|c_1\|&\le3\cdot64M^2t,& \|c_1'\|&\le6\cdot64M^2t,\\
 \|c_2\|&\le64M^2,& \|c_2'\|&\le2\cdot64M^2.
\end{align*}
Consequently, for the seven exact coefficients,
\begin{equation}\label{eq:bilinear-H-normalform-bound}
 \|h_{p,q}\|\le A_H t^{3-p-q},
 \qquad A_H:=11\cdot64^2 M^3.
\end{equation}

\begin{theorem}[Complete bilinear axial-pressure bound]\label{thm:full-H}
For every ordered pair $i,j\ge1$ the whole $H_k$ pressure contribution \eqref{eq:H-composite}, including all mixed-jet and multiplier-commutator terms, satisfies
\begin{equation}\label{eq:full-H-bound}
 \|\mathcal H_{ij}(U_i,U_j)\|_{R,\rho;E}
 \le C_H\|U_i\|_{R,\rho;E}\|U_j\|_{R,\rho;E},
 \quad C_H:=\frac{64^2A_H}{4}(E_0+2E_1+2E_2+2E_3).
\end{equation}
The constant is finite and independent of $i,j$ and all radial Taylor cutoffs.
\end{theorem}
\begin{proof}
Apply the multiplier bound \eqref{eq:bilinear-H-normalform-bound}, the unsigned mixed-derivative majorant \eqref{eq:bilinear-unsigned}, and the first smoothing inequality in \eqref{eq:bilinear-smoothing-bounds} with $d=p+q$. Differentiated multipliers are handled by the same positive Leibniz convolution as multiplication by a radial-degree-zero function; the algebra estimate contributes a second factor $64$. There is one term with $d=0$ and two terms at each of $d=1,2,3$. The pressure prefactor is $1/4$.
\end{proof}

The companion quadratic pressure term is the symmetrized $Q_k$ source. For $r_i=a+1$, $r_j=b+1$ write $W_i=\alpha_i B U_i-G\partial_\eta U_i/r_i$, with $\alpha_i=1-(D+\lambda_i)/r_i$, and similarly for $W_j$. After applying the outer $Z$, the exactly eight derivative-split positions are
\[
\mathcal E_Q=\{(0,0),(1,0),(0,1),(1,1),
 (2,0),(0,2),(2,1),(1,2)\}.
\]
Their coefficients are obtained, for $p,q\in\{0,1\}$, from the terms
\begin{equation}\label{eq:Q-source-normalform}
 (p,q):B\theta g_{p,q}+Gg_{p,q}',\quad
 (p+1,q):Gg_{p,q},\quad (p,q+1):Gg_{p,q},
\end{equation}
where $g_{p,q}=w_{i,p}w_{j,q}$, $w_{i,0}=\alpha_iB$, $w_{i,1}=-G/r_i$, and similarly for $j$. The exact output prefactor is $-1/[8(N+2)^2]$.
The two radial averages satisfy
\begin{equation}\label{eq:quadratic-double-average}
 r_ir_j=(a+1)(b+1)\ge N+1
 \quad\Longrightarrow\quad
 \frac{|\theta|}{r_ir_j}\le2,\qquad
 \frac1{r_ir_j}\le\frac2t.
\end{equation}
These bounds control both the shifted order-two coefficient and the order-three coefficient. Collecting each derivative split and using \eqref{eq:bilinear-multiplier-M} gives
\begin{equation}\label{eq:bilinear-Q-normalform-bound}
 \|q_{p,q}\|\le A_Q t^{2-p-q},
 \qquad A_Q:=32\cdot64^2M^3.
\end{equation}
The exponent is $-1$ for the two third-order channels, as required by the second smoothing inequality.

\begin{corollary}[The complete quadratic pressure source]\label{cor:all-quadratic-pressure}
Both $Q_k$ and $H_k$ in \eqref{eq:quad-reconstruct} are bounded in the common induced analytic space. In particular,
\begin{equation}\label{eq:quadratic-complete-constant}
 \|\mathcal P_{ij}^{\mathrm{quad}}(U_i,U_j)\|
 \le C_{\mathrm{quad,p}}\|U_i\|\|U_j\|,\quad
 C_{\mathrm{quad,p}}=C_H+C_Q,
\end{equation}
where
\begin{equation}\label{eq:CQ}
 C_Q:=\frac{64^2A_Q}{8}(F_0+2F_1+3F_2+2F_3).
\end{equation}
Consequently $\|\mathcal P_{k+1}^{\mathrm{quad}}\|
\le C_{\mathrm{quad,p}}\sum_{i+j=k,\,i,j\ge1}\|U_i\|\|U_j\|$
for every $k\ge2$, with a constant independent of $k$.
\end{corollary}
\begin{proof}
Use the eight exact positions in \eqref{eq:Q-source-normalform},
the bound \eqref{eq:bilinear-Q-normalform-bound}, the unsigned Leibniz
majorant, and the second estimate in \eqref{eq:bilinear-smoothing-bounds}.
The algebra factors give $64^2$, the pressure prefactor is $1/8$,
and there are one, two, three and two channels at derivative orders
$0,1,2,3$. Add Theorem~\ref{thm:full-H} and sum over ordered pairs.
\end{proof}

These constants are deliberately conservative. Unlike the leading-profile theorem, the present bilinear estimate needs no imported-profile regularity hypothesis: $B$ and $G$ are explicit rational functions with the bounds above. It closes the positive-order quadratic pressure contribution of (28), but not the remaining swirl/axial momentum products or the unverified imported leading-profile radii. Thus the complete recurrence (53), termwise physical PDE reconstruction, and Theorem~1.1 are not established here.


\paragraph{The remaining positive-order momentum sources.}
The positive-order convection terms in the swirl and axial equations have the common form
\[
 {\cal S}^{\delta,\beta}_{ij}(U_i,F_j)
 =W_i(D_X+\delta)F_j+
 U_i[G\partial_\eta+B(\beta+\lambda_j-D_X)]F_j,\qquad
 V_i=XW_i,\quad
 W_{i,a}=\alpha_{i,a}B(U_i)_a-\frac G{a+1}\partial_\eta(U_i)_a,
\]
where $\alpha_{i,a}=1-(D+\lambda_i)/(a+1)$,
$B=2\eta/L$, $G=(1-\eta^2)/L$, $i+j=n$,
and $(\delta,\beta)=(1,b)$ in swirl or $(0,c)$ in axial
momentum. Here $b=-A-\frac12$, $c=-A$. For $a\ge i$,
$0\le\alpha_{i,a}\le1$; furthermore $|\beta|<2$.

\begin{theorem}[Uniform positive-order momentum products]
For $0<h<1/100$, $R>0$, $E=[-1,1]$ and $0<\rho\le1/4$,
\[
 \left\|\frac12J_\nu {\cal S}_{ij}^{\delta,\beta}(U_i,F_j)\right\|_{I_{n+1}}
 \le C_{\rm mom}\|U_i\|_{I_i}\|F_j\|_{I_j},
 \qquad C_{\rm mom}=\frac{64^2\,25\,R}{2}
 \left(20+\frac{16}{\rho}\right).
\]
This includes both $\nu=1$ and $\nu=2$, with a constant
independent of the two expansion orders.
\end{theorem}
\begin{proof}
The analytic coefficient algebra has multiplication constant $64$.
The radial-degree-zero functions $B,G$ have norms at most $25$:
on $|\eta|\le3$, $|B|\le300/41$ and $|G|\le500/41$;
Cauchy's inequality on radius-two discs about $[-1,1]$ and
$(m+1)^3/8^m\le1$ also bound $B',G'$ by $25$.
Use the unsigned Leibniz convolution of $U_i,F_j$.
A source with one derivative on either factor is dominated
by the total first-derivative convolution because
$\binom{s}{u}\le\binom{s+1}{u+1}$ and
$\binom{s}{u}\le\binom{s+1}{u}$.
An Euler derivative on one factor costs at most the total
radial degree $N$. The exact regular-singular inverse
multipliers satisfy
\[
 T_0\le 4R/(N+1)^2,\quad T_1\le8R/[\rho(N+1)],
 \quad NT_0\le4R,\quad (N+1)T_1\le8R/\rho.
\]
The bound $|\beta+\lambda_j-a_j|\le
(1+|\beta|)(N+1)$ holds for $a_j\ge j$.
Multiplication by $B$ or $G$ commutes with $J_\nu$ because
they do not depend on $X$, and costs a second factor $64$.
The four source types give, before the factor $1/2$,
$64^2\,25 R[4(1+\delta)+16/\rho+4(1+|\beta|)]$,
which is bounded as stated.
\end{proof}

\paragraph{Complete two-shift delayed momentum forcing.}
For $\beta\in\{b,c\}$, $n\ge2$, $a\ge n-1$, put
$r=a+1$, $\theta=\beta+\lambda_{n-1}-a$, $\theta_2=\theta-D$.
The \emph{full} variable-coefficient normal form is
\[
 Z_{\beta,n-1}^{[2]}F_a=
 (\theta GB'+\theta\theta_2B^2)F_a+
 [GG'+(\theta+\theta_2)GB]F'_a+G^2F''_a.
\]
Since $|\theta|,|\theta_2|\le3r$, its three
radial-normalized coefficient norms are bounded by
$12(64)25^2$, $7(64)25^2$ and $64\,25^2$.
Using $r^2T_0\le4R$, $rT_1\le8R/\rho$ and
$T_2\le8R/\rho^2$ gives, for $\nu=1,2$,
\begin{equation}\label{eq:direct-mom-uniform}
 \left\|\frac12J_\nu Z_{\beta,n-1}^{[2]}F\right\|_{I_n}
 \le\frac{64^2\,25^2R}{2}
 \left(48+\frac{56}{\rho}+\frac8{\rho^2}\right)
 \|F\|_{I_{n-1}}.
\end{equation}
This proves an $n$-independent bound for the complete
\emph{direct} delayed swirl and axial momentum terms.

\paragraph{A pressure-norm obstruction to the stated majorant.}
The preceding theorems bound the velocity forcing. They do
not justify the unchanged normalized amplitude in equation
(52), which includes $\|\Pi_n\|_{I_{n-1}}/(n+1)$.
The isolated delayed-viscosity pressure primitive is not
a bounded map on that same weighted space.

\begin{proposition}[Explicit pressure-primitive counterexample]
Take $h=1/200$, $\rho=1/4$ and $R>0$.
With input $F=X/(1+4i\eta)\in I_1$,
the isolated delayed pressure source produces
$\Pi_3^{\rm del}=X^2P(\eta)$ where
$\|F\|\le8R$ but $\|\Pi_3^{\rm del}\|_{I_2}=\infty$.
\end{proposition}
\begin{proof}
Put $r=2$, $\lambda=2h$,
$W=B(r-D-\lambda)f/r-Gf'/r$,
$f=(1+4i\eta)^{-1}$, and
$Z_q^{(r)}=G\partial_\eta+B(q-r)$.
The exact pressure primitive from (24), (28), (27) is
\[
 P=\frac1{2r}Z_{\lambda-D}^{(r)}Z_\lambda^{(r)}W.
\]
At $\zeta=i/4$ the coefficient of its fourth-order pole is
\[
 \lim_{\eta\to\zeta}(\eta-\zeta)^4P(\eta)
 =-921187500i/4103684801\ne0.
\]
All other coefficient singularities are farther from the
origin. Thus $|P^{(m)}(0)|\asymp4^m m!\,m^3$.
But $w_{2,m}\asymp R^{-2}4^m m!\,m^2$, whereas
$w_{1,m}=R^{-1}4^m m!(m+2)/[8(m+1)]$.
In particular the precise scalar gap is
\[
 \frac{w_{1,m+3}}{4w_{2,m}}
 =\frac{27R}{4\rho^3}(m+5)
       \left(\frac{m+1}{m+4}\right)^2\longrightarrow\infty.
\]
This establishes the claimed unboundedness.
\end{proof}
This operator-level counterexample does not rule out the
specific recursively generated solution. A possible repair is
to eliminate the bare pressure norm from (52), use the
regularised pressure-to-velocity composite throughout the
induction, and reconstruct $\Pi_n$ in a weaker pressure
space, for example using $w_{a,m+1}$ in place of $w_{a,m}$.
The full recurrence (53) and Theorem~1.1 remain unproved
until that repair and all physical-coordinate reconstructions
are independently verified.

\paragraph{Analytic radii of the imported leading profiles.}
The leading fields of the hierarchy are the source profiles from
Appendix~B of \cite{openai2026}, with $Y=\Lambda X$.
Proposition~B.2 constructs these fields on $0\le Y\le4.1$
and a common neighborhood of $[-1,1]$ in $\eta$, once
admissible $\Lambda,C$ and axis data have been chosen. Its
proof selects a radius $S$ with $4.1<S<20$ and a smaller
positive complex $\eta$-neighborhood. This proves qualitative
regularity for fixed admissible parameters; it does not supply
explicit numerical values of $\Lambda,C,S,\delta$, nor a
certified complex supremum for one frozen realization.

\begin{lemma}[Source-radius transfer into the doubled norm]
\label{lem:source-radius-transfer}
Suppose $F(Y,\eta)$ is holomorphic on $|Y|<S$ and in
the $\delta$-neighborhood of $[-1,1]$, with $|F|\le M$.
Then $F_0(X,\eta)=F(\Lambda X,\eta)$ obeys
$\|F_0\|_{\mathcal B}\le M$ provided
$R\le S/(8\Lambda)$ and $\rho\le\delta/8$.
For any fixed integers $p,q\ge0$,
\begin{equation}\label{eq:source-Cauchy-bound}
\|D_X^p\partial_\eta^qF_0\|_{\mathcal B}
\le M\delta^{-q}\sup_{a,b\ge0}
\frac{a^p(a+1)^2(b+1)^2}{8^{a+b}}
\frac{(b+q)!}{b!}<\infty.
\end{equation}
\end{lemma}
\begin{proof}
Cauchy's inequalities give
$|\partial_\eta^{b+q}F_a|\le MS^{-a}(b+q)!\delta^{-b-q}$;
substitution of $Y=\Lambda X$ supplies $\Lambda^a$.
Divide by (37) and use $\binom{2a+b}{b}^{-1}\le1$.
The radial and parameter radius ratios are at most $1/8$.
For $p=q=0$, $(a+1)^2/8^a\le1$ and
$(b+1)^2/8^b\le1$. General fixed derivatives introduce
only polynomial factors absorbed by the geometric terms.
\end{proof}
The source's $U_0,\varphi_0$ and their finite derivatives
therefore have finite multiplier budgets after shrinking the
working radii. The incompressibility formula for $V_0/X$
uses only radial averages, rational factors with $L\ne0$,
and finite parameter derivatives. These operations preserve
holomorphy on a smaller common domain. This is an
\emph{existential} certification, not a numerical enclosure
of any fixed source profile.

\paragraph{The complete current-order swirl-pressure block.}
The two ordered terms $(0,n),(n,0)$ of (27) contribute
\begin{equation}\label{eq:current-pressure-identity}
\Pi_n^{\mathrm{cur}}(F)
=\frac2{C^2}\int_0^X\varphi_0(x,\eta)F(x,\eta)\,dx,
\qquad F=\varphi_n\in I_n.
\end{equation}
They belong in $K_n^{U\varphi}$, not in lower-order forcing.
\begin{lemma}[Quadratic radial gain in current pressure]
\label{lem:current-pressure-quadratic}
For $F\in I_n(\mathcal B)$,
\begin{align}
\left\|\int_0^X F\,dx\right\|_{\mathcal P}
&\le\frac{8R\rho}{(n+1)^2}\|F\|_{\mathcal B},\label{eq:axis-primitive-n2}\\
\|\Pi_n^{\mathrm{cur}}(F)\|_{\mathcal P}
&\le\frac{1024R\rho\|\varphi_0\|_{\mathcal B}}
{|C|^2(n+1)^2}\|F\|_{\mathcal B},\label{eq:current-pressure-n2}\\
\|K_n^{U\varphi}F\|_{\mathcal B}
&\le\frac{1024R\rho C_E\|\varphi_0\|_{\mathcal B}}
{|C|^2(n+1)^2}\|F\|_{\mathcal B},\label{eq:current-block-n2}
\end{align}
where $C_E$ is the pressure-to-velocity constant in
\eqref{eq:P-to-U}.
\end{lemma}
\begin{proof}
Factorial cancellation gives the exact weight ratio
\[
\frac{w_{a,b}}{(a+1)w_{a+1,b+1}}
=R\rho\frac{(a+2)^2(b+2)^2(2a+1)(2a+2)}
{(a+1)^3(b+1)^2(2a+b+1)(2a+b+2)(2a+b+3)}.
\]
Its maximum at fixed $a$ occurs for $b=0$, giving
$4R\rho(a+2)^2/[(a+1)^3(2a+3)]$.
The identity
$2(a+1)(2a+3)-(a+2)^2=3a^2+6a+2\ge0$
bounds this by $8R\rho/(a+1)^2$.
Since $a\ge n$ this proves the first inequality.
The Banach-algebra constant $64$, and the factor two
from the two ordered source pairs, prove the second.
Theorem~\ref{thm:pressure-to-velocity} proves the third.
\end{proof}

\paragraph{Conditional simultaneous pressure convergence.}
The preceding estimate replaces the earlier current-pressure
bound without double-counting $\varphi_0\varphi_n$.
If the three other current-block estimates
\eqref{eq:swirlbound}, \eqref{eq:phiuB} and \eqref{eq:uuB}
hold with a finite combined constant $C_{\rm other}$,
then
\begin{equation}\label{eq:current-K-updated}
\|K_n\|\le\frac{C_{\rm other}}{n+1}
+\frac{1024R\rho C_E\|\varphi_0\|}
{|C|^2(n+1)^2}.
\end{equation}
The Neumann-series inverse has norm at most two for
$n$ larger than twice the sum of the numerator constants.

Write $B_n=\|\varphi_n\|+\|U_n\|$, and suppose
$B_n\le MT^n/(n+1)^2$ for $n\ge1$ with $M,T\ge1$.
Then the complete shifted pressure ledger
\eqref{eq:complete-two-space-P-ledger},
after separating the current ordered pairs, implies
\begin{equation}\label{eq:conditional-pressure-geometric}
\|\Pi_n\|_{\mathcal P}
\le\frac{\Gamma T^n}{(n+1)^2},\quad n\ge3,
\end{equation}
where
\[
\Gamma=2C_{\rm sw}B_0M+4C_{\rm sw}M^2
+\frac{4C_LM}{T}+\frac{4C_DM}{T^2}
+\frac{4C_{\rm quad,P}M^2}{T},
\]
$B_0=\|\varphi_0\|+\|U_0\|$,
$C_L=C_{\rm P,time}+C_{\rm P,lap}+C_{\rm lead,P}$
and $C_D=C_{\rm P,del}$. The proof uses
$\sum_{i=1}^{n-1}[(i+1)^2(n-i+1)^2]^{-1}
\le4/(n+2)^2$ and the elementary adjacent-index
comparisons. The finite startup indices are absorbed by
enlarging $\Gamma$. This implication is not a proof that
the actual velocity sequence obeys the assumed bound.

\paragraph{Physical-coordinate reconstruction, conditional on convergence.}
With $\tau=q(1-\eta^2)$ and $z=q^D\eta$,
\begin{equation}\label{eq:coordinate-inverse-Jacobian}
\left.\frac\partial{\partial q}
[q-z^2q^{1-2D}]\right|_{z\ {\rm fixed}}
=1-2h\eta^2=L(\eta)>0.
\end{equation}
Thus the physical chain-rule coefficients are smooth on
compact subcores with $q>0$. If the velocity and pressure
weighted geometric estimates above hold on a common domain,
the radial valuation and coefficient weights give normal
convergence after each fixed $X,\eta$ derivative when
$T|q^{2h}|<2$ and $|X|<R/8$. Physical differentiation
then passes termwise on compact subcores. Exact recovery
of the unforced equations additionally requires the actual
coefficient hierarchy, including its startup normalization.

\paragraph{The finite-startup radius-selection gap.}
A bound $\|K_n\|\le C_K(\rho)/(n+1)$ may require a cutoff
$N_*(\rho)$. Merely knowing each finite-order coefficient
has some positive analytic radius $\rho_n$ does not
guarantee a common $\rho$ obeying
\begin{equation}\label{eq:radius-startup-fixed-point}
\rho\le\rho_{\rm lead},\qquad
\rho\le\min_{1\le n<N_*(\rho)}\rho_n.
\end{equation}
For example, the abstract choices $\rho_n=e^{-n}$ and
$N_*(\rho)=\lceil1/\rho\rceil+1$ provide positive radii
at every finite order but no $0<\rho\le1$ satisfying
the second inequality: its right-hand side is at most
$e^{-1/\rho}<\rho$. This does not refute the source
Navier--Stokes hierarchy. It shows that a common-radius
induction or noncircular quantitative initialization is
still required before the closed recurrence and
Theorem~1.1 can be claimed.


\paragraph{A fixed-radius inverse for every current-order block.}
The preceding Neumann argument used $\|K_n\|\le C_K/(n+1)$ and therefore treated finitely many small orders separately. That step is unnecessary if one retains radial valuation. Each term of the current-order operator $K_n$ has a final radial inverse $J_1$ or $J_2$; the pressure-feedback channel additionally has an axis primitive. The similarity order $n$ remains fixed while $K_n$ is iterated, but the radial valuation of the input increases.

\begin{theorem}[Factorially convergent current-order resolvent]\label{thm:factorial-resolvent}
Fix a common analytic coefficient space $\mathcal B_{R,\rho}$, and assume the finite leading-profile multiplier norms of Theorem~\ref{thm:fourblock}. Let $K_n$ denote the exact two-by-two operator (47)--(50). A finite constant $C_*$, independent of $n$ and of the input valuation $m$, exists such that, for $m\ge n\ge1$,
\begin{equation}\label{eq:valuation-tame-K}
 K_n(I_m\oplus I_m)\subseteq I_{m+1}\oplus I_{m+1},
 \qquad\|K_nF\|\le\frac{C_*}{m+1}\|F\|.
\end{equation}
Consequently,
\begin{equation}\label{eq:factorial-K}
\|K_n^\ell\|_{I_n^2\to I_{n+\ell}^2}
\le\frac{C_*^\ell n!}{(n+\ell)!}
\le\frac{C_*^\ell}{\ell!},
\end{equation}
and, for every $n\ge1$, $I-K_n$ is invertible on $I_n^2$ with
\begin{equation}\label{eq:fixed-radius-inverse}
 (I-K_n)^{-1}=\sum_{\ell\ge0}K_n^\ell,
 \qquad\|(I-K_n)^{-1}\|\le e^{C_*}.
\end{equation}
There is no radius-dependent exceptional startup order.
\end{theorem}
\begin{proof}
The source identities \eqref{eq:swirlblock}, \eqref{eq:Voperator}, \eqref{eq:phiu}, \eqref{eq:uu}, and \eqref{eq:pressureblock} express each block as a finite sum of complete radial-inverse composites. The estimates \eqref{eq:comp0}--\eqref{eq:compX} hold on arbitrary $I_m$: an undifferentiated term gains $(m+1)^{-2}$ and one $D_X$ or $\partial_\eta$ costs at most one power of $m+1$. The similarity coefficients cause no loss uniform in $m\ge n$, since
\[
 |b+\lambda_n|\le(|b|+2h)(m+1),\quad
 |c+\lambda_n|\le(|c|+2h)(m+1),\quad
 |D+\lambda_n|\le(|D|+2h)(m+1).
\]
The radial-velocity cross blocks include an average and satisfy
$\|\mathcal V_{n,\nu}^{g,q}|_{I_m}\|\le C_V(g,q)/(m+1)^2$.
The pressure block includes an axis primitive, followed by $J_1$;
using \eqref{eq:averageprim} at input valuation $m$ gives
$\|K_n^{U\varphi}|_{I_m}\|\le C_{U\varphi}/(m+1)^2$.
Fixed leading-profile multipliers have nonnegative radial degrees;
the final inverses increase valuation by at least one.
Thus $C_*=1+\max\{C_{\varphi\varphi}+C_{U\varphi},
 C_{UU}+C_V(m_\varphi,q_\varphi)\}$, with the constants
defined in Section~4, is sufficient.

Successively apply the gain at valuations $n,n+1,\ldots,n+\ell-1$.
The factorial inequality in \eqref{eq:factorial-K} follows because
$(n+1)\cdots(n+\ell)\ge\ell!$. The Neumann series converges
absolutely in operator norm, without the assumption $\|K_n\|<1$.
Its partial sums are inverses modulo the remainder $K_n^{N+1}$,
which tends to zero. Summation gives the bound $e^{C_*}$.
\end{proof}

\paragraph{A two-shift convolution majorant.}
The pressure source contains positive-order quadratic sums at
$i+j=n$ and $i+j=n-1$; these cannot be merged without an assumption
on the coefficient sequence.

\begin{proposition}[Constructive majorant with both quadratic shifts]
\label{prop:two-shift-majorant}
Suppose nonnegative $B_n$ are finite at $n=1,2$ and satisfy, for $n\ge3$,
\begin{equation}\label{eq:two-shift-recurrence}
 B_n\le H\left(
 B_{n-1}+B_{n-2}+
 \sum_{\substack{i+j=n\\i,j\ge1}}B_iB_j+
 \sum_{\substack{i+j=n-1\\i,j\ge1}}B_iB_j\right),
\end{equation}
with fixed $H\ge0$. Put $M=[32(1+H)]^{-1}$ and choose
\[
 T\ge\max\left\{2,16H,6\sqrt H,\frac{4B_1}{M},
 3\sqrt{\frac{B_2}{M}}\right\}.
\]
Then $B_n\le MT^n/(n+1)^2$ for all $n\ge1$.
\end{proposition}
\begin{proof}
The startup bounds follow from the choice of $T$.
Apply the inductive hypothesis and \eqref{eq:convolution}.
For $n\ge3$ the right side of \eqref{eq:two-shift-recurrence},
divided by $MT^n/(n+1)^2$, is at most
\[
 \frac{2H}{T}+\frac{4H}{T^2}
 +4HM\left(1+\frac1T\right)<1.
\]
The $n-1$ convolution contributes the final term $4HM/T$;
it is not identified with the $n$ convolution. Induction closes.
\end{proof}

\paragraph{Source-to-majorant reduction on a prescribed strip.}
Put $B_n=\|\varphi_n\|_{\mathcal B}+\|U_n\|_{\mathcal B}$.
Let $C_{\rm del,mom}$ and $C_{\rm mom}$ denote the direct delayed
and bilinear momentum constants above, $C_E$ the shifted
pressure-to-velocity bound, and
$C_L=C_{\rm P,time}+C_{\rm P,lap}+C_{\rm lead,P}$,
$C_D=C_{\rm P,del}$, $C_{\rm sw}$ and $C_{\rm quad,P}$
the shifted-pressure source constants in
\eqref{eq:complete-two-space-P-ledger}. After separating
$\varphi_0\varphi_n$ into the current-order block $K_n^{U\varphi}$,
the assembled lower-order estimates have the form
\begin{equation}\label{eq:fixed-radius-source-budget}
 \|S_nf_n\|\le L_1B_{n-1}+L_2B_{n-2}
 +Q_0\!\sum_{\substack{i+j=n\\i,j\ge1}}B_iB_j
 +Q_1\!\sum_{\substack{i+j=n-1\\i,j\ge1}}B_iB_j,
\end{equation}
with conservative budgets
\[
 L_1=2C_{\rm del,mom}+C_EC_L,\quad
 L_2=C_EC_D,\quad
 Q_0=2C_{\rm mom}+C_EC_{\rm sw},\quad
 Q_1=C_EC_{\rm quad,P}.
\]
They are independent of $n$ on one preselected common analytic
space. The source decomposition
\eqref{eq:lower-swirl}--\eqref{eq:lower-pressure}
and the previous component theorems provide this budget only
upon verifying their exact normalization and leading-profile
hypotheses for the proposed unforced hierarchy.

Under these stated hypotheses the inverse theorem gives
\eqref{eq:two-shift-recurrence} with
\begin{equation}\label{eq:source-majorant-H}
 H=e^{C_*}\max\{1,L_1,L_2,Q_0,Q_1\}.
\end{equation}
Proposition~\ref{prop:two-shift-majorant} then gives geometric
velocity coefficients on $T|q^{2h}|<2$ and $|X|<R/8$.
The shifted pressure ledger implies
$\|\Pi_n\|_{\mathcal P}\le\Gamma T^n/(n+1)^2$
with finite $\Gamma$ as in \eqref{eq:conditional-pressure-geometric}.
The physical coordinate Jacobian is $L(\eta)>0$ on the compact
subcore $q>0$, so each fixed physical derivative passes through
the normally convergent series once the exact startup source
equations and chain rule have been checked.

\paragraph{Scope of this result.}
Theorem~\ref{thm:factorial-resolvent} removes the circular
finite-startup radius cutoff, provided the analytic strip and all
leading multipliers have been fixed beforehand.
Proposition~B.2 of \cite{openai2026} gives an analytic source
leading profile for admissible parameters. The identification
of that particular profile, its pressure datum, and the unforced
hierarchy's startup normalization require independent correspondence
checks. These results are a conditional local reconstruction criterion,
not a proof of the proposed unforced Theorem~1.1 or a global singularity.


\subsection*{Startup coefficients and the remaining source certification}
The fixed-radius inverse and two-space majorant reduce the analytic argument
to a common-domain problem, but they do not by themselves identify the actual
leading datum or verify the first two source equations. Here we audit those
equations against (5.2)--(5.6) of \cite{openai2026}.

\paragraph{Leading data and quantitative radius.}
Appendix B starts with a pressure-axis datum $\Pi_0(0,\eta)$ from its
Lemma A.5 and $U_*(\eta)=4\eta+j_0$, $0<j_0\le0.05$. Its choices
$\delta_*,\sigma_*>0$ and a simply connected neighborhood $\Omega$ obey
nonvanishing conditions for $L$ and $H_*^2+\sigma_*^2$. It sets
\[
\zeta_*=-\frac{LH_*}{H_*^2+\sigma_*^2},\qquad
\varphi_*(\eta)=\exp\left(\Lambda\int_0^\eta\zeta_*(w)\,dw\right),
\]
then chooses $\Lambda$ sufficiently large and $C$ sufficiently large
relative to the complex supremum of $\varphi_*$. Proposition B.2
provides a local analytic solution after those dependent choices.
No particular numerical values for all these objects and their
common complex suprema are supplied. Thus a numerical frozen-profile
certificate cannot be inferred from the existence construction alone.

\paragraph{Exact source startup.}
For $n\ge1$, put $\lambda_n=2hn$ and adopt the source's convention
that negative-index fields vanish. Write
\[
\Pi_n^{\rm cur}=2C^{-2}\int_0^X\varphi_0\varphi_n\,dx,\qquad
\Pi_n^{\rm low}=\int_0^X\left[
 C^{-2}\sum_{i=1}^{n-1}\varphi_i\varphi_{n-i}
 -\frac{\Omega_{n-1}}{2x}\right]dx.
\]
The positive-order gauge is $\Pi_n(0,\eta)=0$, whereas the
order-zero trace $\Pi_0(0,\eta)$ is prescribed separately. The
complete first two radial-momentum sources are
\[
\begin{aligned}
\Omega_0={}&T_{0,0}V_0+
 V_0(V'_0-V_0/(2X))+U_0Z_{0,0}V_0-2XV''_0,\\
\Omega_1={}&T_{0,1}V_1+
 V_0(V'_1-V_1/(2X))+V_1(V'_0-V_0/(2X))\\
 &+U_0Z_{0,1}V_1+U_1Z_{0,0}V_0
 -2XV''_1-Z^{[2]}_{0,0}V_0.
\end{aligned}
\]

\begin{proposition}[Exact first- and second-order forcing]
\label{prop:startup-source-check}
Let $\mathcal K_n$ be the four current-order blocks (47)--(50),
$\mathcal R_\varphi=\tfrac12J_2$, $\mathcal R_U=\tfrac12J_1$
and $w_n=(\varphi_n,U_n)^\top$. Then
\[
w_n=\mathcal K_nw_n+
\begin{pmatrix}\mathcal R_\varphi F_n^\varphi\\
\mathcal R_U F_n^U\end{pmatrix}\qquad(n=1,2),
\]
with exact lower-order sources
\begin{align*}
F_1^\varphi={}&-Z_{b,0}^{[2]}\varphi_0,\\
F_1^U={}&-Z_{c,0}^{[2]}U_0+
 Z_{-2A,1}\left[-\frac12\int_0^X\frac{\Omega_0}{x}\,dx\right],\\
F_2^\varphi={}&-Z_{b,1}^{[2]}\varphi_1+
 V_1(\varphi_1'+\varphi_1/X)+U_1Z_{b,1}\varphi_1,\\
F_2^U={}&-Z_{c,1}^{[2]}U_1+V_1U'_1+U_1Z_{c,1}U_1\\
&+Z_{-2A,2}\int_0^X
  \left(C^{-2}\varphi_1^2-\frac{\Omega_1}{2x}\right)dx .
\end{align*}
In particular $\Omega_0$ has no negative-index delayed viscosity,
$\Omega_1$ includes $-Z^{[2]}_{0,0}V_0$, and the current pressure
factor $2C^{-2}\varphi_0\varphi_n$ is not counted as lower forcing.
\end{proposition}
\begin{proof}
The ordered-pair set $i+j=1$ consists of $(0,1),(1,0)$ and
has no positive-positive pair. At $n=2$ the single positive-positive
pair is $(1,1)$. In both cases $\Omega_{n-1}$ is known at the
beginning of order $n$ and the source convention removes $V_{-1}$.
Separating these pairs in (25)--(28) gives the displayed terms.
The incompressibility formula (24) implies $V_j=O(X)$, so the
pressure integrands are axis-regular. Apply (34)--(35) with zero
positive-order axis integration constants.
\end{proof}

\paragraph{The physical-coordinate identity.}
With $q$ defined implicitly by
$\tau=1-t=q(1-\eta^2)$, $z=q^D\eta$, $X=r^2/(2q)$,
and $L=1-2h\eta^2$, $d=1-\eta^2$, exact differentiation at
fixed physical coordinates yields
\[
\begin{gathered}
q_t=-L^{-1},\quad q_z=2\eta q^{1-D}L^{-1},\quad
\eta_t=D\eta(qL)^{-1},\quad \eta_z=q^{-D}dL^{-1},\\
X_t=X(qL)^{-1},\qquad X_z=-2\eta Xq^{-D}L^{-1}.
\end{gathered}
\]
Thus $\partial_t(q^aF)=q^{a-1}T_aF$ and
$\partial_z(q^aF)=q^{a-D}Z_aF$ exactly, with $L>0$ for
$0<h<1/100$ and $|\eta|\le1$.

\begin{lemma}[Conditional physical residual reconstruction]
\label{lem:conditional-physical-residual}
Assume one fixed leading profile satisfies the order-zero inner
balances, every positive-order coefficient satisfies (23)--(28)
with the exact gauges, and the physical velocity/pressure series
and every derivative entering the Navier--Stokes residual converge
locally uniformly on compact preterminal subcores. Then the
reconstructed field is divergence-free and satisfies the unforced
Navier--Stokes equations pointwise there, including across the axis.
\end{lemma}
\begin{proof}
On each such compact set, $q>0$ and $L>0$ have uniform positive
lower bounds. The displayed exact chain rules permit termwise
physical differentiation. Absolute convergence and the assumed
derivative bounds justify product convolution. At every expansion
order, the physical divergence and momentum residual coefficients
are (23)--(28), and vanish by hypothesis; hence the residual
vanishes. Axis regularity follows from $V_n=O(X)$, the explicit
factor $r$ in $u_\theta$, and analyticity in $X=r^2/(2q)$.
\end{proof}

The source's own higher-order existence lemma permits
order-dependent analytic neighborhoods. Its subsequent smooth
summation uses order-dependent cutoffs and does not assert
convergence of the unmodified infinite series. In addition,
its whole-space theorem has smooth compactly supported forcing.
The outstanding unforced verification requires an explicit common
profile-radius certificate, genuine solutions of the startup
equations above, the complete higher-order convergence estimate,
and a source-specific physical residual check.


\subsection*{An explicit complex fence for the axis data and the first two radial jets}
\label{subsec:axis-fence-jets}
The outstanding source certificate must start with Appendix B.1 of \cite{openai2026}. We isolate a zero-free domain for the rational axis multiplier without replacing the prescribed outer pressure trace by a surrogate.

For $0<h\le1/100$, $0<j_0\le1/20$, $0<\sigma\le1/10$, put
\[
 D=\tfrac12-h,\quad L(z)=1-2hz^2,\quad d(z)=1-z^2,\quad
 U_*(z)=4z+j_0,\quad H_*(z)=Dz+d(z)U_*(z)
\]
and $\mathcal R_\sigma=\{z=x+iy:|x|\le11/10,\ |y|\le\sigma/100\}$.

\begin{proposition}[Explicit zero-free axis-data neighborhood]
\label{prop:axis-data-zero-free}
On $\mathcal R_\sigma$,
\[
 |L(z)|\ge607/625,\qquad |H_*(z)^2+\sigma^2|\ge9\sigma^2/16.
\]
Thus $\zeta_*=-LH_*/(H_*^2+\sigma^2)$ is holomorphic on the interior of $\mathcal R_\sigma$, with
\[
 \sup|\zeta_*|\le25/\sigma^2,\qquad
 \sup_{z\in\mathcal R_\sigma}\left|\exp\left(\Lambda\int_0^z\zeta_*(w)\,dw\right)\right|
 \le\exp(30\Lambda/\sigma^2)
 \quad(\Lambda\ge0).
\]
In particular $C\ge\exp(30\Lambda/\sigma^2)$ suffices for condition (B.16) on this rectangle, but not automatically for the other conditions of Proposition B.2.
\end{proposition}
\begin{proof}
Since $|z|\le6/5$, one has
$|L-1|\le(2/100)(6/5)^2=18/625$. Furthermore
\[
 H_*'=D+4-12z^2-2j_0z,\qquad
 |H_*'|\le219/10<22\quad (|z|\le6/5).
\]
For real $x$, $H_*(x)$ is real. Integrating vertically gives
$|\operatorname{Im}H_*(x+iy)|\le22|y|\le11\sigma/50<\sigma/4$.
Consequently $|H_*\pm i\sigma|\ge3\sigma/4$. The bounds
$|L|\le643/625<103/100$ and $|H_*|\le6217/500<13$
give $|\zeta_*|\le(103/100)13\cdot16/(9\sigma^2)<25/\sigma^2$.
A path from $0$ first horizontally and then vertically has length
less than $6/5$ and stays inside the rectangle. Integrate and exponentiate.
\end{proof}

Cauchy's estimate on parameter disks of radius $\sigma/200$ shows that, if $C$ meets the displayed requirement, the normalized axis datum $g=\varphi_*/C$ satisfies
$|\partial_\eta^m g|\le m!(200/\sigma)^m$ on $[-1,1]$.
The degree-zero doubled coefficient norm therefore obeys
$\|g\|_{\mathcal B}\le1$ for $\rho\le\sigma/1600$.

For the \emph{diagnostic} tuple $h=1/200$, $j_0=1/50$,
$\sigma=1/100$, the polynomial is
\[
 H_*=-4\eta^3-\eta^2/50+(899/200)\eta+1/50.
\]
It obeys $H_*(-1/200)=-99/40000<0$,
$H_*(-1/250)=63123/31250000>0$ and
$\inf_{[-1,1]}(-W_*)\ge14701/5000>2.8$.
The denominator is zero-free on
$|\operatorname{Re}\eta|\le1.1$, $|\operatorname{Im}\eta|\le10^{-4}$,
with $|\zeta_*|\le250000$. This geometric diagnostic is
\emph{not} a source-admissible parameter tuple: condition (B.2)
requires the actual $Z_*$ and pressure trace.

\paragraph{The first two radial jets in the source profile.}
Use the exact rescaled radial equations (B.12)--(B.15). Write
\[
 \Phi=1+p_1Y+p_2Y^2+O(Y^3),\qquad
 u=q_1Y+q_2Y^2+O(Y^3),\qquad
 g=\varphi_*/C,
\]
with $\chi=H_*^2/(H_*^2+\sigma^2)$ and $\zeta_*=-LH_*/(H_*^2+\sigma^2)$. Define
\[
 A_0=W_*+h(1-2\eta U_*),\quad
 B_1=-D\eta q_1-\tfrac12d q_1',
\]
\[
 \mathcal R_{1,1}=
 \frac{(2W_*+h(1-2\eta U_*))p_1+H_*p_1'
       +\Lambda^{-1}(B_1-2h\eta q_1)+dq_1\zeta_*}{L},
\]
\[
 \mathcal R_{2,1}=
 [A(1-4\eta U_*)+dU_*'+W_*]q_1+H_*q_1'
       +d(g^2)'-(4A+2)\eta g^2.
\]
\begin{proposition}[Two exact rescaled radial coefficients]
\label{prop:axis-first-jets}
Every analytic solution with the prescribed Appendix-B data satisfies
\[
 p_1=\frac14\left(-\chi+\frac{A_0}{\Lambda L}\right),\quad
 q_1=-\frac{Z_*}{2L},\quad
 p_2=\frac{-\chi p_1+\Lambda^{-1}\mathcal R_{1,1}}{12},\quad
 q_2=\frac{\mathcal R_{2,1}}{8\Lambda}.
\]
The first pressure coefficients are
\[
 \Pi(Y,\eta)=\Pi_0(\eta)
 +\frac{g^2}{\Lambda}Y+\frac{g^2p_1}{\Lambda}Y^2+O(Y^3).
\]
\end{proposition}
\begin{proof}
Equation (B.14) gives
$W=W_*+\Lambda^{-1}[-2D\eta A_Yu-d\partial_\eta A_Yu]$,
$A_Yu=q_1Y/2+O(Y^2)$, and $H_c=H_*+\Lambda^{-1}du$.
The constant and linear coefficients in
$2(Y\Phi_{YY}+2\Phi_Y)=-\chi\Phi+\Lambda^{-1}R_1$
are $4p_1=-\chi+A_0/(\Lambda L)$ and
$12p_2=-\chi p_1+\Lambda^{-1}\mathcal R_{1,1}$.
In $2(Yu_{YY}+u_Y)=-Z_*/L+\Lambda^{-1}R_2$
they are $2q_1=-Z_*/L$ and
$8q_2=\mathcal R_{2,1}/\Lambda$.
Integrate $\Pi_Y=\Lambda^{-1}g^2\Phi^2$.
\end{proof}

\paragraph{The missing admissibility certificate.}
The source's $\Pi_0$ is the specific exterior-schedule integral
$\Pi_0(\eta)=-\frac12\int_{-\infty}^{\infty}
 E_{\mathrm{id,sched}}(y,\eta)^2\,dy$ of Lemma A.5.
A complete numerical certificate must freeze and interval-enclose
that entire schedule and integral; verify (B.2) for
$\delta_*,\sigma_*$; and provide concrete thresholds for
$\Lambda$, $C$, and the nonlinear contraction ball.
The rational zero-free domain alone proves none of these.
The displayed $Y$-jets of the \emph{leading} profile are not
the distinct similarity-order $n=1,2$ correction fields.
Neither their construction nor the unforced PDE reconstruction
is completed by this note.


\subsection*{The outer pressure datum as a positive measure}
\label{subsec:outer-pressure-measure}
The pressure datum of Lemma A.5 in \cite{openai2026} is defined using
the reference exterior schedule with its pressure-neutral angular
corrections omitted. The proof of that lemma states that each stage
has the form
\begin{equation}\label{eq:measure-profile}
 E_{\mathrm{id,sched}}(y,\eta)
   =c(y)(1+\eta^2)^{-\vartheta(y)},\qquad
 c(y)>0,\quad0\le\vartheta(y)\le1,
\end{equation}
where $c$ and $\vartheta$ are independent of $\eta$. Let $\mu$ be the
pushforward of $c(y)^2dy$ under $y\mapsto\vartheta(y)$, and put
$\mathfrak p=\frac12\mu([0,1])=-\Pi_0(0)>0$.

\begin{proposition}[Exact positive-measure pressure trace]
\label{prop:exact-pressure-measure}
The actual axis datum is exactly
\begin{equation}\label{eq:exact-pressure-measure}
 \Pi_0(\eta)=-\frac12\int_{[0,1]}
 (1+\eta^2)^{-2s}\,d\mu(s),\qquad
 \mu(\{1\})\ge5P_*^2.
\end{equation}
For $|\eta|\le1$,
\begin{equation}\label{eq:pressure-mass-bounds}
 -\mathfrak p\le\Pi_0(\eta)\le
 -\frac52P_*^2(1+\eta^2)^{-2},\qquad
 |\Pi_0'|\le2\mathfrak p,\quad
 |\Pi_0''|\le8\mathfrak p.
\end{equation}
Its first derivative has the strict sign of $\eta\ne0$.
The pressure extends holomorphically to
$|\operatorname{Re}z|\le11/10$,
$|\operatorname{Im}z|\le1/10$, with
\begin{equation}\label{eq:pressure-complex-bound}
 |\Pi_0(z)|\le(100/99)^2\mathfrak p,\qquad
 |\Pi_0^{(m)}(\eta)|\le(100/99)^2\mathfrak p\,m!20^m.
\end{equation}
\end{proposition}
\begin{proof}
The inner interval $y\le0$ has
$c(y)=P_*e^{y/10}$ and $\vartheta=1$; its contribution to
$\mu$ is $5P_*^2\delta_1$. The remaining stages contribute
a finite positive measure. For $w_s=(1+\eta^2)^{-2s}$,
\[
 \frac{w_s'}{w_s}=-\frac{4s\eta}{1+\eta^2},\qquad
 \frac{w_s''}{w_s}
 =-\frac{4s(1-\eta^2)}{(1+\eta^2)^2}
   +\frac{16s^2\eta^2}{(1+\eta^2)^2}.
\]
For $|\eta|\le1$, $|w_s'|\le2$,
$|w_s''|\le8$. Positivity of $\mu$ proves the real estimates;
the inner atom gives strictness of the derivative sign.
For complex $z$ in the stated rectangle,
$\operatorname{Re}(1+z^2)\ge99/100$; the principal complex
power is analytic and has modulus at most $(100/99)^2$.
Integrating over $\mu$ gives the complex bound.
Cauchy's estimate on radius-$1/20$ disks proves the last one.
\end{proof}

\subsection*{Quantitative choice of the B.2 parameters}
The full schedule is not numerically frozen, but the preceding
positivity permits explicit admissibility parameters depending
only on its pressure mass $\mathfrak p$ and the axis offset $j_0$.

\begin{proposition}[A constructive real-interval B.2 bound]
\label{prop:constructive-B2}
Assume the source Lemma A.5, $P_*\ge1$,
$0<h\le1/100$, $0<j_0\le1/20$. Set
\[
 U_*=4\eta+j_0,\quad D=\tfrac12-h,\quad
 A=\tfrac12+h,\quad d=1-\eta^2,\quad
 H_*=D\eta+dU_*,
\]
\[
 Z_*=-A(1-2\eta U_*)U_*-4H_*-d\Pi_0'
       +4A\eta\Pi_0.
\]
Let
\begin{equation}\label{eq:constructive-B2-parameters}
 K=200(1+\mathfrak p),\quad
 \delta_*=\frac{j_0}{12},\quad
 \epsilon=\frac{j_0}{12K},\quad
 \sigma_*=\frac{\epsilon}{30}.
\end{equation}
The unique zero $\eta_0$ of $H_*$ lies in
$(-j_0/4,-j_0/5)$. Moreover on $[-1,1]$,
\begin{equation}\label{eq:B2-certified}
 |Z_*|\le\delta_*
 \quad\Longrightarrow\quad |H_*|>\frac25\epsilon
 \quad\Longrightarrow\quad
 \frac{H_*^2}{H_*^2+\sigma_*^2}\ge\frac{144}{145}
 >\frac{99}{100}.
\end{equation}
These strict inequalities persist on some larger real compact
interval $I\supset[-1,1]$.
\end{proposition}
\begin{proof}
The quotient $H_*/d=D\eta/d+4\eta+j_0$
is strictly increasing on $(-1,1)$, and
$H_*(-j_0/4)=-Dj_0/4<0$ while
$H_*(-j_0/5)=(j_0/5)(1-D-j_0^2/25)>0$.
Write $t=-\eta_0$, so $j_0/5<t<j_0/4$.
At the root, $U_*(\eta_0)=Dt/(1-t^2)<j_0/7$
and $1-2\eta_0U_*(\eta_0)<2$.
The reference inner-tail inequality in
\eqref{eq:pressure-mass-bounds}, together with
$(1+t^2)^{-2}>1/2$, gives
$4A\eta_0\Pi_0(\eta_0)\ge5AtP_*^2\ge Aj_0P_*^2$.
Also $-d\Pi_0'(\eta_0)\ge0$. Therefore
\begin{equation}\label{eq:Z-root-positive}
 Z_*(\eta_0)>Aj_0(P_*^2-2/7)>j_0/3.
\end{equation}
Differentiating the exact formula for $Z_*$ and applying
\eqref{eq:pressure-mass-bounds} gives
$|Z_*'|\le120+20\mathfrak p<K$ on $[-1,1]$.
Indeed $|U_*|\le81/20$, $|H_*'|\le83/5$ and
$A\le51/100$ bound the polynomial terms.
Thus $|\eta-\eta_0|\le\epsilon$ implies
$Z_*(\eta)>j_0/3-K\epsilon=j_0/4>\delta_*$.

Factoring at the zero,
\[
 H_*(\eta)=(\eta-\eta_0)Q(\eta),\quad
 Q=D+4-4(\eta^2+\eta\eta_0+\eta_0^2)
      -j_0(\eta+\eta_0).
\]
The quadratic $Q$ is concave; on $[-1,1]$ its endpoint
values exceed $2/5$ because $D\ge49/100$ and
$|\eta_0|<j_0/4\le1/80$. Hence
$|H_*|>(2/5)|\eta-\eta_0|$. The condition
$|Z_*|\le\delta_*$ excludes $|\eta-\eta_0|\le\epsilon$.
Finally $\sigma_*=\epsilon/30$ yields $\chi\ge144/145$.
Compactness and strictness allow $I$ to be enlarged slightly.
\end{proof}

\paragraph{Unresolved numerical and dynamical data.}
The positive measure $\mu$ and its mass $\mathfrak p$ remain
defined by the actual full outer schedule and its moment matching,
not by an arbitrary model. Numerical evaluation requires freezing
that schedule and certified quadrature. The remaining nonlinear
Appendix-B construction requires effective values of $\Lambda$,
$C$, and a verified contraction ball. The theorem does not solve
the separate similarity-order $n=1,2$ unforced corrections or
establish the physical Navier--Stokes reconstruction.


\subsection*{What can be certified before freezing the exterior schedule}
\label{subsec:stagewise-pressure-enclosure}
The pressure datum of Lemma~A.5 depends on the reference exterior
\emph{after} its moment corrections. Appendix~A.2 selects $M_d$,
$T_d=e^{M_d}+10$, $P_*>e^{T_d}$, followed by sufficiently small
$\lambda,h$; pulse amplitudes and correction coefficients are then
defined by nonlinear moment equations. These prescriptions define
an existence family, not one published numerical pressure mass.

Let $Y_T$ denote the start of the final pure exterior power law
in the global logarithmic radial variable $y$. Omitting only the
pressure-neutral azimuthal corrections prescribed by Lemma~A.5,
write $E(y,\eta)=c(y)f(\eta)^{\vartheta(y)}$ on $0\le y\le Y_T$,
where $f(\eta)=(1+\eta^2)^{-1}$ and $0\le\vartheta\le1$.

\begin{proposition}[Exact source pressure decomposition and strengthened lower bound]
\label{prop:source-pressure-finite-tail}
With $c_T=c(Y_T)$ and $0<h<1/100$,
\begin{align}
\Pi_0(\eta)={}&-\frac52P_*^2f(\eta)^2
-\frac12\int_0^{Y_T}c(y)^2f(\eta)^{2\vartheta(y)}dy
-\frac{c_T^2}{2(1+2h)},\label{eq:pressure-finite-decomp}\\
\mathfrak p:=-\Pi_0(0)={}&\frac52P_*^2+
\frac12\int_0^{Y_T}c(y)^2dy+\frac{c_T^2}{2(1+2h)}.
\label{eq:pressure-mass-decomp}
\end{align}
In particular, the source's first two transitions imply
\begin{equation}\label{eq:source-early-mass}
\mathfrak p\ge (3-\tfrac12e^{-1-T_d})P_*^2
>(3-\tfrac12e^{-1-T_d})e^{2T_d}.
\end{equation}
The positive measure in Proposition~\ref{prop:exact-pressure-measure}
has an atom at exponent one with mass at least
$(6-e^{-1-T_d})P_*^2$.
\end{proposition}
\begin{proof}
The inner reference $y\le0$ has
$E=P_*f(\eta)e^{y/10}$ and thus squared integral
$5P_*^2f^2$. For $y\ge Y_T$, the source prescribes
$\ell=\partial_y\log(\sqrt X E)=-h$ and no $\eta$-dependence.
Thus $E=c_Te^{-(1/2+h)(y-Y_T)}$; integration gives
$c_T^2/(1+2h)$. This proves the decomposition.
On the unit transition following $y=0$,
$\ell=\frac35(1-\sigma(y))\in[0,3/5]$.
For the following axial transition $\ell=0$ over length $T_d$.
Both retain $\vartheta=1$; since
$\partial_y\log c=\ell-\tfrac12\ge-\tfrac12$,
we have $c(y)\ge P_*e^{-y/2}$ on $[0,1+T_d]$.
Its squared integral is at least
$P_*^2(1-e^{-1-T_d})$.
Add the reference inner stage, use positivity of the remaining
integral, and finally use $P_*>e^{T_d}$.
\end{proof}
For the conservative B.2 choice
$\sigma_*=j_0/[72000(1+\mathfrak p)]$,
this proves
\[
\sigma_*<
\frac{j_0}{72000(3-\frac12e^{-1-T_d})e^{2T_d}}.
\]
This is a consequence of that \emph{sufficient} choice, not a
universal upper bound on every admissible $\sigma_*$.

\paragraph{A rigorous stage-envelope algorithm.}
Suppose the actual moment-corrected reference profile has been
frozen and subdivided into finitely many stages
$0=y_0<\cdots<y_s=Y_T$ with lengths $d_j$.
Assume each stage has independently certified bounds
$\ell_j^-\le\ell(y)\le\ell_j^+$.
Let $s_j^\pm=\ell_j^\pm-\tfrac12$ and define
\[
F(s,d)=\int_0^d e^{2st}dt=
\begin{cases}(e^{2sd}-1)/(2s),&s\ne0,\\d,&s=0.\end{cases}
\]
Initialize $c_0^-=c_0^+=P_*$ and propagate
$c_{j+1}^\pm=c_j^\pm e^{s_j^\pm d_j}$.
Comparison of the exact ODE $c'=(\ell-\tfrac12)c$
then gives
\begin{equation}\label{eq:source-stage-enclosure}
\begin{split}
\mathfrak p_-&=\frac52P_*^2+
\frac12\sum_{j=0}^{s-1}(c_j^-)^2F(s_j^-,d_j)
+\frac{(c_s^-)^2}{2(1+2h)},\\
\mathfrak p_+&=\frac52P_*^2+
\frac12\sum_{j=0}^{s-1}(c_j^+)^2F(s_j^+,d_j)
+\frac{(c_s^+)^2}{2(1+2h)},\qquad
\mathfrak p_-\le\mathfrak p\le\mathfrak p_+.
\end{split}
\end{equation}
These interval bounds require genuinely certified inputs for all
stage slopes, matching amplitudes and lengths. The accompanying
outward-rounded program demonstrates the enclosure on a clearly
marked \emph{non-source} constant-slope fixture. It does not
choose or validate the source's missing moment-correction data.

\paragraph{An explicit angular linear inverse.}
The zero-free rectangle of Proposition~\ref{prop:axis-data-zero-free}
gives $|H_*|\le13$ and
$|H_*^2+\sigma_*^2|\ge9\sigma_*^2/16$. Therefore
\[
|\chi|=\left|\frac{H_*^2}{H_*^2+\sigma_*^2}\right|
\le\frac{2704}{9\sigma_*^2}<\frac{301}{\sigma_*^2}.
\]
For disks of radius $\sigma_*/200$ about $[-1,1]$ and
$\rho\le\sigma_*/1600$, Cauchy's inequality yields
$\|\chi\|_{\mathcal B_\rho}\le301/\sigma_*^2$
for a degree-zero coefficient of the source norm (B.4).
The source's product coefficient argument admits the constant
$64$, hence multiplication by $\chi$ has norm
$M_\chi\le19264/\sigma_*^2$.
Its factorial estimate for $T=\frac12J_2\chi$ gives
\begin{equation}\label{eq:source-explicit-linear-resolvent}
\|(1+T)^{-1}\|\le
\sum_{k\ge0}\frac{(40M_\chi)^k}{k!(k+1)!}
\le\exp(770560/\sigma_*^2).
\end{equation}
This very conservative estimate is for the angular
\emph{linear} inverse, not the full nonlinear problem.

For clarity, a sufficient numerical contraction criterion can be
written entirely in terms of the missing remainder norms.
Suppose on a ball of radius $r>0$ the exact remainders
$R_1,R_2$ in (B.15) have bounds $A_1,A_2$ and Lipschitz
constants $L_1,L_2$, uniformly in the remaining parameters.
If $B$ bounds the angular inverse above, the fixed-point map
is invariant and has Lipschitz constant at most $1/4$
provided
\begin{equation}\label{eq:source-Lambda-conditional}
\Lambda\ge
2\max\{40BA_1/r,\ 80A_2/r,\ 40BL_1,\ 80L_2\}.
\end{equation}
Indeed $J_2$ and $J_1$ have norms at most $40$ and $80$
in the source's space. This formula is conditional:
the original text proves the remainder bounds qualitatively
but does not freeze numerical $A_i,L_i$ on a selected ball.

\paragraph{Outstanding physical source data.}
The reference schedule must be fixed with compatible values of
$M_d,P_*,\lambda,h,T_f,c_o,\mathrm{Amp}$ and its exact
moment-correction coefficients, including the root of
(A.19)--(A.20). Without those data, neither
\eqref{eq:source-stage-enclosure} nor
\eqref{eq:source-Lambda-conditional} is a numerical certificate.
The actual similarity-order $n=1,2$ unforced corrections and the
physical residual have not been reconstructed.


\subsection*{An explicit Appendix-A pulse integral and two-moment conditioning}
The exact source step (A.5) obeys $\sigma(1-t)=1-\sigma(t)$. Define
$\phi_b(\xi)=\int_0^\xi \sigma(50v)\,dv$ and
$R_0(\xi)=\phi_b(\xi)(1-\sigma(\xi-10))$, as in
(A.9). The step has mean $1/2$ on $[0,1]$, so
\[
\phi_b(\xi)=\xi-1/100\quad(\xi\ge1/50),
\qquad 0\le\phi_b(\xi)\le1/100\quad(0\le\xi\le1/50).
\]
Also $R_0=0$ for $\xi\ge11$. This permits an exact
rational enclosure of (A.19)'s pulse coefficient:
\begin{proposition}[Pulse integral interval]
\[
0.2450478<K_b:=\int_0^{13}e^{-2\xi}R_0(\xi)^2d\xi<0.245052.
\]
The positive root of $K_bA^2=(1-e^{-26})/4$ is in $(1.009,1.011)$;
its numerical value is about $1.01005$.
\end{proposition}
\begin{proof}
The comparison integral over $[0,\infty)$ with
$R_0$ replaced by $\xi-1/100$ equals $4901/20000$.
The discrepancy in $[0,1/50]$ is at most
$1/500000$. The removed tail past $\xi=10$ is exactly
$(1102901/20000)e^{-20}$. The finite exponential sum
$\sum_{k=0}^{35}20^k/k!>470000000$ yields
$e^{-20}<1/470000000$ and the announced rational interval.
Similarly $\sum_{k=0}^{30}26^k/k!>10^{11}$ brackets the principal
amplitude by rational substitution.
\end{proof}
\begin{proposition}[Residual-to-amplitude certificate]
If the actual source remainder in (A.19) satisfies
$|\mathcal E(A,\eta)|\le1/2500$ and
$|\partial_A\mathcal E|\le1/100$ on
$[1.009,1.011]\times[-1,1]$, then
\[
K_bA^2-(1-e^{-26})/4+\mathcal E(A,\eta)=0
\]
has precisely one solution $A(\eta)\in(1.009,1.011)$ for each $\eta$.
Where differentiable, $|A'(\eta)|\le\sup|\partial_\eta\mathcal E|/0.48$.
\end{proposition}
\begin{proof}
At the left and right endpoints the full residual is respectively
less than $-0.000117$ and greater than $0.000068$.
Furthermore its $A$ derivative is greater than
$2(1.009)(0.2450478)-0.01>0.48$.
Use the intermediate value theorem, monotonicity, and implicit
differentiation.
\end{proof}
The two end corrections of (A.9) have normalized matrix (A.1)
\[
\mathcal B=\begin{pmatrix}b(s_1)&e^{2s_1}b(s_1)\\
b(s_2)&e^{2s_2}b(s_2)\end{pmatrix},
\quad s_1=1/2-\lambda,\quad s_2=1/2-2\lambda,
\quad b(s)=\int e^{sy}\beta(y)\,dy,
\]
where $\beta\ge0$, $\int\beta=1$, $\operatorname{supp}\beta
\subset[-3/20,3/20]$ and the bump centers differ by $2$.
\begin{proposition}[Explicit normalized two-moment inverse]
For $0<\lambda\le1/20$,
\[
|\det\mathcal B|\ge(1369/400)\lambda>3\lambda,
\qquad \|\mathcal B^{-1}\|_\infty<2/\lambda.
\]
Hence a certified normalized moment discrepancy $d$ produces
the unique linear correction $c=\mathcal B^{-1}d$ with
$\|c\|_\infty<2\|d\|_\infty/\lambda$.
\end{proposition}
\begin{proof}
The slopes belong to $[2/5,1/2]$ and
$37/40\le b(s_i)\le40/37<13/12$.
Moreover $e^{2s_2}>2$ and $e^{2\lambda}-1\ge2\lambda$.
Thus $\det\mathcal B$ has absolute value at least
$4(37/40)^2\lambda$. Each adjugate row sum is at most $13/2$.
Divide and use $(13/2)/(1369/400)=2600/1369<2$.
\end{proof}
The true source remainder is only bounded qualitatively by
$C_{\rm pre}\lambda(1+\log(1/\lambda))$, with $C_{\rm pre}$
unfrozen. These results do not determine a source-admissible
$\lambda$, $\mathcal E$, corrected exterior schedule, numerical
$M_0$, nonlinear Appendix-B contraction, or unforced startup
solutions. They furnish explicit verifiable thresholds, not a
completed Navier--Stokes blowup proof.


\subsection*{An exact quadratic form of the Appendix-A pulse remainder}
\label{subsec:exact-A19-quad}
Let $0<\lambda\le1/20$, $X_{\rm end}=X_pe^{13/\lambda}$, and
$E(X_p,\eta)=e_bf(\eta)$, $f=(1+\eta^2)^{-1}$. On the pulse,
$E(X_pe^y,\eta)=e_bf e^{-(1/2+\lambda)y}$ and
$U=E\{AR_0(\lambda y)+c_1\beta_1(y)+c_2\beta_2(y)\}$.
The two translated nonnegative normalized bumps of width $3/10$
are centered at $y_1=13/\lambda-3$ and $y_2=y_1+2$.
Define
\[
P_-=\frac{\lambda S(X_p)}{X_pe_b^2f^2},\quad
D_+=\frac{\lambda}{2X_pe_b^2f^2}\int_{X_{\rm end}}^\infty E^2\,dX,
\quad J_i=\lambda\int_0^{13/\lambda}e^{-2\lambda y}\beta_i(y)^2dy.
\]
The source constructs $E$ after the pulse independently of $A$,
and $U$ vanishes there; hence $P_-,D_+$ are independent of $A$.

\begin{proposition}[Exact remainder and quadratic amplitude equation]
\label{prop:source-exact-A19}
The remainder in (A.19) of \cite{openai2026} is exactly
\begin{equation}\label{eq:A19-exact-remainder}
 \mathcal E(A,\eta)=P_-(\eta)-D_+(\eta)+
              J_1c_1(A,\eta)^2+J_2c_2(A,\eta)^2.
\end{equation}
Let $s_1=1/2-\lambda$, $s_2=1/2-2\lambda$,
and let $\mathcal B$ be the normalized (A.1) two-bump matrix.
Set $m_M=M(X_p)/(X_pe_bf)$ and
$m_J=J(X_p)/(X_pH(X_p)e_bf)$, where $J(X)$ here denotes
the source's accumulated $UH$-moment.
With
\[
g_i=\int_0^{11/\lambda}e^{s_i(y-y_1)}R_0(\lambda y)\,dy,\quad
a=-\mathcal B^{-1}\binom{e^{-s_1y_1}m_M}{e^{-s_2y_1}m_J},
\quad b=-\mathcal B^{-1}\binom{g_1}{g_2},
\]
one has $c_i=a_i+A b_i$ exactly, and $b$ does not depend on $\eta$.
The complete moment equation is
\[
k_2A^2+k_1A+k_0=0,\quad
k_2=K_b+\sum_iJ_ib_i^2>0,\quad k_1=2\sum_iJ_ia_ib_i,
\]
\[
k_0=P_--D_++\sum_iJ_ia_i^2-(1-e^{-26})/4.
\]
In particular $\partial_A^2\mathcal E=2\sum_iJ_ib_i^2\ge0$.
If $k_0<0$, its unique positive root is
$A_*=(-k_1+\sqrt{k_1^2-4k_2k_0})/(2k_2)$.
\end{proposition}
\begin{proof}
The substitution $\xi=\lambda y$, $dX=X_pe^y\,dy$ in the
source definition $S=\int(U^2-E^2/2)\,dX$ gives its
normalized pulse contribution as
$\int_0^{13}e^{-2\xi}\{(AR_0(\xi)+c_1\beta_1(\xi/\lambda)
+c_2\beta_2(\xi/\lambda))^2-1/2\}\,d\xi$.
The principal pulse vanishes for $\xi\ge11$, while the
first bump begins at $\xi\ge13-(63/20)\lambda>11$.
The bumps are disjoint. Thus every pulse--bump and
bump--bump cross term is identically zero. The remaining
pre-pulse and post-pulse contributions give
\eqref{eq:A19-exact-remainder}. The two normalized $M,J$
moments are affine in the pulse amplitude, with weights
$e^{s_1y}$ and $e^{s_2y}$, giving $c=a+Ab$.
Expansion gives the quadratic coefficients and the root criterion.
\end{proof}

\begin{proposition}[Explicit amplitude-sensitive error]
\label{prop:A19-error-budget}
Let $B_2=\int\beta^2$, $D_0=\max(|e^{-s_1y_1}m_M|,
|e^{-s_2y_1}m_J|)$ and $D_1=\max(|g_1|,|g_2|)$.
For $0<\lambda\le1/20$,
\[
D_0\le e^{-(2/5)(13/\lambda-3)}\max(|m_M|,|m_J|),\quad
D_1\le\frac{55}{2}e^{-(2/5)(2/\lambda-3)},\quad
J_1+J_2\le2\lambda B_2 e^{-26+(63/10)\lambda}.
\]
Consequently, for $|A|\le A_{\max}$,
\begin{align*}
0\le\mathcal E-(P_--D_+)&\le
\frac{8B_2}{\lambda}e^{-26+(63/10)\lambda}
(D_0+A_{\max}D_1)^2,\\
|\partial_A\mathcal E|&\le
\frac{16B_2}{\lambda}e^{-26+(63/10)\lambda}
(D_0+A_{\max}D_1)D_1.
\end{align*}
The source-prescribed correction bump is a width-$3/10$ rescaling
of the fixed step derivative $\sigma'$ in (A.5), not a freely chosen
plateau. Its exact analytic envelope is $27/5<B_2\le640/3$;
see Proposition~\ref{prop:true-source-bump-L2}.
If, in addition, the \emph{actual} pre-pulse moments satisfy
$\max(|m_M|,|m_J|)\le1$, then uniformly for
$0<\lambda\le1/20$, $|A|\le1.011$,
\[
0\le\mathcal E-(P_--D_+)<5\cdot10^{-17},
\qquad |\partial_A\mathcal E|<10^{-16}.
\]
\end{proposition}
\begin{proof}
Since $s_i\ge2/5$ and $0\le R_0(\xi)\le11$ is supported
on $[0,11]$,
$|g_i|\le(11/s_i)e^{-s_i(2/\lambda-3)}$.
The left bump endpoint gives the displayed $J_i$ bound.
The inverse estimate \eqref{eq:two-bump-inverse} yields
$\|a\|_\infty\le2D_0/\lambda$ and
$\|b\|_\infty\le2D_1/\lambda$. Substitution proves both bounds.
For the numerical consequence, $t=1/\lambda\ge20$ and
$t e^{-8t/5}$ decreases with $t$. Exact finite rational
exponential series imply $e^{-74/5}<4\cdot10^{-7}$ and
$e^{-5137/200}<10^{-11}$, which prove the asserted
strict numerical inequalities.
\end{proof}

The actual normalized $m_M,m_J$ and the independent balance
$P_--D_+$ are not yet enclosed for one source-admissible
exterior schedule. This exact polynomial reduction therefore
does not certify a numerical amplitude, the nonlinear
Appendix-B profile or any unforced Navier--Stokes singularity.


\subsection*{Moment transport from the Appendix-A reference radius}
The source's reference region (A.7) has
$U=4\eta$, $E=P_*f(\eta)(X/X_R)^{1/10}$,
$f=(1+\eta^2)^{-1}$, and $H=\sqrt{2X}E$.
The cumulative moments are $M=\int_0^XU\,dx$,
$J=\int_0^XUH\,dx$, and
$S=\int_0^X(U^2-E^2/2)\,dx$, as in (4.15).

\begin{proposition}[Exact source moment anchor]
At the first matching radius,
\begin{equation}\label{eq:source-ref-MJS}
\begin{aligned}
M(X_R)&=4\eta X_R,\\
J(X_R)&=\frac52\eta X_R\sqrt{2X_R}\,P_*f,\\
S(X_R)&=X_R\left(16\eta^2-\frac5{12}P_*^2f^2\right).
\end{aligned}
\end{equation}
\end{proposition}
\begin{proof}
Integrate the source's reference power laws. The three needed
integrals over $0\le x=X/X_R\le1$ are
$\int_0^1dx=1$,
$\int_0^1x^{3/5}dx=5/8$ and
$\int_0^1x^{1/5}dx=5/6$.
\end{proof}

Put $y=\log(X/X_R)$,
$a=U/E$, $\ell=\partial_y\log(\sqrt XE)$,
and define the scale-invariant moments
\[
m=\frac{M}{XE},\qquad j=\frac{J}{XHE},\qquad
s=\frac{S}{XE^2},\qquad
d=\frac1{2XE^2}\int_X^\infty E(x,\eta)^2\,dx.
\]
\begin{theorem}[Exact four-moment transport]\label{thm:four-moment-transport}
On every source stage,
\begin{equation}\label{eq:four-moment-transport}
\begin{aligned}
m'&=a-(\ell+\tfrac12)m,&
j'&=a-(2\ell+\tfrac12)j,\\
s'&=a^2-\tfrac12-2\ell s,&
d'&=-\tfrac12-2\ell d.
\end{aligned}
\end{equation}
The forward initial values at $X_R$ are
$m=4\eta/(P_*f)$,
$j=5\eta/(2P_*f)$ and
$s=16\eta^2/(P_*^2f^2)-5/12$.
At the start $X_T$ of the source's pure power-law tail,
$d(X_T)=1/(4h)$.
If $X_p$ is the pulse start and
$X_{\mathrm{end}}=X_pe^{13/\lambda}$ its end, the
amplitude-independent remainder from Proposition~\ref{prop:source-exact-A19}
obeys the exact identity
\begin{equation}\label{eq:source-bias-transport}
P_{\mathrm{pre}}-D_{\mathrm{post}}
=\lambda\bigl(s(X_p)-e^{-26}d(X_{\mathrm{end}})\bigr).
\end{equation}
\end{theorem}
\begin{proof}
The logarithmic derivatives of $XE$, $XHE$, and $XE^2$
are $\ell+1/2$, $2\ell+1/2$, and $2\ell$,
respectively. Use $dM/dy=XU$,
$dJ/dy=XUH$,
$dS/dy=X(U^2-E^2/2)$, and
$d(\frac12\int_X^\infty E^2dx)/dy=-XE^2/2$
to prove the four transport equations.
For $X\ge X_T$, $E=E_T(X/X_T)^{-1/2-h}$;
hence $\frac12\int_{X_T}^\infty E^2dx=X_TE_T^2/(4h)$.
On the pulse,
$XE^2/(X_pe_b^2f^2)=e^{-2\lambda\log(X/X_p)}$;
at $X_{\mathrm{end}}$ this is $e^{-26}$. The bias
identity follows directly from its definitions.
\end{proof}
The first three equations admit forward integrating-factor
formulas; the last is integrated backward from $d(X_T)=1/(4h)$.
An interval calculation need only enclose $a(y,\eta)$ and
$\ell(y,\eta)$ on the \emph{actual} finitely many Appendix-A
stages. It is not necessary to introduce an artificial
semi-infinite pre-pulse profile or an approximate terminal tail.

\subsection*{Correction to the prescribed moment-bump estimate}
Appendix A fixes the positive correction bumps as rescaled copies
of $\sigma'$ from (A.5), on intervals of width $3/10$.
The former use of an alternative normalized plateau bump with
$\int\beta^2\le5$ was not a certificate for this actual
source shape. For the prescribed bump,
\[
\beta(y)=\frac{10}{3}\sigma'\left(
 \frac{10}{3}(y-y_c)+\frac12\right),\qquad
\sigma'(t)=\sigma(t)(1-\sigma(t))
\left(\frac2{t^3}+\frac2{(1-t)^3}\right).
\]
\begin{proposition}[Prescribed source-bump norm and corrected error]
\label{prop:true-source-bump-L2}
The actual normalized smooth-step-derivative bump satisfies
\begin{equation}\label{eq:source-bump-correct}
\frac{27}{5}<B_2:=\int\beta(y)^2\,dy\le\frac{640}{3}.
\end{equation}
If the actual transported source moments satisfy
$\max(|m(X_p)|,|j(X_p)|)\le1$, then for
$0<\lambda\le1/20$, $|A|\le1.011$ the
complete amplitude remainder satisfies
\[
0\le\mathcal E-(P_{\mathrm{pre}}-D_{\mathrm{post}})
<5\cdot10^{-17},\qquad
|\partial_A\mathcal E|<10^{-16}.
\]
\end{proposition}
\begin{proof}
The symmetry $\sigma(1-t)=1-\sigma(t)$ gives
$\sigma(1/4)=1/(1+e^{128/9})<1/20$,
and $\sigma(3/4)>19/20$.
Cauchy--Schwarz on $[1/4,3/4]$ yields
$\int_0^1\sigma'^2>2(9/10)^2$.
On $[1/4,3/4]$ one has $\sigma'\le64$.
For $0<t\le1/4$,
$\sigma'\le e^{2-16}(128+128/27)<1$ using
monotonicity of $t^{-3}e^{-1/t^2}$.
Symmetry handles the remaining quarter. Hence
$\int_0^1\sigma'^2\le64\int_0^1\sigma'=64$,
and scaling by $10/3$ proves the norm bounds.

Insert $B_2\le640/3$ into the already proved
positive remainder inequalities. With $t=1/\lambda\ge20$,
$D_0+1.011D_1\le(57605/2000)e^{6/5-4t/5}$.
The rigorously verified exponential bounds
$e^{-74/5}<4\cdot10^{-7}$,
$e^{-5137/200}<10^{-11}$ and monotonicity of
$t e^{-8t/5+63/(10t)}$ give exact rational upper
bounds $4.531\cdot10^{-17}$ and
$8.652\cdot10^{-17}$ for the two errors.
\end{proof}
A fully source-admissible amplitude interval therefore requires
\[
\max(|m(X_p)|,|j(X_p)|)\le1,\qquad
\lambda\sup_{|\eta|\le1}
|s(X_p,\eta)-e^{-26}d(X_{\mathrm{end}},\eta)|
\le\frac1{2500}-5\cdot10^{-17}.
\]
Neither condition has yet been numerically certified for one
frozen, moment-matched exterior schedule. The corrected
remainder and transport formulas establish no unforced
Navier--Stokes solution or physical residual theorem.


\subsection*{A uniform pulse-entrance envelope from the prescribed source stages}
Write $P=P_*$, $T=T_d=e^{M_d}+10>10$ and $f=(1+\eta^2)^{-1}$. Retain the normalized moment notation
$a=U/E$, $m=M/(XE)$, $j=J/(XHE)$,
$s=S/(XE^2)$ and $d=(2XE^2)^{-1}\int_X^\infty E^2\,dx$.
These satisfy the exact four transport equations
\[
m'=a-(\ell+\tfrac12)m,\quad j'=a-(2\ell+\tfrac12)j,\quad
s'=a^2-\tfrac12-2\ell s,\quad d'=-\tfrac12-2\ell d.
\]
The prime is differentiation in $\log(X/X_R)$.

\begin{proposition}[Source-native pre-pulse moment bound]\label{prop:uniform-source-prepulse}
For the Appendix-A reference stages (A.7)--(A.9), $P>e^T$ and
$0<\lambda\le1/20$ imply
\[
\sup_{|\eta|\le1}\max\{|m(X_p,\eta)|,|j(X_p,\eta)|\}<\frac9{100}.
\]
\end{proposition}
\begin{proof}
The exact reference integrals at $X_R$ yield
$|m|\le8/P$, $|j|\le5/P$. On the first unit interval,
$U=4\eta$, $0\le\ell\le3/5$, and
$E\ge Pf e^{-y/2}$; consequently $|a|<16/P$.
Both moment equations have damping at least $1/2$, giving
$|m|,|j|\le40/P$ at the end.
During the following axial interval of length $T$,
$0\le U/\eta\le4$, $\ell=0$, and
$|a(y)|\le8e^{(1+y)/2}/P$. Integrating the two equations gives
\[
\max(|m|,|j|)
\le\frac{40e^{-T/2}}P+
\frac{8e^{1/2-T/2}}P(e^T-1)
<40e^{-3T/2}+8e^{(1-T)/2}.
\]
Since $T>10$, the last expression is bounded by
$40e^{-15}+8e^{-9/2}<9/100$.
The final inequality uses exact rational exponential-series lower
bounds $e^{15}>10^6$, $e^{9/2}>90$.
Thereafter $U=0$, while $-\lambda\le\ell\le0$ through the
one-unit transition and intermediate power-law interval.
Both damping coefficients remain positive, so neither absolute
moment increases.
\end{proof}

\begin{proposition}[Explicit smallness of the forward $S$-moment]
\label{prop:source-forward-s-explicit}
Let $T_f>0$ be the source's angular interpolation length and
$k=T+T_f+1$. If
\[
0<\lambda\le\frac{10^{-6}}k,
\]
then $\lambda\sup_{|\eta|\le1}|s(X_p,\eta)|<10^{-4}$.
\end{proposition}
\begin{proof}
At $X_R$, $|s|\le5/12+64/P^2$. The first unit interval has
nonnegative damping $2\ell$ and $|a|<16/P$, giving
$|s|<11/12+320/P^2<1$. On the axial interval,
$\int_0^T a^2\,dy\le64e(e^T-1)/P^2<64e^{1-T}<1$,
hence $|s|<T/2+2$. On the one-unit slope transition,
$|s|<e^{2\lambda}(T/2+5/2)<(9/8)(T/2+5/2)$.
For the subsequent constant slope $-\lambda$ lasting
$W=60\log(1/\lambda)$, one has exactly
\[
s(X_p)=q s_{\rm in}-\frac{q-1}{4\lambda},\qquad
q=e^{120\lambda\log(1/\lambda)}.
\]
For $k\ge11$, monotonicity and $e^{17}>11\cdot10^6$ imply
$120\lambda\log(1/\lambda)<1/5000$, so $q<5000/4999$.
Also $(T+5)/k\le15/11$. Therefore
\[
\lambda|s(X_p)|
<10^{-6}\frac{5000}{4999}\frac9{16}\frac{15}{11}
+\frac1{4\cdot4999}<10^{-4}.
\]
\end{proof}

\begin{proposition}[Conditional backward exterior envelope and A.19 root]
\label{prop:source-full-bias-explicit}
Suppose the actual azimuthal moment corrections meet the source's
stated post-pulse slope bounds: $\ell\le0$ on the angular
interpolation and its following $30\log(1/\lambda)$ plateau,
$\ell$ decreases to $-1$ over one further unit interval,
then equals $-1$ for $4\log(1/h)$, and satisfies
$\ell\le-3h/4$ afterward. For $0<h\le1/100$,
\[
0\le d(X_{\rm end})
\le\frac{T_f}{2}+15\log(1/\lambda)+\frac34+\frac{h^7}{3}.
\]
If also $\lambda\le10^{-6}/(T+T_f+1)$, then
\[
\lambda\sup_{|\eta|\le1}
|s(X_p,\eta)-e^{-26}d(X_{\rm end},\eta)|<\frac1{8000}.
\]
Together with the corrected fixed-$\sigma'$ bump estimate and
Proposition~\ref{prop:uniform-source-prepulse}, the true A.19
remainder satisfies $|\mathcal E|<1/2500$ and
$|\partial_A\mathcal E|<1/100$ on the proposed amplitude interval.
Thus the genuine moment equation has a unique root
$1.009<A(\eta)<1.011$ under these simultaneous stage hypotheses.
\end{proposition}
\begin{proof}
Write $R(v)=XE^2/(X_{\rm end}E_{\rm end}^2)$ for
$v=\log(X/X_{\rm end})$. Then $R(0)=1$, $R'=2\ell R$ and
$d(X_{\rm end})=\tfrac12\int_0^\infty R(v)\,dv$.
The first stages contribute at most
$T_f/2+15\log(1/\lambda)+1/2$; the $-1$ hold contributes at most
$1/4$. At its end $R\le h^8$, and the remaining
$\ell\le-3h/4$ tail contributes at most $h^7/3$.
Set $c=10^{-6}$ and $k=T+T_f+1\ge11$.
The elementary inequalities $\lambda T_f\le c$ and
$\lambda\log(1/\lambda)\le17c/11$ yield
\[
\lambda e^{-26}d(X_{\rm end})\le\lambda d(X_{\rm end})
<\frac c2+\frac{255c}{11}
+\frac c{11}\left(\frac34+\frac1{3\cdot100^7}\right)
<\frac1{40000}.
\]
Together with the preceding forward bound this gives
$1/10000+1/40000=1/8000$. The source-specific corrected
bump contribution is $<5\cdot10^{-17}$ and its amplitude
derivative $<10^{-16}$, because $|m|,|j|<9/100<1$.
Apply the previously established A.19 endpoint signs and
strict derivative inequality.
\end{proof}

The estimate is a theorem about stage bounds, not a fully
frozen numerical source realization. In particular, compatibility
of the choice of $\lambda$ with the source's later selection of $T_f$
and the corrected azimuthal moment bumps must be verified. The
exterior pressure mass, Appendix-B nonlinear contraction,
unforced $n=1,2$ corrections and physical Navier--Stokes residual
remain unverified.

\subsection*{Quantitative azimuthal moment restoration at (A.11)}
\label{subsec:exact-A11}
The source's angular correction must be evaluated with the prescribed
translated copies of $\sigma'$, rather than arbitrary moment bumps.
We give an explicit inverse and a quantitative nonlinear correction
on the post-pulse constant-slope interval. This analysis concerns
only the two moment equalities in (A.11); it does not freeze the
preceding stress-cone parameters or the later exterior matching time.

Put $y=\log(X/X_1)$, where $X_1$ is the center of the first
correction bump. The second center is $y=2$, and the unedited
profile has $\ell=-\lambda$ between them. Let
\[
 \beta(t)=\frac{10}{3}\sigma'\!\left(\frac{10t}{3}+\frac12\right),
 \quad \beta_1(y)=\beta(y),\quad\beta_2(y)=\beta(y-2),
 \quad \operatorname{supp}\beta\subset[-3/20,3/20].
\]
Here $\int\beta=1$ and, by the corrected source-bump lemma,
$\int\beta^2\le640/3$. Write
$E_c=E_{\rm base}(1+c_1\beta_1+c_2\beta_2)$.
The source's angular-moment equation and its pressure-neutrality
equation are precisely
\begin{equation}\label{eq:A11-two-exact}
 \begin{split}
 u(c_1+xc_2)&=\Delta_I,\\
 2v(c_1+zc_2)+w(c_1^2+zc_2^2)&=0,
 \end{split}
\end{equation}
where the positive source-bump integrals are
\begin{equation}\label{eq:A11-weights}
\begin{gathered}
 u=\int e^{(1-\lambda)t}\beta(t)\,dt,\qquad
 v=\int e^{-(1+2\lambda)t}\beta(t)\,dt,\qquad
 w=\int e^{-(1+2\lambda)t}\beta(t)^2\,dt,\\
 x=e^{2(1-\lambda)},\qquad z=e^{-2(1+2\lambda)},\qquad
 \Delta_I=\frac1{1-\lambda}-\frac{I(X_1)}{X_1H(X_1)}.
\end{gathered}
\end{equation}
Equivalently, put $\alpha=\Delta_I/u$ and eliminate
$c_1=\alpha-xc_2$. Then the \emph{actual} two-equation
correction is reduced to the scalar quadratic
\[
 \mathsf A c_2^2+\mathsf Bc_2+\mathsf C=0,\quad
 \mathsf A=w(x^2+z),\quad
 \mathsf B=2v(z-x)-2w\alpha x,\quad
 \mathsf C=2v\alpha+w\alpha^2.
\]
On the small branch supplied below $\mathsf B<0$ and the
discriminant is positive. A cancellation-safe evaluation is
$c_2=2\mathsf C/[\sqrt{\mathsf B^2-4\mathsf A\mathsf C}-\mathsf B]$,
followed by $c_1=\alpha-xc_2$. This is an exact specification
of the two source correction coefficients once $\Delta_I$ is known.

The first identity follows by evaluating the change in
$I=\int H\,dX$ and normalizing by $X_1H(X_1)$, including the
fact that $I/(XH)$ approaches $1/(1-\lambda)$ under the
unedited constant-slope flow. The second is exactly
$\int(E_c^2-E_{\rm base}^2)\,dy=0$: disjoint supports
eliminate its mixed quadratic term.

\begin{proposition}[Uniform two-bump inverse and quadratic bound]
\label{prop:A11-matrix-bound}
For $0<\lambda\le1/20$, the Jacobian and quadratic part of
\eqref{eq:A11-two-exact} satisfy
\[
 B_{11}=\begin{pmatrix}u&ux\\2v&2vz\end{pmatrix},
 \qquad Q(c,c)=\binom{0}{w(c_1^2+zc_2^2)},
 \qquad\|B_{11}^{-1}\|_\infty<\frac32,\quad\|Q\|<320.
\]
In particular, its determinant is bounded away from zero
uniformly as $\lambda\downarrow0$.
\end{proposition}
\begin{proof}
Since $|t|\le3/20$ on the support, both $u$ and $v$ belong
to $[4/5,5/4]$ by $e^{-1/5}\ge4/5$ and
$e^{1/5}\le5/4$. Also $6<x<8$ and $0<z<1/7$.
The determinant is $2uv(z-x)$, with magnitude exceeding
$2(4/5)^2(6-1/7)=1312/175$. Row sums of its inverse
are bounded by
\[
 \frac{2(5/4)(1/7)+(5/4)8}{1312/175}
 =\frac{3625}{2624}<\frac32,
 \qquad
 \frac{2(5/4)+5/4}{1312/175}
 =\frac{2625}{5248}<\frac32.
\]
The bound on the prescribed bump implies
$w\le(5/4)(640/3)$, and hence
$\|Q\|\le w(1+z)<6400/21<320$.
\end{proof}

\begin{proposition}[Source moment discrepancy at the first bump]
\label{prop:A11-discrepancy}
Suppose the prescribed pre-pulse and interpolation stages
satisfy their exact source bounds, namely
$\ell\ge-\lambda-1/10$ and
$\ell=-\lambda$ throughout the following
$30\log(1/\lambda)$ interval. Then
\begin{equation}\label{eq:A11-delta-small}
 |\Delta_I|<64\lambda^{28}.
\end{equation}
\end{proposition}
\begin{proof}
For $r=I/(XH)$, the exact source transport law is
$r'+(1+\ell)r=1$. At the reference radius,
$r=5/8$, and during all pre-bump stages
$1+\ell\ge17/20$, so $0\le r\le20/17$.
Under $\ell=-\lambda$, the difference from the fixed
point $1/(1-\lambda)$ decreases at rate $1-\lambda$.
The first bump lies three units before the end of a plateau
of length $30\log(1/\lambda)$. Consequently
\[
 |\Delta_I|<3e^{3(1-\lambda)}\lambda^{30(1-\lambda)}
 \le3e^3\lambda^{57/2}<64\lambda^{28},
\]
since $e^3<21$ and $0<\lambda\le1/20$.
\end{proof}

\begin{theorem}[The two actual (A.11) corrections preserve negative slope]
\label{thm:A11-true-corrections}
For $0<\lambda\le1/20$ the source's two equations
\eqref{eq:A11-two-exact} admit a unique solution in the
small-correction ball with
\begin{equation}\label{eq:A11-c-coeff}
 \|c\|_\infty\le192\lambda^{28}.
\end{equation}
The edited azimuthal field is positive, and on both bump
supports the actual logarithmic slope satisfies
\begin{equation}\label{eq:A11-slope}
 -\frac{3\lambda}{2}\le\ell_c\le-\frac\lambda2.
\end{equation}
Off the supports, $\ell_c=-\lambda$.
\end{theorem}
\begin{proof}
Propositions~\ref{prop:A11-matrix-bound} and
\ref{prop:A11-discrepancy} give
$8\|B_{11}^{-1}\|^2\|Q\||\Delta_I|
 <368640\lambda^{28}<1$. The contraction lemma A.2
of the source therefore supplies a unique small solution
with $\|c\|_\infty\le2(3/2)|\Delta_I|
\le192\lambda^{28}$.
For completeness, the actual step $\sigma$ has derivative
bounds $\|\sigma'\|_\infty\le64$ and
$\|\sigma''\|_\infty\le17152$.
On $[1/4,3/4]$, write
$g=-t^{-2}+(1-t)^{-2}$ and
$\sigma=(1+e^{-g})^{-1}$; then
$|g'|\le256$, $|g''|\le3072$, giving the asserted
second-derivative bound. For $0<t\le1/4$, put $s=1/t\ge4$.
The tail derivatives satisfy
$|\sigma'|\le3s^3e^{-s^2+16/9}<1$ and
$|\sigma''|\le7s^6e^{-s^2+16/9}<1$;
symmetry gives the remaining quarter.
Thus $\|\beta\|_\infty\le640/3$ and
$\|\beta'\|_\infty<200000$.
The two supports are disjoint, so
$\|\sum_i c_i\beta_i\|_\infty
\le(640/3)192\lambda^{28}<1/2$.
Moreover
\[
 |\ell_c+\lambda|
 =\left|\frac{\sum_i c_i\beta_i'}{1+\sum_i c_i\beta_i}\right|
 <76800000\lambda^{28}<\frac\lambda2.
\]
This proves positivity and the slope interval.
\end{proof}

\paragraph{Explicit interpolation and terminal slope parameters.}
One may fix $T_f=512$ in (A.10): since
$0\le J_0(\eta)=\log(1+\eta^2)\le\log2<7/10$
and $\|\sigma'\|_\infty\le64$, its slope satisfies
\begin{equation}\label{eq:A11-Tf-fixed}
 -\lambda-\frac7{80}\le\ell_{\rm interp}\le-\lambda.
\end{equation}
Likewise the terminal cutoff coefficient in (A.12) may be
fixed as $c_o=1/256$, which yields
$0\le f_o'/f_o<h/4$ for $0<h\le1/100$.
For every $T_d>10$, the explicit parameter expressions
\begin{equation}\label{eq:A11-partial-tuple}
 T_f=512,\quad c_o=1/256,\quad
 \lambda=\frac{10^{-6}}{T_d+513},\quad
 h=\frac{\lambda e^{-T_d}}{10^6}
\end{equation}
satisfy all smallness assumptions used in this section and
in the preceding forward/backward pulse-energy inequalities.
This is a \emph{partial} quantitative parameter selection:
it does not certify an admissible choice of $M_d,P_*$,
all later stress-cone constants, or the unique stopping length
of the terminal $Q_s$ matching interval. No full profile,
Appendix-B contraction, unforced similarity corrections or
physical Navier--Stokes solution is thereby established.


\subsection*{Parameter derivatives and explicit terminal matching}
\label{subsec:A11-deriv-stop}
The previous source-native theorem controls the two relative
azimuthal bumps in (A.11). The dependence on the parameter $\eta$
is also tame, and the final constant-slope stopping length can be
determined without a numerical shooting argument.

\begin{proposition}[First-parameter-derivative bounds]
\label{prop:A11-eta-derivative}
Assume the exact Appendix-A stages and Theorem~\ref{thm:A11-true-corrections},
with $T_f=512$ and $0<\lambda\le1/20$. For the actual two correction
coefficients and the corrected azimuthal logarithmic slope,
\[
 \max_i\|\partial_\eta c_i\|_{L^\infty([-1,1])}
 <12\lambda^{28},\qquad
 \|\partial_\eta\ell_c\|_{L^\infty}<5\cdot10^6\lambda^{28}.
\]
The pressure-neutrality and angular-moment identities also
hold after differentiating in $\eta$, since the correction equations
hold identically and the bump integrals have compact support.
\end{proposition}
\begin{proof}
Write $r=I/(XH)$. It satisfies $r'+(1+\ell)r=1$ in logarithmic
radial coordinates. The reference value is $r=5/8$ and is
independent of $\eta$; all preceding unedited slopes are also
$\eta$-independent. During the interpolation (A.10),
\[
 \partial_\eta\ell
 =\frac{\sigma'(y/T_f)}{T_f}\frac{2\eta}{1+\eta^2},
 \quad |\partial_\eta\ell|\le64/T_f=1/8.
\]
Since $0\le r\le20/17$ and $1+\ell\ge17/20$, the
equation for $v=\partial_\eta r$ gives $|v|\le50/289$.
On the following $30\log(1/\lambda)-3$ unedited units,
$v$ decays at rate $1-\lambda$. Thus the moment discrepancy
$\Delta_I$ defined in \eqref{eq:A11-two-exact} obeys
\[
 |\partial_\eta\Delta_I|\le\frac{50}{289}
 e^{3(1-\lambda)}\lambda^{30(1-\lambda)}
 <\frac{1050}{289}\lambda^{28}<4\lambda^{28}.
\]
The exact nonlinear correction system is $B_{11}c+Q(c,c)
=(\Delta_I,0)^T$. By Theorem~\ref{thm:A11-matrix-bound}
and \eqref{eq:A11-c-coeff},
\[
 \|B_{11}^{-1}DQ(c)\|_\infty
 \le(3/2)(640)(192)\lambda^{28}<1/2.
\]
Differentiate the system and invert its Jacobian to obtain
$\|c_\eta\|_\infty\le3|\partial_\eta\Delta_I|
<12\lambda^{28}$.

Put $g=\sum c_i\beta_i$. On the two disjoint supports
$\ell_c=-\lambda+g_y/(1+g)$.
Using $|g|<1/2$, $\|\beta\|_\infty\le640/3$ and
$\|\beta'\|_\infty<200000$ yields
\[
 |\partial_\eta\ell_c|\le
 2(200000)(12)\lambda^{28}
 +4(200000)(192)(640/3)(12)\lambda^{56}
 <5\cdot10^6\lambda^{28}.
\]
\end{proof}

\begin{proposition}[Quantitative terminal stopping length]
\label{prop:A11-terminal-stop}
Assume the actual (A.11) correction and the subsequent
source-prescribed one-unit transition to $\ell=-1$, the hold
$\ell=-1$ for $4\log(1/h)$, and one-unit transition to
$\ell=-h$, with $h<\lambda\le1/20$ and $0<h\le1/100$.
For the terminal cutoff (A.12), fix $c_o=1/256$.
Let $Q_{\rm start}$ denote the value of the source $Q_s$
at the start of the last $\ell=-h$ hold, and $Q_p$ the
cutoff matching value (A.13). Then
\begin{equation}\label{eq:A11-terminal-Q-bounds}
 Q_{\rm start}\ge\frac4e(1-h)\log(1/h)>5,\qquad
 \frac{c_oh}{1-c_oh}\le Q_p
 \le\frac{e^{3(1-h)}c_oh}{1-c_oh}<\frac h{10}.
\end{equation}
In particular the unique positive holding length is
\begin{equation}\label{eq:A11-stop-exact}
 L_{\rm stop}=\frac1{1-h}\log\frac{Q_{\rm start}}{Q_p}>0.
\end{equation}
It is independent of $\eta$.
\end{proposition}
\begin{proof}
The restored angular moment implies
$Q_s=(\lambda-h)/(1-\lambda)>0$ at the exit of (A.11).
The source ODE is $Q_s'+(1+\ell)Q_s=-\ell-h$.
The first transition has nonnegative right-hand side;
on the $\ell=-1$ hold, $Q_s'=1-h$, so its exit value is
at least $4(1-h)\log(1/h)$. During the following unit transition
the right-hand side remains nonnegative and $0\le1+\ell\le1$,
giving the stated lower bound on $Q_{\rm start}$.
For the terminal cutoff $f_o=1-c_oh\psi_o$ of (A.12),
$f_o(0)=1-c_oh$, $f_o(3)=1$, and $f_o'\ge0$.
The integrating factor in (A.13) simplifies exactly to
\[
 Q_p=\frac1{f_o(0)}\int_0^3e^{(1-h)v}f_o'(v)\,dv.
\]
Integration of $f_o'$ gives the displayed interval for $Q_p$.
The numerical inequalities follow from $e^3<21$, $e<3$,
$\log100>4$ and $h\le1/100$. Finally on $\ell=-h$
the equation is $Q_s'=-(1-h)Q_s$; the unique positive length
at which $Q_s$ reaches $Q_p$ is \eqref{eq:A11-stop-exact}.
The exit moment and all later prescribed slopes are
$\eta$-independent, so the length is as well.
\end{proof}

These estimates quantify source constructions that were previously
given qualitatively. They do not certify all of the source's
stress-cone inequalities, the full moment-matched numerical exterior,
Appendix-B contraction or the unforced $n=1,2$ correction systems.


\subsection*{A stress-cone criterion for the two corrected azimuthal stages}

The Appendix-A condition is not merely a sign restriction on
$\ell=\partial_{\log X}\log(\sqrt X E)$. The source's equations
(4.11), (4.20)--(4.22) give
\[
a=2-2\ell,\quad b_s=2D_XU/E,\quad
p_1=XQ_s/L,\quad w=N_s/(EQ_s).
\]
For $a>0$ put $t_s=-b_s/a$,
$v_s=a+b_s^2/a$, $P_c=p_1(1-b_sw/a)$ and
$J_c=p_1(w+b_s/a)$. When $v_s>2$, the full cone reads
$P_c>v_s$ and $(v_s-2)J_c^2<2(P_c-v_s)^2$.

\begin{proposition}[Uniform corrected-support sufficient condition]
\label{prop:A11-cone-source}
On the two genuine (A.11) azimuthal-bump supports,
$U=M=J=0$ and hence $b_s=0$. Assume the corrected slope
$-3\lambda/2\le\ell\le-\lambda/2$, $0<\lambda\le1/20$.
If throughout those supports
\[
Q_s\ge\lambda/4,\quad |N_s|\le E\lambda/4,
\quad X\ge18/\lambda,
\]
then $a>2$, $p_1\ge2a$ and the normalized cone margin obeys
\[
2(1-a/p_1)^2-(a-2)w^2\ge7/20.
\]
\end{proposition}
\begin{proof}
The slope gives $2+\lambda\le a\le2+3\lambda\le43/20$.
As $L\le1$, $p_1=XQ_s/L\ge9/2>2a$.
Also $|w|\le1$. Therefore
$2(1-a/p_1)^2-(a-2)w^2\ge1/2-3/20=7/20$.
Multiply by $p_1^2>0$ and invoke (4.22).
\end{proof}

\begin{proposition}[Conditional source lower bound for $Q_s$]
\label{prop:A11-cone-Qs}
Retain the previously proved A.11 bounds for $c$, $c_\eta$,
$\ell_c+\lambda$, $\partial_\eta\ell_c$ and the initial
first-bump discrepancy:
\[
\|c\|\le192\lambda^{28},\quad
\|c_\eta\|<12\lambda^{28},\quad
|\ell_c+\lambda|<76800000\lambda^{28},\quad
|\ell_{c,\eta}|<5000000\lambda^{28},
\]
and at the first bump center
$|r-(1-\lambda)^{-1}|<64\lambda^{28}$,
$|r_\eta|<4\lambda^{28}$ for $r=I/(XH)$.
If $0<h\le\lambda/100$, then $Q_s\ge\lambda/2$ on both
supports and the intervening constant-slope interval.
\end{proposition}
\begin{proof}
Moving to the left endpoint of the first support enlarges
the initial discrepancies to $80\lambda^{28}$ and
$5\lambda^{28}$. The total interval is shorter than three
logarithmic units. Solve the damped equation
$r'+(1+\ell_c)r=1$, using $0\le r\le20/17$:
\[
|r-(1-\lambda)^{-1}|<3\cdot10^8\lambda^{28},
\qquad |r_\eta|<2\cdot10^7\lambda^{28}.
\]
On either correction support $E=E_{\rm base}(1+c_i\beta_i)$,
with $E_{\rm base}$ independent of $\eta$.
Since $\|\beta_i\|_\infty\le640/3$,
$|H_\eta/H|\le5120\lambda^{28}$, giving
$|I_\eta/(XH)|<2.1\cdot10^7\lambda^{28}$.
The exact identity (4.16) at $U=M=J=0$ is
\[
Q_s=-1+(1-h)r-D\eta I_\eta/(XH).
\]
Hence $Q_s\ge(\lambda-h)/(1-\lambda)-
3.2\cdot10^8\lambda^{28}\ge\lambda/2$, using
$3.2\cdot10^8(1/20)^{27}<49/100$.
\end{proof}

The remaining sufficient support conditions thus reduce to
$|N_s|\le E\lambda/2$ and $X\ge9/\lambda$.
The radial condition can be imposed by increasing $X_R$
after the other source parameters. The first condition is an
unverified source-specific interval requirement on the actual
corrected $N_s$.

From the source's (4.16), on these supports
\[
Q_s=-1+\frac{(1-h)I-D\eta I_\eta}{XH},\quad
N_s=\frac{4h\eta S-(1-\eta^2)S_\eta}{X}
 +4A\eta\Pi-(1-\eta^2)\Pi_\eta .
\]
Thus a genuine interval certificate needs the six cumulative
moment/pressure fields and their indicated $\eta$ derivatives.

\begin{lemma}[Quantitative stability of the A.24 margins]
If both the uncorrected and corrected triples $(a,b_s,w)$
satisfy $a\ge a_0>0$, $|a|\le A_0$, $|b_s|\le B_0$,
$|w|\le W_0$ and corresponding component differences
at most $\epsilon$, then for
$m_1=a-b_sw$ and
$m_2=2-2b_sw-b_s^2/a-(a-2)w^2$,
\[
|\Delta m_1|\le(1+B_0+W_0)\epsilon,
\]
\[
|\Delta m_2|\le
\left[2(B_0+W_0)+2B_0/a_0+B_0^2/a_0^2+
W_0^2+2(A_0+2)W_0\right]\epsilon.
\]
\end{lemma}
\begin{proof}
Expand the products, use
$|\Delta(b_s^2)|\le2B_0\epsilon$,
$|\Delta(a^{-1})|\le\epsilon/a_0^2$, and
$|\Delta(w^2)|\le2W_0\epsilon$.
\end{proof}

On the earlier joining stages the source requires only the
relaxed cone (4.21); the full condition $v_s>2$ applies on
the designated outer intervals. Strict pulse margins from
Appendix A.5 require separate interval bounds on the
source's $o(1)$ errors. The exact checker supplied with this
insert verifies scalar cone implications from externally
certified input intervals, not the actual exterior schedule.
The complete source stress cone and the proposed unforced
Theorem~1.1 remain unproved.


\subsection*{Exterior-tail certification of the A.11 axial cone component}
\label{subsec:A11-actual-Ns-tail}
On the two corrected supports of (A.11), the source has $U=M=J=0$.
After the genuine amplitude equation imposes $S(\infty,\eta)=0$, the source
moment equation (4.15) gives
\begin{equation}\label{eq:A11-tail-SP}
 S(X,\eta)=\frac12\int_X^\infty E(x,\eta)^2\,dx,\qquad
 \Pi(X,\eta)=-\frac12\int_X^\infty\frac{E(x,\eta)^2}{x}\,dx.
\end{equation}
The pressure formula uses the actual condition $\Pi(\infty,\eta)=0$
and the A.11 pressure-neutral moment correction. In particular
the $S$ identity is not valid before the $S(\infty)=0$ condition.

\begin{lemma}[Exact tail form of the axial stress source]
Define $D_s=S/(XE^2)$, $P_s=-\Pi/E^2$, $g=\partial_\eta\log E$,
$d=1-\eta^2$, $A=1/2+h$. Formula (4.16) becomes
\begin{equation}\label{eq:A11-Ns-tail-exact}
 \frac{N_s}{E}=E\left\{4\eta(hD_s-AP_s)
 -d[2g(D_s-P_s)+\partial_\eta D_s-\partial_\eta P_s]\right\}.
\end{equation}
If $\ell=\partial_{\log X}\log(\sqrt X E)\le0$ at all later radii
and $|\partial_\eta\log E|\le\varepsilon$ there, then
\begin{equation}\label{eq:A11-Ns-tail-inequality}
0\le P_s\le\frac12,\qquad
\frac{|N_s|}{E}\le E[4hD_s+2A+2\varepsilon(D_s+1/2)].
\end{equation}
\end{lemma}
\begin{proof}
For $t=\log(x/X)$ write $R(t)=\exp(2\int_0^t\ell(Xe^v)\,dv)$.
Then $D_s=\frac12\int_0^\infty R\,dt$ and
$P_s=\frac12\int_0^\infty e^{-t}R\,dt$. Since $R\le1$, $P_s\le1/2$.
The original tail integrals give
$|S_\eta|/X\le2\varepsilon E^2D_s$ and
$|\Pi_\eta|\le2\varepsilon E^2P_s$. Apply (4.16).
The exact displayed expression follows by differentiation.
\end{proof}

\begin{theorem}[Source-stage sufficient bound for both A.11 supports]
\label{thm:A11-Ns-realized}
Retain the corrected two-bump bounds, $T_f=512$, $T_d>10$,
$P_*>e^{T_d}$ and the post-pulse nonpositive slope and terminal
release bounds used above. Choose
\[
0<\lambda\le\min\left\{
\frac{10^{-6}}{T_d+513},\frac1{10(1+\log P_*)},
\frac1{20}\right\},\qquad0<h\le\lambda/100.
\]
Once the actual source moment matching has imposed $S(\infty)=0$,
the resulting corrected supports satisfy
\begin{equation}\label{eq:A11-Ns-final}
\sup_{\mathrm{A.11\ supports},\,|\eta|\le1}|N_s|/E<\lambda/2.
\end{equation}
Together with $Q_s\ge\lambda/2$, the corrected slope range
$-3\lambda/2\le\ell\le-\lambda/2$, and a choice
$X\ge9/\lambda$, the full strict cone (4.22) follows on these
two supports, with the previous normalized margin $7/20$.
\end{theorem}
\begin{proof}
Past the angular interpolation the unedited source $E$ is
$\eta$-independent. Only the two corrected bumps retain
$\eta$ dependence. The bounds
$\|c_\eta\|_\infty<12\lambda^{28}$,
$\|\beta\|_\infty\le640/3$ and
$|\sum c_i\beta_i|<1/2$ give
$|\partial_\eta\log E|<5120\lambda^{28}=:\varepsilon$
at every radius beyond the first support.
The prescribed later nonpositive slope stages imply
\[
D_s\le B_D:=T_f/2+15\log(1/\lambda)+3+h^7/3
\]
on either support: allow half the remaining interpolation/plateau
length, the two release transitions, the $-1$ hold and its
exponentially suppressed terminal tail.

The first unit transition enlarges $E$ by at most $e^{1/10}$.
The angular flattening has integrated positive contribution at
most $\log2<1$, and the relative correction is less than $3/2$.
The corrected supports lie after a length-$13/\lambda$ pulse
and within a three-unit endpoint reserve. Dropping all other
decay gives
\[
 E\le P_*\exp(4-13/(2\lambda)).
\]
Put $t=1/\lambda\ge20$; since
$t\ge10(1+\log P_*)$ and $\log(4t)\le t$,
\[
\log(4tP_*e^{4-13t/2})\le3-27t/5<0,
\]
hence $E<\lambda/4$.

Also $\lambda T_f\le10^{-6}$, $h\le\lambda/100$,
$\lambda\log(1/\lambda)\le1$ and $\lambda\le1/20$ give
$4hB_D<61/100$, $2A\le1001/1000$, and
$2\varepsilon(B_D+1/2)<1/100$ (use
$\lambda^{27}\le20^{-27}$). Thus the square bracket in
\eqref{eq:A11-Ns-tail-inequality} is strictly less than $2$,
so $|N_s|/E<\lambda/2$. The previously proved
$Q_s$ and slope bounds plus $X\ge9/\lambda$ imply the
source cone inequality at these supports.
\end{proof}

This is an all-parameter \emph{conditional source-stage theorem},
not an enclosure of one independently frozen complete exterior
schedule. Compatibility with every other source-stage stress-cone
margin, pulse $o(1)$ constants, terminal matching, Appendix B
and the unforced coefficient hierarchy still requires separate
verification. The proposed unforced Theorem~1.1 remains open.


\subsection*{Quantitative pulse curvature and the remaining cone error}
The prescribed axial pulse of Appendix A.5 in \cite{openai2026} has
$U=ER_b$, $R_b(y,\eta)=A(\eta)R_0(\lambda y)
+c_1(\eta)\beta_1(y)+c_2(\eta)\beta_2(y)$,
$\beta=1/2-\lambda$ and $\ell=-\lambda$. Here the
$\beta_i$ are the \emph{axial} moment bumps of (A.9), not the
later azimuthal corrections of (A.11).
Put $\phi_b(\xi)=\int_0^\xi\sigma(50v)\,dv$ and
$R_0(\xi)=\phi_b(\xi)[1-\sigma(\xi-10)]$, with the
source's smooth endpoint extensions.

\begin{lemma}[Explicit source convolution remainder]
\label{lem:A5-source-convolution}
The source's fixed smooth step and main pulse satisfy
$\|\sigma'\|_\infty\le64$,
$\|\sigma''\|_\infty\le17152$,
$\|R_0'\|_\infty\le705$ and
$\|R_0''\|_\infty\le192000$.
If $m'+\beta m=A R_0(\lambda y)+\sum_i c_i\beta_i(y)$, then
\begin{equation}\label{eq:A5-explicit-convolution}
\left|m(y)-A\left(\frac{R_0(\lambda y)}{\beta}
-\frac{\lambda R'_0(\lambda y)}{\beta^2}\right)\right|
\le \frac{192000|A|\lambda^2}{\beta^3}
+|m(0)|e^{-\beta y}+\frac{1280}{3\beta}\max_i|c_i|.
\end{equation}
For $|A|\le1011/1000$, the first term is at most
$2{,}131{,}000\lambda^2$. The source's preceding
zero-axial-velocity interval of length $60\log(1/\lambda)$
and the earlier normalized moment bound $|m|<9/100$ imply
$|m(0)|\le(9/100)\lambda^{27}$.
\end{lemma}
\begin{proof}
For $1/4\le t\le3/4$, logarithmic differentiation of
$\sigma$ gives $|\sigma'|\le64$ and
$|\sigma''|\le(256^2+3072)/4=17152$.
For $0<t\le1/4$, the exponentially decaying boundary
expression has the upper bound
$19244e^{-14}<1$; use symmetry at the other endpoint.
The product rule, $|\phi_b'|\le1$,
$|\phi_b''|\le3200$ and $\phi_b\le11$ give
$\|R_0''\|\le3200+128+11(17152)=192000$.
Extend the flat-start pulse by zero and apply Taylor's
formula inside its causal exponential convolution.
The second-order remainder is bounded by
$|A|\lambda^2\|R_0''\|/\beta^3$, since
$\int_0^\infty v^2e^{-\beta v}dv=2/\beta^3$.
Each prescribed bump has supremum at most $640/3$;
the two bump convolutions and the entrance homogeneous
term give the other terms. Use $\beta\ge9/20$ and
exact rational substitution for the final constants.
\end{proof}

\begin{proposition}[Source pulse stress-cone acceptance criterion]
\label{prop:A5-quantified-source-cone}
Assume the genuine A.19 amplitude satisfies
$1.009<A<1.011$ and
\[
0<\lambda\le(2\cdot10^9)^{-1},\quad
0<h\le\lambda/10^4,\quad
\max_i|c_i|\le\lambda/400000,
\]
where $c_i$ are the \emph{A.9 axial} end-bump coefficients.
Writing $R=AR_0$, $c_\eta\ge0$ as in (A.27), and
\[
 C_d=\frac{\lambda(D+c_\eta)(1-\lambda)}
 {\beta^2(\lambda-h+c_\eta)},\qquad
 b_0=-R,\quad w_0=2R-C_dR'_0A,
\]
the source shear is exactly
\[
 b_s=-(1+2\lambda)R_b+2\lambda\,\partial_\xi R_b,\qquad
 |b_s-b_0|<10^{-6}.
\]
If, moreover, the actual (A.29)--(A.30) moment and pressure
calculation yields
$Q_s>0$ and $|w-w_0|\le10^{-6}$ throughout the pulse,
where $w=N_s/(EQ_s)$, then
\begin{equation}\label{eq:A5-quantified-cone}
 b_sw<74/100,\qquad
 2b_sw+b_s^2/a+(a-2)w^2<168/100<2,\quad a=2+2\lambda.
\end{equation}
The full cone follows after the independent radial-scale
criterion of Lemma 4.5 is met.
\end{proposition}
\begin{proof}
The original bump has $\|\beta_i'\|_\infty<200000$.
Thus $|R_b-R|\le(1280/3)\max_i|c_i|$ and
$|\partial_\xi(R_b-R)|\le400000\max_i|c_i|/\lambda$.
Since $|R|<12$ and $|R'|<713$,
$|b_s+R|<1453\lambda<10^{-6}$.
For $h\le\lambda/10^4$ and $c_\eta\ge0$,
$0<C_d\le[(1/2)(10000/9999)]/(1/2-\lambda)^2
<2.001$. Hence $|w_0|<1500$ and $R'\le1.2$.
Completing the squares gives
\[
 \sup_{R\ge0}(-2R^2+2.41R)=58081/80000,\qquad
 \sup_{R\ge0}(-\tfrac72R^2+4.82R)=58081/35000.
\]
If the changes in $b_s$ and $w$ are at most
$\varepsilon=10^{-6}$, the changes in the two cone
polynomials are bounded by
$1512\varepsilon+\varepsilon^2$ and
$3336\varepsilon+4\varepsilon^2$.
The second has in addition the base nonnegative term
$2\lambda(1500)^2\le0.00225$.
The exact rational bounds are strictly less than
$0.74$ and $1.68$ respectively, proving the result.
\end{proof}

The explicit convolution control closes one of the $O(\lambda^2)$
channels in (A.28). The remaining $10^{-6}$ estimate for
$w-w_0$ involves $Q_s,N_s$, the $\eta$-derivatives
of the axial moments, the pressure and the $S$-moment
in (A.29)--(A.30). Its source-dependent constants have not
yet been interval-certified on a single actual parameter
tuple. Neither the other stage cones, terminal matching,
Appendix-B contraction nor an unforced PDE theorem follows.


\subsection*{Exact source residual behind the pulse ratio (A.29)--(A.30)}
\label{subsec:pulse-source-actual-ratio}
The originating manuscript \cite{openai2026} asserts that
$w=N_s/(E Q_s)=2R-C_dR'+o(1)$ on its axial pulse.
Its equations (4.15)--(4.16) give an exact, non-asymptotic
certificate that isolates the errors to be enclosed.
Let $D=1/2-h$, $A=1/2+h$, $\beta=1/2-\lambda$,
$d=1-\eta^2$, $R=A_{\rm pulse}R_0(\lambda y)$ and
$B_{\rm ax}=U/E-R$. On the pulse,
$E=e_b(1+\eta^2)^{-1}e^{-(1/2+\lambda)y}$, so
\[
g=\partial_\eta\log E=-2\eta/(1+\eta^2),\qquad
c_\eta=-D\eta g=2D\eta^2/(1+\eta^2).
\]
Set
\[
q_0=\frac{\lambda-h+c_\eta}{1-\lambda},\quad
C_d=\frac{\lambda(D+c_\eta)(1-\lambda)}
 {\beta^2(\lambda-h+c_\eta)},\quad
w_0=2R-C_dR',\quad m=M/(XE).
\]
Define exact residuals
\[
e_m=m-R/\beta+\lambda R'/\beta^2,\quad
e_Q=Q_s-q_0,\quad
e_N=N_s/E+R-(D+c_\eta)m+D\eta m_\eta.
\]
\begin{proposition}[Exact pulse stress-ratio error identity]
\label{prop:A29A30-exact-residual}
Where $Q_s>0$,
\begin{equation}\label{eq:A29A30-exact-ratio}
\boxed{\displaystyle
w-w_0=\frac{q_0(\lambda/\beta)R+(D+c_\eta)e_m
      -D\eta m_\eta+e_N-w_0e_Q}{q_0+e_Q}.}
\end{equation}
In the source's radial-infinity pressure gauge,
\begin{align}
e_Q={}&-W+
 \frac{(1-h)I-D\eta I_\eta-dJ_\eta
       +2(h-D)\eta J}{XH}-q_0,\label{eq:A29A30-eq}\\
e_N={}&(1-W)R-WB_{\rm ax}
+\frac{4h\eta S-dS_\eta}{XE}
+\frac{4A\eta\Pi-d\Pi_\eta}{E},
\label{eq:A29A30-en}
\end{align}
where $H=\sqrt{2X}E$ and
$W=1-(2D\eta M+dM_\eta)/X$.
\end{proposition}
\begin{proof}
The two residual formulas are immediate from the exact
source identities (4.16) after differentiating $M=XEm$:
$M_\eta=XE(m_\eta+gm)$. The identity
$D+c_\eta-\beta=\lambda-h+c_\eta$ gives
\[
-R+(D+c_\eta)(R/\beta-\lambda R'/\beta^2)
-q_0(2R-C_dR')=q_0(\lambda/\beta)R.
\]
Now write $N_s/E$ using the residual definition, divide
by $Q_s=q_0+e_Q$ and subtract $w_0$.
\end{proof}

\begin{corollary}[Rational sufficient $10^{-6}$ acceptance budget]
\label{cor:A29A30-budget}
Suppose $0<\lambda\le10^{-14}$,
$0<h\le\lambda/10^4$, $|\eta|\le1$,
$|R|\le12$, $|w_0|\le1500$. If on the complete
pulse, including both axial end-bump supports,
\[
|e_Q|/q_0\le10^{-10},\quad
|e_m|/q_0\le10^{-7},\quad
|\eta m_\eta|/q_0\le2\cdot10^{-7},\quad
|e_N|/q_0\le10^{-7},
\]
then $Q_s>0$ and $\sup|w-w_0|<10^{-6}$.
\end{corollary}
\begin{proof}
Since $c_\eta\ge0$,
$q_0\ge(9999/10000)\lambda$,
$0<D<1/2$ and $0<D+c_\eta\le1$.
Formula \eqref{eq:A29A30-exact-ratio} bounds the error
by
\[
\frac{12\lambda/\beta+3\cdot10^{-7}
       +1500\cdot10^{-10}}{1-10^{-10}}
<4.51\cdot10^{-7}<10^{-6}.
\]
\end{proof}

The previous bound $|c_i|\le\lambda/400000$ on the
axial correction bumps suffices for the $b_s$ shear,
but the existing source (A.28) bound permits its
contribution to $|e_m|$ to reach
$1280\max|c_i|/(3\beta)\le(8/3375)\lambda$.
Thus the old assumption does not certify the new
relative $10^{-7}$ error gate.

\begin{lemma}[Conditional stronger axial coefficient estimate]
\label{lem:A29A30-actual-bump}
Assume the actual affine two-moment formula of
Appendix A, $1.009\le A_{\rm pulse}\le1.011$,
the source pre-pulse moment envelope $\max(|m_M|,|m_J|)
\le9/100$ and $|m(0)|\le(9/100)\lambda^{27}$.
For $0<\lambda\le10^{-14}$,
\[
 \max_i|c_i|<10^{-40}\lambda,\qquad
 \sup_{\rm pulse}|e_m|/q_0<3\cdot10^{-8}.
\]
\end{lemma}
\begin{proof}
The source discrepancy formulas and the normalized inverse
estimate $2/\lambda$ give
\[
D_0\le(9/100)e^{-(2/5)(13/\lambda-3)},\qquad
D_1\le(55/2)e^{-(2/5)(2/\lambda-3)}.
\]
Writing $t=1/\lambda$, the bound $e^{6/5}<4$ yields
$\max|c_i|<240t e^{-4t/5}$.
The sixth term of the exponential series implies
\[
 \max|c_i|/\lambda
 \le(10546875/16)t^{-4}<10^{-40}
 \quad(t\ge10^{14}).
\]
The previously established complete (A.28) residual bound is
$|e_m|\le2{,}131{,}000\lambda^2+
(9/100)\lambda^{27}+(1280/(3\beta))\max|c_i|$.
Divide by $q_0\ge(9999/10000)\lambda$ and apply
$\beta\ge9/20$.
\end{proof}

The remaining source verification is therefore concentrated in
$e_Q$, $e_N$, and $\eta m_\eta$ and their \emph{simultaneous}
interval enclosures. The accompanying exact symbolic checker,
rational interval engine and finite Lean scalar gates do not
supply those source enclosures or a complete theorem proof.
The full stress cone on every radial stage, source-admissible
parameter freeze, Appendix B contraction, unforced correction
hierarchy and physical unforced Navier--Stokes reconstruction
remain open.


\section{Normalized source moments and the remaining pulse certificate}
The exact source integrals (4.15)--(4.16) of \cite{openai2026} can be
transported without the very large radial factors appearing in the
unnormalized variables. Write $y=\log(X/X_R)$, $H=\sqrt{2X}E$,
$a=U/E$, $\ell=(\log H)_y$, $g=(\log E)_\eta$, and
\[
 m=\frac{M}{XE},\quad i=\frac{I}{XH},\quad
 j=\frac{J}{XHE},\quad s=\frac{S}{XE^2},\quad
 p=\frac{\Pi}{E^2}.
\]
\begin{proposition}[Ten normalized source transport equations]
\label{prop:norm-source-ten}
On every source stage with $E>0$, the five exact moment equations are
\[
\begin{aligned}
m_y&=a-(\ell+\tfrac12)m,&i_y&=1-(\ell+1)i,\\
j_y&=a-(2\ell+\tfrac12)j,&
s_y&=a^2-\tfrac12-2\ell s,&
p_y&=\tfrac12+(1-2\ell)p.
\end{aligned}
\]
The five parameter-sensitivity equations are
\[
\begin{aligned}
(m_\eta)_y&=a_\eta-(\ell+\tfrac12)m_\eta-\ell_\eta m,\\
(i_\eta)_y&=-(\ell+1)i_\eta-\ell_\eta i,\\
(j_\eta)_y&=a_\eta-(2\ell+\tfrac12)j_\eta-2\ell_\eta j,\\
(s_\eta)_y&=2aa_\eta-2\ell s_\eta-2\ell_\eta s,\\
(p_\eta)_y&=(1-2\ell)p_\eta-2\ell_\eta p.
\end{aligned}
\]
Also $g_y=\ell_\eta$. At the reference radius $X_R$, with
$P=P_*$, the four forward data are
\[
(m,i,j,s)=\left(\frac{4\eta(1+\eta^2)}P,\frac58,
\frac{5\eta(1+\eta^2)}{2P},
\frac{16\eta^2(1+\eta^2)^2}{P^2}-\frac5{12}\right),
\]
with parameter jets
\[
(m_\eta,i_\eta,j_\eta,s_\eta)=
\left(\frac{4(1+3\eta^2)}P,0,
\frac{5(1+3\eta^2)}{2P},
\frac{32\eta(1+\eta^2)(1+3\eta^2)}{P^2}\right).
\]
The pressure uses the source's radial-infinity gauge, not a free
axis constant. At the beginning $X_T$ of its final power-law tail,
\[
 p(X_T)=-\frac{1}{2(1+2h)},\qquad p_\eta(X_T)=0.
\]
Consequently $|\eta m_\eta(X_R,\eta)|\le16/P_*<16e^{-T_d}$.
\end{proposition}
\begin{proof}
The source definitions give
$M_y=XU$, $I_y=XH$, $J_y=XUH$,
$S_y=X(U^2-E^2/2)$, and $\Pi_y=E^2/2$.
Their four denominators have logarithmic derivatives
$\ell+1/2,\ell+1,2\ell+1/2,2\ell$; $E^2$ has logarithmic
derivative $2\ell-1$. Apply the quotient and parameter
product rules. Reference data follow from exact integration of
the source power laws. On the terminal tail
$E=E_T(X/X_T)^{-1/2-h}$, and radial-infinity normalization
gives $\Pi=-E^2/[2(1+2h)]$.
\end{proof}

\begin{proposition}[Scale-free pulse numerator and denominator]
\label{prop:normalized-source-stress}
Put $D=1/2-h$, $A=1/2+h$, $d=1-\eta^2$ and
$W=1-E[2D\eta m+d(m_\eta+gm)]$. Then (4.16) becomes
\[
 Q_s=-W+(1-h)i-D\eta(i_\eta+gi)
 -dE(j_\eta+2gj)+2(h-D)\eta Ej
\]
and
\[
 \frac{N_s}{E}=-Wa+D[m-\eta(m_\eta+gm)]
 +E[4h\eta s-d(s_\eta+2gs)+4A\eta p-d(p_\eta+2gp)].
\]
With $B_{\rm ax}=a-R$, $q_0=(\lambda-h-D\eta g)/(1-\lambda)$,
the exact unclosed residuals are
\[
\begin{aligned}
e_Q={}&-W+(1-h)i-D\eta(i_\eta+gi)
 -dE(j_\eta+2gj)+2(h-D)\eta Ej-q_0,\\
e_N={}&(1-W)R-WB_{\rm ax}
 +E[4h\eta s-d(s_\eta+2gs)+4A\eta p-d(p_\eta+2gp)].
\end{aligned}
\]
The third residual is $\eta m_\eta$ from the sensitivity
system above. These are exact identities, not asymptotic expansions.
\end{proposition}
\begin{proof}
Substitute $M=XEm$, $I=XHi$, $J=XHEj$, $S=XE^2s$ and
$\Pi=E^2p$, with their first parameter derivatives, into (4.16).
For $e_N$, the terms involving $D(m-\eta(m_\eta+gm))$
cancel against $(D-D\eta g)m-D\eta m_\eta$.
\end{proof}

\paragraph{Numerical and formal verification boundary.}
For $\lambda=10^{-14}$ near $\eta=0$, the target
$|e_Q|/q_0\le10^{-10}$ corresponds to an absolute
accuracy of roughly $10^{-24}$, while
$|e_N|/q_0\le10^{-7}$ requires roughly $10^{-21}$.
A credible certificate needs validated arbitrary-precision
source-stage enclosures, including $\ell_\eta,a_\eta$ and the
actual A.9 axial corrections, on a finite exhaustive pulse
cover. The attached exact-rational interval checker accepts only
externally supplied cell bounds and does not certify their
provenance or full coverage. Its default positive and negative
controls are synthetic. Full source realization, strict
stress-cone certification, Appendix-B contraction, unforced
startup solutions and Theorem~1.1 remain open.


\section{Certified transport tubes for the normalized source moments}
\label{sec:source-moment-interval-tubes}
The ten-variable source system \eqref{eq:source-five-transport}--\eqref{eq:source-five-jet-transport} has a triangular split:
\[
z=(m,i,j,s,m_\eta,i_\eta,j_\eta,s_\eta)^\top,\qquad
w=(p,p_\eta)^\top.
\]
Writing $c=(\ell+\tfrac12,\ell+1,2\ell+\tfrac12,2\ell)$,
the eight forward equations are the affine system $z_y=A_fz+b_f$,
\[
A_f=
\begin{pmatrix}
-c_1&0&0&0&0&0&0&0\\
0&-c_2&0&0&0&0&0&0\\
0&0&-c_3&0&0&0&0&0\\
0&0&0&-c_4&0&0&0&0\\
-\ell_\eta&0&0&0&-c_1&0&0&0\\
0&-\ell_\eta&0&0&0&-c_2&0&0\\
0&0&-2\ell_\eta&0&0&0&-c_3&0\\
0&0&0&-2\ell_\eta&0&0&0&-c_4
\end{pmatrix},\quad
b_f=\begin{pmatrix}a\\1\\a\\a^2-\tfrac12\\a_\eta\\0\\a_\eta\\2aa_\eta\end{pmatrix}.
\]
Only the first eight states have reference-radius initial values.
The pressure pair is instead propagated backward from
$y_T$ with $\tau=y_T-y$:
\begin{equation}\label{eq:source-backward-pressure}
w_\tau=
\begin{pmatrix}2\ell-1&0\\2\ell_\eta&2\ell-1\end{pmatrix}w+
\begin{pmatrix}-\tfrac12\\0\end{pmatrix},\quad
w(0)=\begin{pmatrix}-[2(1+2h)]^{-1}\\0\end{pmatrix}.
\end{equation}
\begin{theorem}[Exact-rational conditional affine interval tubes]
\label{thm:source-exact-rational-tube}
Let $z'=A(t)z+b(t)$ on a source cell of length $\Delta>0$,
with independently certified rational interval enclosures
$\mathbf A,\mathbf b$ for every point of that cell, and initial
box $Z_0$. Put $L=\max_i\sum_j\operatorname{mag}\mathbf A_{ij}$,
$B=\max_i\operatorname{mag}\mathbf b_i$,
$M=\max_i\operatorname{mag}(Z_0)_i$.
If $\Delta L<1$, then every solution remains in
$[-S,S]^d$, $S=(M+\Delta B)/(1-\Delta L)$.
The complete trajectory interval may be tightened finitely
by the exact-rational inclusions
\[
Z^{k+1}=Z^k\cap
\operatorname{hull}\bigl(Z_0,Z_0+[0,\Delta](\mathbf A Z^k+\mathbf b)\bigr),
\quad Z^0=[-S,S]^d.
\]
The endpoint is enclosed by
$Z_0+\Delta(\mathbf A Z^k+\mathbf b)$.
\end{theorem}
\begin{proof}
With $V=\sup\|z(t)\|_\infty$, the integral equation gives
$V\le M+\Delta(LV+B)$, hence $V\le S$.
If $Z^k$ encloses every solution value, the integral equation
places every value within the displayed Picard interval; intersection
preserves inclusion. At the endpoint the integration time equals
$\Delta$. All operations are exact rational bounds.
\end{proof}
At $\eta=0$ on the source reference power law,
$a=\ell_\eta=0$, $\ell=3/5$,
$a_\eta=4/E$, and the known terminal reference data are
$(0,5/8,0,-5/12,4/P_*,0,5/(2P_*),0)$.
The pure exterior tail supplies the exact pressure equilibrium
$(-[2(1+2h)]^{-1},0)$.
The accompanying exact-rational verifier replays both checks and
rejects unsafe step sizes.
The certificate is conditional on independently verified
$a,a_\eta,\ell,\ell_\eta$ source bounds, and its short-step version
alone is not a practical discretization of a pulse of length
$13/\lambda$ at extremely small $\lambda$. No complete source
profile or stress-ratio certificate is asserted.

\subsection*{Exact pulse transfers on a fixed shape domain}
Appendix A.9 prescribes on its entire axial pulse, including the
two \emph{axial} moment bumps,
\[
E(y,\eta)=e_b(1+\eta^2)^{-1}e^{-(1/2+\lambda)y},
\quad 0\le y\le13/\lambda.
\]
Consequently $\ell=-\lambda$, $\ell_\eta=0$ and
$g=-2\eta/(1+\eta^2)$ exactly. The axial source is
$a=A(\eta)R_0(\lambda y)+c_1(\eta)\beta_1(y)+c_2(\eta)\beta_2(y)$,
with $a_\eta$ obtained by differentiating only these coefficients
at fixed $y$. The bumps in this formula are not the later,
distinct azimuthal A.11 correction bumps.
\begin{proposition}[Exact source-pulse variation of constants]
\label{prop:source-pulse-exact-transfer}
Let $\alpha=1/2-\lambda$, $\beta=1/2-2\lambda$,
$\gamma=1-\lambda$, $\kappa=1+2\lambda$,
and $y_1=13/\lambda$.
For $q\in\{m,j\}$, with corresponding decay exponent
$d_q\in\{\alpha,\beta\}$,
\[
q(y)=e^{-d_q y}q(0)+\int_0^ye^{-d_q(y-v)}a(v)\,dv,\quad
q_\eta(y)=e^{-d_q y}q_\eta(0)+
\int_0^ye^{-d_q(y-v)}a_\eta(v)\,dv.
\]
Moreover
\[
i(y)=\gamma^{-1}+(i(0)-\gamma^{-1})e^{-\gamma y},
\quad i_\eta(y)=i_\eta(0)e^{-\gamma y},
\]
\[
p(y)=-(2\kappa)^{-1}+
(p(y_1)+(2\kappa)^{-1})e^{-\kappa(y_1-y)},
\quad p_\eta(y)=p_\eta(y_1)e^{-\kappa(y_1-y)}.
\]
The remaining energy states satisfy
\[
s(y)=e^{2\lambda y}s(0)+\int_0^y e^{2\lambda(y-v)}
(a(v)^2-\tfrac12)\,dv,\quad
s_\eta(y)=e^{2\lambda y}s_\eta(0)+
2\int_0^y e^{2\lambda(y-v)}a(v)a_\eta(v)\,dv.
\]
\end{proposition}
\begin{proof}
The pulse has constant $\ell=-\lambda$ and
$\ell_\eta=0$, so each transport equation is solved
exactly by its integrating factor; the pressure is
integrated backward from the post-pulse source value.
\end{proof}
In the fixed coordinate $\xi=\lambda y\in[0,13]$ the
$s,s_\eta$ integrals are over a compact interval with
kernel $e^{2(\xi-v)}$. The $m,j$ kernels are sharply
localized on $\xi$-scale $\lambda$, for which (A.28)'s
controlled convolution method is appropriate.
These exact transfers avoid a prohibitive $1/\lambda$
number of generic ODE cells; certified source amplitudes,
correction coefficients, full pulse coverage and the
$10^{-6}$ stress-ratio proof are still outstanding.

\subsection*{Exact pulse-end axial-moment cancellation}
The axial (A.9) corrections are distinct from the subsequent azimuthal
(A.11) corrections. Their two moment equations admit an exact
terminal and first-parameter-jet simplification.

\begin{theorem}[Pulse-end moments and their parameter jets]
\label{thm:pulse-end-annihilation}
Let $Y=13/\lambda$, $0<\lambda<1/20$, and
$\alpha=\tfrac12-\lambda$, $\beta=\tfrac12-2\lambda$.
On the pulse retain the actual source input
\[
 a(y,\eta)=A_{\rm pulse}(\eta)R_0(\lambda y)
       +c_1(\eta)\beta_1(y)+c_2(\eta)\beta_2(y),
 \qquad \beta_2(y)=\beta_1(y-2).
\]
Assume both nonnegative nontrivial axial bump supports are contained
in $(0,Y)$. Set
\[
 k_\kappa=\int_0^Y e^{-\kappa(Y-v)}\beta_1(v)\,dv>0,\quad
 F_\kappa=\int_0^Y e^{-\kappa(Y-v)}R_0(\lambda v)\,dv,
 \quad \kappa\in\{\alpha,\beta\}.
\]
Then the exact normalized axial moments at pulse exit obey
\begin{equation}\label{eq:pulse-end-affine}
 \binom{m(Y)}{j(Y)}
 =\binom{e^{-\alpha Y}m(0)+A_{\rm pulse}F_\alpha}
             {e^{-\beta Y}j(0)+A_{\rm pulse}F_\beta}
 +\mathcal B\binom{c_1}{c_2},\quad
 \mathcal B=\begin{pmatrix}
 k_\alpha&e^{2\alpha}k_\alpha\\
 k_\beta&e^{2\beta}k_\beta
 \end{pmatrix}.
\end{equation}
The determinant is strictly negative:
\[
 \det\mathcal B=-k_\alpha k_\beta e^{2\beta}
                    (e^{2\lambda}-1)<0.
\]
Solving the two actual moment conditions gives
$m(Y)=j(Y)=m_\eta(Y)=j_\eta(Y)=0$ identically.
These four normalized states remain zero on every subsequent
stage where the source has $U=0$, including the two A.11
azimuthal correction supports.
\end{theorem}
\begin{proof}
The pulse transport equations are $m_y+\alpha m=a$ and
$j_y+\beta j=a$. Apply their exact integrating factors at
$y=Y$. Translating the second bump by two logarithmic units
multiplies its kernel by $e^{2\kappa}$, giving the displayed
matrix and determinant. The latter never vanishes for
$\lambda>0$; hence $c=-\mathcal B^{-1}r$, with $r$ the first
vector in \eqref{eq:pulse-end-affine}, is uniquely determined.
Since $\mathcal B$ is independent of $\eta$, differentiation
gives $c_\eta=-\mathcal B^{-1}r_\eta$. Therefore the two
terminal moments and both first parameter jets vanish
exactly. On later stages $U=0$ implies $M_y=J_y=0$ and
their parameter derivatives vanish; dividing by positive
normalizing denominators preserves zero.
\end{proof}

\begin{corollary}[Six-state post-pulse certification]
\label{cor:postpulse-six-state}
After the pulse, four of the eight forward moment states,
$m,j,m_\eta,j_\eta$, vanish identically. Only
$(i,s,i_\eta,s_\eta)$ require forward interval propagation,
whereas $(p,p_\eta)$ still use the independent infinity-gauge
backward boundary condition. The exact source stresses reduce to
\begin{align}
 Q_s&=-1+(1-h)i-D\eta(i_\eta+gi),\label{eq:pulse-end-Q-reduced}\\
 N_s/E&=E[4h\eta s-d(s_\eta+2gs)+
       4(\tfrac12+h)\eta p-d(p_\eta+2gp)].
 \label{eq:pulse-end-N-reduced}
\end{align}
It is essential to differentiate the actual coefficients $c_i$:
dropping $c_{i,\eta}$ invalidates the jet cancellation.
\end{corollary}
\begin{proof}
Insert the exact terminal zeros and $a=U/E=0$
in \eqref{eq:scale-free-Q}--\eqref{eq:scale-free-N}.
\end{proof}

\begin{proposition}[A conditional positive pulse-exit $Q_s$ margin]
\label{prop:pulse-end-Q-positive}
Suppose $0<\lambda\le1/20$, $0<h\le\lambda/100$ and the
true source pulse entrance has
$|i(0,\eta)|,|i_\eta(0,\eta)|\le100$ on $[-1,1]$.
At $Y=13/\lambda$,
\[
 |Q_s(Y,\eta)-q_0(\eta)|
 \le202e^{-(1-\lambda)Y}
 \le\frac{775680000}{3722098081}\lambda^4
 <\lambda/1000,\quad
 q_0=\frac{\lambda-h-D\eta g}{1-\lambda}.
\]
In particular $Q_s(Y,\eta)>\lambda/2$.
\end{proposition}
\begin{proof}
On the pulse $\ell=-\lambda$, $\ell_\eta=0$,
$i_y+(1-\lambda)i=1$ and
$(i_\eta)_y+(1-\lambda)i_\eta=0$. Consequently
\[
 Q_s(Y)-q_0=e^{-(1-\lambda)Y}
 [(1-h-D\eta g)(i(0)-(1-\lambda)^{-1})
          -D\eta i_\eta(0)].
\]
Since $-D\eta g=2D\eta^2/(1+\eta^2)\in[0,1/2]$,
the bracket has absolute value less than
$\frac32(100+20/19)+50<202$. As
$(1-\lambda)Y\ge247/(20\lambda)$, use
$e^x\ge x^4/24$ and the exact rational comparison
$202\cdot24(20/247)^4\lambda^4<\lambda/1000$.
Finally $q_0\ge99\lambda/100$, proving the claim.
\end{proof}

\paragraph{Formalization and mathematical boundary.}
The independent symbolic tests prove the matrix identities, both
jet cancellations and the reduced stress formulas; a negative
control catches omission of the coefficient derivatives.
The accompanying Lean source proves only the algebraic
$2\times2$ cancellation and rational inequalities, conditional
on nonzero determinant. Neither source integral positivity nor
the pulse-wide A.29--A.30 ratio is Lean-formalized. A single
admissible exterior tuple, full interval cover, all A.11 cone
supports, Appendix-B contraction, actual unforced startup
solutions and the physical zero-force residual remain open.


\subsection*{What a complete pulse certificate must establish}
The translated axial correction bumps impose exact terminal moments, but
the cancellation does not bound source stresses in the pulse interior.
The inequalities in (A.29)--(A.30) are required on
$\mathcal D=[0,13]\times[-1,1]$, with $\xi=\lambda y$.
Point tests and endpoint values are insufficient.

\begin{proposition}[Soundness of a rational rectangular covering]
\label{prop:full-pulse-cover}
Let $\mathcal C$ be finitely many nondegenerate closed rectangles
$C_k=[a_k,b_k]\times[c_k,d_k]\subset\mathcal D$ with rational
endpoints. Form the grid of their two coordinate endpoint sets,
together with $0,13,-1,1$. If every elementary open grid rectangle
lies inside one of the $C_k$, then $\bigcup_k C_k=\mathcal D$.
\end{proposition}
\begin{proof}
Each open elementary rectangle is contained in an input closed
rectangle, so its closure is contained in that rectangle as well.
The closures of all elementary rectangles cover $\mathcal D$.
\end{proof}

\begin{proposition}[Conditional A.29--A.30 interval soundness]
\label{prop:full-pulse-residual-soundness}
Suppose all source field and first-$\eta$ derivative enclosures used
in \eqref{eq:pulse-exact-eQ}--\eqref{eq:pulse-exact-eN} are
independently valid on every rectangle of a complete covering.
Suppose their rational denominator intervals exclude zero and
$X,E,H,Q_s>0$, with the source identity $H^2=2XE^2$.
If exact outward interval arithmetic verifies on every rectangle
\[
|e_Q/q_0|\le10^{-10},\quad |e_m/q_0|\le10^{-7},
\quad |\eta m_\eta/q_0|\le2\cdot10^{-7},\quad
|e_N/q_0|\le10^{-7},\quad |w-w_0|<10^{-6},
\]
then those inequalities hold everywhere on $\mathcal D$.
\end{proposition}
\begin{proof}
The residual formulas are source-native algebraic identities.
Exact rational interval operations enclose their pointwise values
whenever division avoids zero. The finite covering proposition
then transfers the cell bounds to the whole domain.
\end{proof}

\paragraph{A fail-closed source audit.}
The script \texttt{full\_pulse\_cover\_gate.py} checks complete
two-dimensional rational coverage, identical source parameter boxes,
source consistency $H^2=2XE^2$, positive denominators and the existing
A.29--A.30 residual budget on every provided cell. Deliberate
incomplete covers and synthetic invalid inputs are rejected.
The checker requires source and proof hashes, but \emph{hashes
are not mathematical evidence of interval enclosure correctness}.
It therefore always emits the status
\texttt{CONDITIONAL\_SOURCE\_ENCLOSURES\_NOT\_INDEPENDENTLY\_VERIFIED},
even when the arithmetic comparisons pass.

No independently authenticated full-pulse source interval boxes have
been supplied. The two (A.11) cone supports, subsequent transitions,
terminal matching and gluing, Appendix-B fixed point, actual unforced
startup correction equations and reconstructed zero-force
Navier--Stokes residual remain unproved. The originating global
Navier--Stokes construction uses smooth external forcing. The proposed
unforced Theorem~1.1 remains open; elementary Lean scalar gates are not
an end-to-end kernel proof.


\subsection*{Centerline precision and soundness of a pulse certificate}
\label{subsec:centerline-proof-gate}
The source's normalized pulse denominator is
\[
D=\tfrac12-h,\quad\beta=\tfrac12-\lambda,\quad
c(\eta)=\frac{2D\eta^2}{1+\eta^2},\quad
q_0(\eta)=\frac{\lambda-h+c(\eta)}{1-\lambda}.
\]
Thus, for $0<h\le\lambda/10^4$ and $0<\lambda\le10^{-14}$,
\[
\frac{9999}{10000}\lambda\le q_0(0)
=\frac{\lambda-h}{1-\lambda}
\le\frac{\lambda}{1-\lambda},\qquad q_0(\eta)\ge q_0(0)>0.
\]
\begin{proposition}[Necessary absolute source precision]
\label{prop:source-pulse-precision}
If the A.29--A.30 comparison inequalities
$|e_Q/q_0|\le10^{-10}$ and $|e_N/q_0|\le10^{-7}$
hold, then
\[
|e_Q(0)|\le 10^{-10}\frac{\lambda-h}{1-\lambda},
\qquad |e_N(0)|\le10^{-7}\frac{\lambda-h}{1-\lambda}.
\]
At $\lambda=10^{-14}$ and $h=\lambda/10^4$, the
right sides are of order $10^{-24}$ and $10^{-21}$,
respectively. They shrink proportionally for smaller $\lambda$.
\end{proposition}
\begin{proof}
Use $c\ge0$ with equality at $\eta=0$; multiply the
relative inequalities by the exact positive $q_0(0)$.
The endpoint estimates follow by rational arithmetic.
\end{proof}
Thus an absolute error of merely $10^{-16}$ obtained by
subtracting order-one source terms is inadequate. A correct
interval computation must first exploit source cancellations
and use sufficiently accurate outward enclosures.

\paragraph{Cancellation before interval evaluation.}
Set
\[
S=D+c,\quad
C_d=\frac{\lambda S(1-\lambda)}
{\beta^2(\lambda-h+c)},\quad w_0=2R-C_dR',
\]
\[
e_m=m-\frac R\beta+\frac{\lambda R'}{\beta^2},
\quad e_Q=Q_s-q_0,\quad
e_N=\frac{N_s}{E}+R-Sm+D\eta m_\eta.
\]
\begin{lemma}[Exact denominator-cleared pulse identity]
\label{lem:pulse-source-cleared}
Whenever the physical denominators are nonzero,
\begin{align}
\beta^2(N_s/E-w_0Q_s)={}&
q_0\lambda R\beta+S\beta^2 e_m
-D\eta\beta^2m_\eta+\beta^2e_N-\beta^2w_0e_Q.
\label{eq:source-cleared-pulse}
\end{align}
Equivalently,
\[
\frac{N_s}{EQ_s}-w_0=
\frac{q_0(\lambda/\beta)R+Se_m-D\eta m_\eta+
e_N-w_0e_Q}{q_0+e_Q}.
\]
\end{lemma}
\begin{proof}
After multiplying by $\beta^2$, the coefficient of $R$
is $\beta[S-\beta-q_0(2\beta+\lambda)]$; it vanishes since
$S-\beta=q_0(1-\lambda)$ and $2\beta+2\lambda=1$.
The $R'$ coefficient is
$C_dq_0\beta^2-\lambda S=0$. The remaining
terms cancel by definition. Independent symbolic and
rational replays test both identities.
\end{proof}

\paragraph{Authenticating proof inputs rather than labels.}
The existing rational full-cover program checks whether the
provided rectangles cover $[0,13]\times[-1,1]$ and whether
their supplied intervals meet each residual budget. Its
former manifest accepted any syntactically valid SHA-256
string without comparing it to file bytes. The new
\texttt{pulse\_provenance\_replay\_gate.py} binds the three
primary evidence files to recomputed SHA-256 digests, rejects
missing or tampered files and traversal outside the evidence
root, and then runs the exact rectangle and residual checks.
Negative tests also reject empty coverings.

This improves reproducibility but does not prove that the
submitted intervals enclose the actual Appendix-A fields.
Such a theorem must independently validate, on every closed
cell, the source's field construction, derivatives, ODE
comparison and interval arithmetic. No complete source-derived
interval dataset or jointly admissible source tuple is present
in this release. A matching hash is never a substitute for
mathematical enclosure validity.

\paragraph{Lean and physical status.}
\texttt{PulseResidualPrecision.lean} proposes kernel lemmas for
the cleared polynomial identity and centerline rational
scales. Until a Lean job completes successfully, those
statements are not kernel-certified. The full source interval
construction, exterior stress cone, nonlinear Appendix-B
analytic profile, unforced correction hierarchy and exact
physical zero-force Navier--Stokes residual remain unproved.
The originating OpenAI whole-space theorem uses a smooth
external force \cite{openai2026}; it does not prove the
proposed unforced Theorem 1.1.


\subsection*{Actual-source interval enclosures for the fixed pulse}
The full-pulse checker is conditional until the submitted boxes are
proved to enclose the prescribed source functions. We certify a first
genuine layer: the fixed Appendix-A smooth step and pulse, not the
coupled physical moments or stress.

The source's smooth step and fixed pulse are
\[
\sigma(t)=\begin{cases}0&t\le0,\\
\displaystyle\frac{e^{-1/t^2}}{e^{-1/t^2}+e^{-1/(1-t)^2}}&0<t<1,\\
1&t\ge1,\end{cases}\quad
\phi_b(\xi)=\int_0^\xi\sigma(50v)\,dv,\quad
R_0(\xi)=\phi_b(\xi)(1-\sigma(\xi-10)).
\]
The source identity $\sigma(t)+\sigma(1-t)=1$ gives
$\phi_b(\xi)=\xi-1/100$ for $\xi\ge1/50$. Moreover
\[
R_0'=\sigma(50\xi)(1-\sigma(\xi-10))
-\phi_b(\xi)\sigma'(\xi-10),\quad
\partial_\eta R_0=\partial_\eta R_0'=0.
\]
The actual main pulse $A(\eta)R_0(\xi)$ generally has a nonzero
parameter derivative and is not covered by this identity.

\begin{lemma}[Exact exponential and short-primitive enclosures]
For rational $0\le x\le80$ and $x/(N+2)<1$, set
$P_N(x)=\sum_{k=0}^N x^k/k!$. Then
\[
P_N(x)\le e^x\le
P_N(x)+\frac{x^{N+1}}{(N+1)!}\frac1{1-x/(N+2)}.
\]
For $x\ge80$, $e^x>2^{80}$. Also $0\le\sigma'\le64$,
$\sigma$ is nondecreasing, and for rational
$0\le\xi\le1/50$, $h=\xi/N$,
\[
h\sum_{j=0}^{N-1}\sigma(50jh)
\le\phi_b(\xi)\le
h\sum_{j=1}^{N}\sigma(50jh).
\]
\end{lemma}
\begin{proof}
The exponential series has nonnegative terms; successive omitted
terms have ratio no larger than $x/(N+2)$. The inequality $e>2$
follows from the first three Taylor terms. Differentiating the
exact smooth step gives
\[
\sigma'=\sigma(1-\sigma)\left(2/t^3+2/(1-t)^3\right)\ge0.
\]
On $1/4\le t\le3/4$ its value is at most $64$.
On $0<t\le1/4$ it is at most
$(128+128/27)e^{-128/9}<1$; the other end follows by
symmetry. The primitive inequalities are the monotone left and
right Riemann sums.
\end{proof}
\begin{proposition}[Forty actual-source rational pulse boxes]
The source-derived exact-rational implementation generates twenty
$\xi$ cells and two parameter cells covering the closed rectangle
$[0,13]\times[-1,1]$. On all forty cells it encloses $R_0$, $R_0'$,
$\partial_\eta R_0$ and $\partial_\eta R_0'$, as well as $\phi_b$
and $1-\sigma(\xi-10)$. Every interval has a rational containment
guard $2^{-120}$ beyond the derived bounds. On $1\le\xi\le10$,
\[
R_0(\xi)=\xi-1/100\ge99/100,\qquad R_0'(\xi)=1.
\]
\end{proposition}
\begin{proof}
Evaluate source $\sigma$ at rational endpoints with the preceding
Taylor inequalities; bound the initial primitive by monotone
quadrature and the remaining primitive exactly. Propagate
the enclosures through the displayed pulse and derivative identities
using exact rational interval arithmetic. Extend each bound by the
stated guard; test coverage on the induced rational rectangle grid.
The source function is independent of $\eta$. The plateau identity
is immediate from $\sigma(50\xi)=1$ and
$\sigma(\xi-10)=0$ for $1\le\xi\le10$.
\end{proof}

The reproducible source generator and companion verifier recompute
all forty formula layers, check rational covering, and reject
adversarially corrupted data. Their shared arithmetic implementation
is not a substitute for an independent Lean proof of the analytic
enclosure lemma. The strict guard here is a \emph{containment}
margin, not a physical stress-cone margin or a guarantee of adequate
width for the $10^{-24}$ centerline normalized residual.
No source-derived full-cell intervals have yet been established for
$E,U,M,J,S,\Pi,Q_s,N_s$ and their parameter jets under one admissible
Appendix-A parameter tuple. The A.11 correction supports, all-stage
cone, Appendix B contraction, actual unforced $n=1,2$ solutions,
PDE reconstruction and Theorem~1.1 remain open.


\subsection*{A quantitative slack for the azimuthal A.11 supports}
On the two azimuthal correction supports (A.11), suppose the source has
$U=M=J=0$, its correctly normalized terminal moment satisfies
$S(\infty,\eta)=0$, and its pressure has the infinity gauge.
Writing $D_s=S/(XE^2)$, $P_s=-\Pi/E^2$,
$A=1/2+h$, $d=1-\eta^2$ and $g=\partial_\eta\log E$,
the exact source equation (4.16) gives
\[
\frac{N_s}{E}
=E\{4\eta(hD_s-AP_s)-d[
2g(D_s-P_s)+\partial_\eta D_s-\partial_\eta P_s]\}.
\]
If all future logarithmic slopes are nonpositive,
then $0\le P_s\le1/2$. If the future parameter-logarithmic
derivative has modulus at most $\varepsilon$, differentiation of
the positive tail integrals gives
\begin{equation}\label{eq:A11-updated-Ns-tail}
\frac{|N_s|}{E}
\le E[4hD_s+4AP_s+2\varepsilon(D_s+P_s)].
\end{equation}
\begin{proposition}[Strict supportwise $N_s/E$ margin]
Suppose, simultaneously across both complete A.11 supports,
$E\le\lambda/4$, $4hD_s\le61/100$, $P_s\le1/2$,
$2A\le1001/1000$ and $2\varepsilon(D_s+1/2)\le1/100$.
Then
\[
\frac{|N_s|}{E}\le\frac{1621}{4000}\lambda
\quad\text{and}\quad
\frac\lambda2-\frac{|N_s|}{E}\ge\frac{379}{4000}\lambda.
\]
\end{proposition}
\begin{proof}
The square bracket in \eqref{eq:A11-updated-Ns-tail}
is at most $61/100+1001/1000+1/100=1621/1000$.
Multiply by $E\le\lambda/4$ and subtract from $\lambda/2$.
\end{proof}
\begin{corollary}[Improved normalized strict cone]
Suppose additionally that $Q_s\ge\lambda/2$,
$-3\lambda/2\le\ell\le-\lambda/2$, $0<\lambda\le1/20$
and $X\ge9/\lambda$. For $a=2-2\ell$,
$p_1=XQ_s/(1-2h\eta^2)$ and $w=N_s/(EQ_s)$ one has
\[
2(1-a/p_1)^2-(a-2)w^2
\ge\frac{32117077}{80000000}>\frac25.
\]
\end{corollary}
\begin{proof}
The preceding result and $Q_s\ge\lambda/2$ give
$|w|\le1621/2000$. Also
$2+\lambda\le a\le2+3\lambda\le43/20$ and
$p_1\ge9/2>2a$. Therefore the expression is bounded
below by $1/2-(3/20)(1621/2000)^2$.
\end{proof}
The supplied exact-rational checker establishes these arithmetic
implications only. It does not authenticate the analytic bounds
for a frozen Appendix-A source tuple. The originating OpenAI
global theorem includes smooth external forcing; unforced
Theorem~1.1 and full Lean formalization remain open.


\subsection*{Endpoint reduction for the two true A.11 correction supports}
After the angular interpolation the unedited source azimuthal profile
$E_b(y)$, $y=\log X$, is independent of $\eta$. On the interval
$[y_-,y_+]$ between the extreme support endpoints of the two
width-$3/10$ A.11 correction bumps, its logarithmic slope is
$\ell_b=\partial_y\log(\sqrt X E_b)=-\lambda$. Write the actual
corrected profile as
\[
 E_c(y,\eta)=E_b(y)(1+u(y,\eta)),\qquad
 u=\sum_{i=1}^2c_i(\eta)\beta(y-y_i).
\]
The two source bumps have disjoint supports and differ from the
axial moment corrections (A.9).

\begin{proposition}[Full-support corrected-field bounds]
\label{prop:a11-endpoint-jets}
Suppose the genuine moment-correction coefficients obey the source
bounds $|c_i|\le192\lambda^{28}$ and
$|c_{i,\eta}|\le12\lambda^{28}$ and the prescribed smooth-step
derivative satisfies $|\beta|\le640/3$, $|\beta'|\le200000$.
For $0<\lambda\le1/20$ put
$\delta=40960\lambda^{28}<10^{-31}$. Then, simultaneously at
every point of both supports and for every $|\eta|\le1$,
\begin{align*}
 (1-\delta)E_b&\le E_c\le(1+\delta)E_b,\\
 |\partial_\eta\log E_c|&\le
     2560\lambda^{28}/(1-\delta),\\
 |\ell_c+\lambda|&\le38400000\lambda^{28}/(1-\delta)
     <\lambda/2,\\
 |\partial_\eta\ell_c|&\le
     2400000\lambda^{28}/(1-\delta)^2.
\end{align*}
\end{proposition}
\begin{proof}
At most one bump is active at a given radius.
For its coefficient $c$ and bump $\beta$, the exact formulas are
\[
 \partial_\eta\log E_c=\frac{c_\eta\beta}{1+c\beta},
 \quad \ell_c=-\lambda+\frac{c\beta'}{1+c\beta},
 \quad\partial_\eta\ell_c=\frac{c_\eta\beta'}{(1+c\beta)^2}.
\]
The last identity includes an exact numerator cancellation.
Insert the specified source bounds and
$40960/20^{28}<10^{-31}$.
\end{proof}

\begin{theorem}[Positive-tail comparison and endpoint extrema]
\label{thm:a11-endpoint-tails}
Assume the actual terminal $S(\infty)=0$ moment condition and
pressure gauge $\Pi(\infty)=0$, and define
\[
 D_k=\frac1{2XE_k^2}\int_X^\infty E_k(x)^2dx,\qquad
 P_k=\frac1{2E_k^2}\int_X^\infty\frac{E_k(x)^2}{x}dx,
 \quad k=b,c.
\]
For any $y$ on either correction support, with
$\kappa=(1+\delta)/(1-\delta)$,
\[
 \kappa^{-2}D_b(y)\le D_c(y,\eta)\le\kappa^2D_b(y),
 \quad
 \kappa^{-2}P_b(y)\le P_c(y,\eta)\le\kappa^2P_b(y).
\]
Here $\kappa^2\le1+8\delta$. Moreover
\[
 D_b'=2\lambda D_b-\tfrac12,\qquad
 P_b'=(1+2\lambda)P_b-\tfrac12
\]
on all of $[y_-,y_+]$; hence the maximum of each uncorrected
tail moment on both supports occurs at an extreme endpoint.
Writing $D_*=\max(D_b(y_-),D_b(y_+))$ and
$P_*=\max(P_b(y_-),P_b(y_+))$, one obtains
\[
 E_c\le(1+\delta)E_b(y_-),\quad
 D_c\le(1+8\delta)D_*,\quad
 P_c\le(1+8\delta)P_*.
\]
If the \emph{entire corrected future} has $\ell_c\le0$, its
positive pressure kernel gives the independent $P_c\le1/2$.
\end{theorem}
\begin{proof}
The quotient of corrected and uncorrected normalized future
tail kernels is exactly
$\bigl[(1+u(y+t,\eta))/(1+u(y,\eta))\bigr]^2$.
It lies in $[\kappa^{-2},\kappa^2]$ because
$|u|\le\delta$ at every future radius. Positivity of the
integrals gives the comparison. Differentiating the unedited
tails gives $D_b'=-1/2-2\ell_bD_b$ and
$P_b'=-1/2+(1-2\ell_b)P_b$. With $\ell_b=-\lambda$,
their solutions are monotone or constant. Also $E_b$ decreases
strictly on $[y_-,y_+]$ and
$\kappa^2=1+4\delta/(1-\delta)^2\le1+8\delta$.
\end{proof}

\begin{corollary}[A three-value conditional source gate]
\label{cor:a11-three-values}
Let $\varepsilon=2560\lambda^{28}/(1-\delta)$.
If the same actual admissible source tuple and true terminal
normalizations additionally establish
\[
\frac{(1+\delta)E_b(y_-)}{\lambda}\le\frac14,\quad
4h(1+8\delta)D_*\le\frac{61}{100},\quad
2\varepsilon((1+8\delta)D_*+\tfrac12)\le\frac1{100},\quad
2A\le\frac{1001}{1000},
\]
and all corrected future slopes are nonpositive, then
$|N_s|/E_c\le1621\lambda/4000<\lambda/2$ with
exact slack $379\lambda/4000$ on both supports. The further
stress-cone statement needs separate $Q_s$ and radius bounds.
\end{corollary}

\paragraph{Source authentication remains open.}
This reduces the supportwise task to $E_b(y_-)$,
$D_b(y_-),D_b(y_+)$ and the true coefficient and
terminal conditions. It does not supply certified endpoint
values for a single actual Appendix-A tuple, nor prove the
outer stress cone, Appendix-B contraction, distinct unforced
startup solutions or physical zero-force residual. Finite
rational or synthetic interval tests do not fill these gaps.


\subsection*{A source-geometric short-tail estimate for both A.11 supports}
The two azimuthal corrections in source (A.11) are supported in intervals
of width $3/10$ centered three and one logarithmic radial units before
the end of the $-\lambda$ plateau. Write
\[
D(y,\eta)=\frac{1}{2XE(y,\eta)^2}\int_X^\infty E(x,\eta)^2\,dx.
\]
\begin{proposition}[Remaining exterior tail]\label{prop:a11-short-tail-new}
Assume the source's unedited plateau $\ell_b=-\lambda$, followed by
a one-unit transition with $\ell_b\le0$, a $-1$ hold of length
$4\log(1/h)$, and a remaining tail with $\ell_b\le-3h/4$.
For every point of either support,
\[
D_b(y)\le\frac{93}{40}+\frac{h^7}{3}.
\]
If $E_c=E_b(1+u)$ and $|u|\le\delta<1$ at every later radius,
then
\[
D_c\le\left(\frac{1+\delta}{1-\delta}\right)^2
 \left(\frac{93}{40}+\frac{h^7}{3}\right).
\]
For $0<\lambda\le1/20$, $h\le\lambda/100$ and
$\delta=40960\lambda^{28}$ this is strictly less than $233/100$.
\end{proposition}
\begin{proof}
Use $D_b=\frac12\int_0^\infty
\exp(2\int_0^v\ell_b(y+t)dt)\,dv$.
The remainder of the plateau is at most $63/20$ units and
contributes at most $63/40$; the next transition contributes
$1/2$. The $-1$ hold contributes at most $1/4$ and reduces
the energy factor by $h^8$. Its remaining tail contributes
at most $h^8/(3h)=h^7/3$.
The corrected-to-base kernel ratio lies in
$[((1-\delta)/(1+\delta))^2,
((1+\delta)/(1-\delta))^2]$.
Finally $\delta<10^{-31}$ and the explicit rational bound
$(1+8/10000)(93/40+1/1000)<233/100$ completes the proof.
\end{proof}

\begin{theorem}[Improved conditional A.11 axial stress and cone]
\label{thm:a11-short-tail-new}
Assume the genuine source moment and pressure gauges
$S(\infty)=0$, $\Pi(\infty)=0$, together with
$|c_i|\le192\lambda^{28}$,
$|\partial_\eta c_i|\le12\lambda^{28}$,
the source bump sup bounds, nonpositive corrected future slopes,
and $E_c\le\lambda/4$ throughout the two supports.
For $0<\lambda\le1/20$ and $0<h\le\lambda/100$,
\[
\frac{|N_s|}{E_c}<\frac{51}{200}\lambda<\frac{\lambda}{2},
\qquad \frac{\lambda}{2}-\frac{|N_s|}{E_c}>
 \frac{49}{200}\lambda.
\]
If additionally $Q_s\ge\lambda/2$,
$-3\lambda/2\le\ell_c\le-\lambda/2$ and $X\ge9/\lambda$,
then with $a=2-2\ell_c$,
$p_1=XQ_s/(1-2h\eta^2)$ and $w=N_s/(E_cQ_s)$,
\[
2(1-a/p_1)^2-(a-2)w^2\ge
\frac{92197}{200000}>\frac9{20}.
\]
\end{theorem}
\begin{proof}
The exact source stress-tail identity (4.16) gives
\[
\frac{|N_s|}{E_c}\le E_c\{
4hD_c+4AP_c+2\varepsilon(D_c+P_c)\},
\quad A=\tfrac12+h,\quad
\varepsilon=\frac{2560\lambda^{28}}{1-\delta},
\]
where $P_c=-\Pi/E_c^2$. The future slope ensures
$0\le P_c\le1/2$. The proposition gives $D_c<233/100$,
so $4hD_c<1/200$, $4AP_c\le1001/1000$ and
$2\varepsilon(D_c+P_c)<1/200$.
The bracket is smaller than $51/50$, yielding the first
strict margin. For the cone,
$|w|\le51/100$, $a\le43/20$, and $p_1\ge9/2>2a$.
Thus its left side is at least
$1/2-(3/20)(51/100)^2=92197/200000$.
\end{proof}
These are conditional source-stage inequalities, not a certified
numerical realization of an entire Appendix-A parameter tuple.
They do not check the other stress-cone stages, Appendix B's
contraction, unforced startup equations, or a zero-force physical
Navier--Stokes residual. The originating OpenAI theorem includes
smooth external forcing.


\subsection*{Cancellation of the A.11 local-normalization derivatives}
\label{subsec:A11-local-derivative-cancellation}
On and after the two source (A.11) corrections, write
$E_c(y,\eta)=E_b(y)[1+u(y,\eta)]$, where
$u=\sum_{i=1}^2c_i(\eta)\beta(y-y_i)$ and
$E_b$ is independent of $\eta$. The two disjoint corrected supports
have total logarithmic length $3/5$. This permits an exact
cancellation in the axial source (4.16), stronger than the previous
separate estimates for the parameter derivatives of the tails.

\begin{theorem}[Exact future-support formula]\label{thm:A11-local-cancel}
Assume the \emph{genuine} source normalization
$S(\infty,\eta)=\Pi(\infty,\eta)=0$, and let
$D_c=S/(XE_c^2)$, $P_c=-\Pi/E_c^2$,
$g=\partial_\eta\log E_c$,
$R_b(y,t)=\exp[2\int_0^t\ell_b(y+v)\,dv]$.
Then on either A.11 support,
\begin{equation}\label{eq:A11-local-cancel}
2g(D_c-P_c)+\partial_\eta D_c-\partial_\eta P_c
=\frac{\int_0^\infty R_b(y,t)(1-e^{-t})
 [1+u(y+t)]u_\eta(y+t)\,dt}{[1+u(y)]^2}.
\end{equation}
The right-hand integrand vanishes outside the future portions
of the two corrected supports. If $\ell_b\le0$ at all future
radii, $0<\lambda\le1/20$,
$|c_i|\le192\lambda^{28}$,
$|c_{i,\eta}|\le12\lambda^{28}$ and
$\|\beta\|_\infty\le640/3$, then, with
$\delta=40960\lambda^{28}$,
\begin{equation}\label{eq:A11-local-cancel-bound}
|2g(D_c-P_c)+D_{c,\eta}-P_{c,\eta}|
\le1536\lambda^{28}\frac{1+\delta}{(1-\delta)^2}
<\frac1{1000}.
\end{equation}
\end{theorem}
\begin{proof}
Both source tails are positive integrals. At the current radial
position put $v=1+u(y)$ and write their unnormalized numerators
$A_D,A_P$, so $D_c=A_D/(2v^2)$ and
$P_c=A_P/(2v^2)$. Since the base field is $\eta$-independent,
$g=v_\eta/v$. Direct differentiation gives
$\partial_\eta(A/(2v^2))+2g A/(2v^2)=
A_\eta/(2v^2)$, canceling all local-normalization derivatives.
After changing radial integration variable $x=Xe^t$, the
difference $A_D-A_P$ has the kernel $R_b(y,t)(1-e^{-t})$
times $[1+u(y+t)]^2$, proving \eqref{eq:A11-local-cancel}.
Because $\ell_b\le0$, $0\le R_b\le1$; because the supports
are disjoint, $|u_\eta|\le2560\lambda^{28}$ and
their total future length is at most $3/5$. Divide by
$v^2\ge(1-\delta)^2$ to obtain the bound.
\end{proof}

\begin{corollary}[Improved strict stress and cone margins]
\label{cor:A11-local-cancel-cone}
Assume also $D_c<233/100$, $0<h\le\lambda/100$,
$E_c\le\lambda/4$ on both supports, and the future \emph{baseline}
slope is nonpositive. The source axial stress obeys
\begin{equation}\label{eq:A11-new-margin}
\frac{|N_s|}{E_c}<\frac{63}{250}\lambda,\qquad
\frac{\lambda}{2}-\frac{|N_s|}{E_c}>
\frac{31}{125}\lambda.
\end{equation}
This conclusion does not require a separate assumption that the
\emph{corrected} future slope be nonpositive. If additionally
$Q_s\ge\lambda/2$, $-3\lambda/2\le\ell_c\le-\lambda/2$
and $X\ge9/\lambda$, the normalized cone inequality satisfies
\begin{equation}\label{eq:A11-new-cone}
2(1-a/p_1)^2-(a-2)w^2
\ge\frac{144343}{312500}>\frac9{20},
\quad a=2-2\ell_c,\quad
p_1=\frac{XQ_s}{1-2h\eta^2},\quad
w=\frac{N_s}{E_cQ_s}.
\end{equation}
\end{corollary}
\begin{proof}
The future baseline slope gives $P_b\le1/2$; positive tail
comparison gives $P_c\le(1+8\delta)/2$. Thus
$4AP_c<1002/1000$, $4hD_c<1/200$ and the
derivative term is $<1/1000$. The source (4.16) formula
bounds $|N_s|/E_c$ by $E_c$ times their sum,
which is less than $E_c(1008/1000)\le(63/250)\lambda$.
The exact slack is $31\lambda/125$. With the
extra cone hypotheses, $a/p_1\le1/2$,
$a-2\le3/20$ and $|w|\le63/125$. Hence
$2(1-a/p_1)^2-(a-2)w^2
\ge1/2-(3/20)(63/125)^2=144343/312500$.
\end{proof}

The statements are **conditional mathematical implications**, not
a source-authenticated interval for one completely frozen Appendix-A
parameter tuple. The terminal moment/pressure gauges, coefficient
enclosures, complete exterior and pulse stress gates, nonlinear
Appendix-B profile, distinct unforced $n=1,2$ coefficient solutions
and the zero-force physical PDE remain unverified. The originating
OpenAI global Theorem~1.1 includes smooth forcing.

\subsection*{The parameter derivative of the angular moment on both A.11 supports}
\label{subsec:A11-angular-moment-eta-derivative}
The exact source formula (4.16), on the two corrected azimuthal
supports where $U=M=J=0$ and $W=1$, reads
\begin{equation}\label{eq:A11-Q-retain-eta}
Q_s=-1+(1-h)\frac I{XH}-D\eta\frac{I_\eta}{XH},
\qquad D=\tfrac12-h.
\end{equation}
The last term cannot be omitted: within either correction the
angular moment can depend on $\eta$, even though its eventual
exterior target is parameter independent.

Write $E_c=E_b(1+u)$ with
$u=c_1(\eta)\beta(y-y_1)+c_2(\eta)\beta(y-y_1-2)$
and disjoint source-prescribed width-$3/10$ bumps.
Assume $0<\lambda\le1/20$, $0<h\le\lambda/100$, and the
\emph{source-verified} baseline entrance and coefficient-jet bounds
\[
\left|r_b(y_1)-\frac1{1-\lambda}\right|\le64\lambda^{28},
\quad |\partial_\eta r_b(y_1)|\le64\lambda^{28},\qquad
|c_i|\le192\lambda^{28},\quad|c_{i,\eta}|\le12\lambda^{28},
\quad r_b=\frac{I_b}{XH_b}.
\]
The source plateau has $\ell_b=-\lambda$ and $H_b$ independent of
$\eta$ after interpolation. Thus
$r_b'=1-(1-\lambda)r_b$ and
$r_{b,\eta}'=-(1-\lambda)r_{b,\eta}$.
The first support begins $3/20$ units before $y_1$; hence on both
supports
\[
\left|r_b-\frac1{1-\lambda}\right|\le80\lambda^{28},
\qquad |r_{b,\eta}|\le80\lambda^{28},\qquad0<r_b<2.
\]
Put $\delta=40960\lambda^{28}$ and
$\delta_\eta=2560\lambda^{28}$. The prescribed bump bound
$\|\beta\|_\infty\le640/3$ gives $|u|\le\delta<1/2$ and
$|u_\eta|\le\delta_\eta$. Comparing the corrected and baseline
positive moments, including their $\eta$ derivatives, yields
\begin{equation}\label{eq:A11-Q-moment-jet}
\left|\frac{I_c}{XH_c}-r_b\right|
\le\frac{2\delta}{1-\delta}r_b\le8\delta,\qquad
\left|\frac{I_{c,\eta}}{XH_c}\right|
\le\frac{80\lambda^{28}+\delta_\eta r_b}{1-\delta}
\le10400\lambda^{28}.
\end{equation}
The second estimate retains the earlier parameter dependence of
$I_b$; only the \emph{added} bump integral has parameter derivative
$H_bu_\eta$.

\begin{theorem}[Improved A.11 supportwise angular-stress bound]
\label{thm:A11-Q-eta-improved}
Under the above genuine source entrance and correction-jet bounds,
both complete A.11 supports satisfy
\begin{equation}\label{eq:A11-Q-improved}
Q_s\ge\frac{\lambda-h}{1-\lambda}
-350000\lambda^{28}>\frac{49}{50}\lambda.
\end{equation}
If in addition the source-authenticated terminal $S$ and pressure
gauges, nonpositive future baseline slopes, and tail/derivative
bounds of Corollary~\ref{cor:A11-local-cancel-cone} hold, then
\[
\frac{|N_s|}{E_c}<\frac{63}{250}\lambda<\frac\lambda2.
\]
With the same corrected slope bound and $X\ge9/\lambda$, the
companion cone expression obeys
\[
2(1-a/p_1)^2-(a-2)w^2>
\frac{55743}{50000}>\frac9{20},
\quad a=2-2\ell_c,\quad p_1=XQ_s/(1-2h\eta^2),
\quad w=N_s/(E_cQ_s).
\]
\end{theorem}
\begin{proof}
The source identity \eqref{eq:A11-Q-retain-eta} and
\eqref{eq:A11-Q-moment-jet} bound the loss relative to the
equilibrium value $(\lambda-h)/(1-\lambda)$ by
$(80+8\cdot40960+\frac12\,10400)\lambda^{28}
=332960\lambda^{28}<350000\lambda^{28}$.
Since $h\le\lambda/100$ and
$350000\lambda^{27}\le350000/20^{27}<1/100$,
the first bound follows. The $N_s/E$ estimate is the
separately conditional local-normalization cancellation already
proved in Corollary~\ref{cor:A11-local-cancel-cone}.
The improved $Q_s$ bound yields
$|w|<(63/250)/(49/50)=9/35<13/50$.
Moreover $p_1\ge XQ_s\ge441/50$ and
$a\le43/20$, so $a/p_1<1/4$.
Since $0\le a-2\le3/20$,
\[
2(1-a/p_1)^2-(a-2)w^2
> 2(3/4)^2-(3/20)(13/50)^2
=\frac{55743}{50000}.
\]
\end{proof}
This theorem is a \emph{conditional source implication}.
No genuine frozen A.11 entrance $C_\eta^1$ interval, corrected
terminal-gauge/tail dataset, full exterior cone certificate,
Appendix-B contraction, unforced coefficient startup, or
zero-force physical residual has been supplied. The exact-rational
checker tests implications and rejects missing input gates;
its positive fixture does not authenticate source intervals.
The accompanying Lean module formalizes only scalar consequences,
not source integrals or the full unforced Navier--Stokes theorem.

\subsection*{An exact entrance angular-moment parameter derivative}
\label{subsec:a11-entrance-jet-20261021}
The unedited Appendix-A profile has $E(X,\eta)=e(X)(1+\eta^2)^{-1}$
through the axial pulse (the pulse modifies $U$, not $E$).
Thus $I=\int_0^XH\,dx$ and $XH$ share the same angular factor,
so $r_p=I_p/(X_pH_p)$ is independent of $\eta$.
This supplies an exact entrance cancellation previously treated as
a source-dependent interval input.

During the actual source interpolation (A.10), put
$a=1-\lambda$, $J_0=\log(1+\eta^2)$,
$\theta(t)=1-\sigma(t/T_f)$ and
$G(t,\eta)=\exp[(1-\theta(t))(J_0(\eta)-\log2)]$.
The integrated equation $I_y=XH$ gives
\begin{align}
 r_f(\eta)&=\frac{e^{-aT_f}}{G(T_f,\eta)}
   \left(r_p+\int_0^{T_f}e^{at}G(t,\eta)\,dt\right),\\
 r_{f,\eta}(\eta)&=-J'_0(\eta)\frac{e^{-aT_f}}{G(T_f,\eta)}
   \left(r_p+\int_0^{T_f}\theta(t)e^{at}G(t,\eta)\,dt\right).
 \label{eq:a11-new-entrance-eta}
\end{align}
The last formula retains the parameter derivative of the complete
angular moment. In particular, $|r_{f,\eta}|\le r_f\le40/17$
on $|\eta|\le1$ when $0<\lambda\le1/20$ and the actual
pre-interpolation source slope has $1+\ell\ge17/20$.

The unedited source then has $\ell=-\lambda$ on a plateau of
length $30\log(1/\lambda)$. Both entire A.11 correction supports
lie at least $30\log(1/\lambda)-63/20$ after interpolation.
The exact plateau transport, together with
$e^{1197/400}<20$ and the monotonicity of
$e^{(1-\lambda)63/20}\lambda^{30(1-\lambda)-28}$,
gives the **uniform all-support bounds**
\begin{equation}\label{eq:a11-new-rbjet}
  \left|r_b-\frac1{1-\lambda}\right|<18\lambda^{28},
  \qquad |r_{b,\eta}|<12\lambda^{28}.
\end{equation}

For the source's prescribed distinct smooth-step derivative
bumps (A.11), use its exact nonlinear system
$Bc+Q(c)=(\Delta_I,0)^\top$ with
$B=\begin{psmallmatrix}u&ux\\2v&2vz\end{psmallmatrix}$,
$\Delta_I=(1-\lambda)^{-1}-r_b(y_1)$.
Previous source estimates give
$\|B^{-1}\|_\infty<3/2$, $\|Q\|<320$,
$u\ge4/5$, $x\ge6$, $z\le1/7$.
The new $|\Delta_I|<18\lambda^{28}$ and
$|\Delta_{I,\eta}|<12\lambda^{28}$ produce a unique
small correction branch satisfying
\begin{equation}\label{eq:a11-new-cjet}
 \|c\|_\infty\le54\lambda^{28},
 \qquad \|c_\eta\|_\infty<6\lambda^{28}.
\end{equation}
For the derivative, differentiate the *actual* system:
$(B+DQ(c))c_\eta=(\Delta_{I,\eta},0)^\top$.
The first column of $B^{-1}$ is
$(z,-1)^\top/[u(z-x)]$, whose infinity norm is
at most $35/164$. Since
$\|B^{-1}DQ(c)\|<\tfrac12$,
$\|c_\eta\|_\infty\le2(35/164)|\Delta_{I,\eta}|
<6\lambda^{28}$.

The new estimates imply the four previously stated conditional
moment and coefficient $\eta$-jet assumptions in
Theorem~\ref{thm:A11-Q-eta-improved}. The latter's supportwise
$|N_s|/E<63\lambda/250<\lambda/2$ and stress-cone conclusions
**remain conditional** on source-authenticated terminal $S$ and
pressure gauges, normalized future tail bounds, other exterior-stage
bounds and compatible source parameter choices. This is a proof of
a sharper source-specific analytic ingredient, not a complete unforced
Navier--Stokes theorem. The accompanying finite Lean arithmetic gates
do not formalize these integrals or the PDE.


\subsection*{Appendix-A.11: exact entrance derivative and a conditional axial-stress certificate}
The original interpolation (A.10), with \(a=1-\lambda\),
\(J_0=\log(1+\eta^2)\), \(\theta=1-\sigma(t/T_f)\) and
\(G=\exp((1-\theta)(J_0-\log2))\), implies
\[
r_f=\frac{e^{-aT_f}}{G(T_f,\eta)}
 \left[r_p+\int_0^{T_f}e^{at}G(t,\eta)\,dt\right].
\]
All pre-interpolation azimuthal stages retain the common
\((1+\eta^2)^{-1}\) factor. Thus \(\partial_\eta r_p=0\) and
differentiation under the integral yields exactly
\[
\partial_\eta r_f=-J_0'(\eta)\left[
\frac{2r_pe^{-aT_f}}{1+\eta^2}
+\int_0^{T_f}e^{-au}\left(\frac2{1+\eta^2}\right)^{\sigma(u/T_f)}
                  \sigma(u/T_f)\,du\right].
\]
Consequently
\[
|\partial_\eta r_f|\le 2r_pe^{-aT_f}
  +\frac2a[\sigma(1/4)+e^{-aT_f/4}].
\]
For \(T_f\ge512\), \(0<\lambda\le1/20\),
\(r_p\le1/a\), the right side is \(<1.41\cdot10^{-6}\).
The source step has \(\|\sigma'\|_\infty\le64\), so the
interpolation slope belongs to \([-\lambda-1/10,-\lambda]\).

Let the prescribed normalized \(\sigma'\) bumps be centered
at \(L-d_i\), \(L=30\log(1/\lambda)\),
\(d_1=3,d_2=1\), and let \(E=E_0(1+q)\),
\(q=\sum_ic_i(\eta)\beta_i\). Put \(b=1+2\lambda\),
\[
u_i=e^{-ad_i}\int e^{az}\beta(z)\,dz,\quad
v_i=e^{bd_i},\quad
\kappa=\frac{\int e^{-bz}\beta(z)^2\,dz}
             {2\int e^{-bz}\beta(z)\,dz},\quad
\epsilon=(1/a-r_f)e^{-aL}.
\]
The TWO actual A.11 requirements are precisely
\[
u_1c_1+u_2c_2=\epsilon,\qquad
v_1c_1+v_2c_2+\kappa(v_1c_1^2+v_2c_2^2)=0.
\]
The support and exponential-weight estimates give
\[
|\det B|>5,\quad \|B^{-1}\|_\infty<7,\quad
\kappa<150,\quad v_1+v_2<36.
\]
Source comparison gives \(|\epsilon|\le3\lambda^{30a}\),
\(|\partial_\eta\epsilon|\le5\lambda^{30a}\).
On the ball \(\|c\|_\infty\le14|\epsilon|\), the exact
quadratic moment map has Lipschitz constant at most
\(7\cdot10800\cdot14|\epsilon|<1/2\). It has a unique
small solution with
\[
\|c\|_\infty\le42\lambda^{30a},\qquad
\|\partial_\eta c\|_\infty\le70\lambda^{30a}.
\]
The source \(\sigma'\)-bump satisfies \(\|\beta\|_\infty<214\),
\(\|\beta'\|_\infty<56000\); hence
\(\|q\|_\infty<1/4\), \(\|q_\eta\|_\infty<1\) and
\(\|q_y\|_\infty<3\lambda/8\).
Thus the corrected log slope is at most \(-\lambda/2\).

The entrance-moment derivative alone DOES NOT bound \(N_s/E\).
Using the exact source identity (4.16), and explicitly assuming
the completed amplitude constraint \(S(\infty)=0\) and
the pressure gauge \(\Pi(\infty)=0\), one obtains after the
pulse (\(U=M=J=0\)):
\[
N_s(y,\eta)=2\eta\int_y^\infty
(h e^{v-y}-A)E(v,\eta)^2\,dv
+\frac{1-\eta^2}{2}\int_y^\infty
(1-e^{v-y})\partial_\eta E(v,\eta)^2\,dv .
\]
If the remaining source slopes obey \(\ell_0\le0\),
the \(-1\) hold lasts \(4\log(1/h)\), and
\(\ell_0\le-3h/4\) afterward, then the entire future
baseline energy obeys
\[
\int_y^\infty E_0(v)^2\,dv\le E_0(y)^2,\quad
\int_y^\infty e^{v-y}E_0(v)^2\,dv
\le(93/20+2h^7/3)E_0(y)^2
\]
on both A.11 supports. Since their union has length \(3/5\),
the exact \(N_s\) formula and bounds on \(q,q_\eta\) give
\[
\frac{|N_s|}{E}\le4E_0
\le40P_*\lambda^{15}e^{-13/(2\lambda)}.
\]
Thus the sought \(\sup|N_s|/E\le\lambda/2\) follows from
the explicit sufficient scalar gate
\[
80P_*\lambda^{14}e^{-13/(2\lambda)}\le1.
\]
This is a CONDITIONAL SOURCE-FAMILY statement. The source
requires \(S(\infty)=0\), actual post-pulse slope bounds, and
one simultaneously admissible parameter realization; the
diagnostic numbers in the accompanying exact checker do not
certify those missing hypotheses. No unforced Theorem 1.1,
global singularity, or end-to-end Lean proof is asserted.

\subsection*{A stronger local axial-stress margin on the two A.11 supports}
\label{subsec:A11-strong-20261021}
Write $a=1-\lambda$, $b=1+2\lambda$, $T_f=512$ and
$L_p=30\log(1/\lambda)$. Use the actual normalized source
$\sigma'$ bumps $\beta_i$ centered at $L_p-3$ and $L_p-1$.
After the parameter-flattening interpolation the unedited
azimuthal field $E_0$ is independent of $\eta$.
Set $E=E_0(1+\gamma)$,
$\gamma=\sum_{i=1}^{2}c_i(\eta)\beta_i$.
The incoming moment ratio $r_p=I_p/(X_pH_p)$ is independent
of $\eta$ and obeys $0\le r_p\le1/a$.
The exact A.10 entrance formula (proved above) now gives
\[
 0\le r_f\le(a-1/10)^{-1},\qquad
 \sup_{|\eta|\le1}|r_{f,\eta}|<1/700000.
\]
The angular-moment deficit, normalized at the first bump center,
is $\delta=(a^{-1}-r_f)e^{-a(L_p-3)}$. Hence
\begin{equation}\label{eq:A11-strong-delta}
 |\delta|\le50\lambda^{28},\qquad
 |\delta_\eta|\le3\cdot10^{-5}\lambda^{28},
 \qquad0<\lambda\le1/20.
\end{equation}
To retain the actual nonlinear pressure-neutral constraint, define
\[
 u=\int e^{at}\beta(t)\,dt,\quad
 p=\int e^{-bt}\beta(t)\,dt,\quad
 w=\int e^{-bt}\beta(t)^2\,dt,\quad
 x=e^{2a},\quad z=e^{-2b}.
\]
Both exact A.11 equations are
\[
 u(c_1+xc_2)=\delta,\qquad
 2p(c_1+zc_2)+w(c_1^2+zc_2^2)=0.
\]
The source bump satisfies $1/2<u,p<2$, $6<x<8$,
$0<z<1/4$, $w<427$,
$\|\beta\|_\infty\le640/3$ and
$\|\beta'\|_\infty<250000$. Solving the first linear
equation and then the second by contraction on
$\|c\|_\infty\le|\delta|$ gives a map with
\[
 \|\mathcal F(c)\|_\infty
 \le(8/23)|\delta|+743\|c\|_\infty^2,\qquad
 \operatorname{Lip}\mathcal F\le1485|\delta|<1/100.
\]
Thus this ball contains a unique solution, and
differentiation of the \emph{exact} two-dimensional system yields
\[
 \|c\|_\infty\le|\delta|,\quad
 \|c_\eta\|_\infty\le|\delta_\eta|,\quad
 |\gamma|<1/1000,\quad |\gamma_y|<\lambda/4,\quad
 |\partial_\eta\log E|<1/100.
\]
In particular the corrected slope satisfies
$\ell=-\lambda+\partial_y\log(1+\gamma)\le-\lambda/2$.

\begin{theorem}[Improved conditional A.11 axial-stress bound]
\label{thm:A11-support-strong-20261021}
Assume the full source terminal gauges
$S(\infty,\eta)=\Pi(\infty,\eta)=0$ and the prescribed
post-correction stage slopes, including the one-unit
release transition, a $4\log(1/h)$ slope-$-1$ hold,
and $\ell\le-3h/4$ thereafter. Take
\[
 T_f=512,\quad0<\lambda\le
 \min\{1/20,(1+\log P_*)^{-1}\},\qquad
 0<h\le\min\{1/10000,\lambda/100,e^{-T_d}/100\}.
\]
For both full source-corrected A.11 supports,
\begin{equation}\label{eq:A11-strong-support-cone-20261021}
\sup_{\substack{y\in\operatorname{supp}\beta_1
 \cup\operatorname{supp}\beta_2\\|\eta|\le1}}
  \frac{|N_s(y,\eta)|}{E(y,\eta)}
 <\frac{2}{25}\lambda<\frac\lambda2.
\end{equation}
\end{theorem}
\begin{proof}
Because $U=M=J=0$ on these supports, (4.16) and the
genuine terminal gauges give the exact future-energy formula
\[
 N_s=\int_y^\infty\left[
 2\eta(he^{v-y}-A)E(v)^2+
 (1-\eta^2)(1-e^{v-y})E(v)E_\eta(v)\right]dv,
 \qquad A=\tfrac12+h.
\]
Put $\mathfrak d=(2XE^2)^{-1}\int_X^\infty E^2dx$.
The remaining support-to-plateau distance is at most $63/20$.
The release transition, steep hold and final decay imply
$\mathfrak d\le(63/20+1)/2+1/4+h^7/3<3$.
After interpolation the only $\eta$ dependence occurs on the
two bumps, whose combined span is $23/10<3$.
Every later slope is nonpositive, so
\[
 |N_s|/E\le E[4h\mathfrak d+2A+
 3\sup|\partial_\eta\log E|]<2E.
\]
The source reference, pulse and intermediate-stage decay
give, on both corrected supports,
\[
 E<20P_*e^{-13/(2\lambda)}\lambda^{15}.
\]
Since $\lambda\le(1+\log P_*)^{-1}$,
$P_*e^{-13/(2\lambda)}\le e^{-13/2}<1/500$;
therefore $E<\lambda^{15}/25\le\lambda/25$.
Combining the last two inequalities proves the theorem.
\end{proof}

This improves the previous \emph{conditional} A.11 bound
$|N_s|/E<63\lambda/250$ with an explicit sufficient
smallness condition on $\lambda$ relative to $P_*$.
The full global correction parameter order and compatibility
of the terminal gauges remain to be certified; the theorem
does not establish the Appendix-B contraction, unforced
startup solutions or the proposed unforced Navier--Stokes
singularity. The accompanying Lean file proves only finite
scalar inequalities; no complete Lean proof is claimed.


\subsection*{A simultaneous small-parameter choice and a sharper conditional cone margin}
\label{subsec:a11-joint-smallness-20261021}
The local A.11 support estimates require a single order of parameter choices. We distinguish the scalar feasibility of the previously derived smallness inequalities from the stronger assertion that all the source integral, moment and stress conditions hold for the same exterior profile.

\begin{proposition}[Nonempty simultaneous smallness wedge]
\label{prop:a11-smallness-wedge}
Fix $T_d>10$, $P_*>e^{T_d}$ and assume the originating construction permits $T_f=512$. Set
\[
Q=1+T_d+T_f+\log P_*,\qquad
\lambda=\frac{10^{-20}}{Q},\qquad
h=\frac{\lambda e^{-T_d}}{10^6}.
\]
Then, simultaneously,
\[
0<\lambda<
\min\left\{\frac1{20},10^{-14},
\frac{10^{-6}}{T_d+T_f+1},
\frac1{10(1+\log P_*)}\right\},
\]
\[
0<h<\min\left\{\frac{\lambda}{10^4},
\frac{e^{-T_d}}{100},\frac1{10^4}\right\},
\qquad \lambda T_f<10^{-6},\qquad
\lambda\log(1/\lambda)<61\cdot10^{-20}<10^{-18}.
\]
\end{proposition}
\begin{proof}
Since $T_d>10$ and $P_*>e^{T_d}$, $Q>T_d+T_f+1>1$ and $Q>1+\log P_*$. The displayed rational comparisons for $\lambda$ follow. As $0<e^{-T_d}<1$, the bounds on $h$ follow, and $\lambda T_f<10^{-20}$ because $Q>T_f$.
Finally, $\log 10<3$ and $\log Q\le Q-1$, so
\[
\lambda\log(1/\lambda)
=\frac{10^{-20}}Q(20\log10+\log Q)
<10^{-20}\frac{60+Q-1}Q\le60\cdot10^{-20}.
\]
\end{proof}
This proves feasibility of the \emph{listed scalar} inequalities only.
Unquantified source-specific thresholds may require further decreases
of $\lambda$, and no moment, pressure, slope or stress-cone gate is
verified by this construction.

\begin{proposition}[Sharpened conditional A.11 strict cone]
\label{prop:a11-joint-cone}
Assume on both complete A.11 correction supports of the \emph{same}
source profile that
\[
|N_s|/E<2\lambda/25,\qquad
Q_s\ge49\lambda/50,\qquad
-3\lambda/2\le\ell\le-\lambda/2,\qquad
X\ge9/\lambda,
\quad0<\lambda\le1/20,
\]
and $0<L=1-2h\eta^2\le1$, $|\eta|\le1$, $U=M=J=0$.
With $a=2-2\ell$, $p_1=XQ_s/L$ and $w=N_s/(EQ_s)$,
\[
2(1-a/p_1)^2-(a-2)w^2>
\frac{107949}{96040}>1>\frac9{20}.
\]
\end{proposition}
\begin{proof}
The stated source hypotheses give $|w|<4/49$,
$p_1\ge441/50$ and $2<a\le43/20$, so $a/p_1<1/4$.
Since $0<a-2\le3/20$, the expression is bounded strictly
below by neither the source margin nor a weaker inequality, but
\emph{above} by the exact lower estimate
\[
2(3/4)^2-\frac3{20}(4/49)^2=\frac{107949}{96040}.
\]
\end{proof}
The hypotheses $|N_s|/E<2\lambda/25$ and
$Q_s\ge49\lambda/50$ are conditional in the source programme:
their proofs require the actual nonlinear A.11 moment roots and
parameter derivatives, terminal $S$ and pressure gauges,
positive tail integrability, and angular-moment entrance data.
The result is therefore a conditional improvement of the stress
margin on the two correction supports, not a certificate of the
full exterior stress cone.

\paragraph{Joint source and Lean proof obligations.}
The exact fixed $\sigma'$ correction moments $u,p,w$,
the deficit $\delta(\eta)$ and its derivative, and both roots
of
\[
u(c_1+xc_2)=\delta,\qquad
2p(c_1+zc_2)+w(c_1^2+zc_2^2)=0
\]
must be enclosed on the entire source support, with
$x=e^{2(1-\lambda)}$, $z=e^{-2(1+2\lambda)}$.
The small root must be isolated uniformly in $\eta$.
The same frozen source must verify $S(\infty)=\Pi(\infty)=0$,
(A.29)--(A.30) on the entire axial pulse, every exterior
stress-cone stage and the terminal matching conditions.
Subsequently, a source-compatible Appendix B nonlinear profile,
unforced $n=1,2$ corrections, analytic coefficient convergence
and an exact zero-force physical residual are required.
The finite scalar Lean file is not an end-to-end formalization.


\subsection*{A.11: an explicit support-stress certificate}
\label{subsec:A11-support}
In the source Appendix A.11, two width-$3/10$ angular corrections
are inserted three and one logarithmic units before the end of the
uniform plateau. On their supports $U=M=J=0$. Under the terminal
gauges $S(\infty)=\Pi(\infty)=0$, (4.16) becomes
\begin{equation}\label{eq:A11-Ns-exact}
N_s(y,\eta)=\int_y^\infty\left[
2\eta(h e^{v-y}-A)E(v,\eta)^2+
(1-\eta^2)(1-e^{v-y})E(v,\eta)E_\eta(v,\eta)
\right]dv,
\quad A=\tfrac12+h.
\end{equation}
Indeed $S=\frac12\int_X^\infty E^2dx$ and
$\Pi=-\frac12\int_X^\infty E^2/x\,dx$; insert these in
$N_s=(4h\eta S-dS_\eta)/X+4A\eta\Pi-d\Pi_\eta$
and set $x=e^v$.

\begin{lemma}[Piecewise future-energy bound]
Assume $\ell=\partial_y\log(\sqrt X E)\le0$
until the following $\ell=-1$ interval, whose start is at most
$83/20$ units later; that interval lasts $4\log(1/h)$; and
$\ell\le-3h/4$ thereafter, including the terminal transition.
For $0<h\le1/100$,
\begin{equation}\label{eq:A11-future-energy}
E(y)^{-2}\int_y^\infty e^{v-y}E(v)^2\,dv
\le83/20+1/2+2h^7/3<47/10.
\end{equation}
\end{lemma}
\begin{proof}
For $F(t)=e^tE(y+t)^2/E(y)^2$ one has $F'=2\ell F$.
The first part contributes at most $83/20$, the $-1$ hold
at most $1/2$, and its exit factor is at most $h^8$.
The last part contributes at most $h^8 2/(3h)$.
\end{proof}

The actual corrected profile has the form
$E=E_b(1+c_1(\eta)\beta_1+c_2(\eta)\beta_2)$,
where $E_b$ has no $\eta$ dependence on the supports and
$\beta_i(y)=(10/3)\sigma'((10/3)(y-y_i)+1/2)$.
The fixed step has
\begin{equation}\label{eq:A11-bump-C1}
\max\|\beta_i\|_\infty\le640/3,\qquad
\max\|\beta_i'\|_\infty\le2\,000\,000/9.
\end{equation}
For the latter, write
$q=2/t^3+2/(1-t)^3$ and use
$\sigma''=\sigma(1-\sigma)[(1-2\sigma)q^2+q']$.
For $t\in[1/4,3/4]$, its absolute value is at most $17152$.
For $t\le1/4$, let $u=t^{-2}\ge16$ and use
$|\sigma''|\le e^{2-u}(4u^4+26u^2+45)<20000$;
the other endpoint follows by symmetry.

\begin{theorem}[Two-support stress and strict cone, conditional on source coefficients]
Fix $0<\lambda\le1/20$ and
$0<h\le\min(1/100,\lambda/8)$.
Suppose the exact A.11 moment equations, final pressure and $S$
gauges, and the slope hypotheses above hold. Define
$\gamma=\max_i\|c_i\|_{C^0_\eta}$ and
$\gamma_\eta=\max_i\|\partial_\eta c_i\|_{C^0_\eta}$.
If
\begin{equation}\label{eq:A11-certificate-data}
\gamma\le\min\{3/1280,9\lambda/8\,000\,000\},\quad
\gamma_\eta\le3\lambda/10240,\quad
\sup_{\rm supports}E_b\le\lambda/30,
\end{equation}
then $E>0$, $\ell\le-\lambda/2$, and on both supports
\begin{equation}\label{eq:A11-support-Ns}
\boxed{|N_s|/E<\lambda/2}.
\end{equation}
If further $Q_s\ge0$ at the end of interpolation,
there is an unedited constant-slope interval of length at least
$30\log(1/\lambda)-63/20$ before the supports, and
$X\ge35/\lambda$ on the supports, then
\begin{equation}\label{eq:A11-support-cone}
Q_s\ge\lambda/4,\qquad p_1\ge4a,\qquad
2(p_1-a)^2-(a-2)p_2^2\ge21p_1^2/40>0.
\end{equation}
Here $a=2-2\ell$, $b_s=0$,
$p_1=XQ_s/L$, and $p_2=XN_s/(LE)$; thus
the source's admissible stress cone is strict on these
two supports only.
\end{theorem}
\begin{proof}
The disjointness of the bumps gives
$\gamma(640/3)\le1/2$ and
$0<E\le(3/2)E_b\le\lambda/20$.
Moreover
$|\partial_\eta\log E|\le2(640/3)\gamma_\eta
\le\lambda/8\le1$ and
$|\ell+\lambda|\le2(2\,000\,000/9)\gamma\le\lambda/2$.
After the interpolation the $\eta$ derivative is supported only
on these two bumps. Therefore the future-stress formula and
the energy lemma give
\[
|N_s|/E\le E(47/10)[2(A+h)+1]
\le E\,2397/250\le2397\lambda/5000<\lambda/2.
\]
On the unedited plateau
$Q_s'+(1-\lambda)Q_s=\lambda-h$; its nonnegative
initial datum implies $Q_s\ge7\lambda/16>\lambda/4$
at the first support. Across the two corrections
$S_q=-\ell-h-D\eta\,\partial_\eta\log E
\ge5\lambda/16$.
The equation $Q_s'+(1+\ell)Q_s=S_q$
preserves $Q_s\ge\lambda/4$.
Thus $|p_2/p_1|\le2$ and
$2+\lambda\le a\le2+3\lambda$.
By $X\ge35/\lambda$, $p_1\ge35/4>4a$.
Then $(a-2)p_2^2\le12\lambda p_1^2
\le3p_1^2/5$ while
$2(p_1-a)^2\ge9p_1^2/8$, giving the claimed gap.
\end{proof}

\paragraph{The exact nonlinear two-moment calculation.}
Normalize the original A.11 angular-moment defect as
$\Delta_I(\eta)$ at the first bump centre $y_1$.
For the fixed bump $\beta_1$ set
\[
u=\int e^{(1-\lambda)(v-y_1)}\beta_1(v)\,dv,\quad
v_0=\int e^{-(1+2\lambda)(v-y_1)}\beta_1(v)\,dv,\quad
w=\int e^{-(1+2\lambda)(v-y_1)}\beta_1(v)^2\,dv,
\]
and $x=e^{2(1-\lambda)}$, $z=e^{-2(1+2\lambda)}$.
The true A.11 equations are
\[
u(c_1+xc_2)=\Delta_I,\qquad
2v_0(c_1+zc_2)+w(c_1^2+zc_2^2)=0.
\]
Putting $q=\Delta_I/u$, $t=c_2$, $c_1=q-xt$ gives
\[
w(x^2+z)t^2+
[2v_0(z-x)-2wxq]t+(2v_0q+wq^2)=0.
\]
The branch tending to zero with the defect is used.
Its $\eta$ derivative follows from implicit differentiation,
or directly from
\[
t_\eta=
-\frac{(-2wx)t+2v_0+2wq}
 {2w(x^2+z)t+2v_0(z-x)-2wxq}
 \frac{(\Delta_I)_\eta}{u},
\quad(c_1)_\eta=(\Delta_I)_\eta/u-xt_\eta.
\]
This is an explicit interval-root problem, not an
assumed smallness condition. Genuine source-specific interval
bounds on $\Delta_I,(\Delta_I)_\eta$ and $u,v_0,w$ remain to be
supplied. The sufficient baseline bound can be checked from
$E_b\le e_b e^{-13/(2\lambda)}$.

No simultaneous numerical source realization, pulse and exterior
endpoint cone verification, Appendix-B nonlinear contraction,
unforced startup solution, or physical Navier--Stokes residual
follows from these conditional support estimates.

\subsection*{A source-gated final A.11 ratio and an integration-data obstruction}
\label{subsec:a11-source-gated-ratio-20261021}
On both A.11 correction supports the source has $U=M=J=0$.
Its exact stress identity (4.16) consequently reads
\[
N_s=\frac{4h\eta S-(1-\eta^2)S_\eta}{X}
+4A\eta\Pi-(1-\eta^2)\Pi_\eta,\qquad A=\tfrac12+h.
\]
The formulas
$S=\frac12\int_X^\infty E^2\,dx$ and
$\Pi=-\frac12\int_X^\infty E^2\,dx/x$
require, respectively, the true source conditions
$S(\infty)=\Pi(\infty)=0$. They are not implied by the
two A.11 moment-correction equations alone.
With $D_s=S/(XE^2)$, $P_s=-\Pi/E^2$, nonpositive
future logarithmic slopes and
$\sup_{x\ge X}|\partial_\eta\log E(x,\eta)|\le\varepsilon$,
the source identities yield
\[
0\le P_s\le\tfrac12,\qquad
|N_s|/E\le E\{4hD_s+4AP_s+2\varepsilon(D_s+P_s)\}.
\]
\begin{proposition}[Sharpened exact-rational sufficient stress gate]
\label{prop:a11-final-gate-20261021}
Suppose on BOTH complete A.11 supports and every $|\eta|\le1$
that the above genuine source tails hold and
\[
0<\lambda\le\tfrac1{20},\quad0<h\le\lambda/100,\quad
0<E\le\lambda/4,\quad0\le D_s\le233/100,\quad
0\le P_s\le1/2,\quad
\varepsilon\le\frac{2560\lambda^{28}}{1-40960\lambda^{28}}.
\]
Then $|N_s|/E<(503/2000)\lambda<\lambda/2$.
If also $Q_s\ge(49/50)\lambda$, $X\ge9/\lambda$,
$-3\lambda/2\le\ell\le-\lambda/2$ and
$0<L=1-2h\eta^2\le1$, then with
$a=2-2\ell$, $p_1=XQ_s/L$ and $w=N_s/(EQ_s)$,
\[
2(1-a/p_1)^2-(a-2)w^2>
\frac{85676973}{76832000}>\frac9{20}.
\]
\end{proposition}
\begin{proof}
The tail bracket is bounded by
\[
\frac{233}{50000}+1+\frac1{1000}
+\frac{283}{50}\frac{2560\,20^{-28}}{1-40960\,20^{-28}}
<\frac{503}{500}.
\]
Multiplication by $E\le\lambda/4$ proves the first claim.
The other hypotheses give
$|w|<503/1960$, $p_1\ge441/50$ and
$2\le a\le43/20$, so $a/p_1<1/4$.
The displayed cone bound follows from
$2(3/4)^2-(3/20)(503/1960)^2
=85676973/76832000$.
\end{proof}
\begin{proposition}[Why terminal normalization is indispensable]
\label{prop:a11-gauge-obstruction-20261021}
The A.11 bump equations alone do not imply any uniform
bound on $|N_s|/E$ unless the source's accumulated moment
and pressure data are fixed. Formally replacing
$S$ by $S+C_S(\eta)$ and $\Pi$ by $\Pi+C_\Pi(\eta)$
changes $N_s$ by
\[
\Delta N_s=
\frac{4h\eta C_S-(1-\eta^2)C_S'}{X}
+4A\eta C_\Pi-(1-\eta^2)C_\Pi'.
\]
For $C_S=K\eta$, $C_\Pi=0$, this becomes
$\Delta N_s(0)=-K/X$, unbounded as $|K|\to\infty$.
\end{proposition}
This algebraic separation is NOT a counterexample to the
complete source, whose moments and pressure have fixed physical
integral definitions. It proves why genuine terminal gauges and
one independently authenticated A.11 interval covering are
necessary proof inputs. The strict arithmetic gate alone
does not certify any actual source tuple, full exterior cone,
Appendix-B contraction, unforced startup system or
Theorem~1.1. The accompanying finite Lean statements likewise
do not formalize the source integral/PDE arguments.


\subsection*{Exact A.11 axial stress across both bumps and the intervening gap}
\label{subsec:a11-full-hull-20261021}
On the source A.11 corrections, $U=M=J=0$. Let the true terminal
gauges be $S(\infty)=\Pi(\infty)=0$, and let
$E=E_b(1+u)$ where $E_b(y)$ is independent of $\eta$ for
$y\ge y_-$ and $u=\sum_{i=1}^2c_i(\eta)\beta_i(y)$.
The two supports have width $3/10$, their centres differ by
two, and their closed convex hull is $[y_-,y_+]$ of length $23/10$.
Write $F=E^2-E_b^2$. Source (4.16) and the true future-tail
integrals give the exact identity
\[
N_s-N_{s,b}^{\rm tail}
=\int_y^{y_+}\left[
 2\eta(he^{v-y}-A)F(v,\eta)
 +\frac{1-\eta^2}{2}(1-e^{v-y})F_\eta(v,\eta)
\right]\,dv .
\]
This holds at every point of the hull, including the gap; after
$y_+$ the correction vanishes identically. Here $N_{s,b}^{\rm tail}$
is the \emph{auxiliary future-tail functional} of the unedited
$E_b$, not necessarily the physical stress of an unmatched
unedited schedule. In particular the original A.11 constraints
do not by themselves prove $S_b(\infty)=0$.

\begin{proposition}[Conditional full-hull stress gate]
\label{prop:a11-full-hull-20261021}
Suppose the unedited source slope is $\ell_b=-\lambda$ on the
hull, $\ell_b\le0$ afterward, and its normalized future-tail
integral $\sup_{y\in[y_-,y_+]}E_b(y)^{-2}
\int_y^\infty e^{v-y}E_b(v)^2\,dv\le47/10$.
Write $\delta=\sup|u|<1$ and $\epsilon=\sup|u_\eta|$ on the
full hull. Then everywhere, including the gap,
\[
\frac{|N_s|}{E}\le\frac{E_b}{1-\delta}
\left[1+\frac{57}{5}h+
2\left(A+\frac{23}{10}h\right)(2\delta+\delta^2)
+\frac{23}{10}(1+\delta)\epsilon\right].
\]
If additionally $0<\lambda\le1/20$, $0<h\le\lambda/100$,
$E_b\le\lambda/4$, the \emph{actual prescribed} bump
satisfies $\|\beta_i\|_\infty\le640/3$, and
$|c_i|\le54\lambda^{28}$, $|c_{i,\eta}|\le6\lambda^{28}$,
then
\[
\sup_{y\in[y_-,y_+],\ |\eta|\le1}\frac{|N_s|}{E}
\le\frac{101}{400}\lambda<\frac{\lambda}{2}
\]
with a reserve of at least $99\lambda/400$.
\end{proposition}
\begin{proof}
The unedited future-tail contribution is bounded by
$E_b^2[1+2h(1+47/10)]$. Moreover
$F=E_b^2(2u+u^2)$ and
$F_\eta=2E_b^2(1+u)u_\eta$.
On the $23/10$-long hull, the unedited slope
$\ell_b=-\lambda$ gives
$\int_y^{y_+}E_b(v)^2dv\le E_b(y)^2$ and
$\int_y^{y_+}e^{v-y}E_b(v)^2dv\le(23/10)E_b(y)^2$.
Insert these into the exact identity and divide by
$E\ge(1-\delta)E_b$. Source bump disjointness gives
$\delta\le11520\lambda^{28}<10^{-4}$ and
$\epsilon\le1280\lambda^{28}<10^{-4}$.
The resulting worst-case rational bracket is strictly
less than $(101/100)(1-10^{-4})$, and $E_b\le\lambda/4$
finishes the inequality.
\end{proof}
The result is a \emph{conditional extension to the gap};
there is not yet a source-authenticated interval cover of the
corrected A.11 fields, nor a full exterior/terminal cone or a
force-free Navier--Stokes proof. A source-normalization check of
$S(\infty)=\Pi(\infty)=0$ is essential.


\subsection*{Terminal normalization and exact axial parity on A.11}
\label{subsec:A11-parity-terminal-20261021}
The source equation (4.16) reads on the complete A.11 two-bump hull,
where $U=M=J=0$,
\[
N_s={4h\eta S-(1-\eta^2)S_\eta\over X}
 +4A\eta\Pi-(1-\eta^2)\Pi_\eta,\qquad A=\tfrac12+h.
\]
The moment-neutral condition in (A.11) preserves the global pressure
increment $\frac12\int E^2d\log X$. Combined with the actual axis
pressure datum (A.21), it yields $\Pi(\infty,\eta)=0$. The
separate \emph{complete} A.19 amplitude equation, including the
corrected post-pulse field, is needed for $S(\infty,\eta)=0$.
Neither conclusion follows from the local two-bump equations alone.

\begin{lemma}[Exact A.11 equatorial cancellation]
\label{lem:A11-exact-parity-20261021}
Assume the genuine source's unique small branch of the two-bump
system, and the two terminal normalizations above. Its incoming
moment discrepancy is even in $\eta$, and the bump weights are
$\eta$-independent. Hence $c_i(-\eta)=c_i(\eta)$.
The corrected $E,S,\Pi$ are even and $N_s$ is odd, including on
the gap between the two supports. In particular $N_s(X,0)=0$
and $N_s(X,0)/E(X,0)=0$.
\end{lemma}
\begin{proof}
The uncorrected field and the A.10 angular interpolation depend
on $\eta$ through $1+\eta^2$. Reflection in $\eta$ leaves
both moment equations unchanged; uniqueness identifies the reflected
coefficients. Source future-tail integrals preserve evenness.
In the displayed formula the derivative terms are odd, and
all other terms contain an explicit factor $\eta$.
\end{proof}

\begin{theorem}[Sharpened A.11 local parameter-derivative term]
\label{thm:A11-parity-tail-improvement-20261021}
Write $E_c(y,\eta)=E_b(y)(1+u(y,\eta))$ after A.10, where
$u=\sum_i c_i(\eta)\beta_i(y)$ has disjoint supports of combined
logarithmic length $3/5$. Suppose the actual terminal
normalizations hold, $0<\lambda\le1/20$,
$|c_i|\le54\lambda^{28}$,
$|c_{i,\eta}|\le6\lambda^{28}$,
$\|\beta_i\|_\infty\le640/3$, and future baseline
$\ell_b\le0$. Put $\delta=11520\lambda^{28}$ and
$D=S/(XE_c^2)$, $P=-\Pi/E_c^2$,
$g=\partial_\eta\log E_c$. Then
\[
\mathcal T:=2g(D-P)+D_\eta-P_\eta
=\frac{\int_0^\infty R_b(y,t)(1-e^{-t})
(1+u(y+t))u_\eta(y+t)\,dt}{(1+u(y))^2},
\]
where $R_b(y,t)=\exp(2\int_0^t\ell_b(y+v)dv)$, and
\[
|\mathcal T|\le768\lambda^{28}
 \frac{1+\delta}{(1-\delta)^2}<769\lambda^{28}.
\]
If additionally $E_b\le\lambda/[4(1+\delta)]$,
$D_b\le233/100$, $P_b\le1/2$ on the complete
two-support convex hull, and $0<h\le\lambda/100$, then
\[
\frac{|N_s|}{E_c}<\frac{503}{2000}\lambda<\frac\lambda2,
\qquad
\frac\lambda2-\frac{|N_s|}{E_c}>
\frac{497}{2000}\lambda.
\]
\end{theorem}
\begin{proof}
In the future-tail integrals for $D,P$, differentiation
of $(1+u(y))^{-2}$ cancels identically against $2g(D-P)$.
The remaining derivative is supported in the future portions
of the two bumps. Thus $R_b\le1$,
$|u|\le\delta$, $|u_\eta|\le1280\lambda^{28}$
and support length $\le3/5$ prove the $\mathcal T$ bound.
Tail comparison gives $D,P\le\kappa^2 D_b,\kappa^2 P_b$
with $\kappa=(1+\delta)/(1-\delta)$. Source (4.16) implies
$|N_s|/E_c\le E_c(4hD+4AP+|\mathcal T|)$.
The stated result follows from an exact rational comparison
at $\lambda=1/20$, which is the maximum of the positive
upper terms.
\end{proof}

On the uncorrected plateau $\ell_b=-\lambda$,
the normalized tails satisfy
$D_b'=2\lambda D_b-\frac12$ and
$P_b'=(1+2\lambda)P_b-\frac12$.
Both are monotone or constant on the convex hull. Hence the
needed upper bounds reduce to the two outer endpoints,
and $E_b$ is maximal at the first endpoint. This is an
\emph{endpoint reduction}, not an authenticated source interval
certificate. One source-compatible parameter tuple, terminal
moment/pressure integrals, strict $Q_s$ and radial bounds,
the full pulse/exterior cone, Appendix-B contraction and the
unforced coefficient/PDE hierarchy remain to be certified.
The originating OpenAI Theorem 1.1 has smooth external forcing.


\subsection*{Terminal gauges and the A.11 axial stress source}
Source (A.11) preserves the pressure increment, but the separate
equation (A.19) fixes $S(\infty,\eta)=0$.
The pressure terminal condition requires the source's prescribed
axis datum (A.21)--(A.23) and the exact pressure-neutral A.11
correction. They must hold for the \emph{same} corrected profile.
Write $d=1-\eta^2$, $A=1/2+h$, and define the terminal defects
$\sigma=S(\infty,\eta)$ and $\pi=\Pi(\infty,\eta)$.
On both A.11 supports $U=M=J=0$, so (4.15)--(4.16) give
\begin{equation}\label{eq:A11-terminal-gauge-full}
 N_s=N_s^{\rm tail}
 +\frac{4h\eta\sigma-d\sigma_\eta}{X}
 +4A\eta\pi-d\pi_\eta .
\end{equation}
Here $N_s^{\rm tail}$ substitutes
$S_{\rm tail}=\frac12\int_X^\infty E^2dx$ and
$\Pi_{\rm tail}=-\frac12\int_X^\infty E^2dx/x$.
The identity holds before terminal normalization and implies
\[
 \frac{|N_s-N_s^{\rm tail}|}{E}
 \le\frac{4h|\sigma|+|\sigma_\eta|}{XE}
 +\frac{4A|\pi|+|\pi_\eta|}{E}.
\]
It explicitly demonstrates why the two A.11 equations alone
do not certify $|N_s|/E$.
If $\sigma=\pi=0$ and differentiation under the convergent
integrals is justified, then the four tail quantities combine
into one source-exact integral:
\begin{equation}\label{eq:A11-single-future-gate}
 \frac{N_s}{E(X,\eta)}
 =\frac1{E(X,\eta)}\int_X^\infty E(x,\eta)^2
 \left[\frac{2h\eta}{X}-\frac{2A\eta}{x}
 +(1-\eta^2)\partial_\eta\log E(x,\eta)
 \left(\frac1x-\frac1X\right)\right]dx .
\end{equation}
Define $D_s=S_{\rm tail}/(XE^2)$ and
$P_s=-\Pi_{\rm tail}/E^2$. Then $0\le P_s\le D_s$.
If the full future logarithmic slope $\ell\le0$, also
$P_s\le1/2$. If $|\partial_\eta\log E|\le\varepsilon$
at every future radius, the integral gives
\[
 \frac{|N_s|}{E}\le E\left[
 4|\eta|(hD_s+AP_s)
 +2(1-\eta^2)\varepsilon(D_s-P_s)\right].
\]
Thus the jointly certified sufficient source inputs
$0<\lambda\le1/20$, $0<h\le\lambda/100$,
$E\le\lambda/4$, $D_s\le233/100$,
$\varepsilon\le1/1000$, $\ell\le0$, and
$\sigma=\pi=0$ yield the exact-rational implication
\begin{equation}\label{eq:A11-terminal-strict-slack}
 |N_s|/E<253\lambda/1000<\lambda/2 .
\end{equation}
For imperfect gauges, the same result holds if the
two displayed gauge-error ratios sum to less than
$247\lambda/1000$. These are conditional inequalities,
not independently source-authenticated enclosures.

To see the unequal conditioning of the two terminal defects,
write $X=X_pe^y$, $f=(1+\eta^2)^{-1}$,
$E=e_b f e^{-(1/2+\lambda)y}G$ with $G>0$,
$\sigma=X_pe_b^2f^2r/\lambda$, and $\pi=e_b^2f^2p$.
The corresponding normalized corrections to $N_s/E$ are
\[
 \frac{e_bf}{\lambda G}e^{-(1/2-\lambda)y}
 \left[4h\eta r-d\left(2(f_\eta/f)r+r_\eta\right)\right],
\quad
 \frac{e_bf}{G}e^{(1/2+\lambda)y}
 \left[4A\eta p-d\left(2(f_\eta/f)p+p_\eta\right)\right].
\]
The $S$-defect coefficient decays with radius;
the pressure-gauge coefficient grows. This establishes
why the exact terminal roots or authenticated derivative
enclosures are required before applying the stress gate.
It does not close the $Q_s$, complete exterior cone,
Appendix B, or proposed unforced Theorem~1.1.

\subsection*{The A.11 terminal pressure gauge is rigid}
The exact stress identity (4.16) on the two A.11 supports, where
$U=M=J=0$, is
\[
N_s=\frac{4h\eta S-(1-\eta^2)S_\eta}{X}
+4A\eta\Pi-(1-\eta^2)\Pi_\eta,\qquad A=\tfrac12+h.
\]
Writing $\sigma=S(\infty,\eta)$ and $\pi=\Pi(\infty,\eta)$
and substituting the future-tail integrals gives the exact defect
decomposition
\[
N_s=N_s^{\mathrm{tail}}+
\frac{4h\eta\sigma-(1-\eta^2)\sigma_\eta}{X}
+4A\eta\pi-(1-\eta^2)\pi_\eta.
\]
The A.11 edit restores only a pressure \emph{increment}.
The actual source normalization (A.21)--(A.23) must be imposed
for the \emph{same} corrected profile to obtain $\pi=0$.
Independently, the genuine A.19 amplitude equation sets $\sigma=0$.

\begin{proposition}[Pure-power exterior pressure rigidity]
Assume the exact terminal profile is
$E(X,\eta)=c_\infty X^{-A}$, where $c_\infty>0$ is independent
of $\eta$, and $\sigma,\pi\in C^1([-1,1])$.
If $N_s/E$ is bounded as $X\to\infty$ for each
$\eta\in(-1,1)$, then $\pi\equiv0$.
\end{proposition}
\begin{proof}
On the pure tail the future integrals are
$S_t=XE^2/(4h)$ and
$\Pi_t=-E^2/[2(1+2h)]$. Their contribution to
$N_s$ cancels exactly. Hence
\[
\frac{N_s}{E}=
\frac{X^{A-1}}{c_\infty}
[4h\eta\sigma-(1-\eta^2)\sigma_\eta]
+\frac{X^A}{c_\infty}
[4A\eta\pi-(1-\eta^2)\pi_\eta].
\]
Since $A>0$ and $A<1$, boundedness forces
$(1-\eta^2)\pi_\eta=4A\eta\pi$.
Thus $(1-\eta^2)^{2A}\pi$ is constant and, by
continuity at $\eta=\pm1$, the constant vanishes.
\end{proof}

\subsection*{An exact A.19 total-residual stress certificate}
Let $f=(1+\eta^2)^{-1}$ and define the complete A.19
residual along the selected smooth amplitude by
\[
\mathcal F(\eta)=
\frac{\lambda S(\infty,\eta)}{X_p e_b^2 f(\eta)^2}.
\]
Here $\mathcal F_\eta$ is the \emph{total} parameter derivative,
including the derivative of the chosen amplitude.
Under the exact pressure gauge $\pi=0$,
\[
\frac{N_s-N_s^{\mathrm{tail}}}{E}=
K(X,\eta)\left[
\left(4h\eta+\frac{4\eta(1-\eta^2)}{1+\eta^2}\right)
\mathcal F-(1-\eta^2)\mathcal F_\eta\right],
\quad
K=\frac{X_p e_b^2 f^2}{\lambda XE}.
\]
The elementary factorization
$(1+t^2)-2t(1-t^2)=(1-t)^2+2t^3\ge0$
for $0\le t\le1$ gives
$|4\eta(1-\eta^2)/(1+\eta^2)|\le2$.
If an independently certified future-tail bound on
both actual A.11 supports is
$\sup|N_s^{\mathrm{tail}}|/E<253\lambda/1000$,
then the strict sufficient certificate
\[
K_s[(2+4h)\|\mathcal F\|_\infty+
\|\mathcal F_\eta\|_\infty]<247\lambda/1000
\]
implies $\sup|N_s|/E<\lambda/2$.
The old value-only residual check does not suffice.

For $t$ measured from the end of (A.10)'s profile
interpolation, the exact source support has
$X=X_p\exp(13/\lambda+T_f+t)$ and
$E=(e_b/2)\exp[-(1/2+\lambda)(13/\lambda+T_f+t)]B$,
where $B=1+\sum_i c_i(\eta)\beta_i(t)>0$.
Thus
\[
K=\frac{2e_b f^2}{\lambda B}
\exp[-(\tfrac12-\lambda)(13/\lambda+T_f+t)].
\]
Both support intervals satisfy
$t\ge30\log(1/\lambda)-63/20$.
Certified $e_b\le e_*$, $B\ge b_*>0$ imply the explicit
$K_s$ upper bound obtained by replacing $t$ by this
lower endpoint. The actual enclosures remain to be supplied.

\subsection*{Required normalized-tail derivative}
For $D_s=S_t/(XE^2)$, $P_s=-\Pi_t/E^2$, and
$g=\partial_\eta\log E$, the exact local form is
\[
\frac{N_s^{\mathrm{tail}}}{E}
=E[4\eta(hD_s-AP_s)
 -(1-\eta^2)\{2g(D_s-P_s)
 +(D_s)_\eta-(P_s)_\eta\}].
\]
The derivatives $(D_s)_\eta-(P_s)_\eta$ do not vanish
for an arbitrary A.11 correction.
The full future-integral kernel in the previous
A.11 terminal-gauge manuscript remains valid, but a
symbolic replay treating $D_s,P_s$ as parameter constants
is only a special-case check.

\paragraph{Scope.}
The identities above are exact. They do not verify the
A.19 root, the two support/gap interval covers, the
whole future-tail slope and parameter jets, the exterior cone,
Appendix B, the unforced startup recurrence or a full Lean
proof of Theorem 1.1. The originating OpenAI theorem has
smooth forcing.


\subsection*{A.19 exact-root differentiation and a full-hull A.11 future-tail bound}
\label{subsec:A11-root-whole-hull-bridge-20261021}
The current stress proof must distinguish the \emph{total} parameter
derivative from the partial derivative at fixed pulse amplitude.
Write the full source A.19 moment residual, including its genuinely
corrected A.11 azimuthal coefficients, as
\[
 \mathcal F(A,\eta)=\frac{\lambda S(\infty;A,\eta)}
 {X_pe_b^2(1+\eta^2)^{-2}}.
\]
If the exact source moment equation has a $C^1$ root $A_*(\eta)$,
then $\mathcal F(A_*(\eta),\eta)\equiv0$, and hence
\[
 \mathcal F_A(A_*,\eta)A_*'(\eta)
 +\mathcal F_\eta(A_*,\eta)=0.
\]
Thus the total derivative appearing in the stress formula vanishes
\emph{identically}. This uses an exact root, not merely a uniform
small value: $10^{-5}\cos(10^5\eta)$ has supremum at most $10^{-5}$
and derivative supremum one. The source's positive quadratic
amplitude coefficient and its strictly positive derivative margin,
once established for one actual common corrected profile, allow
the implicit function theorem to produce the required $C^1$ branch.

The separate genuine source pressure normalization
$\Pi(\infty,\eta)=0$ must hold for the \emph{same} profile. Under
these two normalizations the terminal-defect part of (4.16)
vanishes, leaving $N_s=N_s^{\rm tail}$. Let
$E_c(y,\eta)=E_b(y)v(y,\eta)$, where $E_b$ is independent of $\eta$
after A.10 and $v=1+u$,
$u=\sum_{i=1}^2 c_i(\eta)\beta_i(y)$.
The two source $\sigma'$ corrections are disjoint, each of
logarithmic support length $3/10$, with centers separated by two.
Their \emph{closed convex hull} contains an intervening gap of
length $17/10$; its total length is $23/10$ while the union
of nonzero supports has length $3/5$.

Define $R_b(y,t)=\exp(2\int_0^t\ell_b(y+s)\,ds)$,
$D_c=S_{\rm tail}/(XE_c^2)$ and $P_c=-\Pi_{\rm tail}/E_c^2$.
The future integrals are exactly
\[
D_c=\frac{\int_0^\infty R_b(y,t)v(y+t)^2\,dt}{2v(y)^2},
\qquad
P_c=\frac{\int_0^\infty e^{-t}R_b(y,t)v(y+t)^2\,dt}
{2v(y)^2}.
\]
The exact normalization cancellation is
\[
 \mathcal T:=2g(D_c-P_c)+(D_c)_\eta-(P_c)_\eta
 =\frac{\int_0^\infty R_b(y,t)(1-e^{-t})
 v(y+t)u_\eta(y+t)\,dt}{v(y)^2},
\]
where $g=\partial_\eta\log E_c$. This identity holds equally
within either corrected support and on their intervening gap;
differentiation under the future integrals requires a
source-justified common dominating envelope.

\begin{theorem}[Conditional complete A.11 hull stress certificate]
\label{thm:A11-root-whole-hull-20261021}
Assume one corrected source profile has the exact $C^1$ A.19
root and the prescribed terminal pressure gauge, $0<\lambda\le1/20$,
$0<h\le\lambda/100$, baseline $\ell_b=-\lambda$ on the
whole hull and $\ell_b\le0$ at all future radii, the source
release/steep-hold/final-tail bound
$D_b\le93/40+h^7/3$ at both hull endpoints, the
authentic coefficient jets
$|c_i|\le54\lambda^{28}$,
$|c_{i,\eta}|\le6\lambda^{28}$,
$|\beta_i|\le640/3$, and $E_c\le\lambda/4$
throughout the hull. Then
\[
 \sup_{\substack{y\text{ in both A.11 supports and their gap}\\|\eta|\le1}}
 \frac{|N_s|}{E_c}
 <\frac{503}{2000}\lambda<\frac\lambda2.
\]
In particular, the strict stress slack is greater than
$497\lambda/2000$.
\end{theorem}
\begin{proof}
Disjointness gives $|u|\le\delta=11520\lambda^{28}$
and $|u_\eta|\le1280\lambda^{28}$. Since $R_b\le1$ and
the future support of $u_\eta$ has length at most $3/5$,
\[
 |\mathcal T|\le\frac35(1280\lambda^{28})
 \frac{1+\delta}{(1-\delta)^2}<10^{-5}.
\]
On the baseline plateau, $D_b'=2\lambda D_b-1/2$ and
$P_b'=(1+2\lambda)P_b-1/2$, so both are monotone or
constant on the \emph{whole} hull, without an unchecked
interior extremum in the gap. Positive-tail comparison yields
\[
D_c<233/100,\qquad
P_c\le\tfrac12[(1+\delta)/(1-\delta)]^2
 <5000001/10000000.
\]
Here $\ell_b\le0$ globally supplies $P_b\le1/2$.
The exact source stress formula then yields
\[
 \frac{|N_s|}{E_c}
 \le E_c\{4hD_c+4(\tfrac12+h)P_c+|\mathcal T|\}
 <\frac\lambda4\frac{503}{500}.
\]
All numerical comparisons are strict rational inequalities
at the source-worst value $\lambda=1/20$.
\end{proof}

This supplies a stronger \emph{logical reduction}, not
an independently authenticated interval realization.
The same-parameter A.19 root and infinity gauge, coefficient
jets, whole future baseline slopes and complete other stress
stages remain to be certified. The accompanying Lean module
proves only scalar/derivative consequences, not the source
integrals or the proposed unforced Navier--Stokes Theorem~1.1.


\subsection*{Independent future-integral comparison on the true A.11 supports}
\label{subsec:a11-direct-future-audit}
This is a second derivation of the existing conditional A.11 stress
gate, useful as an independent check. The stronger full-hull
estimate of the preceding section remains the primary bound.
Let $E=E_0g$, where $E_0$ is the actual parameter-independent
unedited azimuthal profile after interpolation and
$g=1+c_1(\eta)\beta_1(y)+c_2(\eta)\beta_2(y)$.
Set $A=1/2+h$, $d=1-\eta^2$ and
$F_0=\int_y^\infty E(v,\eta)^2dv$,
$F_1=\int_y^\infty e^{v-y}E(v,\eta)^2dv$.
If the same corrected source profile has
$U=M=J=0$, $S(\infty)=0$ and $\Pi(\infty)=0$,
then (4.16) gives
\[
N_s=2\eta(hF_1-AF_0)-\tfrac d2\partial_\eta(F_1-F_0).
\]
The second term is generally nonzero: the A.11 coefficients
$c_i(\eta)$ can depend on $\eta$. It equals
\[
-d\int_y^\infty(e^{v-y}-1)E_0(v)^2
g(v,\eta)\partial_\eta g(v,\eta)\,dv.
\]
For the actual rescaled fixed $\sigma'$ bump, the elementary
bounds $\|\beta\|_\infty\le B_0=640/3$ and
$\|\beta'\|_\infty\le B_1=1715200/9$ are valid.
Given certified $|c_i|\le C_0$, $|c_i'|\le C_\eta$,
and $E_0\le E_*$ on the two supports, put
$q_\pm=1\pm B_0C_0$, $\delta_\eta=B_0C_\eta$
and $D_*=93/20+2h^7/3$. If $q_->0$ then
\[
\frac{|N_s|}{E}\le\frac{E_*}{q_-}
 [\,2q_+^2(A+hD_*)+q_+\delta_\eta D_*\,],
\qquad
\ell_{\rm corrected}\le-\lambda+\frac{B_1 C_0}{q_-}.
\]
The future-integral comparison uses only source stages:
at most $63/20$ units remain on the plateau, then one
transition, the $-1$ hold contributing at most $1/2$,
and the exponentially suppressed tail at most $2h^7/3$.
Thus the two right sides being $\le\lambda/2$ proves both
the requested $N_s/E$ inequality and the exterior slope
condition on BOTH supports.

This is a weaker, independently derived sufficient certificate,
not a new source-authenticated interval cover. An exact Fraction
test at $\lambda=1/20$, $h=10^{-6}$, $E_*=10^{-6}$,
$C_0=C_\eta=10^{-8}$ passes; these input numbers are an explicitly
NON-SOURCE fixture. Neither that test nor the auxiliary scalar
Lean lemmas establishes a true A.11 moment-correction root,
the terminal gauges, or the proposed unforced Theorem~1.1.


\subsection*{A.11 source-root and angular-derivative certificate}
\label{subsec:a11-rootjet-new-20261021}
For the prescribed A.11 smooth-step derivative bumps, the true
normalized correction equations are
\[
 u(c_1+xc_2)=\Delta_I(\eta),\qquad
 2v(c_1+zc_2)+w(c_1^2+zc_2^2)=0,
\]
where $u,v,w>0$, $x=e^{2(1-\lambda)}$,
$z=e^{-2(1+2\lambda)}$ are independent of $\eta$.
Earlier source-specific bounds give
$\|B^{-1}\|_\infty<3/2$ for
$B=\left(\begin{smallmatrix}u&ux\\2v&2vz\end{smallmatrix}\right)$
and $\|Q(c)\|_\infty\le320\|c\|_\infty^2$ for
$Q(c)=(0,w(c_1^2+zc_2^2))$.

\begin{lemma}[Authentic moment-mismatch jet]
\label{lem:a11-rootjet-delta-new}
Assume the source A.10 interpolation $T_f=512$, the
pre-interpolation factorization $E(y,\eta)=c(y)(1+\eta^2)^{-1}$
and the succeeding $-\lambda$ slope plateau of length
$30\log(1/\lambda)$. For $0<\lambda\le1/100$ the
moment mismatch at the first A.11 bump center obeys
\[
 |\Delta_I|<\frac{63}{100}\lambda^{28},\qquad
 |\partial_\eta\Delta_I|<\frac14\lambda^{28}.
\]
\end{lemma}
\begin{proof}
The angular moment ratio $r=I/(XH)$ solves
$r'+(1+\ell)r=1$. It is independent of $\eta$
before A.10, hence $r_\eta=0$ upon interpolation entry.
On A.10, $\ell=-\lambda+\vartheta_f'(\log2-J_0)$,
$J_0=\log(1+\eta^2)$ and
$\vartheta_f=1-\sigma(y/T_f)$, so
$|\ell_\eta|\le|\vartheta_f'|$ and
$\int|\vartheta_f'|=1$.
The source stage bounds give $0\le r\le20/17$ and
$1+\ell\ge17/20$; differentiating the exact ODE shows
$|r_\eta(T_f)|\le20/17$.
The subsequent constant-slope interval has length
$30\log(1/\lambda)-3$ before the first bump.
Both the mismatch from its fixed point $(1-\lambda)^{-1}$
(initially of magnitude below three) and $r_\eta$
are multiplied by
$\exp[-(1-\lambda)(30\log(1/\lambda)-3)]$.
Since $30(1-\lambda)\ge29$, $e^3<21$
and $\lambda\le1/100$, one gets
$|\Delta_I|<63\lambda^{29}\le(63/100)\lambda^{28}$
and $|\Delta_{I,\eta}|<(420/17)\lambda^{29}
<(1/4)\lambda^{28}$.
\end{proof}

\begin{theorem}[Small corrected root and its first parameter jet]
\label{thm:a11-rootjet-new}
Under the lemma, the \emph{actual} nonlinear A.11 equations
have a unique solution in $\|c\|_\infty\le2\lambda^{28}$,
depending $C^1$ on $\eta$, with
\[
 \|c\|_\infty<2\lambda^{28},\qquad
 \|\partial_\eta c\|_\infty<\lambda^{28}.
\]
\end{theorem}
\begin{proof}
On the ball $R_c=2\lambda^{28}$ the map
$c\mapsto B^{-1}((\Delta_I,0)-Q(c))$ has norm at most
$\bigl((3/2)(63/100)+1920\lambda^{28}\bigr)\lambda^{28}
<R_c$ and Lipschitz constant at most
$1920\lambda^{28}<1/100$. It has a unique small fixed
point. The differentiated \emph{source} equations give
$[B+DQ(c)]c_\eta=(\Delta_{I,\eta},0)$.
The inverse norm is at most
$(3/2)/(1-1/100)$, and therefore
$\|c_\eta\|_\infty\le(25/66)\lambda^{28}
<\lambda^{28}$.
\end{proof}

Because the two source bumps are disjoint and
$\|\beta_i\|_\infty\le640/3$, this gives on their
closed convex hull (including the gap)
\[
 \left|\frac{E}{E_b}-1\right|
 \le\frac{1280}{3}\lambda^{28},\qquad
 \left|\partial_\eta(E/E_b)\right|
 \le\frac{640}{3}\lambda^{28}.
\]
It strengthens the previously assumed correction and
parameter-jet constants in the \emph{conditional}
full-hull $|N_s|/E$ theorem. It does not prove the
same profile's exact $S(\infty)=\Pi(\infty)=0$,
the future-tail integral enclosure, $Q_s$ positivity,
all exterior stress-cone stages, the Appendix B fixed
point, or the proposed unforced Theorem~1.1.

The parameter-derivative argument is essential:
$\Delta_I(\eta)=\epsilon\sin(N\eta)$ has a small
uniform value but derivative as large as $N\epsilon$.
The integral/ODE argument above, not the
small-coefficient estimate alone, supplies the
actual source's derivative bound.


\subsection*{A robust rational interval-cell certificate for the complete A.11 hull}
\label{subsec:a11-robust-cell-20261021}
The latest scalar stress and full-hull lemmas still need an authenticated
interval covering of both correction supports and the intervening gap.
We quantify how much independent interval or mesh error they can tolerate.

\begin{theorem}[Error-tolerant strict whole-hull cone]
\label{thm:a11-robust-hull-20261021}
Fix $0<\lambda\le1/20$, $0<h\le\lambda/100$,
$X\ge9/\lambda$, $|\eta|\le1$, and $L=1-2h\eta^2$.
Write $f=N_s/E$, $q=Q_s$ and
$\ell=\partial_y\log(\sqrt X E)$.
Suppose at every point of both prescribed A.11 supports and
their intervening gap there exist reference fields $f_0,q_0,\ell_0$
obeying
\[
 |f_0|\le503\lambda/2000,\qquad q_0\ge49\lambda/50,\qquad
 -3\lambda/2\le\ell_0\le-\lambda/2
\]
and independently certified error bounds
\[
 |f-f_0|\le\lambda/100,\quad
 |q-q_0|\le\lambda/100,\quad
 |\ell-\ell_0|\le\lambda/100.
\]
Then everywhere on the full hull,
\[
 |N_s|/E\le523\lambda/2000<\lambda/2,\qquad
 \lambda/2-|N_s|/E\ge477\lambda/2000.
\]
Moreover, with $a=2-2\ell$, $p_1=Xq/L$,
$w=f/q$, the original normalized cone obeys the
exact strict lower bound
\[
 2(1-a/p_1)^2-(a-2)w^2
 \ge\frac{4192747121}{3763600000}>\frac{11}{10}>\frac9{20}.
\]
\end{theorem}
\begin{proof}
The hypotheses imply
$|f|\le523\lambda/2000$, $q\ge97\lambda/100$
and $-151\lambda/100\le\ell\le-49\lambda/100$.
Thus $0<L\le1$, $p_1\ge873/100$,
$2<a\le2151/1000$, $a/p_1\le239/970<1/4$,
$|w|\le523/1940$ and $a-2\le151/1000$.
Consequently
\[
2(1-a/p_1)^2-(a-2)w^2
\ge2(3/4)^2-\frac{151}{1000}
 \left(\frac{523}{1940}\right)^2
=\frac{4192747121}{3763600000}.
\]
The stress margin is immediate by subtraction. All comparisons
are exact rational inequalities.
\end{proof}

\begin{proposition}[Certified Lipschitz mesh sufficient condition]
\label{prop:a11-full-hull-lipschitz-20261021}
Partition the closed logarithmic A.11 hull of length $23/10$
and the complete parameter interval $[-1,1]$ into finitely many
closed rectangles with disjoint interiors.
For each field $g\in\{f,q,\ell\}$ let $D_{g,y},D_{g,\eta}$
be independently verified bounds on its first derivatives
over each rectangle, and let $\Delta y,\Delta\eta$ be its widths.
Suppose the reference inequalities of the theorem are verified
at its centre, while
\[
 (D_{g,y}\Delta y+D_{g,\eta}\Delta\eta)/2
 \le\lambda/100\quad(g=f,q,\ell).
\]
If also $X\ge9/\lambda$, the theorem holds throughout the
complete hull, including every point in the gap.
\end{proposition}
\begin{proof}
Integrate each partial derivative along the two coordinate
segments from a cell centre to any point of its rectangle.
The errors are bounded by the indicated expression. A finite
closed covering transfers the pointwise rational theorem to
all points and shared boundaries.
\end{proof}

\paragraph{Necessary source-authentication boundary.}
The local (A.11) moment constraints do not by themselves fix
the terminal data $S(\infty,\eta)$ and $\Pi(\infty,\eta)$.
From (4.16), a change of integration data produces
\[
 \Delta N_s=
 \frac{4h\eta C_S-(1-\eta^2)C'_S}{X}
 +4(\tfrac12+h)\eta C_\Pi-(1-\eta^2)C'_\Pi.
\]
At $\eta=0$, $C_S=K\eta$ changes $N_s$ by $-K/X$.
Hence a verified source-root, the same corrected profile's
terminal pressure gauge, full interval/derivative enclosures
and all pulse/exterior stages are indispensable.
The accompanying rational checker deliberately outputs only
CONDITIONAL\_SCALAR\_PASS or FAIL\_CLOSED, never a claim
of a source-authenticated or full Lean certificate.
The original OpenAI theorem constructs a forced whole-space
solution. Appendix-B contraction, unforced $n=1,2$ startups,
all-order convergence and physical residual zero are still open.


\subsection*{Consolidated independent source audit and error-tolerant A.11 hull}
\label{subsec:a11-consolidated-independent-20261021}
A separate 21 October 2026 source audit of the two prescribed
(A.11) corrections starts from $E=E_b(1+c_1(\eta)\beta_1+
c_2(\eta)\beta_2)$ and the *same-profile* gauges
$U=M=J=0$, $S(\infty,\eta)=\Pi(\infty,\eta)=0$.
With $A=1/2+h$ and
\[
 F_0(y,\eta)=\int_y^\infty E(v,\eta)^2\,dv,\qquad
 F_1(y,\eta)=\int_y^\infty e^{v-y}E(v,\eta)^2\,dv ,
\]
its exact independent stress reduction is
\[
N_s=2\eta(hF_1-AF_0)
-\frac{1-\eta^2}{2}\,\partial_\eta(F_1-F_0).
\]
The second term equals
$-(1-\eta^2)\int_y^\infty(e^{v-y}-1)
E_b(v)^2g(v,\eta)\partial_\eta g(v,\eta)\,dv$.
It must be retained when the correction coefficients depend
on the angular variable. The prescribed sigma-prime bump has
the explicit conservative envelopes
$\|\beta_i\|_\infty\le640/3$ and
$\|\beta_i'\|_\infty\le1715200/9$.
These yield a separate \emph{two-support} sufficient estimate
in terms of authenticated $|c_i|$, $|c_{i,\eta}|$ and the
unedited $E_b$ envelope. It does not by itself cover the gap.

The logically stronger hull robustness theorem,
Theorem~\ref{thm:a11-robust-hull-20261021},
covers both supports and the $17/10$ intervening gap.
For $0<\lambda\le1/20$, $0<h\le\lambda/100$ and
$X\ge9/\lambda$, suppose reference bounds
$|f_0|\le503\lambda/2000$, $q_0\ge49\lambda/50$,
$-3\lambda/2\le\ell_0\le-\lambda/2$ have been authenticated
from the corrected source, and the complete actual errors in
$f=N_s/E$, $q=Q_s$ and $\ell$ are each at most $\lambda/100$.
Then the strict stress slack and the original normalized cone obey
\[
\frac{|N_s|}{E}\le\frac{523}{2000}\lambda<\frac\lambda2,
\qquad
\mathcal C\ge\frac{4192747121}{3763600000}>
\frac{11}{10}>\frac9{20}.
\]
The exact rational calculation is reproducible with
\texttt{a11\_robust\_hull\_exact.py}; finite arithmetic alone
does not authenticate the input source profile.

\paragraph{The single-source proof obligation.}
The independent support estimate, the gap-covering robustness
theorem and the more recent A.11 root-jet results must all refer
to one corrected source with the same $(T_f,\lambda,h)$, the same
A.19 amplitude root and infinity gauges, and independently
verified pointwise/derivative intervals. A valid release therefore
requires full closed-cell coverage of the entire $23/10$ hull,
the exterior pulse/transition/terminal stress cones, and
the same-profile $S(\infty)=\Pi(\infty)=0$.
A missing pressure gauge changes $N_s$ by
$4(\frac12+h)\eta C_\Pi-(1-\eta^2)C_\Pi'$, and an unfixed
$S$ datum contributes
$[4h\eta C_S-(1-\eta^2)C_S']/X$.
These defects cannot be hidden in a $\lambda/100$ allowance.
The present implication is proved conditionally; the
authentic full exterior, Appendix-B contraction, unforced
startup hierarchy, force-free Cartesian reconstruction, and a
complete kernel-checked Lean theorem remain unproved.

\subsection*{Backward angular-moment certification from the A.11 endpoint}
\label{subsec:a11-backward-angular-endpoint-20261021}
The source's original stress expression on the corrected azimuthal hull is
\[
Q_s=-1+(1-h)\frac{I}{XH}-D\eta\frac{I_\eta}{XH},
\qquad D=\tfrac12-h.
\]
The last term is essential within a correction support: the moment
and correction coefficients need not be independent of \(\eta\).

On the full closed A.11 hull, of logarithmic length \(23/10\), set
\(E=E_b(1+u)\), \(u=\sum_i c_i(\eta)\beta_i(y)\).
The unedited profile is independent of \(\eta\) there, with
\(\ell_b=-\lambda\). The two smooth-step-derivative bumps are
disjoint and satisfy
\(\|\beta_i\|_\infty\le640/3\) and
\(\|\beta_i'\|_\infty\le1715200/9\).
Assume the authentic root-jet inputs already derived for
\(T_f=512\), \(0<\lambda\le1/100\):
\[
|c_i|\le54\lambda^{28},\qquad
|c_{i,\eta}|\le6\lambda^{28},\qquad
0<h\le\lambda/100.
\]
Let \(y_R\) be the right edge of the second support. The
source's exact A.11 correction makes \(I/(XH)=(1-\lambda)^{-1}\)
at the end of the corrected plateau. From \(y_R\) to that
endpoint there is no edit and \(\ell=-\lambda\). Uniqueness for
\(r_y+(1-\lambda)r=1\), \(r=I/(XH)\), therefore gives the
*backward terminal data*
\[
r(y_R,\eta)=\frac1{1-\lambda},\qquad
r_\eta(y_R,\eta)=0.
\]
This is stronger than merely bounding the entrance moment.

\begin{theorem}[Whole-hull \(Q_s\) without a pointwise mesh]
\label{thm:a11-backward-angular-20261021}
Under the preceding authentic coefficient-jet and terminal-moment
hypotheses, the following hold throughout both correction supports
and the \(17/10\) intervening gap:
\[
\left|r-\frac1{1-\lambda}\right|<6\cdot10^8\lambda^{28},
\quad |r_\eta|<13\cdot10^7\lambda^{28},
\quad \left|\frac{I_\eta}{XH}\right|<131\cdot10^6\lambda^{28},
\]
\[
-\frac{3\lambda}{2}<\ell<-\frac\lambda2,\qquad
Q_s\ge\frac{\lambda-h}{1-\lambda}
-7\cdot10^8\lambda^{28}>\frac{49}{50}\lambda.
\]
\end{theorem}
\begin{proof}
Disjointness and the source bounds imply
\[
|u|\le11520\lambda^{28}<1/100,\quad
|u_y|\le10291200\lambda^{28},\quad
|u_\eta|\le1280\lambda^{28},\quad
|u_{y\eta}|\le(10291200/9)\lambda^{28}.
\]
Because \(\ell+\lambda=u_y/(1+u)\),
\(g:=\partial_\eta\log E=u_\eta/(1+u)\), and
\(\ell_\eta=u_{y\eta}/(1+u)-u_yu_\eta/(1+u)^2\),
exact rational comparisons at \(\lambda=1/100\) give
\[
|\ell+\lambda|<11\cdot10^6\lambda^{28},\quad
|g|<1300\lambda^{28},\quad
|\ell_\eta|<13\cdot10^5\lambda^{28}.
\]
The moment ratio satisfies \(r_y+(1+\ell)r=1\).
Writing \(r_*=(1-\lambda)^{-1}\), integrate this equation
backward from \(r(y_R)=r_*\):
\[
r(y)-r_*
=r_*\int_y^{y_R}\exp\left(\int_y^t(1+\ell(v))\,dv\right)
(\ell(t)+\lambda)\,dt.
\]
The hull length \(23/10\) and \(e^3<21\) imply
\[
|r-r_*|<\frac{100}{99}\frac{23}{10}
21\cdot11\cdot10^6\lambda^{28}
<6\cdot10^8\lambda^{28}.
\]
In particular \(0<r<2\). Differentiating the transport
equation and integrating backward from \(r_\eta(y_R)=0\)
gives
\[
|r_\eta|<2\frac{23}{10}21\cdot13\cdot10^5\lambda^{28}
<13\cdot10^7\lambda^{28}.
\]
Since \(I_\eta/(XH)=r_\eta+rg\), its bound follows.
Substitution into the *full* source formula for \(Q_s\)
gives an error below \(7\cdot10^8\lambda^{28}\).
Finally \(700000000(1/100)^{27}<1/100\), and
\(h\le\lambda/100\). The slope interval follows from
the first derivative bound.
\end{proof}

\begin{corollary}[Strict cone from a separate \(N_s/E\) tail gate]
\label{cor:a11-backward-cone-20261021}
In addition assume *for this same corrected profile* that
the exact \(S(\infty)=\Pi(\infty)=0\) gauges and authenticated
future-integral bounds give \(|N_s|/E\le503\lambda/2000\),
and \(X\ge9/\lambda\) on the hull. Then, with
\(a=2-2\ell\), \(p_1=XQ_s/(1-2h\eta^2)\) and
\(w=N_s/(EQ_s)\),
\[
|N_s|/E<\lambda/2,\qquad
2(1-a/p_1)^2-(a-2)w^2\ge
\frac{85676973}{76832000}>\frac{11}{10}>\frac9{20}.
\]
\end{corollary}
\begin{proof}
The theorem gives \(p_1\ge441/50\),
\(a\le43/20\), \(a/p_1<1/4\),
\(|w|\le503/1960\), and \(a-2\le3/20\).
The displayed value is the exact rational consequence.
\end{proof}

The advance removes the need to separately interval-sample
\(Q_s\) or the corrected slope on the A.11 hull,
*conditional on the same authentic root and terminal moment
gauge*. It does **not** authenticate the A.19 \(S\)-root,
the terminal pressure gauge, the future \(N_s\)-tail,
the full exterior stage covering, Appendix B's nonlinear
contraction, the unforced startup hierarchy, or Theorem~1.1.
The companion Python checker verifies exact rational consequences;
the new Lean module contains only scalar gates and does not
formalize these source integral/ODE/PDE arguments.

\subsection*{Source-local stress without a separate future-tail input (22 October 2026)}
Let \(y\) denote the source's logarithmic radial coordinate and let the unedited azimuthal profile \(E_0(y)\) be independent of \(\eta\) on and after the A.11 hull of length \(23/10\). The actual two-bump correction has
\(E=E_0g\), \(g=1+\sum_{i=1}^2 c_i(\eta)\beta_i(y)\).
Since the prescribed supports are disjoint and
\(\|\beta_i\|_\infty\le640/3\), put
\(a=(640/3)\max_i\|c_i\|_\infty\) and
\(b=(640/3)\max_i\|\partial_\eta c_i\|_\infty\).
Then \(|g-1|\le a\), \(|g_\eta|\le b\), including on the intervening gap.

Under the \emph{exact same-profile} terminal gauges
\(S(\infty,\eta)=\Pi(\infty,\eta)=0\), source (4.16) reads
\[
N_s=2\eta(hF_1-AF_0)-\tfrac{1-\eta^2}{2}
\partial_\eta(F_1-F_0),\quad A=\tfrac12+h,
\]
with \(F_0=\int_y^\infty E(v)^2\,dv\) and
\(F_1=\int_y^\infty e^{v-y}E(v)^2\,dv\).
The prescribed future slopes give
\(\ell_0=\partial_y\log(\sqrt X E_0)\le0\) and
\(F_{1,0}/E_0(y)^2\le93/20+2h^7/3<5\);
also \(F_{0,0}/E_0(y)^2\le1\).
Crucially \(E_0\) is independent of \(\eta\) on the entire future;
thus \(\partial_\eta(F_1-F_0)\) is supported only between
the current point and the right A.11 support edge, and
\[
|\partial_\eta(F_1-F_0)|
\le2(1+a)b E_0(y)^2
\int_0^{23/10}(1-e^{-t})dt
\le(46/10)(1+a)bE_0(y)^2.
\]
Consequently the complete closed hull satisfies
\begin{equation}\label{eq:A11new-tail}
\sup_{y,\eta}\frac{|N_s|}{E}
\le\frac{E_*}{1-a}
\left[2(1+a)^2(A+5h)+(23/10)(1+a)b\right],
\quad E_*=\sup_{\mathrm{hull}}E_0.
\end{equation}
This retains the angular derivative, and requires no isolated
numerical future-tail \(N_s/E\) input.

The source stage order yields
\(E_*\le e^{47/40}P_*e^{-T_d/2}\lambda^{45}
e^{-13/(2\lambda)}
<8P_*e^{-T_d/2}\lambda^{45}e^{-13/(2\lambda)}\).
After choosing the already source-permitted small parameter
\[
0<\lambda\le\min\{1/100,[100(1+\log P_*)]^{-1}\},
\qquad 0<h\le\lambda/100,
\]
this implies \(E_*<\lambda/1000\). If the authentic A.11 root and
parameter derivative have been independently verified to obey
\(|c_i|\le2\lambda^{28}\),
\(|c_{i,\eta}|\le\lambda^{28}\),
then \(a,b<1/1000\) and
\[
\frac{2(1001/1000)^2(2503/5000)
 +(23/10)(1001/1000)(1/1000)}{999/1000}
=\frac{2513764253}{2497500000}<\frac{11}{10}.
\]
Equation \eqref{eq:A11new-tail} therefore proves
\[
\sup_{\mathrm{entire\ A.11\ hull},|\eta|\le1}|N_s|/E
<11\lambda/10000<\lambda/500<\lambda/2.
\]

\paragraph{Dependency boundary.}
This improves the \emph{conditional} whole-hull stress theorem,
but the authentic moment root and exact same-profile infinity
\(S,\Pi\) gauges remain prerequisites. It does not certify all
exterior stress patches, Appendix-B nonlinear contraction,
the unforced similarity-order startup equations or the physical
PDE. The originating OpenAI whole-space construction has forcing.
The complete proof with definitions, integration and source
references is in the standalone 22 October addendum.

\subsection*{An explicit A.11 profile-size certificate from source stages}
\label{subsec:a11-explicit-profile-size-20261022}
The conditional whole-hull stress theorem previously required
$E_c\le\lambda/4$ as an independently supplied source enclosure.
For the exact Appendix-A stages, this requirement follows from
an explicit further smallness choice.

\begin{proposition}[Automatic smallness across both supports and their gap]
\label{prop:a11-source-E-small-20261022}
Assume the source (A.6)--(A.11), $P_*>e^{T_d}$,
$0<\lambda\le1/100$, the A.10 interpolation slope
$\ell\le-\lambda$, and the source-specific small A.11
root $|c_i|\le54\lambda^{28}$. Let $E_b$ be the uniform
precorrection A.11 plateau profile and
$E_c=E_b(1+\sum_i c_i\beta_i)$ for the two disjoint
width-$3/10$ source bumps with
$\|\beta_i\|_\infty\le640/3$.
Then everywhere on the entire closed A.11 two-support
hull (including the gap),
\[
0<E_c\le16P_*\lambda^{45}e^{-13/(2\lambda)}
\le16P_*\lambda^{45}.
\]
Consequently $64P_*\lambda^{44}\le1$ implies
$E_c\le\lambda/4$ uniformly.
\end{proposition}
\begin{proof}
At the inner reference $E\le P_*$. The one-unit initial
transition has $\ell\le3/5$, so multiplies $E$ by at most
$e^{1/10}$. The length-$T_d$ axial stage, one-unit
intermediate transition, and length-$60\log(1/\lambda)$
power-law stage give respectively $e^{-T_d/2}$,
$e^{-1/2}$ and at most $\lambda^{30}$.
Thus $E(X_p,\eta)\le P_*\lambda^{30}$.
The pulse has length $13/\lambda$ and slope $-\lambda$,
giving a factor at most $e^{-13/(2\lambda)}$.
The interpolation cannot increase $E$.
The first bump begins at plateau coordinate
$30\log(1/\lambda)-63/20$. Its baseline factor is at most
\[
e^{-(1/2+\lambda)(30\log(1/\lambda)-63/20)}
\le\lambda^{15}e^{693/400}<8\lambda^{15}.
\]
The last inequality uses $693/400<2$ and $e^2<8$,
with a rational Taylor-tail certificate.
The remaining baseline decreases across the whole hull.
Source bump disjointness gives
$|\sum_i c_i\beta_i|\le11520\lambda^{28}<1$.
Hence $E_c<2E_b$ and the theorem follows.
\end{proof}
\begin{corollary}[A cofinal parameter wedge]
\label{cor:a11-explicit-lambda-20261022}
If $T_f=512$ is permissible for the originating
A.10 profile and $P_*>e^{T_d}>1$, put
\[
Q=1+T_d+512+\log P_*,
\quad
\lambda=\frac{10^{-20}}{(1+P_*)Q},
\quad
h=\frac{\lambda e^{-T_d}}{10^6}.
\]
Then $0<\lambda<1/100$, $0<h<\lambda/100$, and
$64P_*\lambda^{44}<1$.
\end{corollary}
\begin{proof}
As $Q>1$, $\lambda<10^{-20}/(1+P_*)$ and
$64P_*\lambda^{44}
<64\cdot10^{-880}P_* /(1+P_*)^{44}<1$.
The other comparisons are immediate.
\end{proof}
This removes the separately assumed $E_c\le\lambda/4$
from prior conditional A.11 stress inequalities after
the authentic A.11 root bound and source stages are
instantiated. It does \emph{not} authenticate the same
profile's $S(\infty)=\Pi(\infty)=0$, future $D_s,P_s$,
complete exterior stress cone, Appendix B contraction,
unforced coefficient hierarchy, physical zero-force
residual or a trusted end-to-end Lean proof.


\subsection*{A posteriori terminal-gauge robustness for the entire A.11 hull (22 October 2026)}
\label{subsec:a11-gauge-robust-20261022}
The exact source stress identity (4.16) on $U=M=J=0$ is
\[
N_s=\frac{4h\eta S-(1-\eta^2)S_\eta}{X}
+4A\eta\Pi-(1-\eta^2)\Pi_\eta,\qquad A=\tfrac12+h.
\]
Let $C_S(\eta)=S(\infty,\eta)$ and $C_\Pi(\eta)=\Pi(\infty,\eta)$
for the \emph{same corrected profile}. When the integrals converge,
\[
S=C_S+\tfrac12\int_X^\infty E^2\,dx,\qquad
\Pi=C_\Pi-\tfrac12\int_X^\infty E^2\,dx/x.
\]
Hence $N_s=N_s^{(0)}+\Delta N_s$ exactly, with
\[
\Delta N_s=\frac{4h\eta C_S-(1-\eta^2)C'_S}{X}
+4A\eta C_\Pi-(1-\eta^2)C'_\Pi.
\]
This separates the genuine source-tail contribution from any
numerical or analytic terminal-gauge error.

\begin{theorem}[Full-hull stress under finite gauge residuals]
\label{thm:a11-gauge-robust-20261022}
Let $\mathcal H=[y_L,y_L+23/10]\times[-1,1]$ be the complete A.11
two-support hull, including the gap. Assume $0<\lambda\le1/100$ and
$E(y,\eta)=E_0(y)(1+u(y,\eta))$ on the hull, where
$E_0(y)=E_*\exp[-(\frac12+\lambda)(y-y_L)]$,
$E_*>0$ and $|u|\le a\le1/1000$.
Suppose the previous zero-gauge source-local estimate has been
independently established for the actual profile:
\[
\sup_{\mathcal H}|N_s^{(0)}|/E<11\lambda/10000.
\]
Writing $X_L=e^{y_L}$, suppose also
\begin{equation}\label{eq:a11-robust-residual-budget}
\frac{4h\|C_S\|_\infty+\|C_S'\|_\infty}{X_L}
+4A\|C_\Pi\|_\infty+\|C_\Pi'\|_\infty
\le\frac{(1-a)E_*\lambda}{400}.
\end{equation}
Then over both A.11 supports and every point in the gap,
\[
\sup_{\mathcal H}|N_s|/E<111\lambda/10000,\qquad
\lambda/2-\sup_{\mathcal H}|N_s|/E>4889\lambda/10000.
\]
If additionally $0<h\le\lambda/100$,
$Q_s\ge97\lambda/100$,
$-151\lambda/100\le\ell\le-49\lambda/100$,
$X\ge9/\lambda$ and $0<L=1-2h\eta^2\le1$, then
\[
2(1-a_s/p_1)^2-(a_s-2)w^2>1124/1000>9/20,
\]
where $a_s=2-2\ell$, $p_1=XQ_s/L$ and $w=N_s/(EQ_s)$.
\end{theorem}
\begin{proof}
Since $(\frac12+\lambda)(23/10)\le1173/1000<6/5$,
the finite Taylor comparison
$e^{6/5}<1504439/453125<4$ shows that
$E>(1-a)E_*/4$ everywhere on the complete hull.
The terminal contribution therefore obeys
\[
\sup_{\mathcal H}|\Delta N_s|/E\le
\frac4{(1-a)E_*}
\left[\frac{4h\|C_S\|+\|C_S'\|}{X_L}
+4A\|C_\Pi\|+\|C_\Pi'\|\right]\le\lambda/100.
\]
Add the assumed zero-gauge bound and subtract from $\lambda/2$.
For the cone, $p_1\ge873/100$, $a_s/p_1<1/4$,
$0<a_s-2\le151/5000$ and $|w|<111/9700$.
Thus its normalized margin exceeds the exact rational number
\[
2(3/4)^2-\frac{151}{5000}\left(\frac{111}{9700}\right)^2
=\frac{529254389529}{470450000000}>\frac{1124}{1000}.
\]
\end{proof}
\paragraph{Scope.}
The tolerance \eqref{eq:a11-robust-residual-budget} is a new
a posteriori acceptance condition, not an authenticated source
interval. The source's A.11 moment root, exact same-profile terminal
$S$ and pressure gauges, all remaining stress patches, Appendix B,
the force-free correction hierarchy, and the physical zero-force
PDE theorem still require independent proof. The originating
whole-space theorem in \cite{openai2026} has smooth forcing.


\subsection*{The source A.11 support-measure refinement (22 October 2026)}
\label{subsec:a11-support-measure-refinement}
The two source-fixed correction supports have length $3/10$ each,
so their total measure is $3/5$, while their enclosing hull has
length $23/10$. The gap of length $17/10$ does not contribute
to the parameter derivative of the perturbation.

\begin{proposition}[Sharper conditional A.11 full-hull stress]
\label{prop:a11-support-measure-refinement}
Let $E(y,\eta)=E_0(y)(1+u(y,\eta))$ on the A.11 hull,
with $E_0$ independent of $\eta$ along every future ray,
$\partial_y\log(\sqrt X E_0)\le0$, $|u|\le a<1$ and
$|\partial_\eta u|\le b$. Assume $u_\eta$ is supported on
the two disjoint prescribed bumps, of combined measure
at most $3/5$. Let
\[
 F_0=\int_0^\infty E(y+t,\eta)^2\,dt,\qquad
 F_1=\int_0^\infty e^t E(y+t,\eta)^2\,dt.
\]
Assume the \emph{same corrected source profile} has terminal
gauges $S(\infty,\eta)=\Pi(\infty,\eta)=0$ and baseline bounds
$F_{0,0}/E_0^2\le1$, $F_{1,0}/E_0^2\le5$ on the hull.
Then with $A=\frac12+h$ and $E_*=\sup_{\rm hull}E_0$,
\[
 \frac{|N_s|}{E}\le\frac{E_*}{1-a}
 \left[2(1+a)^2(A+5h)+\frac35(1+a)b\right].
\]
If $0<\lambda\le1/100$, $0<h\le\lambda/100$,
$a,b\le1/1000$, and $E_*\le\lambda/1000$, then
\[
 \sup_{\rm hull,\ |\eta|\le1}\frac{|N_s|}{E}
 <\frac{201}{200000}\lambda<\frac\lambda2.
\]
\end{proposition}
\begin{proof}
The baseline slope yields $E_0(y+t)^2\le e^{-t}E_0(y)^2$.
The parameter derivative of the difference of future integrals is
\[
 \partial_\eta(F_1-F_0)=2\int_0^\infty(e^t-1)
 E_0(y+t)^2(1+u(y+t,\eta))u_\eta(y+t,\eta)\,dt.
\]
The kernel $(e^t-1)e^{-t}\le1$ and the support measure
$\le3/5$ give
$|\partial_\eta(F_1-F_0)|\le
\frac65(1+a)bE_0(y)^2$.
The exact source identity, with the specified gauges, is
\[
 N_s=2\eta(hF_1-AF_0)
 -\frac{1-\eta^2}{2}\partial_\eta(F_1-F_0).
\]
Use $F_0\le(1+a)^2E_0^2$ and
$F_1\le5(1+a)^2E_0^2$, and divide by
$E\ge(1-a)E_0$. Under the numerical premises,
$A+5h\le2503/5000$, and exact fraction arithmetic gives
\[
 \frac{2(1001/1000)^2(2503/5000)
 +(3/5)(1001/1000)(1/1000)}{999/1000}
 =\frac{2509510003}{2497500000}<\frac{201}{200}.
\]
Together with $E_*\le\lambda/1000$, this concludes.
\end{proof}

\paragraph{Certification boundary.}
This is a conditional full-hull estimate. The moment-matched
source coefficients and their $\eta$-jets, same-profile terminal
gauges, future-tail assumptions, common parameter selection, and
remaining stress-cone stages have not been certified jointly.
The complete Appendix-B contraction, unforced startup profiles,
physical-coordinate PDE residual and Theorem 1.1 remain open.


\subsection*{Terminal A.11/A.19 residuals as exact stress-gauge defects (22 October 2026)}
\label{subsec:A11-source-terminal-residual-lift}
Let \(\mathcal H\) be the full A.11 two-support hull and intervening gap.
On the source region \(U=M=J=0\), write
\[
 N_s=\frac{4h\eta S-(1-\eta^2)S_\eta}{X}
       +4A\eta\Pi-(1-\eta^2)\Pi_\eta,\quad A=\tfrac12+h.
\]
Fix the \emph{same} corrected source profile throughout and define
\[
 \mathcal F(\eta)=\frac{\lambda S(\infty,\eta)}
  {X_pe_b^2(1+\eta^2)^{-2}},\qquad
 \mathcal G(\eta)=\int_{\text{two A.11 bumps}}(E_c^2-E_b^2)\,dy.
\]
The source baseline pressure trace is chosen to cancel the baseline
terminal pressure. Thus, with
\(K(\eta)=X_pe_b^2/[\lambda(1+\eta^2)^2]\),
the genuine terminal errors are \(C_S=K\mathcal F\),
\(C_\Pi=\mathcal G/2\). Differentiating \emph{totally} in \(\eta\)
along the actual A.19 amplitude and A.11 coefficient branches gives
the exact identity
\begin{equation}\label{eq:a11-source-residual-lift-20261022}
\boxed{\begin{aligned}
 \Delta N_s={}&\frac{K}{X}
  \left[\left(4h\eta+
   \frac{4\eta(1-\eta^2)}{1+\eta^2}\right)\mathcal F
          -(1-\eta^2)\mathcal F_\eta\right]\\
 &+2A\eta\mathcal G-\frac{1-\eta^2}{2}\mathcal G_\eta.
\end{aligned}}
\end{equation}
This follows from \(K_\eta/K=-4\eta/(1+\eta^2)\).
Disjoint A.11 bumps with eta-independent baseline give
\(\mathcal G=\sum_{i=1}^2(2v_ic_i+w_ic_i^2)\) and
\(\mathcal G_\eta=2\sum_i(v_i+w_ic_i)c_{i,\eta}\),
with \(v_i=\int E_b^2\beta_i\,dy\), \(w_i=\int E_b^2\beta_i^2\,dy\).
Consequently, \emph{exact} \(C^1\) simultaneous solutions of the
A.19 and A.11 source equations imply
\(\mathcal F=\mathcal F_\eta=\mathcal G=\mathcal G_\eta=0\);
small residual values alone do not control their total derivatives.

\begin{proposition}[Fail-closed two-support-and-gap stress gate]
Suppose the authentic zero-gauge estimate on the \emph{same}
corrected profile is \(b_0\ge\sup_{\mathcal H}|N_s^{(0)}|/E\).
Assume \(E>0\) and that independently authenticated nonnegative
source bounds satisfy
\[
 \alpha\ge\sup_{\mathcal H}\frac{K}{XE},\qquad
 \beta\ge\sup_{\mathcal H}\frac1E,\qquad
 \epsilon_F\ge\|\mathcal F\|_\infty,\quad
 \epsilon_{F'}\ge\|\mathcal F_\eta\|_\infty,
\]
and analogously \(\epsilon_G,\epsilon_{G'}\).
For \(0<\lambda\le1/20,\;0<h\le\lambda/100\), one has
\[
 \sup_{\mathcal H}\frac{|N_s|}{E}\le b_0+
 \alpha[(2+4h)\epsilon_F+\epsilon_{F'}]
 +\beta[2A\epsilon_G+\epsilon_{G'}/2].
\]
If \(b_0\le\lambda/500\) and the two residual brackets total
strictly less than \(249\lambda/500\), then
\(\sup_{\mathcal H}|N_s|/E<\lambda/2\).
\end{proposition}
\begin{proof}
Use \(|\eta|\le1\), \(0\le1-\eta^2\le1\), and the exact identity
\(2(1+u^2)-4u(1-u^2)=2[(1-u)^2+2u^3]\ge0\)
for \(u=|\eta|\le1\). Hence
\(4|\eta|(1-\eta^2)/(1+\eta^2)\le2\).
Divide the terminal correction by \(E>0\), apply the sup bounds,
and add \(b_0\).
\end{proof}
The code \texttt{a11\_residual\_certificate.py} checks the
rational implication and refuses missing interval inputs.
A positive synthetic test is NOT an authenticated source certificate.
The common A.11/A.19 source roots, full-hull future-tail bound,
other exterior stress cones, Appendix B fixed point, unforced startup
systems, physical PDE residual, and the proposed unforced
Theorem 1.1 remain unproved. Lean code in this tranche covers
only finite scalar identities, not the source integrals or PDE.


\subsection*{A fail-closed full-unforced source-certificate gate (22 October 2026)}
\label{subsec:a11-full-unforced-fail-closed}
The existing conditional A.11 support-measure bound does not itself certify one
actual corrected source profile. On the full two-support hull, the exact
same-profile residual lift is
\[
\Delta N_s=\frac KX\left[
\left(4h\eta+\frac{4\eta(1-\eta^2)}{1+\eta^2}\right)\mathcal F
 -(1-\eta^2)\mathcal F_\eta\right]
+2A\eta\mathcal G-\frac{1-\eta^2}{2}\mathcal G_\eta,
\]
where $\mathcal F$ and $\mathcal G$ are the terminal $S$ and pressure
residuals of the \emph{same} corrected profile and the derivatives are
\emph{total} derivatives along its A.19/A.11 roots.
Let $b_0\ge\sup|N_s^{(0)}|/E$, $\alpha\ge\sup K/(XE)$ and
$\beta\ge\sup1/E$. Then the exact triangle inequality gives
\[
\sup_{\mathcal H}\frac{|N_s|}{E}\le b_0+
\alpha[(2+4h)\|\mathcal F\|_\infty+\|\mathcal F_\eta\|_\infty]
+\beta[(1+2h)\|\mathcal G\|_\infty+
                   \tfrac12\|\mathcal G_\eta\|_\infty].
\]
Therefore $b_0\le\lambda/500$ and a strict total residual budget
$<249\lambda/500$ give $\sup_{\mathcal H}|N_s|/E<\lambda/2$.
This conditional implication is not a new global Navier--Stokes theorem.

The executable \texttt{full\_unforced\_source\_gate.py} uses exact rational
input intervals, checks equality of source-profile hashes across all
stages, rejects absent total $\eta$ derivatives, demands positive
source-backed cone/Qs lower bounds and refuses missing independent
analytic/Lean proof evidence. Fourteen adversarial negative cases must
return \texttt{HOLD}; an artificial arithmetic-positive case is labelled
\texttt{CONDITIONAL} and is \emph{never} source-authenticated.

\paragraph{Missing end-to-end proof.}
One validated parameter tuple must satisfy the actual A.11 and A.19
moment equations simultaneously and certify both terminal gauges.
The full exterior/heat/patch cone, Appendix-B nonlinear leading profile,
the true unforced $n=1,2$ and higher-order source equations,
differentiated physical zero-force residual, and trusted Lean
formalization remain unproved. OpenAI's originating Theorem 1.1
contains smooth compactly supported forcing and cannot be used to
claim an unforced singularity.

\subsection*{Two-coordinate A.11 interval-net certification and a strict rational margin (22 October 2026)}
\label{subsec:a11-two-coordinate-net-and-margin}
This checkpoint supplements the previously derived exact terminal-gauge
residual lift. The full corrected A.11 hull \(\mathcal H\) contains
both azimuthal supports \emph{and} the intervening gap and depends on
the pulse coordinate \(y\) and angular coordinate \(\eta\).
An angular grid alone at fixed \(y\) is insufficient to establish a
uniform estimate on \(\mathcal H\).

Fix the \emph{same} authentic corrected A.11/A.19 profile and define
\[
G_\pm(y,\eta)=\lambda E(y,\eta)\pm 2N_s(y,\eta).
\]
When \(E>0\), the exact desired source stress criterion is equivalent to
\(G_+\ge0\) and \(G_-\ge0\). There is no need to assume that
the physical source integrals are polynomial in \(\eta\).

\begin{proposition}[Two-coordinate interval-net gate]
Let \(V\subseteq\mathcal H\) be finite. Suppose every
\(z=(y,\eta)\in\mathcal H\) has a node \(v=(y_i,\eta_j)\in V\) with
\[
|y-y_i|\le r_y,\qquad|\eta-\eta_j|\le r_\eta.
\]
For each \(f\in\{E,G_+,G_-\}\), suppose independently certified
nonnegative constants \(L_{f,y},L_{f,\eta}\) satisfy
\[
|f(z)-f(v)|\le L_{f,y}|y-y_i|+
                    L_{f,\eta}|\eta-\eta_j|.
\]
If for a common \(\delta>0\) every node satisfies
\[
E(v)\ge\delta+L_{E,y}r_y+L_{E,\eta}r_\eta,\qquad
G_\pm(v)\ge L_{\pm,y}r_y+L_{\pm,\eta}r_\eta,
\]
then \(E(z)\ge\delta>0\) and \(G_\pm(z)\ge0\)
for every \(z\in\mathcal H\). In particular
\(\sup_{\mathcal H}|N_s|/E\le\lambda/2\).
\end{proposition}
\begin{proof}
For each point choose a covering node and subtract the
two coordinatewise Lipschitz-remainder envelopes. The two signed
inequalities imply \(2|N_s|\le\lambda E\), and \(E>0\)
permits division.
\end{proof}
The \(\eta\)-derivative inputs must be \emph{total} derivatives along the
same genuine A.11/A.19 source root, while \(y\)-derivative inputs must
cover support boundaries, the gap, and applicable stage transitions.
A grid scan without authenticated derivative/remainder bounds is
\emph{not} a uniform certificate.

\begin{proposition}[Explicit release margin for the residual lift]
Assume the earlier exact same-profile residual inequality
\(\sup_{\mathcal H}|N_s|/E\le b_0+\mathcal R\),
\(E>0\), and \(\lambda>0\), where
\[
\mathcal R=\alpha[(2+4h)\|\mathcal F\|_\infty+
\|\mathcal F_\eta\|_\infty]+
\beta[(1+2h)\|\mathcal G\|_\infty+
\tfrac12\|\mathcal G_\eta\|_\infty]
\]
with the same definitions and authenticated premises as in the
preceding source-gauge residual section. If
\[
b_0\le\frac{\lambda}{500},\qquad
\mathcal R\le\frac{4979\lambda}{10000},
\]
then the source stress ratio has the strict spare margin
\[
\boxed{\quad
\sup_{\mathcal H}\frac{|N_s|}{E}
\le\frac{\lambda}{2}-\frac{\lambda}{10000}
<\frac{\lambda}{2}.\quad}
\]
\end{proposition}
\begin{proof}
Exactly \(1/500+4979/10000=4999/10000
=1/2-1/10000\). Positivity of \(\lambda\) gives strictness.
\end{proof}

\paragraph{Source authentication and limits.}
Both propositions are complete conditional scalar arguments.
They do not supply the true \((y,\eta)\) intervals or root derivatives.
The companion exact-rational verifier replays
10,000 scalar budgets and 10,000 two-coordinate affine controls and
rejects three negative controls. The Lean source
\texttt{A11RectangularNetScalar.lean} formalizes the corresponding
finite implications subject to independent CI success; it does not
prove the actual PDE-specific source bounds. The authentic joint
parameter tuple and A.11/A.19 roots, full stress cones, Appendix B
contraction, unforced startup hierarchy, physical zero-force
Navier--Stokes residual and the proposed Theorem~1.1 all remain
OPEN/HOLD. The original OpenAI construction has smooth forcing.


\subsection*{Moving-endpoint transport and a quantitative A.11 root-jet reduction (22 October 2026)}
\label{subsec:a11-moving-endpoint-jet-20261022}
For the \emph{actual source-prescribed} translated $\sigma'$ bumps, the angular coefficients satisfy
\[
u(c_1+xc_2)=\Delta_I(\eta),\qquad
2v(c_1+zc_2)+w(c_1^2+zc_2^2)=0.
\]
Here $u,v,w,x,z$ are independent of $\eta$ at fixed $\lambda$.
The exact nonlinear Jacobian and determinant are
\[
J(c)=\begin{pmatrix}u&ux\\2(v+wc_1)&2z(v+wc_2)\end{pmatrix},
\quad
\det J=2u[z(v+wc_2)-x(v+wc_1)].
\]
From the source's $\|B_{11}^{-1}\|_\infty<3/2$,
$\|Q\|<320$ and prior small-root $\|c\|_\infty\le192\lambda^{28}$,
\[
\|B_{11}^{-1}DQ(c)\|_\infty
<184320\lambda^{28}<\frac14
\quad(0<\lambda\le1/20).
\]
Thus $J(c)$ is invertible, and along a genuine $C^1$ root branch,
\begin{equation}\label{eq:a11-jet-reduction-20261022}
\boxed{\ \|c_\eta\|_\infty\le2|\Delta_{I,\eta}|.\ }
\end{equation}

\begin{proposition}[Moving-endpoint transport formula]
Let $r=I/(XH)$ satisfy $r_y+(1+\ell)r=1$.
Assume the source plateau on $[a(\eta),b(\eta)]$ has
$\ell=-\lambda$ and $\lambda$ independent of $\eta$.
For $\alpha=1-\lambda$, $L=b-a$, and differentiable endpoints,
\[
\Delta_I=\frac1\alpha-r(b,\eta)
=-e^{-\alpha L}\left[r(a,\eta)-\frac1\alpha\right],
\]
and the exact \emph{total} angular derivative is
\begin{equation}\label{eq:a11-moving-endpoint-identity-20261022}
\boxed{\ \Delta_{I,\eta}=-e^{-\alpha L}
\left[\frac{d}{d\eta}r(a(\eta),\eta)
-\alpha L_\eta\left(r(a(\eta),\eta)-\frac1\alpha\right)\right].\ }
\end{equation}
\end{proposition}
\begin{proof}
Solve the constant-slope linear ODE on the plateau, then differentiate
the resulting scalar transport identity. The first derivative contains
the moving endpoint term $a_\eta r_y(a,\eta)$; the second
contains $L_\eta=b_\eta-a_\eta$.
\end{proof}

\begin{theorem}[Conditional strengthened A.11 source-root bounds]
Suppose the earlier small-branch existence hypotheses hold, with
$0<\lambda\le1/20$, and for all $|\eta|\le1$,
\[
0\le r(a(\eta),\eta)\le\frac{20}{17},\qquad
L(\eta)\ge30\log(1/\lambda)-3,\quad
\left|\frac{d}{d\eta}r(a(\eta),\eta)\right|\le\frac14,\quad
|L_\eta|\le\frac18.
\]
Then on that true $C^1$ A.11 root branch,
\begin{equation}\label{eq:a11-new-jet-quantitative-20261022}
\boxed{\ 
\|c\|_\infty\le\frac{315}{19}\lambda^{28}<17\lambda^{28},
\qquad
\|c_\eta\|_\infty\le\frac{609}{152}\lambda^{28}<5\lambda^{28}.
\ }
\end{equation}
\end{theorem}
\begin{proof}
Put $d=r(a,\eta)-1/\alpha$. Since
$1\le1/\alpha\le20/19$ and $0\le r(a)\le20/17$,
one has $|d|\le20/19$. The plateau length bound gives
$e^{-\alpha L}\le e^3\lambda^{57/2}
<21\lambda^{57/2}$, and $\sqrt\lambda<1/4$.
Therefore
$|\Delta_I|\le (105/19)\lambda^{28}$ and the previously
proved root bound $\|c\|\le3|\Delta_I|$ gives the
first assertion. Likewise
\[
|\Delta_{I,\eta}|
\le21\lambda^{57/2}
\left(\frac14+\frac{20}{19}\frac18\right)
=\frac{609}{76}\lambda^{57/2}.
\]
Apply \eqref{eq:a11-jet-reduction-20261022} and
$\sqrt\lambda<1/4$ to obtain the jet bound.
\end{proof}

\paragraph{Proof boundary.}
This reduces the earlier $54\lambda^{28}$ and
$6\lambda^{28}$ A.11 input gates to two earlier
\emph{moving-endpoint} source derivative tests.
No real A.10/A.11 profile has yet been certified to satisfy
these $1/4$ and $1/8$ bounds. The separate A.19 root,
same-profile terminal $S$ and $\Pi$ gauges, other cone stages,
Appendix B contraction, unforced startup and zero-force PDE
reconstruction all remain OPEN/HOLD.
The companion rational/symbolic program tests the algebra
and synthetic examples, not source authentication.
The Lean module contains conditional scalar implications only.


\subsection*{A.10 source-fixed interpolation closes the provisional A.11 angular-derivative gates (22 October 2026)}
\label{subsec:a11-A10-fixed-stage-closure}
The original Appendix A.10--A.11 construction fixes all stage boundaries
independently of \(\eta\). For fixed \(T_f=512\) and
\(0<\lambda\le1/100\), the A.11 first bump center follows
the interpolation endpoint by
\(L_1=30\log(1/\lambda)-3\), so \(\partial_\eta L_1=0\).
The moving-endpoint formula of the preceding section is still valid
for hypothetical moving stages, but its \(L_{1,\eta}\) term vanishes
\emph{for this source schedule}.

Write \(a=1-\lambda\), \(J_0=\log(1+\eta^2)\), and
take the source's exact smooth step \(\sigma\) from (A.5).
If the pre-interpolation moment ratio \(r_p=I/(XH)\) has its
source-factorized \(\eta\)-independence, the integrating factor in (A.10)
and \(\sigma(1-t)=1-\sigma(t)\) give
\[
r_f=\frac{2r_p}{1+\eta^2}e^{-aT_f}+
 \int_0^{T_f} e^{-av}
 \left(\frac2{1+\eta^2}\right)^{\sigma(v/T_f)}\,dv.
\]
Differentiation under this genuine finite integral gives
\[
\boxed{
r_{f,\eta}=-\frac{2\eta}{1+\eta^2}
\left[\frac{2r_p}{1+\eta^2}e^{-aT_f}+
\int_0^{T_f}e^{-av}\sigma(v/T_f)
\left(\frac2{1+\eta^2}\right)^{\sigma(v/T_f)}\,dv\right].
}
\]
For \(0\le r_p\le20/17\), \(a\ge99/100\), \(T_f=512\),
the exact step satisfies
\(\sigma(1/4)=1/(1+e^{128/9})<10^{-6}\);
the first 16 positive Taylor terms prove \(e^{128/9}>10^6\).
The first 15 terms give \(e^{30}>10^{10}\), so
\(e^{-aT_f/4},e^{-aT_f}<10^{-10}\).
Splitting the integral at \(T_f/4\) proves
\[
\boxed{|r_{f,\eta}|<
\frac{40}{17}\,10^{-10}
+\frac{200}{99}(10^{-6}+10^{-10})<3\cdot10^{-6}.}
\]
On the source-fixed \(-\lambda\) plateau the exact mismatch
\(\Delta_I=1/a-r(y_1)\) satisfies
\[
\Delta_I=-e^{-aL_1}(r_f-1/a),\qquad
\Delta_{I,\eta}=-e^{-aL_1}r_{f,\eta}.
\]
Assuming the already established \(0\le r_f\le20/17\)
(which follows from its incoming bound and the source slope barrier),
\(e^{-aL_1}\le e^3\lambda^{30a}\le21\lambda^{29}\) yields
\[
|\Delta_I|\le(21/95)\lambda^{28},\quad
|\Delta_{I,\eta}|<63\cdot10^{-6}\lambda^{29}.
\]
Together with the existing genuine A.11 small-branch inverse and
Jacobian gates \(\|c\|_\infty\le3|\Delta_I|\) and
\(\|c_\eta\|_\infty\le2|\Delta_{I,\eta}|\), this gives the
source-conditional sharpened bounds
\[
\boxed{\|c\|_\infty\le(63/95)\lambda^{28}<\lambda^{28},
\qquad\|c_\eta\|_\infty<\lambda^{28}/100000.}
\]
These finite source-stage deductions replace the provisional
\(|D_\eta r(a(\eta),\eta)|\le1/4\) and
\(|L_{1,\eta}|\le1/8\) assumptions of Edition III.
They recover and sharpen earlier A.10 sensitivity analyses
already contained in the cumulative manuscript, not a new
unconditional A.19 theorem.

The original source theorem has \emph{forcing}. The full ordered
parameter selection, same-profile A.19/pressure terminal gauges,
exterior stress cones, Appendix B contraction, unforced startup,
zero-force PDE residual, and complete Lean kernel proof remain
OPEN/HOLD. The finite rational program only audits the displayed
bounds and synthetic controls.


\subsection*{Coupled A.11/A.19 root compatibility and terminal-gauge identity (22 October 2026)}
\label{subsec:a11-a19-coupled-gauge-20261022}
Distinguish the two azimuthal corrections $q_1,q_2$ in (A.11)
from the separate pulse-velocity moment corrections in (A.9)/(A.19).
The source orders the stages so the later azimuthal field
$E_c=E_b(1+q_1\beta_1+q_2\beta_2)$ is independent of the
pulse amplitude $A$, while the pulse changes $U$ and its two
velocity-moment corrections are affine in $A$.
The exact common-profile residuals are
\[
R_1=u(q_1+xq_2)-\Delta_I(\eta),\quad
R_2=2v(q_1+zq_2)+w(q_1^2+zq_2^2),
\]
\[
R_3=k_2(q,\eta)A^2+k_1(q,\eta)A+k_0(q,\eta).
\]
The $k_i$ may depend on the genuinely corrected azimuthal profile.
At fixed $\lambda$, $u,v,w,x,z$ are independent of $\eta$.

\begin{proposition}[Exact block-triangular Jacobian]
The Jacobian in unknowns $(q_1,q_2,A)$ is
\[
\begin{pmatrix}
u&ux&0\\
2(v+wq_1)&2z(v+wq_2)&0\\
\partial_{q_1}R_3&\partial_{q_2}R_3&2k_2A+k_1
\end{pmatrix},
\]
with determinant
\begin{equation}\label{eq:joint-A11-A19-jac-20261022}
\boxed{2u[z(v+wq_2)-x(v+wq_1)](2k_2A+k_1).}
\end{equation}
If the small A.11 root has nonsingular Jacobian, $k_2>0$
and $k_0<0$ at a positive A.19 root, then
\[
2k_2A+k_1=\sqrt{k_1^2-4k_2k_0}>0.
\]
Thus the combined root is locally unique and $C^1$ when the
coefficients are $C^1$. If the needed roots are unique for every
$\eta\in[-1,1]$, the local branches paste to a common $C^1$ family.
\end{proposition}
\begin{proof}
Expand the determinant along its third column. For
$k_2>0,k_0<0$, the quadratic has exactly one positive
root, with a strictly positive derivative given by the square
root of its discriminant. Apply the implicit function theorem.
\end{proof}

\begin{proposition}[Exact pressure-neutrality factor]
On the source plateau,
$E_b(y)^2=E_b(y_1)^2e^{-(1+2\lambda)(y-y_1)}$,
with the two disjoint translated bumps centered two units apart.
Therefore
\begin{equation}\label{eq:joint-A11-pressure-factor-20261022}
\boxed{\ \int(E_c^2-E_b^2)\,dy=E_b(y_1)^2R_2(q,\eta).\ }
\end{equation}
Consequently the exact A.11 second equation forces the pressure
increment and its \emph{total} $\eta$ derivative to vanish identically.
If the prescribed unedited reference already satisfies its
terminal pressure normalization, the same corrected profile
inherits $\Pi(\infty,\eta)=0$. This baseline gauge is a distinct
source assumption; A.11 neutrality alone merely preserves it.
\end{proposition}
\begin{proof}
Expand the square. Disjointness removes all mixed bump terms.
The first support contributes the linear and quadratic weights
$v,w$; translation of the second support multiplies both by
$z=e^{-2(1+2\lambda)}$. This yields exactly $R_2$.
Differentiate along the common $C^1$ root and use the source
pressure convention
$\Pi(\infty)-\Pi^{\rm base}(\infty)=\frac12\int(E_c^2-E_b^2)\,dy$.
\end{proof}

\begin{proposition}[Exact total-amplitude sensitivity]
Along a common $C^1$ positive root of $R_3=0$, let a dot
denote the total $\eta$ derivative through $q_1,q_2$.
Then
\[
\boxed{\quad
 A_\eta=-\frac{\dot k_2A^2+\dot k_1A+\dot k_0}
 {2k_2A+k_1}
 =-\frac{\dot k_2A^2+\dot k_1A+\dot k_0}
 {\sqrt{k_1^2-4k_2k_0}}.
\quad}
\]
A source-certified lower bound $2k_2A+k_1\ge m_A>0$
therefore gives
$|A_\eta|\le
(A^2|\dot k_2|+A|\dot k_1|+|\dot k_0|)/m_A$.
If $R_3$ is the genuine normalized terminal $S$ equation,
then $S(\infty,\eta)$ and its \emph{total} derivative vanish
on this root. With the preceding exact A.11 pressure equation
and a certified baseline gauge, all four terminal-residual variables
$\mathcal F,\mathcal F_\eta,\mathcal G,\mathcal G_\eta$ vanish.
The terminal-gauge contribution to $N_s$ is then identically zero;
the separate future-tail, other stress cones and physical PDE
remain to be proved.
\end{proposition}
\begin{proof}
Differentiate $R_3=0$ along the joint branch and divide
by the positive $A$ derivative. Apply the total derivative
to the two identical-zero terminal residual functions.
\end{proof}

\paragraph{Adversarial negative control and release boundary.}
At $q_1=1/10$, $q_2=0$, $u=v=1$, $x=7$, $z=1/8$,
$w=20$, $\Delta_I=1/10$, and $R_3=A^2-1$, the
values $A=1$, $R_1=R_3=0$ coexist with $R_2=2/5\ne0$.
Thus an A.19 root \emph{alone} does not establish
pressure neutrality. Independent symbolic and 10,000 rational
synthetic checks confirm the displayed finite-dimensional
identities; they do not authenticate a single source-admissible
parameter tuple or the true terminal baseline pressure.
The global originating source theorem has forcing.
Appendix B, unforced similarity startups, exterior stress
certification, zero-force physical PDE residual and proposed
unforced Theorem~1.1 all remain OPEN/HOLD.


\subsection*{Proof Attack 1: exact A.11 pressure normalization (Edition VI, 22 October 2026)}
\label{subsec:a11-pressure-editionVI}
The original source Lemma A.5 (A.21) sets
\[
\Pi_0(\eta)=-\frac12\int_{\mathbb R}E_{\mathrm{id,sched}}(y,\eta)^2\,dy
\]
for the exact reference schedule with the pressure-neutral A.11
corrections omitted. Its pressure equation is
$\Pi_y=E^2/2$, $y=\log(X/X_R)$.
For the actual two translated smooth-step derivative bumps
$\beta_1(y)=\beta(y-y_1)$ and
$\beta_2(y)=\beta(y-y_1-2)$, each supported
on an interval of length $3/10$, write
\[
E_c=E_b(1+q_1\beta_1+q_2\beta_2),\qquad
E_b(y)^2=C(\eta)e^{-b(y-y_1)},\quad b=1+2\lambda.
\]
Define $v=\int e^{-bt}\beta(t)dt$,
$w=\int e^{-bt}\beta(t)^2dt$, $z=e^{-2b}$.
Disjoint support and weighted translation yield
\begin{equation}\label{eq:editionVI-pressure-factor}
\int(E_c^2-E_b^2)dy
=C(\eta)\bigl[2v(q_1+zq_2)+w(q_1^2+zq_2^2)\bigr]
=C(\eta)R_2(\eta).
\end{equation}
Off the two supports the corrected and reference profiles agree.
Hence the actual second A.11 condition $R_2=0$ gives
\begin{equation}\label{eq:editionVI-pressure-zero}
\boxed{\Pi(\infty,\eta)=\Pi_0(\eta)+\frac12
\int E_c^2dy=\frac12 C(\eta)R_2(\eta)=0.}
\end{equation}
On a $C^1$ exact A.11 root branch the \emph{total}
angular derivative of this identically zero function also vanishes.

\paragraph{Provenance and boundary.}
This pressure identity is explicitly part of the originating
source's Lemma A.5, equations (A.21)--(A.23); this is
an independent verification and a resolved redundant
numerical-proof gate, not a novel global theorem.
The existence of one full source-compatible A.11/A.19 branch,
later pressure-preserving edits, other stress-cone patches,
Appendix B, all unforced corrections, the physical zero-force
Navier--Stokes residual, and full Lean formalization remain OPEN.
Independent symbolic, prescribed-bump 68-digit numerical,
and negative-control checks passed; the new Lean code
contains only conditional scalar/integral-equality implications
and requires CI compilation.

\subsection*{Proof Attack 2: heat pressure sign, compensation, and unforced-from-rest obstruction (Edition VII, 22 October 2026)}
\label{subsec:heat-pressure-compensation-VII}
The source heat profile replaces the classical tail
$E_{\rm cl}$ by $E_{\rm new}=E_{\rm cl}[1+\chi_K(y)(\mathcal H(2d/X)-1)]$,
where $d=1-\eta^2$, $0\le\chi_K\le1$, and
\[
\mathcal H(Z)=\Gamma(1+h)^{-1}\int_0^\infty
  e^{-v}v^h(1+Zv)^{-h}\,dv,\quad Z\ge0,\quad h>0.
\]
\begin{proposition}[Strict signed source heat defect]
For every $Z\ge0$, $0<\mathcal H(Z)\le1$ and
$0\le1-\mathcal H(Z)\le h(1+h)Z$.
Set $q=\chi_K(1-\mathcal H(2d/X))$.
Then $0\le q<1$ and
$E_{\rm new}^2-E_{\rm cl}^2=-E_{\rm cl}^2q(2-q)\le0$.
For $\Delta P_{\rm heat}=\frac12\int(E_{\rm new}^2-E_{\rm cl}^2)d\log X$,
\begin{equation}\label{eq:heat-defect-VII}
\boxed{-2h(1+h)d\int\frac{E_{\rm cl}^2\chi_K}{X}d\log X
\le\Delta P_{\rm heat}\le0.}
\end{equation}
It is strictly negative for $|\eta|<1$ if the heat cutoff
is nonzero on a positive-measure set and zero at $\eta=\pm1$.
\end{proposition}
\begin{proof}
The normalized integrand defining $\mathcal H$ is a
probability measure with mean $1+h$.
Integrate $0\le1-(1+Zv)^{-h}\le hZv$.
Then $E_{\rm new}=E_{\rm cl}(1-q)$,
and $q(2-q)\le2q\le4h(1+h)d\chi_K/X$.
\end{proof}
The heat edit alone therefore is NOT pressure-neutral.
As already proved qualitatively in source Proposition A.7,
three ordered nonnegative disjoint E-bumps on the earlier
$U=0$ patch compensate the changes in
$\frac12\int E^2d\log X$,
$-\frac12\int E^2dX$, and
$\int\sqrt{2X}(E-E_{\rm pow})dX$.
The remaining two channels $M,J$ are unchanged since $U=0$.
The normalized linearized radial weights are
\[
f(\eta)x^{-3/2-\lambda},\quad
-f(\eta)x^{-1/2-\lambda},\quad \sqrt{2x},\quad x=X/X_*.
\]
They have distinct exponents, so their generalized
Vandermonde determinant has fixed nonzero sign for
three ordered distinct points. Integration against
ordered disjoint positive bumps preserves that sign;
hence the actual three-moment Jacobian is nonsingular.
\begin{lemma}[Explicit quantitative inverse gate]
For $F(c)=Lc+Q(c)+d$ on $\mathbb R^3$ suppose
$m\ge\|L^{-1}\|_\infty$, $\delta=\|d\|_\infty$,
$\|Q(c)\|_\infty\le K\|c\|_\infty^2$, and
$\|Q(c)-Q(c')\|_\infty\le
K(\|c\|_\infty+\|c'\|_\infty)\|c-c'\|_\infty$.
If $4m^2K\delta<1/2$, $F(c)=0$ has a unique root in
$\|c\|_\infty\le2m\delta$.
\end{lemma}
\begin{proof}
The map $-L^{-1}(d+Q(c))$ takes this ball into itself
and has Lipschitz constant at most $4m^2K\delta<1/2$.
\end{proof}
The real source still needs authenticated $m,K,\delta$
and complete stress-cone intervals. This is an
explicit quantitative version of the local inverse gate,
not a replacement for source Proposition A.7.
\paragraph{Unforced-from-rest obstruction.}
For a smooth finite-energy decaying unforced incompressible
Navier--Stokes solution with $\nu>0$ and $u(\cdot,0)=0$,
the classical identity
\[
\frac12\|u(t)\|_2^2+
\nu\int_0^t\|\nabla u(s)\|_2^2\,ds=0
\]
forces $u\equiv0$. The originating theorem has smooth forcing
and starts from rest; an unforced blow-up claim must not keep
zero initial data under the above regularity assumptions.
\paragraph{Release boundary.}
Exact symbolic, generalized Vandermonde, 10,000 rational and
eight high-precision actual heat-kernel checks passed.
The signed heat deficit, three-bump restoration and no-go
constraint are source-compatible or classical, not a new
global blow-up theorem. The fully frozen corrected source,
complete cone, Appendix B, unforced correction hierarchy
and zero-force physical PDE remain OPEN/HOLD.
Lean verification covers only conditional scalar identities.

\subsection*{Proof Attack 3: exact first-order outer collar obstruction (Edition VIII)}
\label{subsec:n1-outer-collar-obstruction-VIII}
The original forced construction's positive-order similarity equations
(5.2)--(5.6) are linear in each new correction. The later radial
extension supplies a nonzero annular stress $T_1$; merely solving
the inner $n=1$ equations does not make the outer residual vanish.
At the terminal collar where the exterior heat cutoff equals one,
the source gives
\[
u_{\theta,0}=K(r,t)f_o(y),\quad
f_o(y)=1-\rho_o[1-\sigma((y-1)/2)],\quad1<y<3,
\]
with $\rho_o=c_oh>0$ and $K(r,t)>0$, and the only remaining first-order
exterior angular term (with the source's compact-support schedule)
is $-\partial_{zz}u_{\theta,0}$. Its exact physical formula is
\[
\partial_{zz}(Kf_o)=Kq^{-2D}
[a_z(\eta)^2f_o''(y)+b_z(\eta)f_o'(y)],
\]
\[
D=\tfrac12-h,\ d=1-\eta^2,\ L=1-2h\eta^2,\quad
a_z=-2\eta/L,\quad
b_z=(-2D\eta a_z+d\partial_\eta a_z)/L.
\]
\begin{proposition}[Strict first-order slow-field collar residual]
With the genuine source smooth-step
$\sigma(s)=1/(1+\exp[s^{-2}-(1-s)^{-2}])$, on the symmetry plane
$\eta=0$ at the collar midpoint $y=2$,
\begin{equation}\label{eq:n1-outer-collar-obstruction-VIII}
\boxed{-\partial_{zz}(Kf_o)=8Kq^{-2D}\rho_o>0.}
\end{equation}
Thus the prescribed slow-field coefficient sequence cannot satisfy
the unforced $n=1$ angular residual on this entire collar if its
order-one annular stress is simply suppressed and the other source
support assumptions are unchanged.
\end{proposition}
\begin{proof}
At $\eta=0$, $a_z=0$ and $b_z=-2$.
At $s=1/2$, $\sigma=1/2$ and its logistic defining exponent
has derivative $32$, so $\sigma'(1/2)=8$.
Therefore $f_o'(2)=(\rho_o/2)8=4\rho_o$ and the displayed
physical chain rule gives $\partial_{zz}(Kf_o)
=-8Kq^{-2D}\rho_o$, strictly negative.
\end{proof}
\begin{corollary}[Nonzero source annular stress]
The source's equation
$(\partial_R+2/R)T_{1,\theta}=-r_{\theta,1}$
implies that $T_{1,\theta}$ is not identically zero on this
collar under the preceding source stage and support assumptions.
\end{corollary}
This is a local obstruction for the unchanged slow-field profile,
not a nonexistence proof for every unforced alternative.
An oscillatory stress-realization, a redesigned exterior,
or a new order-one correction supported on the collar may modify
the residual. The original paper explicitly retains oscillatory
stress and eventually smooth external forcing. A genuine unforced
theorem would need to prove total force-free cancellation for the
physical reconstructed velocity and nonzero initial data.

\paragraph{Independent proof audit.}
The companion symbolic verifier checks $a_z(0)=0$,
$b_z(0)=-2$, $\sigma'(1/2)=8$, the strict midpoint formula,
and 18 physical-coordinate finite-difference comparisons with
maximum relative error below $10^{-5}$.
Setting $\rho_o=0$ makes the obstruction vanish, a negative control.
The finite Lean candidate proves only exact scalar and positivity
consequences, not the source's full physical-coordinate calculus
or the Navier--Stokes PDE. Kernel compilation remains a separate
GitHub CI gate. The full unforced theorem is OPEN.

\subsection*{Proof Attack 4 (Edition IX): swirl-only collar compatibility and linear Stokes repair}
\label{subsec:n1-swirl-compatibility-edition-IX}
Attack 3 showed that the original terminal collar retains the nonzero
first-order azimuthal viscous residual
$-\partial_{zz}(Kf_o)|_{\eta=0,y=2}
=8Kq^{-2D}\rho_o>0$. We now test a different, limited
repair without promoting it to a solution of the unforced PDE.

\begin{theorem}[Frozen-meridional swirl compatibility no-go]
Let $(u_r,U,u_z)$ be a smooth axisymmetric velocity whose radial
and axial Navier--Stokes momentum equations already vanish on a
connected open region with $r>0$. Replace the azimuthal velocity
$U$ by $U+W$ without changing $u_r,u_z$, and allow a pressure
change $\delta p$. Then the two meridional equations require
\[
\partial_r\delta p=\frac{2UW+W^2}{r},
\qquad\partial_z\delta p=0,
\quad\Longrightarrow\quad
\partial_z\frac{2UW+W^2}{r}=0.
\]
If, for a fixed $(r,t)$, $U(r,z,t)>0$ for all $z$ and
$W(r,\cdot,t)$ is continuous with compact axial support,
these equations force $W(r,z,t)\equiv0$. The same holds on
a connected axial interval containing a zero of $W$,
provided $U>0$ on the interval.
\end{theorem}
\begin{proof}
The swirl contribution to radial convection is $-U^2/r$,
while the axial momentum equation has no swirl term.
Equality of mixed pressure partial derivatives proves the
compatibility equation. Compact support yields a point with
$W=0$; the bracket is independent of $z$ and hence vanishes.
Thus $W(2U+W)=0$. Since $U>0$, the continuous ratio $W/U$
takes values in $\{0,-2\}$ and equals zero at the exterior
point. By connectedness it is identically zero.
\end{proof}
This theorem assumes the original meridional residuals
already vanish. It is NOT a theorem about all modifications
of the source collar; allowing nonzero corrections to $u_r$
or $u_z$, or altering the leading profile, changes the question.

\begin{proposition}[Exact linear solenoidal response]
Let $v_0=U_0(r,z,t)e_\theta$ be smooth axisymmetric swirl and
$\chi$ a smooth axisymmetric cutoff. Put
$g=\nu\chi\partial_{zz}v_0\in L^1([0,T];L^2(\mathbb R^3))$.
Then
\[
w(t)=\int_0^t e^{\nu(t-s)\Delta}g(s)\,ds
\]
is axisymmetric, purely azimuthal, divergence-free, and solves
$\partial_tw-\nu\Delta w=g$, $w(0)=0$, in the mild sense,
with
\[
\|w(t)\|_2\le\int_0^t\|g(s)\|_2\,ds,\qquad
\|\nabla w(t)\|_2\le
\int_0^t\frac{\|g(s)\|_2}{\sqrt{2e\nu(t-s)}}\,ds
\]
when the latter is finite.
\end{proposition}
\begin{proof}
An axisymmetric azimuthal vector field has zero divergence.
The vector heat semigroup commutes with divergence and rotations,
and preserves the azimuthal axisymmetric subspace. Apply the
Duhamel identity and the exact Fourier multiplier maxima
$\sup_{\xi}e^{-\nu\tau|\xi|^2}=1$ and
$\sup_\xi|\xi|e^{-\nu\tau|\xi|^2}=(2e\nu\tau)^{-1/2}$.
\end{proof}
The linear repair does not in general solve the nonlinear PDE:
the projected residual still contains
$\mathbb P[(u_0\cdot\nabla)w+(w\cdot\nabla)u_0+(w\cdot\nabla)w]$
and any uncancelled background residual. It need not retain
compact radial/axial support.

As a known positive control, the exact decaying shear
$u=Ae^{-\nu k^2t}\sin(kz)e_x$ on a periodic domain solves
unforced Navier--Stokes with zero pressure; it is obtained
by correcting stationary $A\sin(kz)e_x$. This is not blowup.
An independent negative control is
$U=re^{-r^2-z^2}$, $W=\varepsilon rz e^{-r^2-z^2}$:
for $F=(2UW+W^2)/r$,
$\partial_z F(r,0)=2\varepsilon r e^{-2r^2}\ne0$ when
$\varepsilon\ne0$, so no pressure with $\partial_z\delta p=0$
can compensate on a neighborhood of that point.

\paragraph{Release boundary.}
The source-specific coupled annular $n=1$ force-free correction,
meridional and pressure equations, all later corrections,
physical zero-force residual and all-order convergence are
OPEN/HOLD. Symbolic verification establishes the elementary
axisymmetric algebra, the negative control and the known shear
positive control. A Lean file tests finite scalar consequences
only, not the full cylindrical PDE or heat-semigroup theory.

\subsection*{Proof Attack 5 (Edition X): compulsory coupled meridional vorticity at order one}
\label{subsec:n1-coupled-vorticity-edition-X}
This is an exact diagnostic for an axisymmetric \emph{pure-swirl}
leading field $u_0=U_0(r,z,t)e_\theta$, and is not a proof that
the originating exterior source is itself a force-free solution.
Put $\mathcal L=\partial_{rr}+r^{-1}\partial_r+\partial_{zz}$,
$\mathcal L_\theta=\mathcal L-r^{-2}$ and perturb by
$V_r e_r+W e_\theta+V_z e_z$, with pressure increment $\pi$.
Denote the residual of the \emph{uncorrected} leading field
by $R^0=(R^0_r,R^0_\theta,R^0_z)$.
Direct substitution into unforced cylindrical Navier--Stokes yields
\begin{align*}
(\partial_t-\nu\mathcal L_\theta)V_r+\pi_r
&=(2U_0W+W^2)/r-V_rV_{r,r}-V_zV_{r,z}-R^0_r,\\
(\partial_t-\nu\mathcal L)V_z+\pi_z
&=-V_rV_{z,r}-V_zV_{z,z}-R^0_z,\\
(\partial_t-\nu\mathcal L_\theta)W
+V_r(U_{0,r}+U_0/r)+V_zU_{0,z}
&=-V_rW_r-V_zW_z-V_rW/r-R^0_\theta,\\
(rV_r)_r+rV_{z,z}&=0.
\end{align*}
At linear order in the perturbation, \emph{if} $R^0$ vanishes
at that order, the first-order meridional vorticity
$\omega_\theta=V_{r,z}-V_{z,r}$ obeys
\begin{equation}\label{eq:n1-coupled-vorticity-X}
\boxed{(\partial_t-\nu\mathcal L_\theta)\omega_\theta
       =\frac2r\partial_z(U_0W),}
\end{equation}
and the pressure obeys
\begin{equation}\label{eq:n1-pressure-poisson-X}
\boxed{\mathcal L\pi=\frac2r\partial_r(U_0W).}
\end{equation}
These follow by taking the meridional curl and divergence,
using the exact vector-Laplacian commutator and solenoidality.
If the leading profile carries a nonzero source residual, its curl
and divergence must appear on the right side at the correct order,
rather than being silently discarded.

With $\psi$ the meridional streamfunction,
$V_r=-\psi_z/r$, $V_z=\psi_r/r$, and
$\omega_\theta=-[\psi_{rr}-\psi_r/r+\psi_{zz}]/r$.
This makes explicit the additional elliptic axis and matching
boundary problem. In the homogeneous linearized system, for a smooth
base and vanishing energy boundary fluxes, writing
$\mathsf E=\tfrac12(\|V\|_2^2+\|W\|_2^2)$ gives
\[
\mathsf E'(t)+\nu\mathsf D(t)
=\int W[V_r(U_0/r-U_{0,r})-V_zU_{0,z}]\,dx,
\quad\mathsf D\ge0.
\]
With $M(t)=\|(U_0/r-U_{0,r},-U_{0,z})\|_\infty$,
$\mathsf E'\le M\mathsf E$; homogeneous data $\mathsf E(0)=0$
therefore imply $V=W=0$. This is a conditional standard energy
estimate and does not prohibit \emph{inhomogeneous} collar correction.

Independent SymPy replay passed the cylindrical divergence,
streamfunction, Laplacian-commutator, pressure-Poisson, curl
and energy cancellation identities for explicit Gaussian manufactured
fields, together with 10,000 exact rational Young-inequality tests.
The manufactured field has $\partial_z(U_0W)|_{z=0}
=r^2e^{-2r^2-t}\ne0$, a constructive coupling positive control.
No source-specific true order-one force-free correction,
nonlinear fixed point, all-order convergence, or physical
zero-forcing theorem has been established.
The associated Lean file verifies finite algebraic scalar
consequences only, subject to actual kernel compilation.

\subsection*{Proof Attack 6 (Edition XI, 23 October 2026): pressure-free collar witness and necessary correction budget}
\label{subsec:n1-pressure-free-dual-witness-XI}
The prescribed source's terminal-collar residual computed in Edition VIII
obeys, at a physical symmetry-plane point $r=r_0>0,z=0,t=t_0<1$,
\[
R_{\theta,1}(r_0,0,t_0)
=8K(r_0,t_0)q(t_0)^{-2D}\rho_o=:R_*>0,
\qquad \rho_o=c_oh>0.
\]
This is an exact local residual under the original slow-field
support and coefficient decomposition; it is \emph{not} automatically
the complete residual of every modified velocity field.

\begin{proposition}[Nonvanishing solenoidal collar witness]
Suppose $R_{\theta,1}(r,z,t_0)$ is continuous on a neighborhood of
$(r_0,0)$. There is a compactly supported smooth
nonnegative bump $\zeta(r,z)$, not identically zero and supported
away from $r=0$, such that $\Phi=\zeta(r,z)e_\theta$ satisfies
\[
\nabla\cdot\Phi=0,\qquad
\int_{\mathbb R^3}\nabla p\cdot\Phi\,dx=0
\quad\text{for every sufficiently regular pressure }p,
\]
but
\begin{equation}\label{eq:XI-positive-pairing}
B_\Phi=\int_{\mathbb R^3}R_1\cdot\Phi\,dx
=2\pi\int R_{\theta,1}\zeta r\,dr\,dz>0.
\end{equation}
In particular no pressure gradient can cancel the specified
source first-order residual on that neighborhood.
\end{proposition}
\begin{proof}
Continuity permits an annular neighborhood where
$R_{\theta,1}\ge R_*/2>0$. Choose $\zeta\ge0$ with
nonempty interior support there. An axisymmetric azimuthal field
is divergence-free. The pressure pairing vanishes by integration
by parts; the residual pairing is strictly positive.
\end{proof}
The corresponding Leray-projected residual has the strict
$H^{-1}$ lower bound
\[
\|\mathbb P R_1(t_0)\|_{H^{-1}}
\ge B_\Phi/\|\Phi\|_{H^1}>0,\qquad
\|\Phi\|_{H^1}^2=2\pi\int
(\zeta^2+\zeta_r^2+\zeta_z^2+\zeta^2/r^2)\,r\,dr\,dz.
\]
Thus the obstruction is not a pressure-gauge artifact.

\begin{proposition}[Exact necessary nonlinear correction budget]
For a smooth divergence-free background $u_B$ of physical unforced
momentum residual
$R_B=\partial_tu_B-\nu\Delta u_B+
(u_B\cdot\nabla)u_B+\nabla p_B$,
every truly unforced correction $u=u_B+w$ with $\nabla\cdot w=0$
must obey the exact equation
\[
\partial_tw-\nu\Delta w
+(u_B\cdot\nabla)w+(w\cdot\nabla)u_B
+(w\cdot\nabla)w+\nabla\pi=-R_B.
\]
For any compact smooth solenoidal test $\Phi$, assuming all norms exist,
\begin{align*}
|\langle R_B,\Phi\rangle|
\le{}&
\|\partial_tw\|_{H^{-1}}\|\Phi\|_{H^1}
+\nu\|\nabla w\|_2\|\nabla\Phi\|_2\\
&+\|u_B\|_\infty\|w\|_2\|\nabla\Phi\|_2
+\|\nabla u_B\|_\infty\|w\|_2\|\Phi\|_2
+\|w\|_4^2\|\nabla\Phi\|_2.
\end{align*}
\end{proposition}
\begin{proof}
Subtract the two Navier--Stokes systems, test by $\Phi$,
and use divergence-free integration by parts, Holder's inequality
and $H^{-1}$--$H^1$ duality. Pressure drops out exactly.
\end{proof}
The numerical test uses a source-\emph{shaped} local fixture
$K(r)=re^{-r^2}$ solely as a positive control, not the
actual source amplitude. It independently checks Cartesian
divergence, the smooth-step midpoint, a positive integral pairing,
121 sample points and 10,000 exact rational positivity controls.
The actual source physical residual must still be authenticated
globally, and the missing matched coupled order-one corrector,
nonzero initial data, all-order convergence and zero-force
Navier--Stokes identities remain OPEN/HOLD.
The Lean module only addresses finite scalar consequences, not
the continuum witness or the full PDE.

\subsection*{Proof Attack 7 (Edition XII, 23 October 2026): a nonflat source-correction startup theorem}
\label{subsec:n1-nonflat-start-XII}
Let $u_B$ be any smooth divergence-free background with its
\emph{full physical} unforced momentum residual
\[
R_B=\partial_tu_B-\nu\Delta u_B+
(u_B\cdot\nabla)u_B+\nabla p_B.
\]
For a genuinely force-free corrected velocity $u=u_B+w$,
the exact nonlinear perturbation identity is
\begin{equation}\label{eq:XII-full-correction}
\partial_tw-\nu\Delta w+(u_B\cdot\nabla)w+
(w\cdot\nabla)u_B+(w\cdot\nabla)w+\nabla\pi=-R_B,
\qquad \nabla\cdot w=0.
\end{equation}
\begin{theorem}[Nonflat response to an active solenoidal residual]
Fix $t_0$ and suppose a compactly supported smooth
divergence-free $\Phi$ satisfies
$B_\Phi:=\langle R_B(t_0),\Phi\rangle>0$.
If $w(t_0,\cdot)=0$ and the corrected field solves the
unforced equations with sufficient regularity, then
\begin{equation}\label{eq:XII-no-flat-start}
\boxed{\langle\partial_tw(t_0),\Phi\rangle=-B_\Phi<0.}
\end{equation}
If $w:[t_0,t_0+\epsilon)\to L^2$ is $C^1$,
there exists $\delta>0$ such that
\begin{equation}\label{eq:XII-quantitative-start}
\boxed{\|w(t)\|_2\ge\frac{B_\Phi}{2\|\Phi\|_2}(t-t_0),
\qquad t_0<t<t_0+\delta.}
\end{equation}
Thus an order-one correction that is flat in time at
$t_0$ cannot cancel an active source residual.
\end{theorem}
\begin{proof}
At $t_0$, $w$ and all its spatial derivatives vanish.
Pair \eqref{eq:XII-full-correction} with $\Phi$;
all convection and viscosity terms vanish at $t_0$.
Pressure vanishes by solenoidality and integration by parts.
This yields \eqref{eq:XII-no-flat-start}.
Set $F(t)=\langle w(t),\Phi\rangle$; then
$F(t_0)=0$ and $F'(t_0)=-B_\Phi$. By continuity of
$F'$, $F'(t)\le-B_\Phi/2$ for short positive times.
Integrate and use Cauchy--Schwarz to prove
\eqref{eq:XII-quantitative-start}.
\end{proof}
\paragraph{Application to the terminal collar.}
Edition XI found an exactly positive azimuthal first-order
source residual at the collar midpoint and, by continuity,
a compactly supported solenoidal witness $\Phi=\zeta e_\theta$.
\emph{If} that prescribed source residual remains the complete
pressure-free order-one residual of the final selected background
at a matching time $t_0$, the above theorem forces a nonzero
instantaneous correction derivative. This does not prove
the existence or nonexistence of a fully matched correction.
No appeal to pressure adjustment can remove the witness.

\begin{proposition}[Exact manufactured smooth positive control]
For the globally smooth finite-energy divergence-free
background $u_B=r e^{-r^2-z^2}e_\theta$, $p_B=0$, $\nu=1$,
\[
R_{B,\theta}=
2r(5-2r^2-2z^2)e^{-r^2-z^2}.
\]
On $1/2\le r\le1$, $|z|\le1/2$ this is
at least $(5/2)e^{-5/4}>0$.
With
$\zeta=(r-1/2)^4(1-r)^4(1/4-z^2)^4$
on this rectangle, zero outside, the axisymmetric test
$\Phi=\zeta e_\theta$ is divergence-free and has
\[
\int\zeta r\,dr\,dz=\frac{1}{270950400},\qquad
\boxed{B_\Phi\ge
\frac{5\pi e^{-5/4}}{270950400}>0.}
\]
The displayed polynomial bump is $C^3$, sufficient for the
weak test; it may be approximated by a positive $C_c^\infty$
bump without losing strict positivity.
\end{proposition}
\paragraph{Adversarial global-neutrality control.}
A local positive residual is not a global net-torque
obstruction: for $g(z)=e^{-z^2}$,
$-g''(0)=2>0$ but
$\int_{\mathbb R}-g''(z)\,dz=0$ exactly.
Thus source-specific matching cannot be excluded using a
local sign alone.

\paragraph{Verification boundary.}
Independent symbolic Gaussian vector-Laplacian and
divergence tests, exact rational support mass, high-precision
positive pairing with a strict analytic lower bound, a
zero-torque negative control and $10{,}000$ exact-rational
scalar checks passed. The true source-specific coupled
inhomogeneous annular problem, physical residual authentication,
all-order convergence, nonzero smooth data and complete
unforced Navier--Stokes theorem remain OPEN/HOLD.
The companion Lean candidate formalizes \emph{finite scalar
consequences only}, not the field-theoretic arguments.

\subsection*{Proof Attack 8 (Edition XIII): exact first projected-Stokes correction and its nonlinear residual}
\label{subsec:n1-projected-stokes-XIII}
A first-order correction must cancel the \emph{whole physical} momentum residual, not just the azimuthal viscous channel. The following theorem is a standard local semigroup construction carried out on the full space; it is not an authentic globally matched correction for the prescribed source exterior.
Let $t_0<T$, $\nu>0$, and let $u_B$ be a smooth divergence-free background in $\mathbb R^3$, with a chosen background pressure $p_B$. Define
\[
R_B=\partial_tu_B-\nu\Delta u_B+(u_B\cdot\nabla)u_B+\nabla p_B,
\qquad \mathbb P=I-\nabla\Delta^{-1}\nabla\cdot.
\]
The inverse Laplacian is the standard free-space Fourier multiplier; the projected field is independent of the arbitrary additive constant in pressure.
\begin{theorem}[Exact first-step correction and its residual]
Assume $R_B\in C([t_0,T];H^1(\mathbb R^3))$, $u_B$ is divergence-free and $u_B,\nabla u_B$ are bounded. Define
\begin{equation}\label{eq:XIII-first-corrector}
 w_1(t)=-\int_{t_0}^t e^{\nu(t-s)\Delta}\mathbb P R_B(s)\,ds,
\quad \delta p_1=-\Delta^{-1}\nabla\cdot R_B.
\end{equation}
Then $w_1$ is divergence-free, $w_1(t_0)=0$, and
\begin{equation}\label{eq:XIII-corrected-jet}
(\partial_t-\nu\Delta)w_1=-\mathbb P R_B,
\qquad \partial_tw_1(t_0)=-\mathbb P R_B(t_0).
\end{equation}
The unforced Navier--Stokes momentum residual of the \emph{corrected} pair $(u_B+w_1,p_B+\delta p_1)$ is exactly
\begin{equation}\label{eq:XIII-exact-leftover}
\boxed{\mathcal R_1=(u_B\cdot\nabla)w_1+(w_1\cdot\nabla)u_B+(w_1\cdot\nabla)w_1.}
\end{equation}
Set $a=t-t_0$, $M=\sup_{s\in[t_0,T]}\|R_B(s)\|_{H^1}$, and $U=\sup_s(\|u_B(s)\|_{\infty}+\|\nabla u_B(s)\|_{\infty})$. For an absolute $H^1(\mathbb R^3)\hookrightarrow L^4$ embedding constant $C_4$,
\begin{equation}\label{eq:XIII-time-remainder-bound}
\boxed{\|w_1(t)\|_{H^1}\le aM,\qquad
\|\mathcal R_1(t)\|_{H^{-1}}\le UaM+C_4^2a^2M^2.}
\end{equation}
For each compact solenoidal test field $\Phi$, the startup response automatically meets the Edition XII condition
$\langle\partial_tw_1(t_0),\Phi\rangle=-\langle R_B(t_0),\Phi\rangle$.
\end{theorem}
\begin{proof}
The Leray projector and heat semigroup commute, are contractions on $H^1$, and preserve divergence-free fields. Duhamel differentiation proves \eqref{eq:XIII-corrected-jet}, including its right derivative at $t_0$. Since $R_B=\mathbb PR_B+\nabla\Delta^{-1}\nabla\cdot R_B$ and $\nabla\delta p_1=-(I-\mathbb P)R_B$, expanding the corrected momentum equation cancels \emph{all} of $R_B$ against the linear Stokes response and the pressure adjustment, leaving exactly \eqref{eq:XIII-exact-leftover}. Contraction yields $\|w_1\|_{H^1}\le aM$. For an $H^1$ test $\Phi$, integrate $u_B\cdot\nabla w_1$ and $w_1\cdot\nabla w_1$ by parts, using solenoidality. Their pairings are bounded by $\|u_B\|_\infty\|w_1\|_2\|\nabla\Phi\|_2$ and $\|w_1\|_4^2\|\nabla\Phi\|_2$ respectively. The remaining cross term is bounded by $\|\nabla u_B\|_\infty\|w_1\|_2\|\Phi\|_2$. Sobolev embedding proves \eqref{eq:XIII-time-remainder-bound}. At startup the pressure gradient is orthogonal to every solenoidal test, proving the final identity.
\end{proof}
\begin{proposition}[The meridional response is already compulsory for smooth Gaussian swirl]
For the stationary, smooth finite-energy, unforced-equation \emph{non-solution} $u_B=Ue_\theta$, $U=r\exp[-(r^2+z^2)]$, with $p_B=0$ and $\nu=1$, the \emph{complete} physical residual has
\[
R_{B,r}=-r e^{-2(r^2+z^2)},\quad R_{B,z}=0,\quad
R_{B,\theta}=2r(5-2r^2-2z^2)e^{-(r^2+z^2)}.
\]
In particular,
\begin{equation}\label{eq:XIII-gaussian-curl}
\boxed{\partial_zR_{B,r}-\partial_rR_{B,z}
=4rz\,e^{-2(r^2+z^2)}\ne0\quad(r>0,\ z\ne0).}
\end{equation}
As a gradient has zero curl, $\mathbb PR_B$ has a nontrivial meridional component. Therefore the exact first response $w_{1,t}(t_0)=-\mathbb PR_B(t_0)$ cannot be a purely azimuthal correction.
\end{proposition}
\paragraph{Exact global positive control.} On a periodic domain, the stationary shear $u_B=A\sin(kz)e_x$ has $R_B=\nu k^2u_B$ and its projected first step gives $w_1=(e^{-\nu k^2(t-t_0)}-1)u_B$. Here all quadratic convection terms vanish, and the \emph{fully corrected} field $u=Ae^{-\nu k^2(t-t_0)}\sin(kz)e_x$ is an exact smooth \emph{unforced} Navier--Stokes solution. This is well-known and not a blow-up construction.
\paragraph{Scope and outstanding theorem.} The theorem requires uniform $H^1$ residual and bounded smooth source data over a short physical time interval, with no imposed radial/terminal compact-support or prescribed similarity/oscillatory matching. The projected heat semigroup is nonlocal and may destroy strict annular localization. Its remainder \eqref{eq:XIII-exact-leftover} is generally nonzero. Proving a \emph{second nonlinear correction}, original source-stage matching, a terminal blow-up rate, nonzero smooth initial data, complete zero-force physical NSE reconstruction, and full Lean formalization remains OPEN/HOLD. The companion Lean checks are finite scalar implications only.

\subsection*{Proof Attack 9 (Edition XIV): nonlinear projected Picard closure and terminal scaling}
\label{subsec:n1-XIV-picard}
Let $u_B\in C([t_0,t_0+T];H^3_\sigma(\mathbb R^3))$ be a
smooth solenoidal background, with its \emph{complete physical}
momentum residual $R_B=\partial_tu_B-\nu\Delta u_B+
(u_B\cdot\nabla)u_B+\nabla p_B\in C_tH^3$.
Define
$A=\sup\|u_B\|_{H^3}$, $F=\sup\|\mathbb PR_B\|_{H^3}$.
Choose a valid tensor-product constant $C_3\ge1$ with
$\|u\otimes v\|_{H^3}\le C_3\|u\|_{H^3}\|v\|_{H^3}$.
For the free-space Leray projection $\mathbb P$,
\[
\|e^{\nu\tau\Delta}\mathbb P\nabla\cdot Q\|_{H^3}
\le(2e\nu\tau)^{-1/2}\|Q\|_{H^3}.
\]
Set $M=2TF$ and $\gamma=2C_3\sqrt{T/(2e\nu)}$.
\begin{theorem}[Exact short-time nonlinear unforced closure]
If
\begin{equation}\label{eq:XIV-contraction-condition}
\boxed{\gamma(2A+2M)=
2C_3\sqrt{\frac{T}{2e\nu}}(2A+4TF)\le\frac12,}
\end{equation}
then the projected mild map
\[
\mathcal Tw(t)=
-\int_{t_0}^te^{\nu(t-s)\Delta}\mathbb P R_B(s)\,ds
-\int_{t_0}^te^{\nu(t-s)\Delta}\mathbb P\nabla\cdot
(u_B\otimes w+w\otimes u_B+w\otimes w)(s)\,ds
\]
has a unique fixed point in the ball
$\|w\|_{C_tH^3}\le M$, with $w(t_0)=0$.
The field $u=u_B+w$ solves \emph{the full unforced}
Navier--Stokes equations in the mild sense on this interval,
with initial datum $u_B(t_0)$, after reconstructing
\[
\nabla\pi=-(I-\mathbb P)
[R_B+\nabla\cdot(u_B\otimes w+w\otimes u_B+w\otimes w)].
\]
For smooth compatible data this is classical for positive times.
This theorem guarantees no terminal blow-up or exterior matching.
\end{theorem}
\begin{proof}
The heat semigroup and $\mathbb P$ contract on $H^3$, and
$\int_0^T(2e\nu s)^{-1/2}ds=2\sqrt{T/(2e\nu)}$.
For $\|w\|\le M$,
$\|\mathcal Tw\|\le TF+\gamma(2AM+M^2)\le M$.
For $w,v$ in the same ball,
$\|\mathcal Tw-\mathcal Tv\|
\le\gamma(2A+2M)\|w-v\|\le\frac12\|w-v\|$.
Banach's theorem yields the fixed point. The complementary
gradient projection determines $\pi$ and eliminates every
remaining physical momentum-residual component.
\end{proof}
\begin{proposition}[Terminal-scale boundary of this sufficient test]
With $\varepsilon=1-t_0$, $T=\theta\varepsilon$
for fixed $0<\theta<1$, and authentic two-sided envelopes
$A_\varepsilon\asymp\varepsilon^{-\alpha}$,
$F_\varepsilon\asymp\varepsilon^{-\beta}$, the two
nonnegative contributions to \eqref{eq:XIV-contraction-condition}
scale as $\varepsilon^{1/2-\alpha}$ and
$\varepsilon^{3/2-\beta}$. For $\alpha>1/2$ or
$\beta>3/2$, this particular fixed-fraction sufficient
criterion eventually fails. Failure of a sufficient criterion
is \emph{not} nonexistence of an unforced solution.
\end{proposition}
The result is a quantitative statement of \emph{classical local
Navier--Stokes well-posedness}. The Leray projector and heat kernel
are nonlocal and do not certify retention of the prescribed compact
collar, pressure/axis matching or source symmetry. Even if
$\|u_B(t)\|_{H^3}\to\infty$, its singularity is not inherited without
a single matched correction with
$\|w(t)\|_{H^3}=o(\|u_B(t)\|_{H^3})$, or an equivalent
lower bound through $t\uparrow1$. No such result is proved.
Exact symbolic algebra, 10,000 rational contraction-envelope
tests, an intentional failed-gate negative control and a
known decaying-shear positive control passed. The Lean file
checks only scalar implications; infinite-dimensional semigroup
and PDE formalization remain OPEN/HOLD.

\subsection*{Proof Attack 10 (Edition XV): exact anisotropic $H^3$ scaling and a terminal contraction obstruction}
\label{subsec:XV-anisotropic-H3}
Write $\tau=1-t$, $D=\tfrac12-h$ and $A=\tfrac12+h$, $0<h<1/100$. The originating \emph{forced} construction uses
\[
\tau=q(1-\eta^2),\quad z=q^D\eta,\quad X=\frac{r^2}{2q},
\quad u_\theta^{(0)}=q^{-A}E(X,\eta).
\]
Its core has characteristic radial scale $\tau^{1/2}$, axial scale $\tau^D$ and nonzero swirl of size $\tau^{-A}$. Amplitude claims alone do not supply the derivative lower bound below.

\begin{proposition}[Exact rescaled-chart lower bound]
Write $R=(x_1,x_2)/\sqrt{\tau}$, $Z=z/\tau^D$. On a fixed compact similarity chart away from $|\eta|=1$,
\[
q=\tau Q(Z),\qquad Q-Z^2Q^{2h}=1,\qquad
\eta=Z/Q^D,\quad X=|R|^2/(2Q).
\]
The leading Cartesian swirl velocity consequently has form
$u_{\rm swirl}^{(0)}(x,t)=\tau^{-A}\mathcal U(R,Z)$ with $\mathcal U$ independent of $\tau$.
Suppose on a fixed compact chart $\mathcal K$ some third \emph{radial Cartesian} derivative obeys
$C_*^2=\int_{\mathcal K}|\partial_{R_i}^3\mathcal U_j|^2\,dR\,dZ>0$.
For any complete background agreeing with the leading field there,
\begin{equation}\label{eq:XV-aniso-H3}
\boxed{\|u_B(t)\|_{H^3(\mathbb R^3)}\ge C_*\tau^{-\alpha},
\qquad \alpha=A+\frac32-\frac{1+D}{2}
=\frac54+\frac{3h}{2}.}
\end{equation}
This cannot be promoted to a bound for an unknown corrected full source until its actual nondegeneracy and later modifications are verified.
\end{proposition}
\begin{proof}
The similarity identities yield $1=Q-Z^2Q^{2h}$. The spatial Jacobian is $\tau^{1+D}$ and three radial derivatives introduce $\tau^{-3/2}$. The chosen squared derivative integral is therefore exactly
$C_*^2\tau^{-2A-3+1+D}$.
\end{proof}

\begin{proposition}[Fixed-fraction Picard gate cannot remain certified]
Let $T=\theta\tau$, $0<\theta<1$ fixed. For the Edition XIV sufficient contraction gate, let
$A_{H^3}=\sup_{[t_0,t_0+T]}\|u_B\|_{H^3}$ and $F\ge0$ its projected residual bound. Under the preceding nondegeneracy assumption,
\[
2C_3\sqrt{\frac{T}{2e\nu}}(2A_{H^3}+4TF)
\ge4C_3C_*\sqrt{\frac{\theta}{2e\nu}}\,
\tau^{-3/4-3h/2}\longrightarrow\infty.
\]
Thus the sufficient condition that the left side be at most $1/2$ fails on late fixed-fraction intervals. The \emph{background contribution alone} forces failure, irrespective of $F$.
For adaptive $T=c\tau^p$, the background term tends to zero if $p>2\alpha=\frac52+3h$ (at equality a small enough constant may suffice). This is a failure of one proof technique, not a nonexistence result for force-free corrections.
\end{proposition}

\begin{proposition}[Explicit divergence-free Gaussian model and raw residual]
On $\mathbb R^3$, write $(X,Y,Z)=(x_1/\sqrt{\tau},x_2/\sqrt{\tau},z/\tau^D)$ and set
\[
u_G=\tau^{-A}(-Y,X,0)e^{-(X^2+Y^2+Z^2)},\qquad p_G=0.
\]
This is smooth for every $t<1$, finite energy, axisymmetric and solenoidal. Its \emph{exact} third-derivative norm is
\[
\|\partial_{x_1}^3(u_G)_2\|_2^2
=\frac{105}{4}\left(\frac\pi2\right)^{3/2}\tau^{-2\alpha}.
\]
For the complete \emph{unprojected} physical Navier--Stokes residual $R_G$, the radial centrifugal component gives
\[
\|\partial_{x_1}^3(R_G)_1\|_2^2
=105\left(\frac\pi4\right)^{3/2}\tau^{-2\beta},
\qquad \beta=\frac94+\frac{5h}{2}.
\]
These are model identities, not authenticated bounds for the source's \emph{Leray-projected} residual.
\end{proposition}
\begin{proof}
The third derivative of $Xe^{-aX^2}$ equals
$(-6a+24a^2X^2-8a^3X^4)e^{-aX^2}$.
The squared Gaussian polynomial expectations equal $105/4$ for $a=1$ and $105$ for $a=2$. The centrifugal component is
$(R_G)_1=-\tau^{-2A-1/2}X e^{-2(X^2+Y^2+Z^2)}$.
Apply the exact anisotropic Jacobian.
\end{proof}
\paragraph{Independent tests and HOLD boundary.}
The companion exact verifier replays the source-chart exponent algebra, both Gaussian constants, and $10{,}000$ rational choices of $0<h<1/100$, all passing. The actual source-profile positive derivative constant $C_*$, its complete time-dependent $H^3$ and projected-residual constants, external matching, physical zero-forcing identity, survival of singularity after correction, and infinite-dimensional Lean theorem remain \textbf{OPEN/HOLD}. The originating OpenAI construction has nonzero external forcing.

\subsection*{Edition XVI: source-native third Cartesian jet of the analytic leading swirl (23 October 2026)}
\label{subsec:editionXVI-source-axis-jet}
Unlike the Gaussian surrogate in Edition XV, this calculation uses the originating forced construction's \emph{actual} Appendix B.1--B.3 axis profile. Let $0<h\le1/100$, $D=\frac12-h$, $A=\frac12+h$, and $j_0,\sigma_*,\Lambda,C>0$. Its prescribed data are
\[
U_*=4\eta+j_0,\quad H_*=D\eta+(1-\eta^2)U_*,
\quad W_*=1-(1-\eta^2)U_*'-2D\eta U_*,
\]
\[
L=1-2h\eta^2,\quad\chi=\frac{H_*^2}{H_*^2+\sigma_*^2},
\quad E(X,\eta)=\sqrt{2X}\,\varphi(X,\eta)/C,\quad
\varphi=\varphi_*(\eta)\Phi(\Lambda X,\eta).
\]
The exact Appendix B radial equation at $Y=0$ implies
\[
p_1(\eta)=\partial_Y\Phi(0,\eta)
=\frac14\left[-\chi(\eta)+
\frac{W_*(\eta)+h(1-2\eta U_*(\eta))}{\Lambda L(\eta)}\right].
\]
This same first coefficient was independently derived in the cumulative axis-jet section.

\begin{theorem}[Genuine source-leading swirl nondegeneracy]
Let $u_{\rm swirl}^{(0)}(x,t)=\tau^{-A}\mathcal U(x_1/\sqrt\tau,x_2/\sqrt\tau,z/\tau^D)$, $\tau=1-t$. Then the source's exact axis field obeys
\begin{equation}\label{eq:XVI-source-axis-jet}
\boxed{\partial_{R_1}^3\mathcal U_2(0,0,0)=
-\frac3{4C}\left(\Lambda\frac{j_0^2}{j_0^2+\sigma_*^2}+3-h\right)<0,}
\end{equation}
and in particular
\[
\boxed{|\partial_{R_1}^3\mathcal U_2(0)|\ge897/(400C)>0.}
\]
Consequently the actual analytic \emph{leading} profile has a fixed positive-volume compact chart $\mathcal K$ about the similarity axis with $C_*=\|\partial_{R_1}^3\mathcal U_2\|_{L^2(\mathcal K)}>0$.
\end{theorem}
\begin{proof}
At $\eta=0$ the source has $U_*(0)=H_*(0)=j_0$, $W_*(0)=1-4=-3$, $L(0)=1$, $\varphi_*(0)=1$. Hence $\Lambda p_1(0)=-(\Lambda j_0^2/(j_0^2+\sigma_*^2)+3-h)/4$. Along $Z=R_2=0$, the exact similarity variables give $q=\tau$, $X=R_1^2/2$, and
$\mathcal U_2(R_1,0,0)=R_1\Phi(\Lambda R_1^2/2,0)/C$.
The cubic coefficient is $\Lambda p_1(0)/(2C)$, whose third derivative is $3\Lambda p_1(0)/C$. Because $3-h\ge299/100$, the strict lower bound follows. Continuity gives the positive local $L^2$ coefficient.
\end{proof}

\begin{corollary}[Explicit missing modulus and final-correction gates]
If $M_4$ bounds the gradient of the third derivative on a similarity ball $B_{r_0}$, $a=897/(400C)$ and $0<\delta\le r_0$ satisfies $M_4\delta\le a/2$, then
\[
C_*\ge(a/2)\sqrt{4\pi\delta^3/3}.
\]
For the \emph{complete} velocity $u=u_{\rm swirl}^{(0)}+v$, define $V_\tau(R,Z)=\tau^A v(\sqrt\tau R,\tau^D Z,1-\tau)$. If
$\|\partial_{R_1}^3(V_\tau)_2\|_{L^2(\mathcal K)}\le C_*/2$,
the triangular inequality yields
\[
\|u(t)\|_{H^3}\ge(C_*/2)\tau^{-5/4-3h/2}.
\]
The needed $M_4$ and \emph{full oscillatory/higher-order correction} estimate are not certified here.
\end{corollary}
\paragraph{Verification and scope.}
Independent symbolic Cartesian Taylor calculations, 10,000 exact-rational source-parameter sign/lower comparisons, and 75-digit diagnostic third finite differences passed. This establishes the missing \emph{leading-profile} nondegeneracy assumption from Edition XV. It is a deduction from prior-source Appendix B equations, not a new singularity theorem. Actual source-admissible parameter calibration, later correction stability, the global physical zero-forcing equation and full Lean proof are still \textbf{OPEN/HOLD}.

\subsection*{Edition XVII: effective Cauchy-radius certificate for the genuine source leading third jet}
\label{subsec:XVII-cauchy}
Edition XVI already proves from the source Appendix B axis equation that
$f(0):=\partial_{R_1}^3\mathcal U_2(0)$ satisfies
$|f(0)|\ge a:=897/(400C)>0$.
Assume $f$ has a holomorphic extension to the polydisc
$|Z_j|<\rho$ ($j=1,2,3$) and
$\sup|f|\le B$, $B\ge a$.
Put $\delta=\min\{\rho/2,a\rho/(12B)\}$.
By the one-variable Cauchy derivative bound on each coordinate line
within the polydisc, for $R\in[-\delta,\delta]^3$,
\[
|f(R)-f(0)|\le\sum_{j=1}^3\frac{B|R_j|}{\rho-\delta}
\le\frac{6B\delta}{\rho}\le a/2.
\]
Consequently
\[
\boxed{\|\partial_{R_1}^3\mathcal U_2\|_{L^2([-\delta,\delta]^3)}
\ge (a/2)(2\delta)^{3/2}>0.}
\]
The proof is conditional on a source-valid holomorphic radius
$\rho$ and bound $B$; neither is presently certified for a
frozen parameter tuple. Nor is the third-derivative norm of
all later oscillatory and higher-order corrections on this cube.
This converts the existential neighborhood in Edition XVI into
an explicit \emph{certification target}; it is not a new
unforced Navier--Stokes singularity theorem.

\subsection*{Proof Attack 13 (Edition XVIII): source-exact exterior third radial derivative}
\label{subsec:editionXVIII-exterior-derivative}
Unlike the Gaussian surrogate in Edition XV, the original OpenAI
\emph{forced} source gives an exact exterior heat profile in
Theorem 4.6(v), equation (4.29). Put $\tau=1-t$,
$A=1/2+h$, $D=1/2-h$, $0<h<1/100$, $X_b<\infty$ and $c_\infty>0$.
For $X=r^2/(2q)\ge X_b$ the source asserts
\[
E(X,\eta)=c_\infty X^{-A}\mathcal H(2(1-\eta^2)/X),\quad
\mathcal H(Z)=\frac1{\Gamma(1+h)}
\int_0^\infty e^{-v}v^h(1+Zv)^{-h}dv.
\]
Using $\tau=q(1-\eta^2)$ yields the physical axially independent
exterior swirl
\[
u_\theta^{(0)}=c_\infty(r^2/2)^{-A}\mathcal H(4\tau/r^2).
\]
With $R=r/\sqrt\tau$ and $f(R)=\tau^Au_\theta^{(0)}$,
\begin{equation}\label{eq:XVIII-mixture}
\boxed{ f(R)=c_\infty2^A\int_0^\infty
\frac{e^{-v}v^h}{\Gamma(1+h)}
\frac{dv}{R(R^2+4v)^h}. }
\end{equation}
\begin{proposition}[Strict source-exterior third radial derivative]
For $a\ge0$, $h>0$, $R>0$ and $g_{a,h}(R)=R^{-1}(R^2+a)^{-h}$,
\[
-g'''_{a,h}(R)=\frac{P_h(a,R)}{R^4(R^2+a)^{h+3}}>0,
\]
\[
P_h=6a^3+(6h+18)a^2R^2+(12h+18)aR^4+
(1+2h)(2+2h)(3+2h)R^6.
\]
In particular $f'''(R)<0$ at every exterior rescaled radius.
\end{proposition}
\begin{proof}
Three direct differentiations give the displayed polynomial.
All its coefficients are positive, and the gamma density in
\eqref{eq:XVIII-mixture} is positive. Dominated differentiation
on fixed compact $R$ intervals gives strict negativity.
\end{proof}
\begin{proposition}[Explicit source-parameter $H^3$ lower bound]
Fix $\zeta=1/4$ and $b>a\ge2\sqrt{X_b}$, and put
\[
m_{a,b}=
\frac{c_\infty2^Ae^{-1}}{\Gamma(h+2)}
\frac{(1+2h)(2+2h)(3+2h)a^2}{(b^2+4)^{h+3}}>0.
\]
On the physical exterior slab
$a\le R=r/\sqrt\tau\le b$,
$|Z|=|z|/\tau^D\le\zeta$,
$-f'''(R)\ge m_{a,b}$. The genuine \emph{leading field}
therefore satisfies
\begin{equation}\label{eq:XVIII-source-H3-lower}
\boxed{\|u^{(0)}(t)\|_{H^3(\mathbb R^3)}
\ge m_{a,b}\sqrt{\frac{\pi(b^2-a^2)}{12}}\,
\tau^{-5/4-3h/2}.}
\end{equation}
\end{proposition}
\begin{proof}
The source relation for $Q=q/\tau$ is
$Q-Z^2Q^{2h}=1$; on $|Z|\le1/4$ it gives $1\le Q<2$.
Hence $X=R^2/(2Q)\ge R^2/4\ge X_b$.
Retain only the last positive term in $P_h$ and integrate
\[
\frac1{\Gamma(1+h)}\int_0^1e^{-v}v^h\,dv
\ge \frac{e^{-1}}{\Gamma(h+2)}.
\]
The radial third derivative of physical swirl is
$\tau^{-A-3/2}f'''(R)e_\theta$; it is a unit-directional
contraction of the third Cartesian derivative tensor.
Its square is at most six times the standard sum of
third-order Cartesian derivative squares. Integrate over
the slab using $r\,dr\,d\theta\,dz=
\tau^{1+D}R\,dR\,d\theta\,dZ$ to obtain
\eqref{eq:XVIII-source-H3-lower}.
\end{proof}
This removes the derivative-nondegeneracy \emph{premise} of
Edition XV for the actual \emph{leading exterior field}.
No numerical freeze of $h,c_\infty,X_b$ is asserted;
the illustrative $h=1/200,c_\infty=1,X_b=2,a=4,b=5$
gives $m\approx0.00205353$ and $C_{\rm ext}\approx0.00315216$
only as a mathematical example, not authentic source constants.
The positive exterior lower bound makes Edition XIV's
fixed-fraction \emph{sufficient} isotropic $H^3$ Picard gate fail
near $\tau=0$ for the unmodified leading exterior profile.
It does not prove that later source oscillations, cutoffs, a
new unforced correction or the final physical field preserve
this bound. The complete Leray-projected residual and
force-free PDE remain OPEN/HOLD, as does the end-to-end Lean proof.
Exact symbolic polynomial, 10,000 rational sign tests, and
independent high-precision gamma-integral checks passed.
The candidate Lean file formalizes only finite scalar consequences.


\section{A second-order Leray correction and a meridional dual witness}
\label{sec:second-order-leray}
This section concerns a smooth background and does not establish unforced
finite-time blowup. Its short-time scheme is an explicit instance of standard
Leray--Stokes perturbation theory, not an originality claim.
Let $s>5/2$, $\nu>0$, and let $u_B$ be a smooth divergence-free field on
$\mathbb T^3$, with complete residual
\[
R_B=\partial_tu_B-\nu\Delta u_B+(u_B\cdot\nabla)u_B+\nabla p_B.
\]
Let $\mathbb P$ denote the Leray projection,
$S(t)=e^{\nu t\Delta}$ and
$B(a,b)=\mathbb P[(a\cdot\nabla)b]$.
For a slab $[t_0,t_0+T]$ put
$A=\sup_t\|u_B(t)\|_{H^s}$,
$F=\sup_t\|\mathbb P R_B(t)\|_{H^s}$.
Define
\begin{align*}
w_1(t)&=-\int_{t_0}^tS(t-v)\mathbb P R_B(v)\,dv,\\
N_1&=B(u_B,w_1)+B(w_1,u_B)+B(w_1,w_1),\\
w_2(t)&=-\int_{t_0}^t S(t-v)N_1(v)\,dv.
\end{align*}
\begin{proposition}[Exact residual after two projected corrections]
With $w=w_1+w_2$, the projected unforced residual is precisely
\[
\mathcal E_2=B(u_B,w_2)+B(w_2,u_B)+B(w_1,w_2)
+B(w_2,w_1)+B(w_2,w_2).
\]
If $\|B(a,b)\|_{H^{s-1}}\le C_s\|a\|_{H^s}\|b\|_{H^s}$,
and
\[
H_\nu(T)=T+\sqrt{\frac{2T}{e\nu}},\quad
W_1=TF,\quad W_2=C_sH_\nu(T)(2AW_1+W_1^2),
\]
then
\[
\|w_1\|_{CH^s}\le W_1,\quad
\|w_2\|_{CH^s}\le W_2,\quad
\|\mathcal E_2\|_{CH^{s-1}}
\le C_s(2AW_2+2W_1W_2+W_2^2).
\]
The full unprojected residual equals $\mathcal E_2$ upon taking
\[
\pi=-\Delta^{-1}\nabla\cdot
\{R_B+(u_B\cdot\nabla)w+(w\cdot\nabla)u_B+(w\cdot\nabla)w\}
\]
as the correction pressure, with zero spatial mean.
\end{proposition}
\begin{proof}
The defining Duhamel equations give
$(\partial_t-\nu\Delta)w_1=-\mathbb P R_B$ and
$(\partial_t-\nu\Delta)w_2=-N_1$. Expand the full quadratic
nonlinearity and cancel the terms containing $w_1$ alone.
The Fourier multiplier inequality
$\sqrt{1+|k|^2}e^{-\nu t|k|^2}
\le1+(2e\nu t)^{-1/2}$
and time integration prove the stated $H_\nu(T)$ bound.
The bilinear estimate then yields all three inequalities.
Finally, the Poisson pressure removes the gradient component of the
total forcing; crucially, it retains $\nabla\cdot R_B$.
\end{proof}
For fixed $A,F$ the resulting residual bound is $O(T^{3/2})$;
there is no claim of exact zero residual or terminal continuation.

\begin{proposition}[Explicit meridional dual witness on $\mathbb R^3$]
Let $u_B=r e^{-(r^2+z^2)}e_\theta$ be a stationary Gaussian
pure-swirl background. Define
$\chi(x)=(x-\frac12)^4(1-x)^4$ for $x\in[\frac12,1]$ and
$\chi(x)=0$ otherwise, and put
$\psi(r,z)=\chi(r)\chi(z)$, $\Phi=\nabla\times(\psi e_\theta)$.
Then $\Phi$ is compactly supported and divergence-free, and
\[
\langle R_B,\Phi\rangle
=8\pi\int_{1/2}^1\!\int_{1/2}^1
r^2 z e^{-2(r^2+z^2)}\chi(r)\chi(z)\,dr\,dz
>\frac{\pi e^{-4}}{322560^2}>0.
\]
Consequently a smooth unforced correction vanishing at $t_0$
must satisfy
$\langle\partial_t w(t_0),\Phi\rangle
<-\pi e^{-4}/322560^2$.
\end{proposition}
\begin{proof}
The meridional centrifugal residual has
$(R_B)_r=-r e^{-2(r^2+z^2)}$, $(R_B)_z=0$
modulo any pressure gradient. Thus
$(\nabla\times R_B)_\theta=4rz e^{-2(r^2+z^2)}$.
Use integration by parts, the cylindrical volume element
$2\pi r\,dr\,dz$, and the exact beta integral
$\int_{1/2}^1\chi(x)\,dx=1/322560$.
The initial-time identity follows by testing the correction equation
with the solenoidal $\Phi$.
\end{proof}
The Gaussian witness is a pressure-independent \emph{positive control},
not a computation on the source's fully corrected field.



\section{An exact heat identity for the source-leading exterior}
\label{sec:source-heat-exterior}
The preceding source-leading exterior analysis can be strengthened by an
exact integration-by-parts identity, which also makes clear why a local
exterior calculation does \emph{not} give smooth global unforced blowup.
For $h>0$, $A=\frac12+h$, $c_\infty>0$, let
\[
g(R,v)=\frac{1}{R(R^2+4v)^h},\qquad
p_h(v)=\frac{e^{-v}v^h}{\Gamma(1+h)},\qquad
f(R)=c_\infty 2^A\int_0^\infty p_h(v)g(R,v)\,dv .
\]
This is the gamma-mixture expression for the published heat-exterior
\emph{leading} profile, not the fully matched physical field.
\begin{theorem}[Exact exterior radial heat identity]
For $R>0$,
\[
\boxed{f''(R)+\frac{1}{R}f'(R)-\frac{1}{R^2}f(R)
=\left(\frac12+h\right)f(R)+\frac R2 f'(R).}
\]
\end{theorem}
\begin{proof}
Put $\mathcal L=\partial_R^2+R^{-1}\partial_R-R^{-2}$
and $S=R^2+4v$. Direct differentiation yields
\begin{align*}
\frac{(\mathcal L-A-\frac R2\partial_R)g}{g}
&=\frac{4h\{R^2(1+h-v)-4v^2\}}{S^2}\\
&=4h\left\{\frac{1+h-v}{S}
-\frac{4(1+h)v}{S^2}\right\}.
\end{align*}
Because $p_h'/p_h=h/v-1$ and
$\partial_v g/g=-4h/S$, the previous expression equals
\[
(\mathcal L-A-\tfrac R2\partial_R)g
=\frac{1}{p_h(v)}
\partial_v\left[p_h(v)\frac{4hv}{S}\,g(R,v)\right].
\]
For fixed $R>0$ the boundary bracket vanishes at $v=0$
as $v^{1+h}$ and at $v=\infty$ exponentially.
All required differentiated integrands are integrable.
Integrating in $v$ proves the theorem.
\end{proof}
For $\tau=1-t$, the axially independent leading swirl
$U(r,t)=\tau^{-A} f(r/\sqrt{\tau})$ consequently solves
$\partial_tU=(\partial_r^2+r^{-1}\partial_r-r^{-2})U$
for viscosity $\nu=1$. On any annulus separated from the axis,
$u=Ue_\theta$ also solves the \emph{unforced} three-dimensional
Navier--Stokes equations \emph{locally}, taking
$p(r,t)=\int_{r_*}^{r}U(s,t)^2/s\,ds$ for its centrifugal pressure.
The identity does not assert that the genuine source's
meridional, oscillatory, transition, collar, and gluing corrections
satisfy the force-free equations.

\begin{proposition}[Exterior-only axis singularity]
For every $h>0$,
\[
\lim_{R\downarrow0} R f(R)
=\frac{c_\infty 2^{1/2-h}}{\Gamma(1+h)}>0.
\]
Hence for each fixed $\tau>0$,
\[
U(r,t)\sim \frac{c_\infty 2^{1/2-h}}{\Gamma(1+h)}
\frac{\tau^{-h}}{r}\quad(r\downarrow0).
\]
Thus the exterior-only flow is \emph{already singular before}
the proposed terminal time and is not a globally smooth
finite-energy solution from admissible initial data.
\end{proposition}
\begin{proof}
Dominated convergence in the gamma mixture gives
$\mathbb E[V^{-h}]=\Gamma(1)/\Gamma(1+h)$.
The displayed limit and its similarity consequence follow.
\end{proof}
The genuine remaining proof obligation is therefore exact
force-free \emph{matching} across the axis replacement, collars,
meridional velocity and all transition regions, together with
uniform terminal control, not a repair of the idealized
exterior's heat equation.



\section{A positive-time analyticity obstruction to exact exterior matching}
\label{sec:analytic-no-exact-gluing}
The preceding exterior profile is real-analytic away from the symmetry
axis, but has an unavoidable $r^{-1}$ singularity on that axis.
Combining this fact with the established positive-time spatial
analyticity of strong unforced Navier--Stokes solutions
\cite{zhang2026analyticity} gives a useful restriction on exact gluing.

\begin{theorem}[No exact annular matching of the singular exterior]
Let $h>0$, $c_\infty>0$, $0<t_0<T$ and $\tau=1-t_0>0$.
Suppose $u$ is a finite-energy, $H^s$ strong unforced solution
on $\mathbb R^3$ with $s>5/2$, smooth up to time $t_0$ and
spatially analytic at $t_0$. Define, on the connected domain
$\Omega=\{(x,y,z):r=\sqrt{x^2+y^2}>0\}$, the exterior comparison field
\[
u_{\mathrm{ext}}(x,y,z)=
\tau^{-(1/2+h)} f(r/\sqrt\tau)\,e_\theta,
\]
where $f$ is the gamma-mixture in Section~\ref{sec:source-heat-exterior}.
Then $u(t_0,\cdot)$ cannot agree identically with
$u_{\mathrm{ext}}$ on any nonempty open subset of $\Omega$.
\end{theorem}
\begin{proof}
For $R>0$ the gamma-mixture is real-analytic in $R$:
locally in $R$, a complex neighbourhood avoids the zeros of
$R^2+4v$ for all $v\ge0$ and the exponentially weighted integral
is locally uniformly convergent with its derivatives. The Cartesian
factor $e_\theta=(-y/r,x/r,0)$ is analytic on $\Omega$.
Thus both $u(t_0)$ and $u_{\mathrm{ext}}$ are real-analytic
on the connected domain $\Omega$.
Equality on one nonempty open set forces equality on $\Omega$.
But Proposition in Section~\ref{sec:source-heat-exterior} shows
$|u_{\mathrm{ext}}|\sim C\tau^{-h}/r$ for $C>0$ near the axis.
Its squared norm has divergent radial integral
$\int_0^\epsilon r^{-1}\,dr$ on every bounded axial interval,
contradicting $u(t_0)\in L^2(\mathbb R^3)$.
\end{proof}

\noindent\textbf{Precise scope.}
This rules out \emph{exact} equality with the specified
singular analytic leading exterior on a nonempty open patch
at a positive time. It does not rule out a nonzero analytic
matching tail, an approximate exterior profile, or other
unforced singularity constructions. The analyticity theorem is
classical prior art; its application here is a restricted
source-specific diagnostic, not a resolution of Navier--Stokes
regularity. The actual full matched source profile has not
been proved equal to $u_{\mathrm{ext}}$ on an open set.



\section{Circulation maximum principle and a necessary leading-order cancellation}
\label{sec:circulation-max-gate}
The following obstruction applies to a \emph{smooth axisymmetric unforced}
whole-space solution with bounded initial circulation. It does not apply
as stated to arbitrary nonaxisymmetric solutions or directly to the
separate periodic-torus closure problem. The scalar maximum principle
used here is established prior art \cite{circulationMP2023}.

For an axisymmetric velocity $(u_r,U,u_z)$ let $\Gamma=rU$. The
azimuthal momentum equation implies
\begin{equation}\label{eq:Gamma-axis-XXII}
 (\partial_t+u_r\partial_r+u_z\partial_z)\Gamma
 =\nu(\partial_{rr}-r^{-1}\partial_r+\partial_{zz})\Gamma.
\end{equation}
With smooth axis-compatible velocities and the usual spatial
comparison-principle hypotheses, and $M=\|\Gamma(t_0)\|_\infty<\infty$,
the maximum principle gives
$\|\Gamma(t)\|_\infty\le M$ for $t_0\le t<T$.

\begin{theorem}[Mandatory leading swirl cancellation in the axisymmetric route]
\label{thm:circulation-cancel-XXII}
Let $0<h$, $A=\frac12+h$, $\tau=1-t$ and fix an exterior similarity
radius $R>0$ where the Edition XX source-leading profile obeys $f(R)>0$.
Let $r_\tau=R\sqrt{\tau}$ and choose any compatible $z_\tau$.
Assume the fully chosen smooth axisymmetric background has
\[
 (u_B)_\theta(r_\tau,z_\tau,1-\tau)
 =\tau^{-A}[f(R)+\delta(\tau)],\qquad \delta(\tau)\longrightarrow0.
\]
If $u=u_B+w$ is a smooth axisymmetric \emph{unforced} solution with
bounded initial circulation and the above maximum-principle hypotheses,
then
\begin{equation}\label{eq:circulation-cancellation-XXII}
 \left|\tau^A w_\theta(r_\tau,z_\tau,1-\tau)+f(R)\right|
 \le |\delta(\tau)|+\frac{M}{R}\tau^h.
\end{equation}
Consequently
$\lim_{\tau\downarrow0}\tau^A w_\theta(r_\tau,z_\tau,1-\tau)
=-f(R)\ne0$.
No correction $w$ that is subleading in this \emph{rescaled azimuthal
pointwise amplitude} can yield such an axisymmetric unforced solution.
\end{theorem}
\begin{proof}
The maximum principle gives
$|r_\tau u_\theta(r_\tau,z_\tau,1-\tau)|\le M$,
so $|\tau^A u_\theta|\le M\tau^{A-1/2}/R=M\tau^h/R$.
Insert $u_\theta=(u_B)_\theta+w_\theta$ and use the assumed source
asymptotic and the triangle inequality. This proves
\eqref{eq:circulation-cancellation-XXII}.
\end{proof}

\begin{corollary}[Necessary forcing budget if the circulation survives]
For an axisymmetric \emph{forced} solution with azimuthal force $F_\theta$,
the right side of \eqref{eq:Gamma-axis-XXII} gains $rF_\theta$.
If $\int_{t_0}^{t}\|rF_\theta(s)\|_\infty\,ds<\infty$ for every
$t<T$, and the actual forced velocity obeys
$u_\theta(r_\tau,z_\tau,1-\tau)\ge\frac12\tau^{-A}f(R)$,
then
\begin{equation}\label{eq:force-budget-XXII}
 \int_{t_0}^{1-\tau}\|rF_\theta(s)\|_\infty\,ds
 \ge\frac12 Rf(R)\tau^{-h}-M.
\end{equation}
Thus such retained circulation requires a divergent \emph{cumulative}
weighted azimuthal force as $\tau\downarrow0$. This is not a measured
or certified bound for the complete forced source construction.
\end{corollary}
\begin{proof}
The inhomogeneous maximum principle implies
$\|\Gamma(t)\|_\infty\le
 M+\int_{t_0}^t\|rF_\theta(s)\|_\infty\,ds$.
Evaluate at $(r_\tau,z_\tau)$ and rearrange.
\end{proof}

\paragraph{Nonaxisymmetric boundary of the argument.}
For general smooth cylindrical fields, direct calculation for $r>0$
gives
\begin{equation}\label{eq:nonaxis-circulation-XXII}
\begin{aligned}
 &(\partial_t+u_r\partial_r+(u_\theta/r)\partial_\theta
     +u_z\partial_z)\Gamma\\
 &\quad=\nu(\Delta-2r^{-1}\partial_r)\Gamma
      +2\nu r^{-1}\partial_\theta u_r-\partial_\theta p
      +rF_\theta.
\end{aligned}
\end{equation}
The angular meridional and pressure terms are not negligible by
symmetry in a nonaxisymmetric correction. The separated term
$r^{-1}\partial_\theta u_r$ may fail to be uniformly bounded near
the axis even for smooth Cartesian velocities, so we do not infer
an unchecked global $L^\infty$ estimate from
\eqref{eq:nonaxis-circulation-XXII}. A nonaxisymmetric route remains
open, as does any distinct axisymmetric singularity mechanism that
does not retain the positive growing source-leading circulation.

\paragraph{Audit and scope.}
The chain rules, viscosity cross term and similarity power were
independently checked with SymPy; 50,000 exact rational finite budget
tests passed. A separate Lean file certifies only finite scalar
implications, not the axisymmetric maximum principle or its physical
hypotheses. The source's leading exterior identity does not by itself
prove that the \emph{fully matched background} satisfies the asymptotic
assumed in Theorem~\ref{thm:circulation-cancel-XXII}. This is a
source-specific obstruction, not an unforced blow-up proof.



\section{The first nonaxisymmetric Fourier pair: pressure and harmonic leakage}
\label{sec:nonaxis-Fourier-XXIII}
This section studies a possible nonaxisymmetric escape from the
circulation obstruction. Fourier-mode selection, the cylindrical
vector Laplacian and Reynolds averaging are established prior art.
The source Navier--Stokes singularity is \emph{forced}; we do not
claim that a force-free construction has been completed.

Let $u_B=U(r,z,t)e_\theta$ be an axisymmetric pure-swirl background
(not assumed force-free). For
$v_m=e^{im\theta}(a e_r+b e_\theta+c e_z)$ put
$D_m=\partial_{rr}+r^{-1}\partial_r-m^2r^{-2}+\partial_{zz}$
and $L_m=\partial_t-\nu D_m$. Linearizing, with a
modewise right-hand side $H_m$, gives
\begin{align}
L_m a+\nu(a+2im b)/r^2+(imU/r)a-2Ub/r+\pi_{m,r}&=H_{m,r},\\
L_m b+\nu(b-2im a)/r^2+(imU/r)b+(U_r+U/r)a+U_zc
  +(im/r)\pi_m&=H_{m,\theta},\\
L_m c+(imU/r)c+\pi_{m,z}&=H_{m,z},\\
a_r+a/r+im b/r+c_z&=0.
\label{eq:nonaxis-system-XXIII}
\end{align}
The off-diagonal viscosity coefficients reflect the vector
Laplacian in its rotating basis and cannot be omitted.
Divergence and solenoidality give the necessary pressure equation
\begin{equation}\label{eq:nonaxis-pressure-XXIII}
D_m\pi_m=\operatorname{div}_mH_m
-\frac{2im}{r}(U_ra+U_zc)+\frac2r\partial_r(Ub),
\end{equation}
where
$\operatorname{div}_mH_m=H_{m,r,r}+r^{-1}H_{m,r}
+im r^{-1}H_{m,\theta}+H_{m,z,z}$.
This formula does not establish existence of an axis-regular
Poisson inverse with the original physical boundary data.

\begin{proposition}[A smooth first-harmonic source of mean angular torque]
\label{prop:nonaxis-torque-XXIII}
On $\mathbb R^2\times\mathbb S^1$ with axial period $2\pi$,
put $f(r)=e^{-r^2}$ and let $\lambda\in\mathbb R$.
For $m=+1$ set
\[
a=f\cos z,\qquad b=ia+\lambda r^2f\cos z,\qquad
c=-f'\sin z-i\lambda rf\sin z,
\]
and define the real velocity
$v=e^{i\theta}(a e_r+b e_\theta+c e_z)
+e^{-i\theta}(\bar a e_r+\bar b e_\theta+\bar c e_z)$.
It is smooth, finite-energy and divergence-free; in Cartesian form,
\[
v=(2f\cos z(1-\lambda xy),\,2\lambda x^2f\cos z,\,
       (4x+2\lambda y)f\sin z).
\]
Its azimuthal Reynolds tensor $Q=\langle v\otimes v\rangle_\theta$
has
\[
Q_{r\theta}=2\lambda r^2e^{-2r^2}\cos^2z,\qquad
Q_{z\theta}=\lambda r(2r^2-1)e^{-2r^2}\sin2z,
\]
and so the averaged angular convection satisfies
\begin{equation}\label{eq:nonaxis-az-force-XXIII}
g_\theta=\frac1{r^2}\partial_r(r^2Q_{r\theta})
             +\partial_zQ_{z\theta}
=2\lambda r e^{-2r^2}(2-2r^2+\cos2z).
\end{equation}
Its radial-slice torque is exactly
\begin{equation}\label{eq:nonaxis-slice-XXIII}
\int_0^\infty r^2g_\theta(r,z)\,dr
=\frac{\lambda}{4}\cos2z.
\end{equation}
At each point, $Q$ is positive semidefinite and has rank at most
two.
\end{proposition}
\begin{proof}
The modal divergence vanishes directly:
$a_r+a/r+ib/r+c_z=0$. Independently, the displayed Cartesian
field has zero divergence and is regular across the axis.
The angular average gives
$Q_{ij}=2\operatorname{Re}(v_{+,i}\overline{v_{+,j}})$,
a sum of two positive rank-one real tensors.
Substitution in the cylindrical tensor-divergence formula gives
\eqref{eq:nonaxis-az-force-XXIII}; exact Gaussian integration gives
\eqref{eq:nonaxis-slice-XXIII}. Direct real-angle integration
independently verifies both Reynolds tensor components.
\end{proof}

\begin{proposition}[Non-gradient second-harmonic residual]
\label{prop:nonaxis-mtwo-XXIII}
For the field in Proposition~\ref{prop:nonaxis-torque-XXIII}, the
$m=+2$ coefficient $N=(v_+\cdot\nabla)v_+$ of self-convection
has, at $\lambda=1$, $r=1$, $z=0$,
\begin{equation}\label{eq:nonaxis-mtwo-curl-XXIII}
(\nabla\times N)^{(+2)}_z=2(-2+5i)e^{-2}\ne0.
\end{equation}
Thus this \emph{particular} correction cannot be closed at
$m=0,\pm1$ by choosing scalar pressure.
\end{proposition}
\begin{proof}
In the rotating cylindrical basis, the $m=+2$ coefficients are
\[
N_r=aa_r+iba/r+ca_z-b^2/r,\quad
N_\theta=ab_r+ib^2/r+cb_z+ab/r,\quad
N_z=ac_r+ibc/r+cc_z.
\]
The axial curl coefficient is
$r^{-1}\partial_r(rN_\theta)-2iN_r/r$.
Exact differentiation gives \eqref{eq:nonaxis-mtwo-curl-XXIII}.
A scalar pressure gradient has zero curl.
\end{proof}

\paragraph{Connection to the existing H8 shell programme.}
The H8 inversion gate requires suitably small \emph{radial-slice}
torque in a slender shell, not merely zero total torque.
Equation~\eqref{eq:nonaxis-slice-XXIII} has zero average over $z$
but an $O(|\lambda|)$ nonzero slice torque. Any source-compatible
use of this mechanism requires $\lambda$ small at the right
slenderness scale or an independent cancelling profile.
The present manufactured domain is not the H8 physical source.

\paragraph{The full unforced and Lean gates.}
If the base residual lies in $m=0$, pure $m=\pm1$ perturbations
cannot cancel it linearly. Quadratic cross terms can generate
$m=0$ but the present witness generates $m=\pm2$ as well.
An exact zero-force theorem needs source-matched A.11/A.19
residuals and collars, a fully coupled pressure-compatible
harmonic inverse, nonlinear norm convergence and survival of
the terminal singularity. Independent SymPy and Cartesian
reconstructions, direct angular averages and 12,000 exact
rational finite tests pass. Lean establishes only finite
mode-index and Gram-minor algebra, not these PDE theorems.



\subsection{Exact two-shell compensation of the first-harmonic slice torque}
\label{subsec:nonaxis-two-shell-XXIII}
The preceding Gaussian $m=\pm1$ model produces a nonzero radial-slice
torque. Its scalar obstruction can be removed constructively at the
level of smooth, solenoidal manufactured fields, without claiming
that the authentic physical source has been reconstructed.

\begin{proposition}[Two-shell exact radial-slice torque cancellation]
\label{prop:two-shell-XXIII}
Let $f_1,f_2$ be nonzero real $C^\infty$ compactly supported
radial functions whose supports lie in disjoint annuli in $(0,\infty)$.
Put $I_j=\int_0^\infty r^3f_j(r)^2\,dr>0$.
For each $j$ use the $m=+1$ amplitudes
\[
a_j=f_j\cos z,\quad
b_j=ia_j+\lambda_j r^2f_j\cos z,\quad
c_j=-f_j'\sin z-i\lambda_j rf_j\sin z
\]
and add their conjugate $m=-1$ fields.
If $\lambda_1\ne0$ and
\[
\boxed{\lambda_2=-\lambda_1 I_1/I_2},
\]
the resulting real velocity is smooth, divergence-free, finite-energy,
and its averaged nonlinear angular source satisfies
\[
\boxed{\int_0^\infty r^2g_\theta(r,z)\,dr=0}
\qquad\text{for every }z,
\]
while $g_\theta$ does not vanish identically on either of the two
annular supports.
\end{proposition}
\begin{proof}
Each modal divergence vanishes by direct differentiation.
The disjoint radial supports eliminate all cross products of
the two profiles. For an arbitrary one of the radial amplitudes
$f=f_j$, Fourier averaging gives
\[
Q_{z\theta}
=-2\lambda(r^2ff'+rf^2)\sin z\cos z.
\]
The cylindrical torque identity is
$r^2g_\theta=\partial_r(r^2Q_{r\theta})
+\partial_z(r^2Q_{z\theta})$.
The radial boundary term is zero because $f$ is compactly supported.
Integration by parts yields
\[
\int_0^\infty (r^4ff'+r^3f^2)\,dr
=-\int_0^\infty r^3f^2\,dr=-I.
\]
Thus the first shell contributes $2\lambda_1I_1\cos2z$
and the second $2\lambda_2I_2\cos2z$.
The displayed choice cancels their sum exactly.
If the local source were identically zero on a support, setting
$z=\pi/2$ would force $r^2ff'+rf^2=0$ throughout the
support, hence $f\propto r^{-1}$ wherever $f\ne0$,
contradicting smooth compact support of a nonzero annular profile.
\end{proof}

\paragraph{Independent rational fixture and scope.}
For the piecewise-polynomial comparison profiles
$f_1=(r-1)^4(2-r)^4$ on $[1,2]$ and
$f_2=(r-3)^4(4-r)^4$ on $[3,4]$ (zero outside),
exact integration gives
\[
 I_1=\frac{29}{1847560},\qquad
 I_2=\frac{3269}{16628040},\qquad
 \lambda_1=1,\quad \lambda_2=-\frac{261}{3269}.
\]
These specific test profiles are $C^3$, not $C^\infty$; the
preceding theorem concerns genuine smooth compact bumps.
An independent Fraction/SymPy calculation reproduces all rational
values and exactly zero slice torque. No claim is made that the
two-shell model solves the full original A.11/A.19 source, the
seven-ray H8 cone, the $m=\pm2$ nonlinear residual, the infinite
Fourier hierarchy, or the unforced Navier--Stokes blow-up problem.



\section{Weighted angular-flux feasibility before source-matched stress closure}
\label{sec:weighted-flux-XXIV}
This section supplies a source-facing \emph{necessary} test for the
nonaxisymmetric correction route of Section~\ref{sec:nonaxis-Fourier-XXIII}.
The argument is elementary positivity and integration by parts; we do
not claim novelty for the covariance inequality or Navier--Stokes
energy identity. The proposed physical source and full unforced
blow-up theorem remain uncertified.

\begin{theorem}[Weighted covariance and two-slice torque obstruction]
\label{thm:weighted-flux-XXIV}
Let $v$ be a real smooth field, compactly supported in an annulus
$0<r<b$, periodic in $z$, and let
$Q_{ij}(r,z)=\langle v_i v_j\rangle_\theta$.
Define the azimuthal nonlinear mean momentum flux, its radial-slice
torque, and the two weighted diagonal moments by
\[
g_\theta=\frac1{r^2}\partial_r(r^2Q_{r\theta})
               +\partial_zQ_{z\theta},\quad
m(z)=\int_0^\infty r^2g_\theta(r,z)\,dr,
\]
\[
T(z)=\int_0^\infty r^2Q_{\theta z}(r,z)\,dr,\qquad
E_j(z)=\int_0^\infty r^2Q_{jj}(r,z)\,dr
\quad(j=\theta,z).
\]
Then, for every $z$,
\begin{equation}\label{eq:weighted-CS-XXIV}
m(z)=T'(z),\qquad
T(z)^2\le E_\theta(z)E_z(z),\qquad
2|T(z)|\le E_\theta(z)+E_z(z).
\end{equation}
For any primitive $M'(z)=m(z)$ and any interval $I$ in $z$,
\begin{equation}\label{eq:torque-osc-XXIV}
\boxed{\ \sup_{z\in I}(E_\theta(z)+E_z(z))
       \ge\operatorname{osc}_{z\in I}M(z).\ }
\end{equation}
If the radial support lies in $r\le b$ and
$\mathcal E_{\theta z}(z)=2\pi\int_0^\infty
r\,\langle v_\theta^2+v_z^2\rangle_\theta\,dr$ is the
non-halved directional kinetic-energy per axial slice, then
\begin{equation}\label{eq:physical-energy-XXIV}
\boxed{\ \sup_{z\in I}\mathcal E_{\theta z}(z)
 \ge \frac{2\pi}{b}\operatorname{osc}_{z\in I}M(z).\ }
\end{equation}
The bound is necessary for the \emph{portion of the actual source
assigned to the quadratic Reynolds stress}. In a full source
comparison all other linear, mean and pressure terms must first be
separated; $m$ cannot be identified with the entire source residual
without that decomposition.
\end{theorem}
\begin{proof}
Smoothness and radial compactness allow integration by parts in
$r$ and differentiation in $z$, giving $m=T'$.
The angular covariance matrix is positive semidefinite:
for any real $a,b$, $\langle(av_\theta+bv_z)^2\rangle_\theta\ge0$.
The resulting pointwise Cauchy--Schwarz bound, followed by weighted
integral Cauchy--Schwarz, gives $T^2\le E_\theta E_z$.
By arithmetic--geometric mean,
$2|T|\le E_\theta+E_z$.
For any two axial locations $z_+,z_-$,
\[
|M(z_+)-M(z_-)|=|T(z_+)-T(z_-)|
 \le |T(z_+)|+|T(z_-)|
 \le \sup_{z\in I}(E_\theta+E_z).
\]
Take the supremum over pairs, proving \eqref{eq:torque-osc-XXIV}.
Finally $E_\theta+E_z\le
b\mathcal E_{\theta z}/(2\pi)$ on $r\le b$, proving
\eqref{eq:physical-energy-XXIV}.
\end{proof}

\begin{corollary}[Exact Gaussian fixture and two-shell interpretation]
\label{cor:Gaussian-XXIV}
For the smooth $m=\pm1$ Gaussian model of
Proposition~\ref{prop:nonaxis-torque-XXIII},
\[
T(z)=\frac{\lambda}{8}\sin(2z),\qquad
m(z)=\frac{\lambda}{4}\cos(2z),\qquad
\operatorname{osc}_{[0,2\pi]}T=\frac{|\lambda|}{4}.
\]
Thus any stress realizing this \emph{entire} slice-torque target
must have $\sup_z(E_\theta+E_z)\ge|\lambda|/4$.
Conversely, the two-shell compensation of
Proposition~\ref{prop:two-shell-XXIII} makes $T\equiv0$ and
$m\equiv0$; it cannot reproduce a target with nonzero radial-slice
torque, even though its local angular source need not vanish.
\end{corollary}

\begin{proposition}[Zero-initial-energy exclusion for unforced smooth flow]
\label{prop:zero-data-XXIV}
If $u$ is a smooth finite-energy unforced solution on $\mathbb R^3$
or a periodic torus, and $u(0)=0$, then $u\equiv0$ on its entire
smooth lifespan. Therefore removing the force from the published
OpenAI construction \emph{while retaining its zero initial datum}
cannot prove a nontrivial unforced singularity.
\end{proposition}
\begin{proof}
The standard smooth energy identity yields
\[
\frac12\|u(t)\|_2^2+
\nu\int_0^t\|\nabla u(s)\|_2^2\,ds
=\frac12\|u(0)\|_2^2=0.
\]
Both left terms are nonnegative, so $u(t)=0$.
The originating forced theorem explicitly has $u(0)=0$
\cite{openai2026}. An unforced singularity attempt must instead
exhibit admissible \emph{nonzero} smooth initial data.
\end{proof}

\paragraph{Proof and release status.}
Independent Gaussian integration, $30{,}000$ rational positive
covariance tests and $30{,}000$ two-slice budget tests pass; a
separate scalar Lean module states the finite covariance inequality,
two-slice bound, and zero-data energy arithmetic. It does not
formalize angular averaging, Bochner integration, energy equality,
the source or the full PDE. The existing source certificate remains
\texttt{FAIL\_CLOSED\_MISSING\_SOURCE\_DATA}: source A.11/A.19
profiles, actual mean forcing, full cone, pressure, and zero-force
PDE residual are not authenticated. The seven-ray H8 realizability
and its spatially varying dual obstructions remain separate,
stronger gates. \emph{No unforced blow-up or complete Lean proof
is established.}



\subsection{A rank defect of one angular Fourier pair}
\label{subsec:angular-rank-defect-XXIV}
The preceding slice-flux budget does not guarantee feasibility of
a positive-definite H8 stress. A second, independent algebraic
gate is the rank of an angularly averaged covariance.

\begin{proposition}[One-pair rank obstruction]
\label{prop:single-pair-rank-XXIV}
At a fixed $(r,z)$, let a real cylindrical velocity consist
of one angular Fourier pair,
\[
v(\theta)=e^{i\theta}c+e^{-i\theta}\overline c,
\qquad c=a+ib\in\mathbb C^3,\quad a,b\in\mathbb R^3,
\]
with vectors expressed in the co-rotating orthonormal
$(e_r,e_\theta,e_z)$ frame. Then its angularly averaged covariance
is
\[
\boxed{Q=\langle v\otimes v\rangle_\theta
=2(aa^{\mathsf T}+bb^{\mathsf T}),\qquad
\operatorname{rank}Q\le2,\quad\det Q=0.}
\]
Thus on a region with $\mu\rho>0$ the \emph{direct stress
realization} $Q\succeq\mu\rho I_3$ is impossible using only
this one angular pair. Adding two disjoint radial shells at the
\emph{same} Fourier frequency does not remove this obstruction:
at each point their combined mode is still one complex
three-vector.
\end{proposition}
\begin{proof}
The angular mean removes the $m=\pm2$ exponentials in
$v\otimes v$ and retains $c\bar c^{\mathsf T}
+\bar c c^{\mathsf T}=2(aa^{\mathsf T}+bb^{\mathsf T})$.
Its range lies in $\operatorname{span}\{a,b\}$.
When $a,b$ are independent, the nonzero vector $a\times b$
is in its kernel; when dependent, its rank is at most one.
Every strictly positive-definite $3\times3$ matrix has rank three.
\end{proof}

\begin{proposition}[Two-frequency algebraic repair]
\label{prop:two-frequency-PSD-XXIV}
For independent angular frequencies $m=1$ and $m=2$, a
\emph{pointwise algebraic} covariance can be strictly positive:
take $c_1=(1,i,0)$ and $c_2=(0,0,1)$. Then
\[
v(\theta)=2\operatorname{Re}(e^{i\theta}c_1+
e^{2i\theta}c_2),\qquad
\boxed{\langle v\otimes v\rangle_\theta=2I_3.}
\]
\end{proposition}
\begin{proof}
Orthogonality of the different angular Fourier frequencies
eliminates their cross terms. The real and imaginary
polarizations of $c_1$ provide the first two Cartesian
coordinates and $c_2$ the third. Each contributes $2$ to its
diagonal entry, giving the identity.
\end{proof}

\paragraph{What is and is not repaired.}
The two-frequency example is \emph{finite-dimensional algebra},
not a globally smooth, incompressible, axis-regular wave with a
source-compatible heat, transport or pressure equation.
Furthermore, the H8 target is constrained by the projected
divergence and pressure-gauge equivalence; a pointwise rank
obstruction to \emph{direct} covariance equality does not rule
out all projected-stress representations with gauge freedom.
The seven-ray strict cone, virtual-work dual inequalities,
angular-mode coupling and $m=3,4$ leakage remain to be
verified for an actual flow. Independent symbolic determinant,
cross-product null-vector and $5{,}000$ exact-rational tests
pass; the finite Lean scalar cross-product certificate is not
a PDE or full cone proof.



\subsection{Smooth solenoidal two-frequency stress with a strict local floor}
\label{subsec:two-frequency-solenoidal-XXIV}
The algebraic rank-three example above can in fact be realized by
a smooth finite-energy \emph{divergence-free} field, though it is not
an unforced Navier--Stokes solution.

\begin{proposition}[Annular divergence-free positive covariance]
\label{prop:two-frequency-solenoidal-XXIV}
Let $f,g\in C_c^\infty((0,\infty);\mathbb R)$ have supports away
from $r=0$, and set $r=\sqrt{x^2+y^2}$.
On $\mathbb R^2\times(\mathbb R/2\pi\mathbb Z)$ define
\[
v_x=2f(r)\cos z,\quad v_y=0,\quad
v_z=-2f'(r)\frac{x}{r}\sin z+
          2g(r)\frac{x^2-y^2}{r^2}.
\]
This field is smooth, compactly supported in the transverse plane,
finite-energy on the stated domain, and exactly divergence-free.
In the rotating cylindrical frame it has angular frequencies
$m=\pm1,\pm2$ only. Its azimuthally averaged co-rotating stress is
\[
Q=2\begin{pmatrix}
 A^2&0&-AB\\0&A^2&0\\-AB&0&B^2+C^2
\end{pmatrix},
\quad A=f(r)\cos z,\quad B=f'(r)\sin z,\quad C=g(r).
\]
In particular,
\[
\boxed{\det Q=8f(r)^4 g(r)^2\cos^4z.}
\]
Choose $f=g=1$ on a nonempty annular plateau and zero outside
a slightly larger annulus. On that plateau and
$|z|\le\pi/4$, the stress satisfies the \emph{explicit} floor
\[
\boxed{Q=2\operatorname{diag}(\cos^2z,\cos^2z,1)
\succeq I_3.}
\]
\end{proposition}
\begin{proof}
Since $v_y=0$ and the last term of $v_z$ is $z$-independent,
\[
\partial_xv_x
=2f'(r)\frac{x}{r}\cos z,\qquad
\partial_zv_z=-2f'(r)\frac{x}{r}\cos z,
\]
so $\nabla\cdot v=0$ identically. Compact radial supports away
from $r=0$ ensure smoothness in Cartesian coordinates.
The cylindrical components are
\[
v_r=2f\cos\theta\cos z,\quad
v_\theta=-2f\sin\theta\cos z,\quad
v_z=-2f'\cos\theta\sin z+2g\cos2\theta .
\]
Their nonzero Fourier amplitudes are
$c_1=(f\cos z,\,if\cos z,\,-f'\sin z)$
and $c_2=(0,0,g)$.
Angular orthogonality yields the displayed covariance.
Its determinant is
$8A^4C^2$ by direct expansion; on the plateau $B=0$.
Finally $2\cos^2z\ge1$ on $|z|\le\pi/4$.
\end{proof}

\paragraph{Physical and formal boundaries.}
This theorem removes the \emph{local strict positive-definite
covariance} obstruction for a two-frequency divergence-free
manufactured field. It does not establish membership in the
specific seven-ray H8 coefficient cone with a positive floor,
nor the H8 virtual-work dual inequalities, A.11/A.19 source
matching, compact anti-divergence, a genuinely unforced time
evolution, or singularity preservation. Quadratic interaction
of the $m=\pm1,\pm2$ modes creates frequencies through
$m=\pm4$, requiring a larger and ultimately convergent
harmonic hierarchy. An independent Cartesian divergence
calculation, symbolic determinant factorization and a scalar
Lean determinant identity were checked; full PDE formalization
remains open.



\subsection{Simultaneous local strict covariance and nonzero slice torque}
\label{subsec:two-mode-torque-XXIV}
The preceding two-frequency field has zero mean angular torque in
the special choice $\lambda=0$. A small additional first-harmonic
polarization produces an explicitly nonzero torque without losing
the positive covariance floor.

\begin{proposition}[Divergence-free two-harmonic torque shaper]
\label{prop:two-mode-torque-XXIV}
Fix real $f,g\in C_c^\infty((1/2,3/2))$ such that
$f=g=1$ on $[3/4,5/4]$. Let $|\lambda|\le1/16$.
Use the $m=+1$ cylindrical amplitudes
\[
a=f(r)\cos z,\quad
b=ia+\lambda r^2f(r)\cos z,\quad
c=-f'(r)\sin z-i\lambda rf(r)\sin z
\]
and the independent $m=+2$ amplitude
$(a_2,b_2,c_2)=(0,0,g(r))$; add complex-conjugate modes
to make the field real. On $\mathbb R^2\times\mathbb S^1$
of axial period $2\pi$, the resulting smooth finite-energy
velocity is divergence-free. Its co-rotating azimuthal stress
$Q=\langle v\otimes v\rangle_\theta$ satisfies
\[
\boxed{\det Q=8 f(r)^4g(r)^2\cos^4z.}
\]
On the plateau
$3/4\le r\le5/4,\ |z|\le\pi/6$,
\begin{equation}\label{eq:strict-positive-torque-XXIV}
\boxed{\quad Q\succeq\frac{653}{512}I_3
             \succ I_3.\quad}
\end{equation}
Moreover, the radial-slice angular flux and torque are
\[
\boxed{
T(z)=\lambda I_f\sin 2z,\qquad
m(z)=2\lambda I_f\cos2z,\qquad
I_f=\int_0^\infty r^3 f(r)^2\,dr>0.
}
\]
In particular, the torque is nonzero when $\lambda\ne0$.
\end{proposition}
\begin{proof}
The first-harmonic divergence vanishes because
$a_r+a/r+ib/r+c_z=0$; the second mode has axial
velocity independent of $z$, and therefore also has zero
divergence. Compact annular support gives Cartesian
smoothness and finite energy.

Writing the real and imaginary first-harmonic polarizations as
\[
p=(f\cos z,\,\lambda r^2 f\cos z,\,-f'\sin z),\qquad
q=(0,\,f\cos z,\,-\lambda r f\sin z)
\]
and the independent second-harmonic polarization as
$s=(0,0,g)$ gives $Q=2(pp^{\mathsf T}+qq^{\mathsf T}
+ss^{\mathsf T})$. Since the determinant of the
three-column matrix $(p,q,s)$ is $f^2g\cos^2z$,
$\det Q=8f^4g^2\cos^4z$, independently of $\lambda$.

On the plateau the diagonal entries satisfy
$Q_{11}=2\cos^2z,\ Q_{22}\ge2\cos^2z,\ Q_{33}\ge2$.
The only nonzero off-diagonal terms are
$Q_{12}=2\lambda r^2\cos^2z$ and
$Q_{23}=-2\lambda r\sin z\cos z$.
Using $2|xy|\le x^2+y^2$ in the quadratic form gives a
uniform lower bound equal to the smallest of
\[
2\cos^2z(1-|\lambda|r^2),\quad
2\cos^2z(1-|\lambda|r^2)-2|\lambda|r|\sin z\cos z|,
\quad
2-2|\lambda|r|\sin z\cos z|.
\]
The hypotheses imply these are respectively bounded below
by $693/512$, $653/512$ and $123/64$, giving the stated
positive floor.

Angular Fourier orthogonality removes all mixed $m=1,m=2$
terms in $Q_{\theta z}$. The first pair contributes
\[
Q_{\theta z}=-\lambda r(rff'+f^2)\sin 2z.
\]
Integrating $r^2Q_{\theta z}$ in $r$ and using
$\int r^4ff'=-2\int r^3f^2$ proves
$T=\lambda I_f\sin2z$; differentiating yields the torque.
\end{proof}

\paragraph{Release boundary.}
The construction simultaneously passes three \emph{local}
kinematic gates---solenoidality, strict covariance positivity
and nonzero slice torque---but remains a manufactured field,
not a solution of the time-dependent unforced Navier--Stokes
equations. In particular it is not certified to lie in the
specific H8 seven-ray positive-coefficient cone, is not matched
to the A.11/A.19 source, and does not cancel $m=3,4$
quadratic mode leakage. Strict local covariance cannot
replace the global anti-divergence, pressure, source and
singularity-preservation proof obligations.
Independent symbolic determinant, full modal divergence,
angular torque and rational covariance-floor calculations pass.



\section{Exact seven-ray cone completion and the compulsory carrier}
\label{sec:seven-ray-cone-XXV}
We now test the smooth two-frequency covariance of Edition XXIV
against the actual H8 seven-ray geometry. These are exact
source-independent tensor identities in a local co-rotating frame,
not an unforced Navier--Stokes construction.

\begin{theorem}[Direct seven-ray cone obstruction]
\label{thm:ray-direct-XXV}
Let $t=1/4$, and in the local frame
$(e_1,e_2,e_3)=(e_r,e_\theta,e_z)$ put
\[
v_1=e_1,\qquad v_{2,3}=e_1\pm te_2,\qquad
v_{4,5}=e_1\pm te_3,\qquad
v_{6,7}=e_1+t(e_2\pm e_3).
\]
For $d_j\ge\mu>0$, the stress
$S=\sum_{j=1}^7d_jv_j\otimes v_j$ obeys
\[
\boxed{S_{rr}-S_{\theta\theta}\ge 105\mu/16>0.}
\]
On the Edition XXIV annular plateau, its manufactured two-frequency
Reynolds covariance has
$Q_{rr}-Q_{\theta\theta}=-2\lambda^2r^4\cos^2z\le0$.
Hence $Q+\pi I$ is not directly in the strict seven-ray cone for
any scalar gauge $\pi$.
\end{theorem}
\begin{proof}
Each ray satisfies $(v_j)_r=1$ and
$|(v_j)_\theta|\le1/4$. Consequently
$S_{rr}-S_{\theta\theta}
=\sum_jd_j(1-(v_j)_\theta^2)
\ge(15/16)\sum_jd_j\ge105\mu/16$.
The manufactured $Q$ has $Q_{rr}=2\cos^2z$ and
$Q_{\theta\theta}=2\cos^2z(1+\lambda^2r^4)$.
An isotropic gauge does not change their difference.
\end{proof}

\begin{theorem}[Strict seven-ray completion with a positive carrier]
\label{thm:ray-carrier-XXV}
On $3/4\le r\le5/4$, $|z|\le\pi/6$ and
$|\lambda|\le1/16$, write $s=\sin z$ and $c=\cos z$.
Define
\begin{align*}
d_1&=1,&
d_2&=1+4\lambda r^2c^2,&
d_3&=3-4\lambda r^2c^2,\\
d_4&=2+16s^2+16\lambda^2(r^2s^2-r^4c^2)
  +16\lambda rcs,\\
d_5&=2+16s^2+16\lambda^2(r^2s^2-r^4c^2)
  -16\lambda rcs,\\
d_6&=1-16\lambda rcs,&
d_7&=1+16\lambda rcs,\\
\pi&=\frac38-2c^2(1+\lambda^2r^4),\\
K&=\frac{85}{8}+32s^2+32\lambda^2r^2s^2
                     -30\lambda^2r^4c^2.
\end{align*}
For the complete co-rotating covariance tensor $Q$ of the
two-frequency field in Edition XXIV, the following holds exactly:
\begin{equation}\label{eq:ray-completion-XXV}
\boxed{Q+K e_r\otimes e_r+\pi I
=\sum_{j=1}^7 d_jv_j\otimes v_j,\qquad
d_j\ge\frac38>0,\quad K>0.}
\end{equation}
\end{theorem}
\begin{proof}
The covariance in the plateau is
\[
Q=2\begin{pmatrix}
c^2&\lambda r^2c^2&0\\
\lambda r^2c^2&c^2(1+\lambda^2r^4)&-\lambda rcs\\
0&-\lambda rcs&1+\lambda^2r^2s^2
\end{pmatrix}.
\]
The seven-ray synthesis identities are
\begin{align*}
S_{rr}&=\sum_jd_j,&
S_{\theta\theta}&=(d_2+d_3+d_6+d_7)/16,\\
S_{zz}&=(d_4+d_5+d_6+d_7)/16,&
S_{r\theta}&=(d_2-d_3+d_6+d_7)/4,\\
S_{rz}&=(d_4-d_5+d_6-d_7)/4,&
S_{\theta z}&=(d_6-d_7)/16.
\end{align*}
Substitution verifies all six symmetric tensor components.
Using $r\le5/4$, $|s|\le1/2$, $|c|\le1$ and
$|\lambda|\le1/16$ yields
\[
d_1=1,\quad d_2\ge39/64,\quad d_3\ge167/64,\quad
d_{4,5}\ge5007/4096,\quad d_{6,7}\ge3/8.
\]
Also $K\ge85/8-30(1/256)(625/256)>0$.
\end{proof}

\begin{proposition}[Non-gauge carrier curl]
\label{prop:ray-carrier-curl-XXV}
The spatially varying radial carrier in
\eqref{eq:ray-completion-XXV} has
\[
\nabla\cdot(K e_r\otimes e_r)=(K_r+K/r)e_r,
\]
and
\[
\boxed{\bigl[\nabla\times\nabla\cdot
(K e_r\otimes e_r)\bigr]_\theta
=\left(\frac{64}{r}+192\lambda^2r+300\lambda^2r^3\right)
\sin z\cos z.}
\]
The curl is nonzero on open subpatches with $0<|z|<\pi/6$.
Thus this carrier is not a pressure-gauge adjustment.
\end{proposition}
\begin{proof}
The cylindrical divergence of an axisymmetric pure radial
tensor gives $(K_r+K/r)e_r$.
Take its axial derivative and use the displayed formula for $K$.
A gradient has zero curl.
\end{proof}

\paragraph{Physical and independent-verification limits.}
This is an exact local positive cone completion, not the full
H8 four-field stress or a source-matched unforced PDE solution.
A positive radial carrier has a physical projected-divergence
effect and its construction must be justified; a pressure gauge
alone cannot supply it. The co-rotating ray frame also generates
spatial derivative terms. Source A.11/A.19, collar matching,
the seven-ray virtual-work dual, time evolution, all-mode
convergence and singularity inheritance remain unproved.
Independent symbolic reconstruction of all six tensor entries,
the rank-six tensor map with coefficient kernel
$(2,-2,0,-1,-1,1,1)$, five thousand exact rational parameter
replays, and an independent carrier-curl derivative audit pass.
The accompanying Lean declarations cover only finite scalar
identities, not the full PDE theorem.



\section{A global annular seven-ray quotient-cone completion and the axial localization obstruction}
\label{sec:xxvi-global-ray}

The exact seven-ray synthesis of Edition XXV has a stronger
source-independent consequence on a periodic-axial cylinder. It
also has a sharp limitation when one tries to use a compact
three-dimensional H8 source shell. Neither statement proves a
zero-force Navier--Stokes singularity.

\begin{theorem}[Global smooth positive seven-ray completion modulo pressure]
\label{thm:xxvi-global-cone}
Work on $\Omega=\mathbb R^2\times\mathbb S^1$ in cylindrical
coordinates $(r,\theta,z)$, or inside an embedded transverse
disk of $\mathbb T^3$. Let $Q(r,z)$ be any smooth symmetric
$3\times3$ co-rotating tensor, independent of $\theta$, with
its radial support contained in a compact annulus
$0<r_0\le r\le r_1<\infty$.
Choose $\chi\in C_c^\infty((0,\infty);[0,1])$ equal to one
in a neighborhood of the radial projection of $\operatorname{supp}Q$.
Let
\[
M=\max_{i,j}\|Q_{ij}\|_\infty,\qquad
A=16M+1,\qquad B=29A+64M+1.
\]
For the same seven rays $v_j$ of
Theorem~\ref{thm:ray-direct-XXV} (with $t=1/4$), define
\begin{align*}
d_1&=Q_{rr}-16(Q_{\theta\theta}+Q_{zz})+(B-29A)\chi,\\
d_2&=8Q_{\theta\theta}+6A\chi+2Q_{r\theta},\\
d_3&=8Q_{\theta\theta}+8A\chi-2Q_{r\theta},\\
d_4&=8Q_{zz}+7A\chi+2Q_{rz}-8Q_{\theta z},\\
d_5&=8Q_{zz}+7A\chi-2Q_{rz}+8Q_{\theta z},\\
d_6&=A\chi+8Q_{\theta z},\\
d_7&=A\chi-8Q_{\theta z}.
\end{align*}
Then all coefficients are smooth and radially compactly supported,
\begin{equation}\label{eq:xxvi-cone}
\boxed{\quad d_j\ge\chi\ge0\ (1\le j\le7),\qquad
\sum_{j=1}^7d_jv_j\otimes v_j
=Q+A\chi I+B\chi e_r\otimes e_r.\quad}
\end{equation}
The resulting tensor $S$ is globally smooth and radially compactly
supported, and its physical Cartesian tensor divergence satisfies
\begin{equation}\label{eq:xxvi-pressure-quotient}
\boxed{\quad
\nabla\cdot S=\nabla\cdot Q+\nabla\Phi,\qquad
\Phi(r)=(A+B)\chi(r)+B\int_0^r\frac{\chi(s)}s\,ds.\quad}
\end{equation}
In particular $\mathbb P\nabla\cdot S
=\mathbb P\nabla\cdot Q$ in the periodic-axial or periodic-torus
Leray quotient.
\end{theorem}
\begin{proof}
For the seven rays of Edition XXV the six components satisfy
\[
S_{rr}=\sum d_j,\quad
16S_{\theta\theta}=d_2+d_3+d_6+d_7,\quad
16S_{zz}=d_4+d_5+d_6+d_7,
\]
\[
4S_{r\theta}=d_2-d_3+d_6+d_7,\quad
4S_{rz}=d_4-d_5+d_6-d_7,\quad
16S_{\theta z}=d_6-d_7.
\]
Direct substitution of the displayed coefficients yields exactly
\eqref{eq:xxvi-cone} in all six components.

Where $Q\ne0$ one has $\chi=1$.
Using $|Q_{ij}|\le M$ gives
\[
d_1\ge31M+1,\quad d_2\ge86M+6,\quad
d_3\ge118M+8,\quad
d_{4,5}\ge94M+7,\quad
d_{6,7}\ge8M+1.
\]
Where $\chi\ne1$, $Q=0$, and the seven coefficients are
$(64M+1)\chi$, $6A\chi$, $8A\chi$, $7A\chi$,
$7A\chi$, $A\chi$, $A\chi$, respectively.
Every coefficient is therefore at least $\chi$.
The co-rotating frame causes no axis singularity because $\chi$
and $Q$ vanish near $r=0$.

Finally, on the full cylindrical domain
\[
\nabla\cdot(B\chi e_r\otimes e_r)
=B(\chi'+\chi/r)e_r
=\nabla\left(B\chi+B\int_0^r\chi(s)/s\,ds\right).
\]
The isotropic term contributes $\nabla(A\chi)$.
The scalar potential is smooth at the axis, becomes constant
outside the outer radial support, and patches to a periodic
potential on the embedded-disk three-torus.
This proves \eqref{eq:xxvi-pressure-quotient}.
\end{proof}

\begin{proposition}[Necessary carrier amplitude in any strict representation]
\label{prop:xxvi-gauge-cost}
If, at a point with $\chi>0$,
\[
Q+A\chi I+B\chi e_r\otimes e_r
=\sum_{j=1}^7d_jv_j\otimes v_j,
\qquad d_j\ge\mu\chi,
\]
then
\begin{equation}\label{eq:xxvi-gauge-cost}
\boxed{\quad B\chi\ge Q_{\theta\theta}-Q_{rr}
                  +\frac{105}{16}\mu\chi.\quad}
\end{equation}
Thus the free choice of large $A,B$ in
Theorem~\ref{thm:xxvi-global-cone} cannot be replaced by
an asserted source-compatible H8 gain or coefficient-size budget.
\end{proposition}
\begin{proof}
Each ray satisfies $(v_j)_r=1$ and
$|(v_j)_\theta|\le1/4$. Therefore
$S_{rr}-S_{\theta\theta}\ge(15/16)\sum_jd_j
\ge(105/16)\mu\chi$.
The isotropic gauge cancels from the left side, leaving
$Q_{rr}-Q_{\theta\theta}+B\chi$.
\end{proof}

\begin{proposition}[No compact axially localized pure-radial gauge on $\mathbb R^3$]
\label{prop:xxvi-axial-no-go}
Let $k(r,z)$ be smooth and compactly supported in both $r$
(away from the axis) and $z$ on $\mathbb R^3$.
If $\nabla\cdot(k e_r\otimes e_r)$ is curl-free, then
$k\equiv0$. In particular, simply multiplying the gradient-neutral
carrier of Theorem~\ref{thm:xxvi-global-cone} by a compact axial
cutoff introduces a non-gradient force unless it vanishes.
\end{proposition}
\begin{proof}
The carrier divergence is $(k_r+k/r)e_r$ and the
azimuthal component of its curl is
$\partial_z(k_r+k/r)$.
Curl-freeness and compact axial support imply
$k_r+k/r=0$ for every $(r,z)$.
Thus $\partial_r(rk)=0$, so $rk=C(z)$.
Compact radial support forces $C(z)=0$.
\end{proof}

\paragraph{Independent double verification and release boundary.}
Six general tensor component identities were recomputed symbolically
and $10{,}000$ exact rational tests included arbitrary symmetric
tensors (not only covariance matrices), the active support,
and the cutoff collar. An independent Cartesian derivative check
verified a radial carrier's divergence is a gradient. A separate
axial-cutoff calculation found an explicit nonzero curl and
reproduced Proposition~\ref{prop:xxvi-axial-no-go}.
This global quotient-cone representation is mathematically exact
but uses large unconstrained gauges: $B=29(16M+1)+64M+1$.
It does \emph{not} give the small coefficient/stress norms required
in the real H8 source, prove the physical A.11/A.19 residual,
produce a time-dependent unforced solution, or verify a terminal
singularity. The original forced Navier--Stokes construction
cannot be made unforced by this algebra alone.



\section{Quantitatively calibrated seven-ray completion}
\label{sec:xxvii-calibrated}

Edition XXVI proves pressure-quotient realizability using large
source-independent gauge constants. Here we substantially reduce
those constants and prove a sharp zero-source floor. All statements
are exact on a periodic-axial cylinder, or on an embedded annular
disk of a three-torus; no assertion of a compact axial gauge on
$\mathbb R^3$ is made.

\begin{theorem}[Improved universal gauge with an arbitrary cone floor]
\label{thm:xxvii-universal}
Let $Q(r,z)$ be a smooth symmetric co-rotating tensor, radially
compactly supported away from the axis, and let $\chi(r)$ be a
smooth radial cutoff equal to one throughout the radial
projection of $\operatorname{supp}Q$. Write
\[
M=\max_{i,j}\|Q_{ij}\|_{L^\infty},\qquad
\mu>0,\qquad
a=\frac{3}{8}\mu+\frac{13}{4}M,\qquad
b=\frac{85}{8}\mu+\frac{471}{4}M,\qquad
h=\mu+8M.
\]
Define the seven coefficients (ray aperture $1/4$)
\begin{align*}
d_1&=Q_{rr}-16(Q_{\theta\theta}+Q_{zz})
          +(\mu+33M)\chi,\\
d_2&=8Q_{\theta\theta}+2Q_{r\theta}
          +(\mu+10M)\chi,\\
d_3&=8Q_{\theta\theta}-2Q_{r\theta}
          +(3\mu+26M)\chi,\\
d_4&=8Q_{zz}+2Q_{rz}-8Q_{\theta z}
          +(2\mu+18M)\chi,\\
d_5&=8Q_{zz}-2Q_{rz}+8Q_{\theta z}
          +(2\mu+18M)\chi,\\
d_6&=8Q_{\theta z}+(\mu+8M)\chi,\\
d_7&=-8Q_{\theta z}+(\mu+8M)\chi.
\end{align*}
Then, everywhere,
\begin{equation}\label{eq:xxvii-completion}
\boxed{\quad
d_j\ge\mu\chi,\qquad
\sum_{j=1}^7d_jv_j\otimes v_j
=Q+a\chi I+b\chi e_r\otimes e_r.\quad}
\end{equation}
The carrier is globally pressure-neutral in the same periodic
geometries as Edition XXVI:
\[
\nabla\cdot\bigl(\sum_jd_jv_j\otimes v_j-Q\bigr)
=\nabla\left((a+b)\chi+
  b\int_0^r\frac{\chi(s)}s\,ds\right).
\]
In particular, every matrix entry of the constructed stress obeys
\begin{equation}\label{eq:xxvii-size}
\boxed{\quad
\left\|\sum_jd_jv_j\otimes v_j\right\|_{\ell^\infty_{\rm entry}}
\le122M+11\mu.\quad}
\end{equation}
\end{theorem}
\begin{proof}
The seven ray tensor identities are those in
Section~\ref{sec:xxvi-global-ray}. For a general
seven-coefficient inverse with middle parameter $h\chi$,
the off-diagonal and diagonal equations give
\[
\begin{aligned}
d_6&=h\chi+8Q_{\theta z},&d_7&=h\chi-8Q_{\theta z},\\
d_2&=8Q_{\theta\theta}+2Q_{r\theta}
      +(8a-2h)\chi,&
d_3&=8Q_{\theta\theta}-2Q_{r\theta}+8a\chi,\\
d_4&=8Q_{zz}+2Q_{rz}-8Q_{\theta z}
      +(8a-h)\chi,&
d_5&=8Q_{zz}-2Q_{rz}+8Q_{\theta z}
      +(8a-h)\chi,\\
d_1&=Q_{rr}-16(Q_{\theta\theta}+Q_{zz})
      +(b-31a+2h)\chi.
\end{aligned}
\]
The chosen values satisfy exactly
\[
8a-2h=\mu+10M,\quad 8a=3\mu+26M,\quad
8a-h=2\mu+18M,\quad
b-31a+2h=\mu+33M.
\]
Where $Q\ne0$, $\chi=1$ and $|Q_{ij}|\le M$.
The worst-case signed contributions to $d_1,\ldots,d_7$
are respectively $-33M,-10M,-10M,-18M,-18M,-8M,-8M$.
Every coefficient therefore exceeds $\mu$. In a cutoff
collar $Q=0$, the displayed positive constants multiply $\chi$,
so $d_j\ge\mu\chi$ globally. The tensor identity follows by
substitution in the six synthesis equations.

The carrier divergence is a gradient as in Edition XXVI.
The only potentially largest entry of
$S=Q+a\chi I+b\chi e_r\otimes e_r$ is the radial diagonal:
\[
|S_{rr}|\le M+a+b=M+121M+11\mu=122M+11\mu.
\]
The other diagonals are at most
$(17/4)M+(3/8)\mu$, and the off-diagonals at most $M$.
\end{proof}

\begin{proposition}[Sharp zero-source cost, and minimax within a fixed-splitting inverse]
\label{prop:xxvii-sharp-floor}
For $Q=0$, $\chi=1$, any strict seven-ray representation
\[
aI+b e_r\otimes e_r=\sum_jd_jv_j\otimes v_j,\qquad
d_j\ge\mu>0
\]
requires
\begin{equation}\label{eq:xxvii-sharp-floor}
\boxed{\quad b\ge\frac{85}{8}\mu.\quad}
\end{equation}
Equality is attained by
$(d_1,\ldots,d_7)=(\mu,\mu,3\mu,2\mu,2\mu,\mu,\mu)$.

Moreover, within the source-independent fixed-$h$ inverse
displayed in the proof of Theorem~\ref{thm:xxvii-universal},
with fixed scalar $a,b,h$ and arbitrary symmetric
$|Q_{ij}|\le M$, one necessarily has
\[
h\ge\mu+8M,\qquad
a\ge\frac{\mu+2h+10M}{8},\qquad
b\ge\mu+31a-2h+33M,
\]
which imply
\begin{equation}\label{eq:xxvii-fixed-minimax}
\boxed{\quad
b\ge\frac{85}{8}\mu+\frac{471}{4}M.\quad}
\end{equation}
The construction of Theorem~\ref{thm:xxvii-universal}
attains this bound within the stated fixed-inverse class.
This is \emph{not} a claim that the pressure quotient has
the same optimal constant among all nonlinear
source-dependent coefficient selections.
\end{proposition}
\begin{proof}
At $Q=0$, the three off-diagonal tensor constraints give
$d_6=d_7=h$, $d_3=d_2+2h$, $d_4=d_5$.
Equality of the two transverse diagonal entries gives
$d_4=d_5=d_2+h$. Hence
\[
a=(d_2+2h)/8,\qquad
b=d_1+\frac{31}{8}d_2+\frac{23}{4}h.
\]
Since $d_1,d_2,h\ge\mu$, this gives the sharp $85\mu/8$
bound and its equality case.

For the fixed inverse in the theorem, testing all signed
corners of the independent box $|Q_{ij}|\le M$ gives
the three inequalities displayed, as well as
$8a\ge\mu+10M$ and $8a-h\ge\mu+18M$.
The first three dominate the latter two for
$h\ge\mu+8M$. Eliminating $a$ and then $h$ yields
\[
b\ge\frac{39}{8}\mu+\frac{287}{4}M
                 +\frac{23}{4}h
\ge\frac{85}{8}\mu+\frac{471}{4}M.
\]
The displayed explicit coefficients attain equality.
\end{proof}

\paragraph{Independent audit, baseline, and real source limitation.}
A distinct symbolic reconstruction checks every one of the
six symmetric tensor components, and $20{,}000$ exact-rational
box tests verify all seven coefficient floors.
An independent linear-programming calculation finds the
fixed-inverse optimum $(a,h,b)=(29/8,9,1027/8)$ at
$(\mu,M)=(1,1)$, and independently confirms the
sharp zero-source constant $85/8$.
The $M$ slope of the carrier decreases from the source-independent
Edition XXVI value $528$ to $471/4=117.75$.
The original A.11/A.19 source remains unauthenticated;
no H8 smallness threshold or all-mode Navier--Stokes
closure follows from a universal pressure-quotient identity.
The true end-to-end unforced theorem and Lean PDE proof remain open.



\section{An optimal universal radial carrier by source-dependent seven-ray splitting}
\label{sec:xxviii-sharp-source-dependent}
The linear-programming optimum of Edition XXVII concerns a
source-independent \emph{fixed} coefficient splitting. Allowing the
isotropic gauge and the redundant seven-ray coefficient to depend
smoothly on the source tensor yields a smaller universal radial-carrier
cost. This result is purely a symmetric-tensor pressure-quotient
lemma on a periodic-axial annulus, not a source-matched unforced
Navier--Stokes theorem.

\begin{theorem}[Sharp universal annular carrier and smooth seven-ray selection]
\label{thm:xxviii-sharp-universal}
Let $Q(r,z)$ be a smooth symmetric co-rotating $3\times3$ tensor
with compact radial support away from the symmetry axis on
$\mathbb R^2\times\mathbb S^1_z$, or in an embedded transverse
disk on $\mathbb T^3$. Let
\[
M=\max_{i,j}\|Q_{ij}\|_\infty,\qquad \mu>0,
\]
and take a smooth radial cutoff $0\le\chi\le1$ equal to one
on a neighborhood of the full radial support of $Q$.
There exist \emph{explicit smooth} coefficients $d_1,\ldots,d_7$
and a smooth scalar $a=a(r,z)$, all radially compactly
supported, such that
\begin{equation}\label{eq:xxviii-exact-stress}
\boxed{
d_j\ge\mu\chi,\qquad
S:=\sum_{j=1}^7d_jv_j\otimes v_j
=Q+aI+B\chi e_r\otimes e_r,\quad
B=\frac{85}{8}\mu+\frac{343}{4}M.
}
\end{equation}
The rays have aperture $1/4$, exactly as in Editions XXV--XXVII.
The reconstructed stress satisfies the explicit bound
\begin{equation}\label{eq:xxviii-supbound}
\boxed{\quad
\|S\|_{L^\infty,\mathrm{entry}}\le90M+11\mu.\quad}
\end{equation}
Moreover,
\begin{equation}\label{eq:xxviii-projected}
\nabla\cdot(S-Q)
=\nabla\left[a+B\chi+B\int_0^r\frac{\chi(s)}s\,ds\right],
\qquad
\mathbb P\nabla\cdot S=\mathbb P\nabla\cdot Q.
\end{equation}
The carrier coefficient $343/4$ is \emph{sharp} for the universal
independent-entry box $|Q_{ij}|\le M$, when the radial coefficient
$B$ is uniform over the source family and the isotropic scalar
$a(Q)$ and all seven ray coefficients may depend on $Q$.
It is not claimed sharp for the smaller class of positive
semidefinite Reynolds covariance matrices or for a fixed
physical A.11/A.19 source.
\end{theorem}

\begin{proof}[Constructive exact-rational proof]
Represent the independent symmetric entries of $Q$ by
$q=(Q_{rr},Q_{\theta\theta},Q_{zz},
Q_{r\theta},Q_{rz},Q_{\theta z})$.
Write $s=(s_1,\ldots,s_6)\in\{-1,+1\}^6$, with the
$64$ vertices ordered lexicographically.
For each $s$ the accompanying frozen table
\texttt{XXVIII\_64\_CORNER\_EXACT\_CERTIFICATES.txt}
gives rational numbers
$(a_s,d_{s,1},\ldots,d_{s,7})$ with
$d_{s,j}\ge0$, $a_s\le13/4$, and the six \emph{exact} identities
\begin{equation}\label{eq:xxviii-vertices}
\sum_{j=1}^7d_{s,j}v_j\otimes v_j
=Q_s+a_sI+\frac{343}{4}e_r\otimes e_r,
\end{equation}
where $Q_s$ is the symmetric matrix with independent
entries $s_i$. These are not floating-point certificates:
all $64$ rows are independently checked with rational
arithmetic by \texttt{verify\_xxviii\_sharp\_cone.py}.

For $M>0$ put $x_i=q_i/M\in[-1,1]$ and use the nonnegative
multilinear barycentric polynomials
\[
w_s(x)=2^{-6}\prod_{i=1}^6(1+s_i x_i),\qquad
\sum_s w_s=1,\qquad
\sum_s w_s s_i=x_i.
\]
Let
$d^0=(1,1,3,2,2,1,1)$ be the source-free ray
coefficient vector. It obeys
\[
\sum_j d_j^0v_j\otimes v_j
=\frac38 I+\frac{85}{8}e_r\otimes e_r.
\]
Define
\begin{align*}
a(r,z)&=\chi(r)\left[\frac38\mu+
 M\sum_s w_s(q/M)a_s\right],\\
d_j(r,z)&=\chi(r)\left[\mu d_j^0+
 M\sum_s w_s(q/M)d_{s,j}\right].
\end{align*}
All functions are smooth, because each $w_s$ is a polynomial
and $Q$ is smooth, and the radial cutoff is identically one
where $Q\ne0$. Weighted addition of
\eqref{eq:xxviii-vertices} proves
\eqref{eq:xxviii-exact-stress} everywhere, including the
radial collar where $Q=0$ and $0<\chi<1$.
The nonnegative weights and certificates prove $d_j\ge\mu\chi$.
For $M=0$, use only the source-free formula. Since
$0\le a_s\le13/4$, one has
$0\le a\le(3\mu/8+13M/4)\chi$.
Combining with $B$ proves
\eqref{eq:xxviii-supbound}; all off-diagonal entries
of $S$ equal those of $Q$. The pressure-equivalence
formula follows by direct tensor differentiation of the
radial-only carrier and the isotropic tensor.

For \emph{sharpness}, define a dual linear functional on
symmetric matrices by
\[
L(H)=-H_{rr}-15H_{\theta\theta}+16H_{zz}
-\frac{31}{4}H_{r\theta}-46H_{\theta z}.
\]
Then $L(I)=0$, $L(e_r\otimes e_r)=-1$, and the exact
seven ray values are
\[
L(v_j\otimes v_j)
=\left(-1,-\frac{31}{8},0,0,0,-\frac{23}{4},0\right).
\]
Hence every strict ray tensor with $d_j\ge\mu$
satisfies $L(S)\le-85\mu/8$. If
$S=Q+aI+B e_r\otimes e_r$, it follows that
\[
B\ge L(Q)+\frac{85}{8}\mu.
\]
Take the admissible signed box corner
$q=(-M,-M,M,-M,-M,-M)$: then
$L(Q)=343M/4$. This proves that the displayed
$B$ is the smallest possible universal radial coefficient,
with a matching feasible exact rational corner certificate.
\end{proof}

\paragraph{Interpretation and remaining source theorem.}
The improvement is from the Edition XXVII $471M/4$
carrier slope to $343M/4$---a $27.2\%$ reduction---and
from $122M+11\mu$ to $90M+11\mu$ in the displayed
entrywise stress bound. Unlike a single local LP solution,
the multilinear formula constructs smooth globally
cutoff-compatible coefficients. However, the original
OpenAI source's actual A.11/A.19 profiles, physical forcing,
three-dimensional compact shell virtual-work conditions,
space--time stress budget, nonlinear all-harmonic correction
and zero-force singularity theorem remain \emph{unproved}.
A finite Lean scalar dual is supplied; the 64-corner
interpolation and physical PDE are not yet formalized in
Lean. The fully unforced Navier--Stokes problem remains open.



\section{The authentic source residual and the necessary nonzero-data gate}
\label{sec:source-startup-XXIX}

This section corrects the target of the unforcing programme.
The source moment identities (A.11) and (A.19) control leading-profile
matching, not the physical forcing after cutoff. In the original
forced Navier--Stokes source \cite{openai2026},
Proposition 10.1, equation (10.4), forms the compactly supported
velocity $u_B=\operatorname{curl}(c{\mathcal A})+c{\mathcal B}e_\theta$,
with a temporal cutoff giving $u_B(0)=0$. Equation (10.5)
then defines the actual force
\[
f_B=\partial_t u_B+(u_B\cdot\nabla)u_B-\nu\Delta u_B+\nabla p_B.
\]
Appendix A is not a certificate that this vector residual,
or its divergence-free Leray projection, is zero.

\begin{theorem}[Projected residual and compulsory unforced startup]
\label{thm:xxix-startup}
Let $\Omega=\mathbb R^3$ or $\mathbb T^3$, $\nu>0$, and
$u_B$ be any smooth divergence-free finite-energy field
on $[0,T)$ with $u_B(0)=0$. Let $w$ be smooth
and divergence-free and put $u=u_B+w$.
For an arbitrary smooth pressure $p$, define the residual
\[
R=\partial_tu+(u\cdot\nabla)u-\nu\Delta u+\nabla p,
\]
and assume $\mathbb P R\in L^1([0,t];L^2(\Omega))$
for each $t<T$. Then
\begin{equation}\label{eq:xxix-source-budget}
\boxed{\quad
\|u_B(t)\|_2\le
\|w(0)\|_2+
\int_0^t\|\mathbb P R(s)\|_2\,ds+
\|w(t)\|_2.\quad}
\end{equation}
In particular, if $w(0)=0$ and the completed field solves
the \emph{unforced} Navier--Stokes equations, i.e.
$\mathbb P R\equiv0$ with the pressure reconstructed
appropriately, then $u\equiv0$ and necessarily
\begin{equation}\label{eq:xxix-no-go}
\boxed{\qquad w=-u_B\qquad\text{on }[0,T).}
\end{equation}
Consequently a convergent retarded perturbation series
with zero initial correction cannot remove the entire
forced source while retaining any nonzero part of
the forced velocity. This does not rule out a nonzero
unforced initial datum.
\end{theorem}
\begin{proof}
Apply the orthogonal Leray projection to the residual
equation. Testing against $u$ and integrating by parts
over space yields
\[
\frac12\frac{d}{dt}\|u(t)\|_2^2+
\nu\|\nabla u(t)\|_2^2
=\langle\mathbb P R(t),u(t)\rangle.
\]
The usual regularization of $\|u(t)\|_2$ at its zeros
and Cauchy--Schwarz imply
\[
\|u(t)\|_2\le\|u(0)\|_2+
\int_0^t\|\mathbb P R(s)\|_2\,ds.
\]
Since $u(0)=w(0)$ and
$u_B(t)=u(t)-w(t)$, the triangle inequality gives
\eqref{eq:xxix-source-budget}.
If $w(0)=0$ and $\mathbb P R=0$, the energy identity
gives $\|u(t)\|_2=0$ for every $t<T$. This proves
\eqref{eq:xxix-no-go}.
\end{proof}

\begin{corollary}[A quantitative non-gauge forcing floor for the authentic forced velocity]
\label{cor:xxix-real-force}
For the nontrivial compact forced source solution $u_B$
of Theorem 1.1 and Proposition 10.1 in \cite{openai2026},
at any fixed $t_*\in(0,1)$ at which
$\|u_B(t_*)\|_2=m_*>0$, its actual
Section 10 residual satisfies
\begin{equation}\label{eq:xxix-real-source-floor}
\boxed{\quad
\int_0^{t_*}\|\mathbb P f_B(s)\|_2\,ds\ge m_*>0.
\quad}
\end{equation}
Thus the entire physical forcing cannot be a pressure
gradient. Pure pressure-equivalent seven-ray modifications
of Editions XXV--XXVIII cannot remove this projected
force while leaving $u_B$ unchanged.
\end{corollary}
\begin{proof}
Apply Theorem~\ref{thm:xxix-startup} with $u=u_B$,
$w=0$, and $R=f_B$. Existence of a fixed positive
$t_*$ follows from the source theorem's nontrivial
late-time velocity, although its numerical $L^2$
value is not supplied by the Appendix A moments.
\end{proof}

\begin{proposition}[Exact three-dimensional periodic shear negative control]
\label{prop:xxix-shear}
On the spatial torus of period $2\pi$ in every coordinate,
take $\nu>0$ and, for $0\le t\le t_*$,
\[
u_B(x,y,z,t)=(t\sin y,0,0),\quad p_B=0.
\]
It is divergence-free, its convection vanishes identically,
and the complete physical residual is
\[
f_B=(1+\nu t)\sin y\,e_x.
\]
Moreover
\[
\|u_B(t)\|_2=2\pi^{3/2}t,\quad
\int_0^t\|f_B(s)\|_2\,ds
=2\pi^{3/2}(t+\nu t^2/2),
\]
and the exact kinetic-energy equality is
\[
\|u_B(t)\|_2^2+
2\nu\int_0^t\|\nabla u_B(s)\|_2^2\,ds
=2\int_0^t\langle f_B(s),u_B(s)\rangle\,ds.
\]
The retarded correction $w=-u_B$ begins from zero and
makes the completed field the trivial unforced solution,
not a singular flow.
\end{proposition}
\begin{proof}
Differentiate the sine mode explicitly:
$\Delta u_B=-t\sin y\,e_x$,
$(u_B\cdot\nabla)u_B=0$. The spatial integral
$\int_{\mathbb T^3}\sin^2y\,dx\,dy\,dz=4\pi^3$
gives the displayed norms and equality.
\end{proof}

\paragraph{Corrected theorem dependency and formalization boundary.}
The next actual PDE input is the complete Section 10
cutoff residual $f_B$ and its Leray projection,
not only the A.11/A.19 moment-root data.
A proposed smooth unforced correction must start with
nonzero solenoidal initial velocity $w(0)$, must satisfy
the exact projected equation throughout the early-time
cutoff interval and terminal region, and must preserve
the intended singularity despite the energy budget.
The finite Lean component accompanying this edition
formalizes only scalar and normed-space inequalities;
it does not formalize the Navier--Stokes energy identity,
Leray projection, source expressions or unforced
singularity. The exact three-dimensional manufactured
shear and independent 50,000 rational budget fixtures
pass their numerical/symbolic audit.
Nothing in this section proves the Millennium problem.


\section{Heat-semigroup spectral barriers for nonzero-data unforcing}
\label{sec:xxx-spectral-barrier}

The forced velocity of Section~10 of \cite{openai2026} vanishes
initially. Edition~XXIX proved that a correction also starting from
zero cannot produce a nontrivial smooth unforced flow. We now develop
an independently checkable, frequency-resolved \emph{necessary}
condition for a correction starting from nonzero data. These estimates
are elementary Fourier--Duhamel consequences of the energy inequality,
not a new regularity theorem.

\begin{theorem}[Exact periodic Fourier energy barrier]
\label{thm:xxx-periodic}
Let $u$ be a smooth unforced three-dimensional incompressible
Navier--Stokes solution with viscosity $\nu>0$ on the torus
$\mathbb T^3=(\mathbb R/2\pi\mathbb Z)^3$. Normalize spatial volume
to one, put $E_0=\|u(0)\|_{L^2}^2$ and
$\widehat u_k=(2\pi)^{-3}\int_{\mathbb T^3}u(x)e^{-ik\cdot x}\,dx$.
For every nonzero $k\in\mathbb Z^3$ and $0\le t<T$,
\begin{equation}\label{eq:xxx-standard-mode}
\boxed{
|\widehat u_k(t)|\le e^{-\nu|k|^2t}|\widehat u_k(0)|+
\frac{1-e^{-\nu|k|^2t}}{\nu|k|}E_0.}
\end{equation}
If $\int_{\mathbb T^3}u(0)\,dx=0$, the stronger estimate holds:
\begin{equation}\label{eq:xxx-mean-zero-mode}
\boxed{|\widehat u_k(t)|\le
 e^{-\nu|k|^2t}|\widehat u_k(0)|+E_0G_{\nu,k}(t),}
\end{equation}
with
\[
G_{\nu,k}(t)=\begin{cases}
\displaystyle |k|\frac{e^{-2\nu t}-e^{-\nu|k|^2t}}
{\nu(|k|^2-2)},&|k|^2\ne2,\\[5pt]
\displaystyle |k|t e^{-2\nu t},&|k|^2=2.
\end{cases}
\]
Both formulas remain valid at all finite times before any alleged
singularity. They impose necessary conditions, not sufficient ones.
\end{theorem}
\begin{proof}
The projected Fourier equation reads
\[
\frac{d}{dt}\widehat u_k+\nu|k|^2\widehat u_k
=-i\mathbb P_k\sum_{\ell\in\mathbb Z^3}
 (k\cdot\widehat u_\ell)\widehat u_{k-\ell}.
\]
The right side has modulus at most
$|k|\sum_\ell|\widehat u_\ell||\widehat u_{k-\ell}|
\le |k|\sum_\ell|\widehat u_\ell|^2=|k|E(t)$,
by Cauchy--Schwarz and Parseval. The smooth energy identity gives
$E(t)\le E_0$. Integrating against the heat multiplier proves
\eqref{eq:xxx-standard-mode}. For mean-zero data, the mean stays
zero, Poincar\'e's inequality on the $2\pi$ torus gives
$E'(t)\le-2\nu E(t)$, hence $E(t)\le E_0e^{-2\nu t}$.
Integrating $|k|E_0e^{-\nu|k|^2(t-s)}e^{-2\nu s}$
proves the second formula, including the resonant
$|k|^2=2$ limit.
\end{proof}

\begin{theorem}[Whole-space dyadic shell and residual tradeoff]
\label{thm:xxx-whole-space}
Let $u=u_B+w$ be a smooth divergence-free finite-energy field on
$\mathbb R^3\times[0,T)$, with $u_B(0)=0$, $\nu>0$, and physical
residual
$R=\partial_tu+(u\cdot\nabla)u-\nu\Delta u+\nabla p$.
For $K>0$ let $\mathcal B_K$ be the $L^2$ Fourier projector onto
$\{\xi:K\le|\xi|\le2K\}$. Put
\[
x=\|w(0)\|_2,\quad \mathcal R_t=\int_0^t\|\mathbb P R(s)\|_2ds,
\quad b_K(t)=\|\mathcal B_Ku_B(t)\|_2,
\quad\delta_K(t)=\|\mathcal B_Kw(t)\|_2.
\]
Assume the standard smooth-solution energy identity and appropriate
spatial decay/integrability. Then
\begin{equation}\label{eq:xxx-physical-tradeoff}
\boxed{\quad
[b_K(t)-\delta_K(t)]_+\le
 e^{-\nu K^2t}x+
 \frac{\sqrt K}{\sqrt2\,\pi\nu}(x+\mathcal R_t)^2+
 \mathcal R_t.\quad}
\end{equation}
In particular, for an \emph{exactly unforced} completed flow
($\mathcal R_t=0$), writing
$a=e^{-\nu K^2t}$, $c=\sqrt K/(\sqrt2\pi\nu)$ and
$h=[b_K(t)-\delta_K(t)]_+$, a necessary explicit lower bound is
\begin{equation}\label{eq:xxx-nonzero-startup-floor}
\boxed{\quad
\|w(0)\|_2\ge
\frac{2h}{a+\sqrt{a^2+4ch}}\qquad(h>0).\quad}
\end{equation}
This lower bound applies directly to the authentic Section 10
background once $b_K$ has been sourced and independently certified.
It does not require pointwise evaluation of a whole-space Fourier
transform.
\end{theorem}
\begin{proof}
Use the unitary whole-space Fourier transform
$\widehat u(\xi)=(2\pi)^{-3/2}\int e^{-ix\cdot\xi}u(x)dx$.
For the projected quadratic nonlinearity,
\[
\left|\widehat{\mathbb P\operatorname{div}(u\otimes u)}(\xi)\right|
\le(2\pi)^{-3/2}|\xi|\|u\|_2^2,
\]
because $|\xi\cdot u|\,|u|\le|\xi|\,|u|^2$ under the
Fourier integral and $\|\mathbb P_\xi\|_{\mathrm{op}}=1$.
The projected residual equation and smooth energy inequality imply
$\sup_{s\le t}\|u(s)\|_2\le x+\mathcal R_t$.
Duhamel's formula, Minkowski, and Plancherel now yield
\[
\|\mathcal B_Ku(t)\|_2\le
 e^{-\nu K^2t}x+\mathcal R_t+
 \frac{(2\pi)^{-3/2}}\nu(x+\mathcal R_t)^2
 \left(\int_{K\le|\xi|\le2K}|\xi|^{-2}d\xi\right)^{1/2}.
\]
The shell integral is exactly $4\pi K$, giving
$\sqrt K/(\sqrt2\pi\nu)$. The triangle inequality
$\|\mathcal B_Ku_B\|_2\le\|\mathcal B_Ku\|_2+
\|\mathcal B_Kw\|_2$ proves
\eqref{eq:xxx-physical-tradeoff}. For $\mathcal R_t=0$,
solve $cx^2+ax-h\ge0$ for $x\ge0$ and rationalize the
positive root to obtain \eqref{eq:xxx-nonzero-startup-floor}.
\end{proof}

\begin{proposition}[Exact shear controls and limits of inference]
\label{prop:xxx-controls}
For each integer $n\ge1$, the periodic shear
$u=Ae^{-\nu n^2t}\sin(ny)e_x$ is an exact unforced
solution and obeys the purely linear part of the Fourier bound
with no quadratic transfer. In contrast
$u_B=t\sin(ny)e_x$ has projected physical force
$(1+\nu n^2t)\sin(ny)e_x$ and begins from zero; it is not an
unforced solution. Further, no pointwise whole-space bound of the
form $|\widehat u_0(\xi_0)|\le C\|u_0\|_2$ with a universal
constant $C$ is available for arbitrary Schwartz data; one must
use Fourier shells or a specified test function.
\end{proposition}
\begin{proof}
The shear fields are divergence-free and their convection term
vanishes identically. The heat equation gives the first field;
straight differentiation gives the stated force for the second.
For the Fourier pointwise claim, normalized smooth frequency bumps
supported in balls of radius $\varepsilon$ have unit $L^2$
norm but pointwise height of order $\varepsilon^{-3/2}$.
\end{proof}

\paragraph{Adversarial verification and limitations.}
An independently written exact SymPy replay checked the heat
convolution at $|k|^2=2$, the periodic shear equations, and the
unitary Fourier-shell constant. Forty thousand rational Cauchy
and energy-budget fixtures passed; a second independent audit
tested 1,200 divergence-free finite Fourier fields at four
wavevectors. This is a necessary condition; the constant in the
whole-space shell bound grows like $\sqrt K$, so it does not alone
exclude a shrinking-core singularity. No actual Fourier-shell
mass $b_K(t)$ or correction estimate $\delta_K(t)$ has yet been
certified for the complete authentic A.11/A.19-corrected source.
The end-to-end Navier--Stokes PDE and Lean formalization remain open.


\section{Authentic Fourier-shell witness and independent correction energy bound}
\label{sec:xxxi-shell}
The estimates below are for OpenAI's actual Section~10 localized forced velocity
\cite{openai2026}, not an invented profile. Its source is the complete
force in (10.5), including every space--time cutoff derivative.

\begin{lemma}[Zero mean of the authentic source]
For every $0\le t<1$, the velocity in (10.4),
$u_B=\operatorname{curl}(c\mathcal A)+c\mathcal B e_\theta$,
is compactly supported, divergence-free and has exactly zero spatial mean.
Its support is contained in $\{r<r_0,\ |z|<z_0\}$.
One may take $V=2\pi r_0^2z_0$ and $R=\sqrt{r_0^2+z_0^2}$.
\end{lemma}
\begin{proof}
The compact curl integrates to zero. The azimuthal scalar factor is
independent of $\theta$; integrating $e_\theta=(-\sin\theta,\cos\theta,0)$
over one rotation also gives zero. The support follows from the
explicit cutoff conditions immediately before (10.4).
\end{proof}

\begin{theorem}[Quantitative guaranteed Fourier shell for the actual source]
Fix $t<1$ at which $m=\|u_B(t)\|_2>0$, and put
$G=\|\nabla u_B(t)\|_2$. For unitary Fourier transform define
\[
L=\left(\frac{5\pi^2}{4R^2V}\right)^{1/5},\quad
H=\max\{2L,4G/m\},\quad N=\lceil\log_2(H/L)\rceil,\quad K_j=2^jL.
\]
For at least one $0\le j<N$ the whole-space Fourier shell satisfies
\begin{equation}\label{eq:xxxi-lower}
\boxed{\quad \|\mathbf1_{K_j\le|\xi|\le2K_j}\widehat u_B(t)\|_2
\ge m\sqrt{\frac{13}{16N}}>0.\quad}
\end{equation}
Moreover, from the actual source asymptotic (10.21), for sufficiently
small $\tau>0$ define $J_\tau=\frac12e_0\tau^{-1/2-h}$ and
$D_\tau=\|\nabla u_B(1-\tau)\|_\infty$. Then
\begin{equation}\label{eq:xxxi-source-mass}
\boxed{\quad
m(1-\tau)\ge\frac{J_\tau}{2}\sqrt{\frac{4\pi}{3}}
\left(\frac{J_\tau}{2D_\tau}\right)^{3/2}>0.\quad}
\end{equation}
These bounds are explicit in source-defined quantities, not numerical
certificates of a uniquely fixed profile.
\end{theorem}
\begin{proof}
Zero mean and compact support imply
$|\widehat u_B(\xi)|\le(2\pi)^{-3/2}|\xi|R\sqrt V\,m$.
Integrating over $|\xi|<L$ gives at most
$R^2VL^5m^2/(10\pi^2)=m^2/8$.
Plancherel bounds the energy above $H$ by $G^2/H^2\le m^2/16$.
Hence the disjoint dyadic shells from $L$ through $2^NL$ carry at
least $13m^2/16$, and pigeonholing proves the claim.
The source's positive azimuthal amplitude at $x_\tau$ is at least
$J_\tau$; the global Lipschitz constant $D_\tau$ keeps its
magnitude above $J_\tau/2$ in a ball of radius
$J_\tau/(2D_\tau)$. Integrating over this ball gives the last bound.
\end{proof}

\begin{theorem}[Independent physical correction shell upper bound]
Let $u=u_B+w$ be a smooth finite-energy exactly unforced
Navier--Stokes comparison solution on $[0,t]$ and $\nu>0$.
Write $x=\|w(0)\|_2$, $f_B$ for the full source (10.5),
\[
\begin{gathered}
U=\sup_{s\le t}\|u_B(s)\|_2,\quad
A=\int_0^t\|\nabla u_B(s)\|_\infty ds,\quad
F=\int_0^t\|\mathbb P f_B(s)\|_2ds,\quad W=e^A(x+F),\\
F_{K,t}=\int_0^t e^{-\nu K^2(t-s)}
\|\mathbf1_{K\le|\xi|\le2K}\widehat{\mathbb P f_B}(s)\|_2ds,\quad
a_K=e^{-\nu K^2t},\quad c_K=\frac{\sqrt K}{\sqrt2\pi\nu}.
\end{gathered}
\]
Then
\begin{equation}\label{eq:xxxi-upper}
\boxed{\quad
\sup_{s\le t}\|w(s)\|_2\le W,\qquad
\delta_K(t)\le a_Kx+F_{K,t}+c_K(2UW+W^2).\quad}
\end{equation}
Combining with the Edition~XXX unforced shell estimate yields
\begin{equation}\label{eq:xxxi-tradeoff}
\boxed{\quad
b_K(t)\le2a_Kx+c_Kx^2+F_{K,t}+c_K(2UW+W^2).\quad}
\end{equation}
\end{theorem}
\begin{proof}
Subtract the forced PDE for $u_B$ from the unforced PDE for $u$.
Testing the correction equation against $w$ eliminates the
$(u\cdot\nabla)w$ term and bounds the cross-advection
by $\|\nabla u_B\|_\infty\|w\|_2^2$.
Gronwall proves $\|w\|_2\le W$.
Projected Duhamel applied to the three tensors
$u_B\otimes w$, $w\otimes u_B$, $w\otimes w$ yields
\[
\|\mathcal B_Kw(t)\|_2\le a_Kx+F_{K,t}
+\frac{2UW+W^2}{(2\pi)^{3/2}\nu}
\left(\int_{K\le|\xi|\le2K}|\xi|^{-2}d\xi\right)^{1/2}.
\]
The annular integral is $4\pi K$. The last claim follows
from the triangle inequality and Edition~XXX.
\end{proof}

\begin{corollary}[Conditional retention-based initial-data test]
Assume in addition that the correction actually preserves a
fraction of the selected shell:
$\delta_{K_j}(t)\le\kappa b_{K_j}(t)$ for some $\kappa<1$.
Put $B=m\sqrt{13/(16N)}$. Then every unforced comparison flow
satisfies
\begin{equation}\label{eq:xxxi-retention}
\boxed{\quad x\ge\min_{0\le j<N}
\frac{2(1-\kappa)B}
{a_{K_j}+\sqrt{a_{K_j}^2+4c_{K_j}(1-\kappa)B}}.\quad}
\end{equation}
\end{corollary}
\begin{proof}
At the shell selected by the lower theorem, retention and the
Edition~XXX unforced bound give $(1-\kappa)B\le a_{K_j}x+c_{K_j}x^2$.
Solve this quadratic and take the minimum over the finite list.
\end{proof}

\paragraph{Limitations and release gate.}
The authentic source gives neither frozen numerical values of
$(e_0,r_0,z_0,D_\tau,G,F_{K,t})$ nor the independent spectral
retention inequality. The unconditional correction upper bound
may be very weak near the singular time; in particular $x=0$,
$u=0$, $w=-u_B$ satisfies it while cancelling the entire blowup.
Thus the source-backed shell lower bound and independently proved
correction upper bound do not by themselves reject a specific
nonzero initial-data budget. They are a test protocol, not a
complete unforced Navier--Stokes construction or Lean PDE proof.



\section{The missing numerical source and a rigorous spectral-retention no-go}
\label{sec:xxxii-source-retention}
The actual 166-page OpenAI source \cite{openai2026}, Section~10,
(10.3)--(10.5), and Appendix~B, Proposition~B.2, specifies
existence parameters and arbitrarily chosen smooth cutoffs,
not a single frozen numerical $u_B$. For example, with
$D=1/2-h$, the source requires
$0<h<1/100$, $q_*,C_0>0$, and
$C_0(\tau_0+z_0^{1/D})<q_*/2$, where $0<\tau_0<1/2$.
The numerical $h,q_*,C_0$, the nonlinear profile, all cutoffs,
and the complete Section~10 residual are not certified as a
single independent numerical source in this programme.

\begin{proposition}[Conditional admissible geometric choices]
For any externally authenticated $q_*,C_0>0$ and admissible $h$,
\[
r_0=1,\quad \tau_0=\min\{1/4,q_*/(16C_0)\},\quad
z_0=(q_*/(16C_0))^{1/2-h}
\]
satisfy $C_0(\tau_0+z_0^{1/D})\le q_*/8<q_*/2$.
This is an admissible \emph{symbolic cutoff geometry only},
not a numerical realization of the source's velocity.
\end{proposition}
\begin{proof} Substitute the definitions; $\tau_0<1/2$
holds directly. \end{proof}

\begin{proposition}[Exact interval-friendly Fourier-shell certificate]
\label{prop:xxxii-interval}
Suppose externally certified bounds for the actual zero-mean,
compact Section~10 velocity at an explicit $t_*\in(0,1)$ are
\[
0<m_-\le\|u_B(t_*)\|_2,\quad
R^2\le R_+^2,\quad V\le V_+,\quad
\|\nabla u_B(t_*)\|_2/\|u_B(t_*)\|_2\le Q_+.
\]
If rational $L>0$ and integer $N\ge1$ obey
$R_+^2V_+L^5\le45/4$ and $2^NL\ge4Q_+$,
then some $0\le j<N$ satisfies
\[
\boxed{\quad b_{2^jL}(t_*)^2\ge
\frac{13}{16N}m_-^2.\quad}
\]
\end{proposition}
\begin{proof}
Zero mean and compact support yield low-frequency
energy at most $R_+^2V_+L^5\|u_B\|_2^2/(10\pi^2)
\le\|u_B\|_2^2/8$ using $\pi^2>9$.
The gradient ratio gives high-frequency energy above
$2^NL$ at most $\|u_B\|_2^2/16$.
The intermediate $N$ shells thus contain at least
$13\|u_B\|_2^2/16$; pigeonhole.
\end{proof}

\begin{theorem}[No automatic shell retention from nonzero initial data]
\label{thm:xxxii-noret}
Fix $0<t_*<1$ and any Fourier shell with
$b_K(t_*):=\|\mathcal B_Ku_B(t_*)\|_2>0$.
For every $0\le\kappa<1$ there are arbitrarily small
\emph{nonzero} smooth compactly supported divergence-free
initial data with a global smooth exactly unforced
Navier--Stokes solution $v_\varepsilon$ whose correction
$w_\varepsilon=v_\varepsilon-u_B$ violates retention:
\[
\boxed{\quad
\|\mathcal B_Kw_\varepsilon(t_*)\|_2>
\kappa b_K(t_*).\quad}
\]
\end{theorem}
\begin{proof}
Choose any $\phi\in C_c^\infty(\mathbb R^3;\mathbb R^3)$,
$\operatorname{div}\phi=0$, $\|\phi\|_2=1$.
Global smooth small-data well-posedness of three-dimensional
unforced Navier--Stokes applies to
$v_\varepsilon(0)=\varepsilon\phi$ for $\varepsilon>0$
sufficiently small. The unforced energy inequality gives
$\|v_\varepsilon(t_*)\|_2\le\varepsilon$.
Since $\mathcal B_K$ is an $L^2$ contraction,
\[
\|\mathcal B_Kw_\varepsilon(t_*)\|_2
\ge b_K(t_*)-\varepsilon.
\]
Choose $\varepsilon<(1-\kappa)b_K(t_*)$.
This corrector cancels the singular background in the
completed globally smooth solution. It disproves
universal retention, not the possibility of some specially
constructed singularity-preserving corrector.
\end{proof}

\begin{proposition}[One-time shell retention is not blowup persistence]
On the $2\pi$ periodic three-torus with $\nu=1$, the forced
background $u_B=t\sin y\,e_x$ and the smooth, exactly unforced
solution $v=\tfrac14e^{1/2-t}\sin y\,e_x$ satisfy
$w:=v-u_B=-u_B/2$ at $t_*=1/2$.
Thus $\|\mathcal B_1 w(t_*)\|_2=
\|\mathcal B_1u_B(t_*)\|_2/2$, despite $v$ being
globally smooth.
\end{proposition}
\begin{proof}
The transport nonlinearity vanishes for both shears,
and differentiation establishes their claimed equations.
\end{proof}

\begin{proposition}[A sufficient pointwise singularity-retention gate]
For the actual source path
$x_\tau=(\sqrt{2X_{\rm in}\tau},0,0)$ of (10.20),
the source (10.21) gives
$u_{B,\theta}(x_\tau,1-\tau)=
\tau^{-A}(e_0+O(\tau^{2h}))$,
$A=1/2+h$, $e_0>0$.
If a nonzero-data zero-force corrector satisfied
$|w_\theta(x_\tau,1-\tau)|\le
(e_0/4)\tau^{-A}$ for all sufficiently small $\tau$,
the completed velocity would blow up:
\[
|u_\theta(x_\tau,1-\tau)|
\ge(e_0/4)\tau^{-A}\longrightarrow\infty.
\]
\end{proposition}
\begin{proof}
For small $\tau$, the background azimuthal component is at
least $(e_0/2)\tau^{-A}$; apply the triangle inequality.
\end{proof}

\paragraph{Exact verification and release boundary.}
The attached rational interval program rejects the current
source ledger because the source-specific bounds have not
been frozen; its positive acceptance fixture is clearly
synthetic. Another checker reproduces 25,000 exact rational
small-nonzero-data counterexamples and an explicit periodic
smooth retention control. These are mathematics and test
infrastructure, \emph{not} a numerical evaluation of the
actual OpenAI velocity. A physically admissible singularity-
preserving unforced correction, and its full Lean PDE proof,
remain unproved.


\section{A shrinking-core gradient and frequency barrier for unforced corrections}
\label{sec:xxxiii-shrinking}
In the authenticated forced construction of OpenAI \cite{openai2026},
Section~10, equations (10.20)--(10.21), along
$x_\tau=(\sqrt{2X_{\rm in}\tau},0,0)$,
\[
u_{B,\theta}(x_\tau,1-\tau)
=\tau^{-A}(e_0+O(\tau^{2h})),
\quad A=1/2+h,\quad 0<h<1/100,\quad e_0>0.
\]
The background also obeys a uniform whole-space kinetic-energy
bound $B=\sup_{0\le t<1}\|u_B(t)\|_2<\infty$.
These statements follow from the published Section~10 source,
without inventing a numerical cutoff profile.

\begin{lemma}[Two explicit three-dimensional interpolation bounds]
\label{lem:xxxiii-agn}
For a smooth vector field $v$ on $\mathbb R^3$ with
$v\in L^2$ and $\nabla v\in L^\infty$,
\[
\boxed{\quad\|v\|_\infty^5\le
\frac{24}{\pi}\|v\|_2^2\|\nabla v\|_\infty^3.\quad}
\]
For $v\in L^2$ with $\Delta v\in L^2$,
\[
\boxed{\quad\|v\|_\infty\le
C_2\|v\|_2^{1/4}\|\Delta v\|_2^{3/4},
\quad C_2=\frac4{\pi\sqrt2\,3^{7/8}}.\quad}
\]
Here $\|\nabla v\|_\infty$ denotes the operator-norm Lipschitz
constant, and $C_2$ is an admissible explicit constant,
not a claim of best possible interpolation constant.
\end{lemma}
\begin{proof}
At a point $|v(x_0)|=a>0$, the Lipschitz estimate gives
$|v|\ge a/2$ on $B(x_0,a/(2L))$, where
$L=\|\nabla v\|_\infty$. Integrating the squared norm
gives $\|v\|_2^2\ge\pi a^5/(24L^3)$ and then the
first bound by taking suprema.
For the second, use the unitary Fourier transform,
split at $|\xi|=K$, and use Cauchy--Schwarz:
\[
\|v\|_\infty\le
\frac{K^{3/2}}{\pi\sqrt6}\|v\|_2+
\frac{K^{-1/2}}{\pi\sqrt2}\|\Delta v\|_2.
\]
The exact radial integrals are
$\int_{|\xi|\le K}1\,d\xi=4\pi K^3/3$ and
$\int_{|\xi|>K}|\xi|^{-4}d\xi=4\pi/K$.
Choose $K^2=\|\Delta v\|_2/(\sqrt3\|v\|_2)$;
the two terms sum to the stated constant.
\end{proof}

\begin{theorem}[Necessary cancellation cost at the authentic singular trajectory]
\label{thm:xxxiii-cancellation}
Let $u$ be any smooth finite-energy, exactly unforced
three-dimensional Navier--Stokes comparison solution on
$[0,1)$ with $\|u(0)\|_2=x<\infty$.
Write $w=u-u_B$, $M=x+B$. The energy inequality and
triangle inequality give $\|w(t)\|_2\le M$ for all $t<1$.
Suppose that for all sufficiently small $\tau$ the
completed azimuthal velocity is no greater than
$(e_0/4)\tau^{-A}$ along the source path:
\[
|u_\theta(x_\tau,1-\tau)|\le
(e_0/4)\tau^{-A}.
\]
Then both the following lower bounds are \emph{necessary}:
\begin{align}
\|\nabla w(1-\tau)\|_\infty
&\ge\left(\frac{\pi(e_0/4)^5}{24M^2}\right)^{1/3}
\tau^{-5A/3},
\label{eq:xxxiii-gradient-cost}\\
\|\Delta w(1-\tau)\|_2
&\ge\left(\frac{e_0}{4C_2M^{1/4}}\right)^{4/3}
\tau^{-4A/3}.
\label{eq:xxxiii-h2-cost}
\end{align}
In particular, either reversed upper bound with the
same coefficient suffices to preserve a divergent
completed azimuthal velocity along $x_\tau$.
\end{theorem}
\begin{proof}
The source (10.21) is at least $(e_0/2)\tau^{-A}$
for all sufficiently small $\tau$. By the assumed
cancellation, $|w(x_\tau,1-\tau)|\ge
|w_\theta(x_\tau,1-\tau)|\ge(e_0/4)\tau^{-A}$.
Apply Lemma~\ref{lem:xxxiii-agn} to the correction,
using its energy bound $M$, to obtain both inequalities.
If either upper bound holds, the lemma instead bounds
$|w|\le(e_0/4)\tau^{-A}$ and the source's velocity
growth survives by the reverse triangle inequality.
\end{proof}

\begin{theorem}[Necessary high-frequency tail for cancellation]
\label{thm:xxxiii-tail}
With the same notation, define a moving Fourier cutoff
\[
K_\tau=\left(\frac{\pi\sqrt6\,e_0}{8M}\right)^{2/3}
\tau^{-2A/3},\qquad
P_{>K_\tau}=\mathbf 1_{\{|\xi|>K_\tau\}}(D).
\]
Under the cancellation assumption of
Theorem~\ref{thm:xxxiii-cancellation}, the correction
necessarily satisfies
\[
\boxed{\quad
|(P_{>K_\tau}w)(x_\tau,1-\tau)|
\ge(e_0/8)\tau^{-A}.\quad}
\]
Hence $|(P_{>K_\tau}w)(x_\tau,1-\tau)|=o(\tau^{-A})$
is a separate, sufficient spectral criterion for
preservation of the source's singularity.
The threshold scales as
$K_\tau\asymp\tau^{-1/3-2h/3}$, strictly slower
than the inverse radial similarity scale $\tau^{-1/2}$.
\end{theorem}
\begin{proof}
Unitary Fourier inversion and the volume of a ball give
$\|P_{\le K}w\|_\infty\le
MK^{3/2}/(\pi\sqrt6)$.
At $K=K_\tau$ this upper bound equals
$(e_0/8)\tau^{-A}$. The correction must have magnitude
at least $(e_0/4)\tau^{-A}$ at the source point;
the high-frequency piece therefore contributes at least
$(e_0/8)\tau^{-A}$. The exponent comparison follows
from $2A/3=1/3+2h/3<1/2$.
\end{proof}

\paragraph{Scope and release boundary.}
The source-specific ingredients used here are only the
actual Section~10 uniform energy and asymptotic growth,
not unverified numeric A.11/A.19 parameters.
The two Fourier constants, interpolation optimizer,
Gaussian solenoidal positive control, and 30,000 rational
dichotomy fixtures were independently reproduced.
These are necessary cancellation costs and sufficient
conditional retention inequalities, not an existence
proof for a nonzero-data unforced correction.
The exact correction $w=-u_B$ corresponding to the
trivial unforced solution $u=0$ is not excluded and
may satisfy the compulsory growth costs.
The full unforced Navier--Stokes singularity theorem
and kernel-checked Lean PDE proof remain open.


\begin{thebibliography}{9}
\bibitem{newton2026}
T.~Newton, \emph{Exact Local Unforced Terminal-Core Dynamics in Three-Dimensional Navier--Stokes: Anisotropic Concentration and Swirl Reversals}, Zenodo (2026). \url{https://doi.org/10.5281/zenodo.23474915}.
\bibitem{openai2026}
OpenAI, \emph{Finite Time Blowup for Navier--Stokes}, research manuscript (2026), Appendix B.
\url{https://cdn.openai.com/pdf/32d9f210-8b73-45e0-91bc-82a30aef8a9a/navier-stokes.pdf}.
\bibitem{zhang2026analyticity}
P.~Zhang, \emph{3-D Anisotropic vs. Classical Incompressible Navier--Stokes Equations}, Proceedings of the International Congress of Mathematicians 2026, Vol.~5, pp.~388--407 (2026). \url{https://doi.org/10.1137/25M1801980}.
\bibitem{circulationMP2023}
\emph{Quantitative Control of Solutions to the Axisymmetric Navier--Stokes Equations in Terms of the Weak $L^3$ Norm},
Annals of PDE (2023), Section 5, equation (49).
\url{https://doi.org/10.1007/s40818-023-00156-7}.
\end{thebibliography}
\end{document}
